\documentclass[11pt]{article}
\usepackage[review]{acl}
\usepackage{times}
\usepackage{latexsym}
\usepackage[T1]{fontenc}
\usepackage[utf8]{inputenc}
\usepackage{microtype}
\usepackage{inconsolata}
\usepackage{amsmath,amssymb,amsthm,booktabs,graphicx}
\usepackage{algorithm,algorithmic}
\usepackage{placeins}
\newtheorem{theorem}{Theorem}
\newtheorem{proposition}{Proposition}
\newtheorem{lemma}{Lemma}
\newtheorem{corollary}{Corollary}
\newcommand{\one}{\mathbf{1}}
\newcommand{\SD}{\mathbin{\triangle}}
\newcommand{\rbZeroCov}{80.9}
\newcommand{\rbOptCov}{99.6}
\newcommand{\rbActive}{61.4}
\newcommand{\wmtZeroCov}{84.9}
\newcommand{\wmtOptCov}{99.5}
\newcommand{\wmtActive}{45.1}
\newcommand{\exclZeroCov}{99.2}
\newcommand{\exclOptCov}{99.9}
\newcommand{\exclActive}{3.8}

\newcommand{\rbLowActive}{72.9}
\newcommand{\rbMiddleActive}{88.7}
\newcommand{\rbHighActive}{0.5}
\newcommand{\wmtLowActive}{83.5}
\newcommand{\wmtMiddleActive}{89.7}
\newcommand{\wmtHighActive}{0.0}
\newcommand{\exclLowActive}{77.9}
\newcommand{\exclMiddleActive}{44.4}
\newcommand{\exclHighActive}{0.0}
\newcommand{\deMedianActive}{1.3}
\newcommand{\deMeanActive}{83.3}
\newcommand{\esMedianActive}{1.1}
\newcommand{\esMeanActive}{82.1}
\newcommand{\zhMedianActive}{11.3}
\newcommand{\zhMeanActive}{85.9}
\newcommand{\scoreRbMedianCov}{96.4}
\newcommand{\scoreRbMeanCov}{98.8}
\newcommand{\compWmtMedianActive}{2.5}
\newcommand{\compWmtMedianLow}{0.0}
\newcommand{\compWmtMedianHigh}{67.8}

\newcommand{\granRbColumns}{4.93}
\newcommand{\granRbMinimum}{70.4}
\newcommand{\granRbOnlineColumns}{5.63}
\newcommand{\granWmtMinimum}{45.6}

\newcommand{\trueRowWmtMeanQueries}{1,892.9}

\newcommand{\fullMeanProven}{1,920}
\newcommand{\fullMeanUnclosed}{0}
\newcommand{\fullMeanUnionProven}{1,920}
\newcommand{\fullMeanStrictGain}{89}
\newcommand{\fullMeanJointSaving}{6,269}
\newcommand{\frontierQueries}{5210.6}
\newcommand{\frontierRefQueries}{5286.5}
\newcommand{\frontierWidthQueries}{5349.7}
\newcommand{\frontierWins}{403}
\newcommand{\frontierLosses}{554}
\newcommand{\frontierSafetyQueries}{5207.6}
\newcommand{\frontierExclusionQueries}{5343.7}

\newcommand{\dStarClosed}{1,879}
\newcommand{\dStarUnclosed}{41}
\newcommand{\dStarZero}{1,024}
\newcommand{\dStarPositive}{896}
\newcommand{\dStarPositiveNoGain}{807}
\newcommand{\dStarUndetermined}{0}
\newcommand{\dStarTighter}{1,207}

\newcommand{\toleranceRBten}{77.9}
\newcommand{\toleranceWMTten}{137.4}

\title{How Many Evaluator Scores Certify Output Retention Under a Churn Budget?}
\author{Anonymous ACL submission}
\begin{document}
\maketitle
\begin{abstract}
Replacing an evaluator can change which language outputs pass an ensemble cutoff. We study how many candidate-evaluator scores certify a retained subset of selected outputs. Its replacement-average churn must meet a budget, and its size must lie within a stated tolerance of the largest feasible retained count. Normalization, cutoff and a finite family of single-evaluator replacements are fixed; guarantees hold for every assignment of unseen scores consistent with observations. We characterize the smallest certificates under complete score knowledge and lower-bound the queries needed to find one. For odd median pools, safety and exclusion evidence are disjoint, yielding closed-form minima. Mean and even-median pools allow sharing. Under exact certification from empty history, we construct mean-pooling families whose joint certificate approaches half the best union of separately minimum certificates as pools grow. We solved 1,920 targets numerically from 960 cached mean-pooling instances of RewardBench~2 and five WMT panels. Joint gains of at most 2.61\% beyond the best separate union appeared in 89 targets. Whole-column targets showed no joint gain. Finding a two-entry certificate can require queries linear in batch size, even at positive tolerance. Certificate size, discovery cost and inference time require separate evaluation; retention stability does not certify output quality.
\end{abstract}

\section{Introduction}
Reward models, translation metrics and LLM judges score language outputs \citep{malik2026rewardbench2,freitag2023wmt23,zheng2023judging}. Replacing one evaluator can move previously accepted outputs across the same ensemble cutoff. Before rescoring the whole batch, an analyst may retain part of the original list and defer the remaining decisions for review. Deferring every decision trivially meets a budget on these membership changes (churn). Useful certification must also bound the retained-count shortfall relative to the budget-constrained optimum. Two questions follow: how many observed scores suffice to certify near-optimal retention, and how many queries are needed to find them?

We fix the score-normalization maps, the cutoff and the finite family of candidate replacements. A score completion assigns unseen scores consistently with the observations. A certificate bounds replacement-average churn and retained-count loss for every completion. The formulation caps average, not per-replacement, churn and imposes no return-size or per-prompt quotas. Stability differs from quality: evaluator agreement can coexist with shared errors and reward hacking \citep{eisenstein2024herding}.

A certificate shows two things. The retained set meets the churn budget (safety), and no consistent completion admits a feasible set larger by more than the tolerance (exclusion). The central question is whether separately minimum safety and exclusion witnesses, optimally deduplicated over the same safe-count targets, combine into a minimum certificate. Changing the retained subset can lengthen its safety proof while shortening the combined evidence. Odd medians use disjoint agreements and oppositions, yielding closed-form minima. Mean and even-median observations can support both tasks. We characterize joint minima. Under exact stopping from empty history, mean families approach a factor-two advantage over the best separate union as pools grow. Short certificates can require discovery cost linear in batch size even at positive tolerance. On 1,920 mean targets from NLP panels, numerical joint gains occurred in 89 targets and were modest.

\section{Related Work}
\label{sec:related}
\paragraph{Risk control.}
Conformal risk control (CRC) bounds expected monotone loss and risk-controlling prediction sets give high-probability bounds \citep{angelopoulos2024crc,yeh2025conformal,bates2021rcps}. SCoRE controls general bounded losses in selective deployment \citep{bai2026score}. SCOPE and multi-expert CRC address pairwise acceptance and multiple judges \citep{badshah2026scope,cheng2026multiexpert}, and replicable calibration stabilizes thresholds across samples \citep{papamichalis2026replicableconformalprediction}. These guarantees hold under their sampling assumptions. Ours holds for every completion of one fixed batch under declared evaluator replacements; input and order perturbations, which change LLM-judge verdicts \citep{wang2024fair,li2026llmsreliablyjudgeyet}, lie outside it. Ensemble agreement need not exclude common error \citep{afrin2026evaluatorensemblesrewardhacking}.

\paragraph{Verification and selection.}
Priced information and matroid verification distinguish proof size from discovery cost \citep{charikar2002priced,diwan2026basis}. Matroid verification also selects the certified basis \citep{diwan2026basis}. Boolean evaluation studies adaptive queries \citep{deshpande2016sbfe,gkenosis2018ssclass}. Robust selection and knapsack choose queries before uncertain parameters are resolved \citep{chen2026selection,goerigk2015robustqueries}. Uncertain-profit knapsack certifies good packings and excludes better ones under resource augmentation \citep{schloter2025knapsack}. Here replacement-average flip costs couple two proof tasks, safety and count exclusion. Odd medians keep their evidence disjoint, whereas mean pools can beat optimally deduplicated separate minima. Unlike black-box certification \citep{blanc2022certification,gupta2022monotone}, we query a fixed matrix under known pooling. Recoverable selection, covering and query-competition models study later exchanges, sharing and discovery \citep{chassein2017recoverabletwostagerobustselection,megow2023set,erlebach2015survey,erlebach2021parallel}.

\paragraph{Ranking and NLP scoring.}
Ranking robustness covers weight changes, preference completions and word attacks \citep{asudeh2018stable,imber2020topk,wu2022certdr}. PAC selection samples rewards \citep{kalyanakrishnan2012pac}. Uncertain top-$k$ combines rankings with possible worlds or averages completions \citep{soliman2007topk,amarilli2017topk}; our safe and optimistic counts play the role of necessary and possible answers. Col-Bandit queries MaxSim entries \citep{pony2026colbanditquerytimetopkestimation}; budget-adaptive ranking uses joint confidence intervals \citep{oettershagen2026topkbudgetadaptiveranking}. Sorted-list aggregation supports sorted and random accesses \citep{fagin2003optimal}; our completion certificates assume no sorted oracle. Selective ensembles address prediction consistency \citep{black2022selective}; pruning and early exit reduce scoring, and Best-of-Evidence queries factors within a fixed pool \citep{li2022pruning,inoue2019ensemble,busolin2021exit,zhang2026bestofevidencebestofnselectionpartial}. MT filtering \citep{perrella-etal-2024-beyond} and metric biases toward adequacy \citep{shayegh-etal-2025-feeding} motivate separate quality assessment. Predictive churn concerns model updates and training-sample changes \citep{watsondaniels2025predictive,prastalo2026reducing}; fixed-threshold routing can change even with matching class predictions \citep{kim2026lottery}. We instead certify a retained subset within a count tolerance of the largest feasible count and use minimum certificates to benchmark acquisition.

\section{Selection under Evaluator Replacements}
\label{sec:setting}
Certification concerns one partly observed score matrix: rows index outputs and columns index evaluators. Pools $P,R$ contain $M\ge1$ base and $T\ge1$ candidate evaluators, respectively. We fit normalization maps and cutoff $0<t_0\le1$ on a reference split of scored outputs. Reference denotes this fitting split, not a human translation. Write $s_e(i)\in[0,1]$ for normalized scores. Replacing base member $j$ with candidate $r$ gives $P^{j\to r}$; all $H=MT$ replacements have uniform weight. The partly observed matrix, maps and cutoff remain fixed. Base scores are initially visible; queries reveal candidate test scores.

The base list is $S_0(P)=\{i:c_P(i)\ge t_0\}$, where $c_P$ is median or mean consensus. For even-sized pools, the median averages the two central scores; the closed-form median witness results require odd $M$. Cutoff ties are included. For a batch of $N\ge1$ outputs, choose a nominal selection fraction $q\in(0,1]$ and budget $\alpha\ge0$. Let $k=\max(1,\operatorname{round}(qN))$ be the nominal count, rounding half-integers to even. This count scales risk without imposing a return quota. Write $\sigma_Q(i)=\one[i\in S_0(Q)]$. Known base scores and one candidate score determine all $M$ row replacement outcomes. Define their flip count $f_{ir}$ and row total $F_i$:
\begin{align}
 f_{ir}&=\sum_{j=1}^{M}\one[\sigma_P(i)\ne\sigma_{P^{j\to r}}(i)],\nonumber\\
 F_i&=\sum_{r=1}^{T}f_{ir}.
 \label{eq:cost}
\end{align}
The commitment mask $C$ records outputs whose baseline accept or reject decisions are finalized without review; the returned set is $S_0(P)\cap C$. Its average normalized churn over the declared replacements is
\begin{align}
 R(C)&=\frac1H\sum_{j,r}\frac{|(S_0(P)\SD S_0(P^{j\to r}))\cap C|}{2k}\nonumber\\
 &=\frac{\sum_{i\in C}F_i}{2kH}.
 \label{eq:risk}
\end{align}
Decisions outside $C$ are referred for review. Allowing $\rho$ average membership changes per replacement gives $\alpha=\rho/(2k)$. We maximize returned count subject to $R(C)\le\alpha$. Only baseline-selected outputs receive entry queries; an unselected output enters $C$ only when no declared replacement can change its membership. A whole-evaluator query still reveals all $N$ rows.

The integer flip budget is $b_H=\lfloor2k\alpha H\rfloor$. With all costs $F_i$ observed, sort baseline-selected outputs by increasing $F_i$. The cardinality of the longest prefix within $b_H$ is the full-information optimum $V^*$.

\section{Certificates from Partial Scores}
\label{sec:partial}
A certificate uses visible base and queried candidate scores to verify churn feasibility and bound count shortfall. Unobserved scores range freely over $[0,1]$; this is a completion set, not a statistical independence assumption.

Let $f_i(x)$ count flips across the $M$ replacements when a candidate assigns score $x$ to output $i$. Thus $f_{ir}=f_i(s_r(i))$ and $p_i=\max_{x\in[0,1]}f_i(x)$. Under either pooling rule, replacement consensus is nondecreasing in $x$. For selected outputs, $f_i(x)$ is nonincreasing, with its maximum at $x=0$ and zero at $x=1$. For unselected outputs, it is nondecreasing, with its maximum at $x=1$ and zero at $x=0$. Visible base scores determine $p_i$ exactly.

After $n_i$ queries on row $i$, let $\ell_i$ sum their observed flip contributions. Bounding each unobserved count between zero and $p_i$ gives
\begin{equation}
 \ell_i\le F_i\le u_i=\ell_i+(T-n_i)p_i.
 \label{eq:interval}
\end{equation}
For nonnegative costs $w$, define retention capacity
\[
 V(w)=\max\{|C\cap S_0(P)|:\sum_{i\in C}w_i\le b_H\}.
\]
Sorting selected-output costs computes $V(w)$. Maximizing retention under $u$ yields a mask, feasible for every consistent completion, returning $V^-=V(u)$ outputs. The optimistic count $V^+=V(\ell)$ is attainable in some completion.

Keeping observed values fixed, assigning every unobserved entry its zero-flip endpoint realizes $\ell$; assigning these entries their maximum-flip endpoints realizes $u$. Known caps on a candidate's deviation from consensus, or Lipschitz bounds in base-score space, keep both endpoint vectors attainable with narrower entry ranges.
\begin{proposition}[Safe retention and stopping]
\label{prop:stop}
At any adaptive query history, let $C^-$ maximize retained count subject to the upper-cost budget. Then
\[
 R(C^-)\le\alpha,\qquad V^-\le V^*\le V^+.
\]
If $V^+-V^-\le\varepsilon$, its returned count is at most $\varepsilon$ outputs below $V^*$. Under the declared replacement family and completion model, these guarantees hold for every acquisition order and stopping time.
\end{proposition}
\begin{proof}
Equation~\ref{eq:interval} bounds each true cost, so upper-cost feasibility implies true and then lower-cost feasibility. Maximizing cardinality gives the count bounds, and the returned count $V^-$ falls at most the gap below $V^*$.
\end{proof}

Integer $\varepsilon\ge0$ bounds count shortfall; zero requires count optimality. Full revelation gives $V^-=V^+=V^*$.

An output is structurally invariant ($p_i=0$) if no candidate score can change its membership. Zero-query retention uses upper costs $Tp_i$ and can also admit noninvariant outputs. Strategies start there and recompute masks after queries; the safe count never decreases, though retained sets need not be nested. Let $m_0$ and $m$ count invariant and active ($p_i>0$) selected rows; only active rows are queryable. For $Q$ entry queries, active inspection is $Q/(mT)$, or zero when $m=0$ (Appendix~\ref{app:protocol}).

\section{Minimum Certificate Complexity}
\label{sec:median}
To measure how few revealed entries suffice, an offline constructor uses the true matrix to choose entries. The verifier sees only revealed entries and uses bounds valid for every consistent completion. Each entry has unit cost.

\paragraph{Odd medians: disjoint evidence.}
For $M=2h+1$ with nonnegative integer $h$, median selection requires $h+1$ votes at or above $t_0$. A single replacement changes at most one vote. Selected rows with at least $h+2$ votes are invariant interior rows; rows with exactly $h+1$ are boundary rows. A below-cutoff candidate excludes a boundary row precisely when replacing one of its $h+1$ voting members. Thus $p_i=h+1$ at the selected boundary and zero in the interior.

Here the active selected rows form the boundary set $\mathcal B$, and the $m_0$ invariant selected rows are interior. For $i\in\mathcal B$, define opposition bits and their true row sums by
\[
 y_{ir}=\one[s_r(i)<t_0],\qquad z_i=\sum_{r=1}^{T}y_{ir}.
\]
Each opposition causes $h+1$ flips, so $F_i=(h+1)z_i$. Measuring boundary costs in opposition units, feasibility is exactly $\sum_{i\in C\cap\mathcal B}z_i\le b$, where
\begin{equation}
 b=\left\lfloor\frac{b_H}{h+1}\right\rfloor=\left\lfloor\frac{2k\alpha MT}{h+1}\right\rfloor.
 \label{eq:booleanbudget}
\end{equation}
Thus $\alpha$ budgets normalized churn, $b_H$ counts flips, and $b$ counts oppositions. For $M=7$, $T=3$, $k=2$ and $\alpha=.10$, these budgets are $b_H=8$ and $b=2$ (Appendix~\ref{app:worked-example}). Query counts measure observations. On a selected boundary row, an agreement (candidate score at or above the cutoff) reduces the upper opposition bound by one. A below-cutoff opposition increases the lower bound by one. With $a_i$ agreements and $o_i$ oppositions observed, $o_i\le z_i\le T-a_i$.

\begin{theorem}[Safety witness]
\label{thm:safe}
If the true matrix admits a feasible set of $L$ boundary rows, the minimum number of observations needed to certify feasibility of an $L$-row set under every consistent completion is
\[
 Q_{\rm safe}(L)=\max\{0,LT-b\}.
\]
If no $L$-row set is feasible under true costs, no such certificate exists.
\end{theorem}
An $L$-row set starts with upper cost $LT$; only agreements reduce it, one unit each, proving necessity. Any truly feasible $L$-row set contains enough agreements to attain the bound. Appendix~\ref{app:proofs} gives full proofs, endpoints and a worked example.

Feasibility need not establish optimality: a larger retained set may still be possible. Let $z_{(1)}\le\cdots\le z_{(m)}$ and let $L$ be the largest prefix size whose sum is at most $b$. The full retained optimum is $m_0+L$. If $L<m$, then $\sum_{i=1}^{L+1}z_{(i)}>b$, so the threshold $t^*$ below exists.
\begin{theorem}[Exact optimality witness]
\label{thm:exact}
If $L=m$, the minimum certificate proving exact optimality is $\max(0,mT-b)$. Otherwise set $p=L+1$, $B_0=b+1$ and $D(t)=\sum_{i=1}^{p}\min(z_{(i)},t)$. Define
\begin{equation}
 t^*=\min\{t\in\{0,\ldots,T\}:D(t)\ge B_0\}.
 \label{eq:level}
\end{equation}
The exact offline certificate size is
\begin{equation}
 Q_{\rm exact}=\max(0,LT-b)+B_0+(m-p)t^*.
 \label{eq:exact}
\end{equation}
\end{theorem}
The opposition cap $t^*$ differs from cutoff $t_0$. For $L<m$, order rows by true cost. Reveal $b+1$ oppositions across the first $p$ rows, capped at $\min(z_{(i)},t^*)$ in row $i$, and $t^*$ oppositions in every remaining row. These cells are disjoint from safety agreements, so their costs add. For $L=m$, safety alone suffices (Appendix~\ref{app:proofs}).



Let $s$ be the certified safe boundary count. Gap tolerance $\varepsilon$ requires optimistic boundary count at most $s+\varepsilon$, so every $(s+\varepsilon+1)$-row boundary set must be excluded unless $s+\varepsilon+1>m$.
\begin{corollary}[Minimum gap witness]
\label{cor:gap-witness}
Let $L^*$ be the true optimal boundary count and $\varepsilon\ge0$ an integer. For integer $p$ with $L^*<p\le m$, let $D_p(t)=\sum_{i=1}^{p}\min(z_{(i)},t)$ and define
\[
 t_p=\min\{t\in\{0,\ldots,T\}:D_p(t)>b\},
\]
with exclusion cost $E(p)=b+1+(m-p)t_p$. For $p>m$, set $E(p)=0$. The minimum entry certificate size for gap at most $\varepsilon$ under this completion model is
\begin{equation}
 \begin{aligned}
 Q_\varepsilon=\min_{s\in\mathcal S_\varepsilon}\big[&\max(0,sT-b)\\
 &+E(s+\varepsilon+1)\big],
 \end{aligned}
 \label{eq:gap-witness}
\end{equation}
where $\mathcal S_\varepsilon=\{\max(0,L^*-\varepsilon),\ldots,L^*\}$. In particular, $Q_0=Q_{\rm exact}$.
\end{corollary}
These minima lower-bound entry queries under the same instance, budget, model, tolerance and empty history (Appendix~\ref{app:gap-witness}); residual witnesses account for existing observations (Appendix~\ref{app:residual}).


\paragraph{Joint evidence: mean and even medians.}
An entry with true flip contribution $f_{ir}$ raises its row's lower cost by $f_{ir}$ and reduces its upper cost by $p_i-f_{ir}$. Mean pooling can have integer $0<f_{ir}<p_i$, so one query contributes to both proofs. Table~\ref{tab:joint-example} shows why: the five-entry safety witness retains row 2, and the best separate union needs twelve entries. If row 1 is retained instead, its six safety entries also support exclusion, and one more entry completes a seven-entry joint certificate. Appendix~\ref{app:deduplicated} realizes these mean scores and proves minimality.

\begin{table}[t]
\centering\small
\begin{tabular}{@{}lccr@{}}
\toprule
Evidence & Row 1 & Row 2 & Entries\\
\midrule
Minimum safety & \texttt{......} & \texttt{.00000} & 5\\
Minimum exclusion & \texttt{111111} & \texttt{1.....} & 7\\
Best separate union & \texttt{111111} & \texttt{100000} & 12\\
Joint & \texttt{111111} & \texttt{1.....} & 7\\
\bottomrule
\end{tabular}
\caption{Mean evidence with caps $(7,6)$ and flip budget six. Digits are revealed flip counts; dots are unqueried entries. Joint evidence retains row~1 (upper cost six). The lower costs of rows 1 and 2 sum to seven, exceeding the budget and excluding retention of both rows.}
\label{tab:joint-example}
\end{table}

For history $\mathcal O$, let $\ell_i^0,u_i^0$ be its row bounds and $U_i$ its unqueried candidates. Using true contributions offline, define $\mathcal D_i=\{(|J|,\sum_{r\in J}f_{ir}):J\subseteq U_i\}$. Here $n_i$ counts new queries and $v_i$ sums their contributions. A state $(n_i,v_i)$ gives $\ell_i=\ell_i^0+v_i$, $u_i=u_i^0-n_ip_i+v_i$. Let $L^*$ be the true optimal active count.
\begin{theorem}[Joint entry witness]
\label{thm:joint}
For the fixed true matrix, history $\mathcal O$ and integer $\varepsilon\ge0$, the minimum additional entry count $Q_{\rm joint}$ is $\min\sum_i n_i$, over $s\in\{\max(0,L^*-\varepsilon),\ldots,L^*\}$, row states in $\mathcal D_i$ and a safe set $C$ of $s$ active rows satisfying $\sum_{i\in C}u_i\le b_H$. If $p=s+\varepsilon+1\le m$, also require an integer $0\le t\le TM$ with
\[
 \sum_i(t-\ell_i)_+\le pt-b_H-1,
\]
where $(x)_+=\max(0,x)$. For $p>m$, exclusion is automatic.
\end{theorem}
Every successful query set induces reachable row states and a safe count $s$. Lemma~\ref{lem:order-threshold} (Appendix~\ref{app:joint-evidence}) equates exclusion with the $p$ smallest lower costs summing to at least $b_H+1$. Conversely, feasible row states certify at least $s$ safe active rows and optimistic active count at most $s+\varepsilon$. Only the latter condition is automatic when $p>m$. Every successful set is represented, and every feasible state choice is attained by a query set of the same size, proving minimality (Appendix~\ref{app:joint-evidence}). The offline dynamic program is polynomial in explicit $m,T,M$, with a potentially large budget-indexed grid. An equivalent integer program groups equal within-row outcomes into $G$ row--outcome groups, using $O(G+m)$ variables and constraints per safe-count target without guaranteeing polynomial solver time. Appendix~\ref{app:compact-joint} gives the formulations and even-median extension.

\paragraph{What joint optimization adds.}
Let $Q_{\rm split}$ minimize the sum of the separately minimum safety and exclusion sizes over admissible safe counts. For each such count, consider all tied minimum safety and exclusion witnesses. Let $Q_{\cup,\rm sep}$ minimize their union at the same count, then over admissible counts. For the same instance, history and tolerance,
\[
 Q_{\rm joint}\le Q_{\cup,\rm sep}\le Q_{\rm split}\le2Q_{\rm joint}.
\]
The chain holds for any nonnegative replacement weights, with $Q_{\rm joint}$ defined by the weighted integer program (Appendix~\ref{app:model-scope}). Observations only tighten both bounds, and a joint witness serves each task. One realizable mean family attains $Q_{\cup,\rm sep}/Q_{\rm joint}\to2$; a second, with $M=7$ and $\alpha=.10$ fixed, has limit $53/30$. Both assume empty history and exact stopping (Appendix~\ref{app:deduplicated}). A strict gain requires every minimum joint witness to contain an entry with $0<f_{ir}<p_i$ (Corollary~\ref{cor:overlap-budget}).

Whole-evaluator queries reveal all $N$ scores in a column, so their minima depend on evidence positions (Appendix~\ref{app:granularity}). Every history admits a completion requiring no further exclusion evidence. No history therefore guarantees a strict remaining joint advantage (Proposition~\ref{prop:residual-no-uniform}); predictions can still guide discovery.

\section{Acquiring a Certificate}
\label{sec:acquisition}
Offline minima show how much evidence suffices, not how to find it. Baselines are fixed order, random order and largest width $u_i-\ell_i$ (Appendix~\ref{app:protocol}). Reference priority fits one ridge model per candidate to reference-row flip counts $f_{ir}$. Features are pool scores, cutoff margin, score range and $p_i/M$; predictions are clipped to $[0,p_i]$ (Appendix~\ref{app:protocol}).

Among query-eligible rows, reference priority selects the row with the smallest estimated total cost: observed flips plus predictions for remaining entries. Within the chosen row it chooses the unqueried candidate with largest predicted upper-cost reduction $p_i-\hat f_{ir}$. Frontier priority computes the active capacity under observed costs plus predictions for unqueried entries and clips it to $[V^- -m_0+1,V^+-m_0]$, obtaining target $s$. It scores predicted progress towards safety for $s$ and exclusion of $s+1$. It omits exclusion when $s+1>m$ or its deficit is nonpositive (Appendix~\ref{app:frontier}); safety-only frontier always omits it. Its exclusion target $s+1$ ignores the stopping tolerance, which only the verifier uses. Predictions guide order, not bounds; none of these rules has a query-optimality guarantee.

\begin{algorithm}[t]
\caption{Score inspection for certified retention}
\label{alg:inspect}
\begin{algorithmic}[1]
\REQUIRE Base scores, pooling rule, $T$, $t_0$, $k$, $\alpha$, integer $\varepsilon\ge0$, policy $\pi$ with its reference-fitted models and seeded tie state.
\STATE Compute membership, caps $p_i$ and $b_H$; set $\ell_i=0$, $u_i=Tp_i$ on active selected rows.
\STATE Maximize retained count under upper costs to obtain $C^-$; compute $V^-=V(u)$ and $V^+=V(\ell)$.
\WHILE{$V^+-V^->\varepsilon$}
\STATE Use $\pi$ to choose an unqueried active candidate--output entry from initially visible information and query history.
\STATE Query $s_r(i)$ and compute $f_{ir}$ from the base scores; set $\ell_i\gets\ell_i+f_{ir}$, $u_i\gets u_i-p_i+f_{ir}$; recompute $C^-$, $V^-$ and $V^+$.
\ENDWHILE
\STATE Return $S_0(P)\cap C^-$, its churn upper bound and retained-count gap.
\end{algorithmic}
\end{algorithm}

Sorting computes each capacity bound in $O(m\log m)$ time; invariant selected rows have zero churn cost.
\begin{proposition}[Discovery can be expensive]
\label{prop:positive-discovery}
For every integer $m\ge2$, a family of mean-pooling instances with $M=T=7$, $q=.05$, $\alpha=.10$, $|S_0|=m$ and $N=\Theta(m)$ has two-entry minimum certificates at $\varepsilon=\lceil.02m\rceil$. Every deterministic policy valid under the completion model needs at least $m$ queries on some instance of the family. Every randomized always-valid policy has expected cost at least $(m+1)/2$ on some fixed instance.
\end{proposition}
Hence no competitive factor independent of batch size holds at this tolerance (proof in Appendix~\ref{app:deduplicated}).

\section{NLP Score-Inspection Experiments}
\label{sec:experiments}
We tested these bounds and acquisition rules on cached NLP evaluator panels.
\begin{table*}[t]
\centering\small
\begin{tabular}{llrrrrrr}
\toprule
Panel & Pooling & Z (\%) & Fixed (\%) & Random (\%) & Width (\%) & Ref. (\%) & Cov.\ (\%) \\
\midrule
RewardBench 2 & median & 0.0 & 63.8 & 64.0 & 63.0 & 61.4 & 97.3 \\
WMT23, 13 streams & median & 10.0 & 57.9 & 58.4 & 65.0 & 45.1 & 97.7 \\
WMT23, 12 streams & median & 92.5 & 3.9 & 3.8 & 3.9 & 3.8 & 99.6 \\
RewardBench 2 & mean & 0.0 & 90.9 & 91.5 & 91.7 & 86.5 & 97.5 \\
WMT23, 13 streams & mean & 0.0 & 91.3 & 92.2 & 92.6 & 85.7 & 97.5 \\
WMT23, 12 streams & mean & 0.0 & 89.5 & 90.2 & 89.9 & 85.4 & 97.0 \\
\bottomrule
\end{tabular}
\caption{$q=.05$, $\alpha=.10$, gap $V^+-V^-\le\lceil.02|S_0|\rceil$. Z: percentage of 40 splits stopping without queries. Order columns: mean active inspection ($100Q/(mT)$), zero for empty active sets. Ref.: reference priority; Cov.: its mean selected coverage ($100K/|S_0|$).}
\label{tab:stopping}
\end{table*}


\subsection{Panels and protocol}
Our RewardBench~2 panel contains 14 reward models scoring 1,723 four-completion prompts \citep{malik2026rewardbench2}. Our WMT23 en-de panel contains 557 source segments, 13 output streams and 24 reference-based metrics \citep{freitag2023wmt23}. Removing \texttt{synthetic\_ref} gives a twelve-stream variant with refitted maps and cutoff. WMT24 en-de, en-es and ja-zh panels use document splits, exclude human-reference outputs and keep four depths, 40 splits and $\alpha=.10$ \citep{freitag2024wmt24}.

In RewardBench and WMT23, the first seven evaluators of a frozen ordering form the base pool, leaving $T=7$ RewardBench or $T=17$ WMT23 candidates (Appendix~\ref{app:protocol}). Forty seeded group-preserving splits assign 574/574/575 RewardBench prompts or 185/185/187 WMT23 source segments to reference, calibration and test sets. Only CRC and grouped SCoRE use the calibration split. Each evaluator uses a reference-fitted right empirical cumulative distribution function (ECDF). For $n_{\rm ref}$ reference outputs, the cutoff is the $\max(1,\operatorname{round}(qn_{\rm ref}))$-th highest reference consensus (Appendix~\ref{app:protocol}). RewardBench and WMT23 use both pooling rules, $q\in\{.05,.10,.20,.30\}$ and $\alpha\in\{.02,.05,.10,.15,.20\}$ (primary $.10$).

Primary stopping uses $V^+-V^-\le\varepsilon=\lceil.02|S_0(P)|\rceil$, returning $K=V^-$ and hence $V^*-K\le\varepsilon$; exact stopping uses zero.

Main-panel queries access cached scores, so counts are logical accesses, not inference savings; reference and base scoring are charged separately (Appendix~\ref{app:protocol}).

\subsection{Retention baselines and risk guarantees}
Selected coverage is $K/|S_0|$. At $q=.05$, $\alpha=.10$ and median pooling, zero-query coverage spanned 80.9--99.2\% across RewardBench and the two WMT23 panels, versus 99.5--99.9\% under full information (Appendix Table~\ref{tab:baselines}).

Of 15,360 certified masks at $\alpha=.10$, 6,985 had a replacement with normalized churn above $.10$, though each met the average budget (Appendix~\ref{app:individual-replacements}). Capping every replacement, equivalently guarding against adversarial replacement weights, makes full-information retention NP-hard (Proposition~\ref{prop:per-replacement-hard}). At $q=.05$ on the six mean panels, such caps raised reference-priority inspection at the primary stop from 82.1--86.5\% to 95.3--99.5\% of active entries (Appendix~\ref{app:per-replacement}).

Envelope CRC \citep{yeh2025conformal} and grouped SCoRE target the same churn loss in expectation under group exchangeability, so neither ensures the budget on every batch. Grouped SCoRE uses candidate calibration scores but no candidate test scores. Its RewardBench coverage at $q=.05$ was \scoreRbMedianCov\% under median and \scoreRbMeanCov\% under mean (Appendices~\ref{app:baselines} and~\ref{app:review-score}).

\subsection{Inspection and certificate size}

At nominal $q=.05$ (Table~\ref{tab:stopping}), with median pooling, reference priority averaged 61.4\% active-entry inspection on RewardBench and 45.1\% on thirteen-stream WMT23. Mean pooling averaged 85.4--86.5\% across the three panels.

\begin{table*}[t]
\centering\small
\begin{tabular}{lrrrrr}
\toprule
Panel & Num. closed & Primary $Q$ & Primary Ref./$Q$ & Exact Ref./$Q$ & Gain/num. closed \\
\midrule
RewardBench 2 & 320/320 & 1216.9 & 1.205 & 1.186 & 0/320 \\
WMT23, 13 & 320/320 & 2012.1 & 1.394 & 1.318 & 7/320 \\
WMT23, 12 & 320/320 & 2024.3 & 1.291 & 1.258 & 13/320 \\
WMT24 en-de & 320/320 & 6922.7 & 1.185 & 1.169 & 3/320 \\
WMT24 en-es & 320/320 & 5968.4 & 1.187 & 1.174 & 3/320 \\
WMT24 ja-zh & 320/320 & 8717.0 & 1.119 & 1.107 & 63/320 \\
\bottomrule
\end{tabular}
\caption{Full-test mean evidence at $\alpha=.10$: four depths, 40 splits, and two count-gap targets per panel. Primary uses $\varepsilon=\lceil .02|S_0|\rceil$; exact uses zero. Primary $Q$ averages the minimum joint entry count over 160 primary targets per panel; each ratio uses its matching stopping target. All joint and union minima are numerically closed. Ratios average splitwise reference queries divided by numerically closed positive matching minima, excluding zero minima. Strict gain compares joint evidence with the best union of separately optimal witnesses. Appendix~\ref{app:full-mean-grid} reports all depths.}
\label{tab:full-mean-main}
\end{table*}

\paragraph{Matching minima.} At $q=.05$ and $\alpha=.10$, splitwise medians of exact-stop reference queries over the exact odd-median minimum (Theorem~\ref{thm:exact}) were 1.968 for RewardBench and 1.662 for thirteen-stream WMT23. Both medians use splits with $Q_{\rm exact}>0$ (Appendix Table~\ref{tab:entry-overhead-summary}).

On the six panels, four depths and 40 splits yield 960 mean instances and 1,920 targets across the two stops. For every target, the saved solver records contain integer-checked joint and best separate-union witnesses whose sizes equal numerical global lower bounds. Table~\ref{tab:full-mean-main} reports these closed records. A nominal 60-second rerun of the outcome MILP, including construction, closed 931/960 primary and 945/960 exact targets (Table~\ref{tab:resource-strata}). Numerical closure is not an exact-arithmetic proof; Appendix~\ref{app:full-mean-grid} reports one rejected conflicting bound.

\paragraph{Beyond deduplication.} Joint optimization improved on $Q_{\cup,\rm sep}$ in \fullMeanStrictGain{} of the 1,920 targets, saving \fullMeanJointSaving{} entries in total (Appendix~\ref{app:full-mean-grid}). Of 57 instances unable to retain the full baseline list, all 57 exact targets and 32 primary targets showed strict gain. Full retention makes exclusion automatic in the other 903 instances (1,806 targets). In the remaining 25 no-gain primary targets, a saved optimal witness certifies at least $|S_0|-\varepsilon$ outputs safe, making exclusion automatic (Table~\ref{tab:explanatory-overlap}). Across the 960 targets for each stopping rule, median savings were zero. Mean savings were 5.34 entries for primary stopping and 1.19 for exact stopping. Maximum reductions relative to $Q_{\cup,\rm sep}$ were 2.61\% and .87\%, respectively (Table~\ref{tab:gain-distribution}). Reference-predicted active cost relative to the budget is available before candidate scoring. In a post hoc check, it ranked gain targets above no-gain targets (AUC .971 primary, .959 exact; Appendix~\ref{app:joint-mechanism}).

\paragraph{Online discovery.} Across 960 primary mean targets, frontier priority averaged \frontierQueries{} queries versus \frontierRefQueries{} for reference and \frontierWidthQueries{} for width. It beat reference on \frontierWins{} targets and lost on \frontierLosses{}, with three ties. Its wins averaged 420 entries and its losses 174. Safety-only frontier had a slightly lower overall mean. On the 32 targets needing exclusion, frontier priority saved 20.28 entries over safety-only frontier. Under full retention it added 3.72. A tolerance-aware target, $s+\varepsilon+1$, did not help these 32 targets (Appendix~\ref{app:frontier}).

\subsection{Budget, composition, external panels and quality}
Inspection need not decrease with the churn budget: safety can become easier while exclusion becomes harder (Appendix Figure~\ref{fig:review-budgets}). Twenty random pools per RewardBench or WMT23 panel showed sensitivity to composition, which changes the replacement family and selected set (Appendix~\ref{app:review-composition}). Under median pooling at $q=.05$, WMT24 average zero-query coverage was 98.1--99.8\% (Appendix~\ref{app:review-quality}).

\paragraph{Summarization evaluation and measured scoring cost.}
A frozen SummEval \citep{fabbri-etal-2021-summeval} protocol covers all ten three-base/two-candidate compositions. On its 48 non-full mean targets, joint certificates were smaller than the best separate union on 27 and equal on 21. With mean pooling, frontier priority saved two queries on average relative to safety-only frontier when $V^*+\varepsilon<|S_0|$ and added queries under full retention. A fresh five-metric panel gave 192 runs on one A40, 24 per method. Primary reference priority averaged 89.97 s over 24 runs versus 108.18 s for full scoring, both including common setup. All query policies took 89.51--90.43 s, so this experiment showed no runtime advantage of joint frontier (Appendix~\ref{app:summeval-confirmation}).

\paragraph{Output quality.} Against label-blind, same-consensus top-$K$ at matched count, RewardBench preferred-label precision differed by $-0.46$ to $-0.03$ percentage points across depths under median reference priority. Average returned-set Multidimensional Quality Metrics (MQM) scores, higher meaning better, differed in either direction at matched annotated count. Neither comparator is churn-constrained (Appendices~\ref{app:quality} and~\ref{app:review-quality}).

\subsection{Reference mapping and query granularity}
\label{sec:query-granularity}
On thirteen-stream WMT23 with median pooling at nominal $q=.05$, switching from right to midrank ECDF raised zero-query selected coverage from 84.9\% to 99.8\%. Average active-entry inspection at primary stopping fell from 45.1\% to 0.3\%. Test groups stayed fixed, but baseline-set Jaccard similarity was 0.77: the saving concerned a changed selection task (Appendix~\ref{app:mapping-column}).

At $q=.05$ and $\alpha=.10$, median RewardBench reference priority revealed 1.77\% of all $NT$ candidate-test entries by entry queries. Reference-ordered column queries, which charge $N$ entries each, revealed 71.8\%. All 1,920 mean column targets on the six-panel grid had equal joint and best separate-union minima, including the column counterparts of the 89 entry-gain targets. Other realizable instances permit strict column gains (Appendix~\ref{app:mean-even-columns}).

\section{Discussion}
Unlike matroid feasibility in basis verification \citep{diwan2026basis}, retention feasibility depends on unobserved flip costs, which couple safety and exclusion evidence. A strict gain arises when no pair of tied separate minima reaches the joint minimum after deduplication; overlap alone does not decide this. Two matched instances share component minima, joint minimum and intermediate counts, yet their best separate unions are eight and seven (Appendix~\ref{app:deduplicated}).

A gain also needs both proof tasks to remain, but this is not sufficient. Every no-gain target on the six-panel mean grid had an optimal joint witness with automatic exclusion. Yet 13 of the 40 SummEval mean targets requiring exclusion showed no gain. Whole-column access removed the remaining room: every non-full column target needed all columns.

Discovery adds its own cost. On the same 960 primary targets, reference queries exceeded $Q_{\rm joint}$ by 809.58 entries on average, whereas $Q_{\cup,\rm sep}-Q_{\rm joint}$ averaged 5.34. Proposition~\ref{prop:positive-discovery} bounds only the worst case, and prediction error alone does not explain this overhead. True-score priorities were not query-optimal either (Appendix~\ref{app:query-interventions}).
\FloatBarrier
\section{Conclusion}
Evaluator replacement makes retention a choice of outputs and the evidence justifying them. Under unit-cost entry access, odd medians separate safety and exclusion evidence. Mean pools can beat optimal separate unions, which needs both proof tasks; NLP gains were modest. Small certificates ensure neither cheap discovery nor cheap inference. Strict-gain conditions under general histories and batching-aware acquisition remain open.

\section*{Limitations}
The certificates cover single-member replacements from the declared finite family with fixed weights, maps and cutoff. The core experiments, closed forms and Theorem~\ref{thm:joint} use uniform weights; Appendix~\ref{app:model-scope} also tests fixed nonuniform weights. They control average churn. Per-replacement caps define an NP-hard program whose optimum is never above the average-budget optimum (Appendix~\ref{app:per-replacement}). Evaluators outside the family are not covered. Adding an unretained output of $S_0(P)$ changes the mask and requires rechecking the mask's budget, potentially using the existing observations. Returning an output outside $S_0(P)$ requires a different selection rule. Newly generated outputs also change the batch being certified. Nominal $2k$ normalization does not constrain the returned count to $k$. The odd-median closed-form witnesses assume unit-cost candidate--output queries and Boolean opposition completions without cross-entry constraints. Deviation and Lipschitz bounds are valid only when their constants hold for the true scores. The appendix experiment uses oracle constants from the test batches, and constants estimated on reference data often fail on test batches (Appendix~\ref{app:constrained}). The near-full compatibility criterion of Proposition~\ref{prop:compatibility} assumes empty history, exact stopping and restricted exclusion-witness ties. It covers no exact target on the six-panel mean grid. The joint dynamic program has a large budget-indexed state space. Outcome grouping reduces the integer model size without guaranteeing polynomial solver time. Forty splits overlap and the two WMT23 variants share outputs, so aggregate comparisons are descriptive. True-score priorities are diagnostic, not deployable policies; improvements beyond optimal deduplication are modest in solved comparisons.

RewardBench and WMT23 were used during method development. WMT24 results had not been inspected when the risk objective and panel-inclusion criteria were fixed. After an exact-replay failure, the numerical query priorities were corrected; this affected query order, not the certificate definition. This protocol freeze is partial and does not constitute a fully preregistered evaluation. Independence between evaluator training data and the benchmarks is also not established, so training-data overlap remains possible. Some released meta-metrics additionally use human-score calibration. The rationale for the original frozen evaluator ordering is not established by the records. Composition sensitivity does not make this ordering representative, and sampled pools do not cover all future evaluators. Reference priority has no online competitive guarantee. Incomplete MQM coverage restricts quality diagnostics to annotated outputs. The SummEval timing study measures five fixed scorers on one A40, with batch-eight full scoring and sequential entry queries. It does not optimize online microbatching or model deployment, and its three repetitions measure runtime variability, not independent task replication.

\section*{Ethical Considerations}
The experiments reuse evaluator panels and existing labels or generate controlled numerical scores; they collect no new human labels or language responses. Source reconstruction matches existing benchmark text to cached scores and output identities locally. Available records do not document a systematic audit for personally identifiable information or offensive content. Stable selection can preserve evaluator errors, so a certificate does not endorse retained content. Appendices~\ref{app:protocol} and~\ref{app:review-quality} document data sources and acquisition. We have not established redistribution permission for every separate RewardBench result or contributed WMT score. Benchmark and toolkit licenses do not resolve the status of every asset. Code, derived records and selected query traces are available at \url{https://anonymous.4open.science/r/acl-arr-B388}. Third-party texts, full score arrays and some large or unreleased traces are excluded. The repository documents how to rebuild them or gives their hashes and summaries. Intended use is research on retention certification. Historical compute accounting is incomplete; the repository records available details.

\section*{Acknowledgments}
Generative AI tools supported problem framing, proof discussion, coding, analysis and manuscript drafting. Human authors remain responsible for claims, code, results, citations and final text.

\bibliography{references}
\clearpage
\appendix
\setcounter{dbltopnumber}{3}
\renewcommand{\dbltopfraction}{.9}
\renewcommand{\dblfloatpagefraction}{.6}
\raggedbottom
\paragraph{Appendix guide.}
Repeated baseline, budget and grouped-SCoRE grids are provided in \nolinkurl{supplement/result_grids.csv}. Full-depth mapping, column, tolerance and discovery grids are in \nolinkurl{supplement/extended_result_grids.csv}. Companion JSON files preserve headers, captions and table sources. Matched minimum-evidence and query costs for both pooling rules are in \nolinkurl{supplement/matched_evidence_summary.csv}, with per-instance records in \nolinkurl{supplement/matched_evidence_records.csv}. Full proofs, replay specifications and additional diagnostics are provided below.
Proofs begin in Appendix~\ref{app:proofs} (p.~\pageref{app:proofs}); replay specification begins in Appendix~\ref{app:protocol} (p.~\pageref{app:protocol}). Appendix~\ref{app:baselines} specifies baselines; D--J collect retention, quality and risk comparisons; K--L examine selection protocols and loss definitions. Minimum-evidence extensions are in Appendix~\ref{app:granularity} (p.~\pageref{app:granularity}), with a compact joint mean formulation in Appendix~\ref{app:compact-joint} (p.~\pageref{app:compact-joint}) and full-batch evidence in Appendix~\ref{app:full-mean-grid} (p.~\pageref{app:full-mean-grid}). N--O discuss tolerance and operational scope. Prediction interventions and frontier acquisition appear in Appendix~\ref{app:acquisition-diagnostics} (p.~\pageref{app:acquisition-diagnostics}). Appendix~\ref{app:summeval-confirmation} reports SummEval and measured scoring costs.
\FloatBarrier
\section{Proofs and Completion Model}
\label{app:proofs}
\subsection{Notation and units}
Table~\ref{tab:notation} collects the main quantities; Section~\ref{sec:median} and Appendix~\ref{app:worked-example} give worked examples.
\begin{table}[t]
\centering\small
\begin{tabular}{@{}p{.23\columnwidth}p{.69\columnwidth}@{}}
\toprule
Symbol & Meaning and unit\\
\midrule
$N,M,T$ & Outputs, base evaluators, candidate evaluators.\\
$k,S_0$ & Nominal count and baseline-selected set; $|S_0|$ need not equal $k$.\\
$H=MT$ & Number of single-member replacement scenarios.\\
$f_{ir},p_i$ & Entry's flip count and its row's maximum flip count.\\
$F_i$ & Full row cost $\sum_r f_{ir}$, in flips.\\
$b_H$ & Integer flip budget $\lfloor2k\alpha H\rfloor$.\\
$\ell_i,u_i$ & Lower and upper row costs, in flips.\\
$V^-,V^+,V^*$ & Safe, optimistic, and true optimal retained counts.\\
$\varepsilon$ & Allowed certified retained-count gap.\\
$Q_{\rm joint}$ & Minimum joint evidence under the stated query unit.\\
$Q_{\cup,\rm sep}$ & Minimum union over separately minimum proofs.\\
$Q_{\rm split}$ & Minimum sum of separate proof sizes.\\
\bottomrule
\end{tabular}
\caption{Core notation. In odd-median Boolean formulas, dividing by the nonzero flip cap gives opposition-count costs and budget $b$. A column query reveals and is charged for all $N$ score entries.}
\label{tab:notation}
\end{table}

\subsection{Flip geometry and sharp intervals}
Our cutoff rule includes scores equal to the threshold. Under the nonempty reference ECDF conventions used here, every score in the reference split maps into $(0,1]$, so a cutoff fitted from reference consensus satisfies $0<t_0\le1$. For median pooling with $M=2h+1$, let $v_i$ be the number of base scores at least $t_0$. If $v_i\ge h+2$, removing one vote leaves at least $h+1$ votes. If $v_i\le h-1$, adding one vote leaves at most $h$ votes. These rows are invariant. When $v_i=h+1$, a below-cutoff candidate changes membership exactly when it replaces one of the $h+1$ voting members. When $v_i=h$, an at-or-above-cutoff candidate changes membership exactly when it replaces one of the $h+1$ nonvoting members. Thus the maximum flip count is $h+1$ on either boundary and zero elsewhere.

For mean pooling, each replacement mean is nondecreasing in candidate score $x$. On a selected row, increasing $x$ cannot create an exclusion, so $f_i(x)$ is nonincreasing, with maximum $f_i(0)$. At $x=1$, the candidate score is at least every base score it can replace; all replacement means remain at or above the cutoff, giving $f_i(1)=0$. On an unselected row, $f_i(x)$ is nondecreasing, with maximum $f_i(1)$ and $f_i(0)=0$. The same argument applies to the average of the two central scores in an even-sized median pool: increasing a candidate score cannot decrease either order statistic. It also covers ties because membership always uses $\ge$.

Each row\textquotesingle s cost sums observed flip counts and unobserved counts bounded by zero and $p_i$. With no cross-entry constraints, joint endpoint assignments realize both $\ell$ and $u$. Consequently, $V(\ell)$ is the largest full-information retention capacity among consistent completions, and $V(u)$ the smallest. These extrema are sharp in the stated model. The declared completion set does not encode additional constraints from evaluator functions, score correlations or raw-score support. Sharpness is relative to this set; it does not claim every endpoint assignment is realizable by a deployed evaluator.

\subsection{Safety witness}
Suppose a certificate proves some boundary mask $C$ of cardinality $L$ safe. For each $i\in C$, let $a_i$ count its observed agreements. In the worst-case completion its opposition cost is $\sum_{i\in C}(T-a_i)$. Feasibility requires
\[
 \sum_{i\in C}a_i\ge LT-b.
\]
Each query reveals at most one agreement, establishing the lower bound $\max(0,LT-b)$.

A truly feasible set $C$ of $L$ boundary rows contains $LT-\sum_{i\in C}z_i\ge LT-b$ agreement cells. Revealing any $\max(0,LT-b)$ of them makes its upper cost at most $b$, attaining the lower bound. If no truly feasible $L$-row set exists, no observations can certify one; this infeasible target is represented by infinite certificate cost. For $L=0$ the empty mask requires no information. When $b\ge LT$, every $L$-mask is safe before querying.

\subsection{Exact optimality witness}
Let $L$ be the true optimal boundary count. When $L=m$, certifying the full boundary safe certifies optimality automatically; the safety witness gives the claimed value. Consider $L<m$, so $p=L+1$ and the sum of the $p$ smallest true row costs exceeds $b$. Therefore the level $t^*$ in Equation~\ref{eq:level} exists.

A certificate establishing infeasibility of every $p$-row mask under every consistent completion must reveal opposition counts $x_i$ satisfying
\[
 0\le x_i\le z_i,\qquad\sum_{i=1}^{p}x_{(i)}\ge B_0=b+1.
\]
Otherwise, choose the $p$ rows with the smallest observed opposition counts and complete their unknown cells with agreements. This makes the set feasible in a consistent completion. Rowwise inequalities $x_i\le z_i$ imply order-statistic inequalities $x_{(j)}\le z_{(j)}$ for every $j$. Let $t=x_{(p)}$, the $p$-th smallest observed count. Then
\[
 B_0\le\sum_{i=1}^{p}x_{(i)}\le\sum_{i=1}^{p}\min(z_{(i)},t),
\]
so $t\ge t^*$. Every one of the remaining $m-p$ observed row counts is at least $t$. The number of observed opposition cells is consequently at least $B_0+(m-p)t^*$.

For attainment, order the rows by their true costs. Assign integer observed counts to the first $p$ rows between zero and $\min(z_{(i)},t^*)$, with sum exactly $B_0$. This is possible because the sum of the caps is at least $B_0$. If all were below $t^*$, their sum would be at most $D(t^*-1)<B_0$; at least one must therefore equal $t^*$. The level also satisfies $t^*\le z_{(p)}$, since the cap at $z_{(p)}$ leaves the first $p$ true row counts unchanged. Assign observed count $t^*$ to every subsequent row. Each is within its row's true opposition capacity. The $p$ smallest observed counts then sum to $B_0$, and the total number of revealed opposition cells is exactly $B_0+(m-p)t^*$.

Choose safety agreements from any truly optimal set of $L$ rows. Their cells are disjoint from the exclusion witness\textquotesingle s opposition cells. Revealing the union certifies feasibility and rules out any feasible $p$-row set under every completion, at the sum of the two costs. Every exact certificate must meet both lower bounds, proving Equation~\ref{eq:exact}.

For $L=0<m$, $p=1$ and $t^*=b+1$; exact exclusion requires $b+1$ oppositions on every row. For $m=0$, full optimality is immediate and the certificate size is zero. The formula is not an online competitive bound and does not apply as a lower bound to a positive-$\varepsilon$ stop. Nonuniform replacement weights and unequal entry costs require a different characterization. Appendix~\ref{app:granularity} treats whole-column revelation separately.

\subsection{Online difficulty and acquisition objectives}
For one row and $0\le b<T$, feasibility asks whether its opposition count is at most $b$. An adversary can arrange the first $T-1$ replies to contain exactly $b$ oppositions in total. The final reply then distinguishes cost $b$ from cost $b+1$. Any deterministic exact strategy has a worst case of $T$ queries. An offline feasible certificate needs only $T-b$ agreements; an infeasible certificate needs $b+1$ oppositions.

At $b=T-1$, a feasible row with one agreement has a one-entry offline certificate, whereas an adversary can force $T$ queries. Proposition~\ref{prop:fixed-discovery} strengthens this obstruction: with $M=T=7$, $q=.05$, $\alpha=.10$ and $|S_0|=k$, exact discovery cost can grow with batch size while the offline minimum remains one.

The integer retention gap need not show diminishing returns. For one row with $T=2$ and $b=0$, the first observed agreement leaves the gap at one, whereas the second reduces it to zero. The second observation therefore gives a larger marginal reduction than the first, so adaptive-submodular guarantees need further conditions before they apply to this objective. Stochastic Boolean query methods analyze objectives and distributions with their own conditions \citep{deshpande2016sbfe,gkenosis2018ssclass}.

\subsection{Minimum witnesses at positive retention tolerance}
\label{app:gap-witness}
The same agreement and opposition geometry characterizes the smallest certificate at any integer tolerance. The implemented target is the observable frontier gap $V^+-V^-\le\varepsilon$. The realized regret satisfies $0\le V^*-V^-\le V^+-V^-$ and can be smaller than the gap. The gap criterion certifies shortfall for every consistent completion; checking realized regret alone on the complete matrix would not supply this guarantee. Appendix~\ref{app:worked-example} gives a state with regret zero and gap one.

Corollary~\ref{cor:gap-witness} in Section~\ref{sec:median} states the minimum witness size.

\begin{proof}
Selected interior rows add the same $m_0$ to both count bounds and cancel from the gap. Consider any valid certificate and let $s$ be its safe boundary count, so $V^-=m_0+s$. Since $V^-\le m_0+L^*\le V^+$ and $V^+-V^-\le\varepsilon$, this count lies between $\max(0,L^*-\varepsilon)$ and $L^*$. Safety requires at least $\max(0,sT-b)$ agreements by Theorem~\ref{thm:safe}. If $p=s+\varepsilon+1\le m$, observed opposition counts must rule out every $p$-row mask, making the optimistic boundary count at most $p-1$. The order-statistic argument for Theorem~\ref{thm:exact} applies to this $p$, yielding at least $E(p)$ oppositions. Their cells are disjoint, so every certificate has at least the corresponding sum in Equation~\ref{eq:gap-witness}.

Conversely, choose any $s$ in this range and a true-feasible $s$-row mask. Reveal a minimum safety witness in it. When $p\le m$, reveal the capped opposition witness of size $E(p)$; it exists because $p>L^*$, so the $p$ smallest true costs exceed $b$. The union certifies at least $s$ safe rows and at most $s+\varepsilon$ feasible rows in any completion. When $p>m$, the upper count is already at most $m\le s+\varepsilon$, and no opposition witness is needed. Thus the minimum is attained. The empty boundary and tolerances at least $m$ give zero observations.
\end{proof}

Increasing $\varepsilon$ enlarges the set of permitted certificates, so $Q_\varepsilon$ is nonincreasing. The offline constructor uses the true matrix to locate the witness; it does not prescribe an online query order. Mean costs need not be Boolean opposition sums, and whole-column queries couple the available cells, so neither model directly inherits this median entry formula. Theorem~\ref{thm:joint} treats joint mean evidence, while Appendix~\ref{app:column-witness} treats exact median columns.

\subsection{A complete certificate example}
\label{app:worked-example}
Consider median pooling with $N=20$, $M=7$, $T=3$, $k=2$, $\alpha=.10$ and cutoff $t_0=.5$ fixed by construction. Each of the three selected rows has four base scores $.8$ and three $.2$; the remaining seventeen rows have all seven scores $.2$ and remain unselected under a single replacement. Each candidate--output entry assigns $.1$ for opposition and $.9$ for agreement. Their opposition matrix is
\[
 Y=\begin{pmatrix}0&0&0\\1&0&0\\1&1&0\end{pmatrix}.
\]
Each opposition reverses membership under four member replacements. Hence true flip costs are $(0,4,8)$ and the integer flip budget is $b_H=\lfloor2k\alpha H\rfloor=8$, equivalent to the bit budget $b=2$. Table~\ref{tab:worked-example} reports $\ell/4$ and $u/4$ in opposition-count units.

\begin{table}[t]
\centering\small
\begin{tabular}{rccrr}
\toprule
Queries & $\ell/4$ & $u/4$ & $V^-$ & $V^+$\\
\midrule
0 & $(0,0,0)$ & $(3,3,3)$ & 0 & 3\\
4 & $(0,0,0)$ & $(0,2,3)$ & 2 & 3\\
7 & $(0,1,2)$ & $(0,2,3)$ & 2 & 2\\
\bottomrule
\end{tabular}
\caption{Certificate updates for the three-row example with opposition budget $b=2$. The first four observations are row~1's three agreements and row~2's second entry. The final three are row~2's first opposition and row~3's first two oppositions.}
\label{tab:worked-example}
\end{table}

After the first four queries, returning rows~1 and~2 has certified risk at most $8/84=2/21\le.10$, while its actual risk is $4/84=1/21$. After four queries the frontier gap is one, although realized regret is already zero. Enumerating all $2^9=512$ subsets of the nine observable entries in this fixed matrix gives $Q_0=7$ and $Q_1=4$. Thus a lower query count at positive tolerance can reflect a smaller required witness as well as the acquisition order.

\FloatBarrier
\section{Reproducible Experimental Specification}
\label{app:protocol}
Code, derived records and query traces are in the anonymous repository \url{https://anonymous.4open.science/r/acl-arr-B388}; scripts there rebuild the score panels from pinned public sources.
\paragraph{Baseline query orders.} Fixed order uses seeded row priority followed by candidate index; random order permutes active entries. Largest width maximizes $u_i-\ell_i$ over eligible rows, then uses candidate index.

\paragraph{Panel construction.}
The frozen panels contain numerical scores and identifiers. For each RewardBench evaluator, we exclude records in the \texttt{Ties} subset. After unwrapping singleton score lists, the initial eligibility filter requires four numeric values. Finiteness is checked separately after panel construction and is not part of this type-and-length filter. We retain prompt IDs eligible for every evaluator; each prompt forms a four-output group. For each evaluator, alignment checks match the native prompt and every completion to the corresponding model-formatted input, preserving evaluator-specific formatting. Aligned records retain native order: the chosen completion followed by three rejected completions. Selection ties use the separate label-blind priorities below. 

The \href{https://huggingface.co/datasets/allenai/reward-bench-2-results}{allenai/reward-bench-2-results} snapshot is pinned at \texttt{b19c7033e964187d12e74}\allowbreak\texttt{a43a07f2d727a3d37e5}. Within each model JSON, the builder traverses the parallel \texttt{id}, \texttt{subset} and \texttt{scores} arrays in their native order and retains the last eligible occurrence of a repeated identifier. Common prompt IDs are then sorted numerically. The model/file mapping is \nolinkurl{evidence/source-model-files.json}; \nolinkurl{ACQUISITION_INDEX.md} gives pinned-source acquisition and reconstruction commands. 

For WMT23, we retain reference-based metrics under the common reference-A configuration and remove duplicate raw \texttt{refA.seg.score} files by SHA-256 of their byte contents before parsing scores and aligning systems.

Before assigning splits, the WMT23 builder checks the Pearson correlation of each retained metric with BLEURT-20 \citep{sellam-etal-2020-bleurt,pu-etal-2021-learning} over all 557 source segments and 13 output streams. All correlations are positive, with a minimum of 0.18425. A nonpositive correlation would reject the build. Since all are positive, this check changes neither score signs nor metric inclusion. It provides no human-quality labels and lies outside the reported post-freeze acquisition counts. A post hoc reference-only check over the same 40 splits also finds positive correlations for all 24 metrics: the minima are 0.09446 for thirteen streams and 0.07931 after removing \texttt{synthetic\_ref}. This post hoc diagnostic leaves scores unchanged; the original preprocessing used the full panel.

RewardBench splits use seed $20260815+r$ and WMT23 uses $87000+r$, for $r=0,\ldots,39$. Each seed shuffles group indices with Python \texttt{random.Random(seed).shuffle}. RewardBench reference and calibration sets each contain 574 prompts, leaving 575 for testing. WMT23 reference and calibration sets each contain 185 source segments, leaving 187 for testing. Let $n_{\rm ref}$ count reference outputs, including all outputs within each reference group. For each evaluator $e$, we separately fit
\[
 \widehat F_e(s)=\frac{1}{n_{\rm ref}}\sum_{\ell=1}^{n_{\rm ref}}\one[s^{\rm ref}_{e,\ell}\le s].
\]
All reference scores tied at $s$ enter the cumulative count. The cutoff is the reference consensus at position $n_{\rm ref}-\max(1,\operatorname{round}(qn_{\rm ref}))$ in ascending order, with zero-based indexing. The implementation uses Python's nearest-integer \texttt{round}, which resolves exact half-integer ties to the even integer, for both this index and the nominal test count $k$. Cutoff membership includes scores equal to the boundary. Throughout the tables, a top-$100q$\% setting denotes the nominal selection fraction $q$, not an exact returned fraction.

The WMT twelve-stream panel excludes the entry named \texttt{synthetic\_ref} before CDF fitting, reference consensus and all subsequent calculations. It has 12 outputs per segment instead of 13. RewardBench has 2,300 test outputs, thirteen-stream WMT23 has 2,431, and the twelve-stream panel has 2,244. The family contains 49 replacements for RewardBench and 119 for each WMT23 panel. 

The ordered RewardBench evaluators are Tulu-3-8B, QRM-Gemma-27B, WorldPM-72B, Skywork-V2-8B, DeBERTa-v3, InternLM2-20B, Starling-34B, UltraRM-13B, URM-LLaMa-8B, OffsetBias-8B, GRM-Llama3-8B, Skywork-V2-3B, LDL-Gemma-27B and BTRM-Qwen2-7B. The first seven constitute the pool and the last seven are candidates. This is the frozen original ordering; its selection rationale is not established by the available records. Analysis files record the source-model identity mapping. A results-dataset revision identifies the cached-score snapshot, not necessarily the model-weight checkpoints. Replaying scores requires the former; reproducing scorer execution additionally requires verified model revisions and interfaces.

The ordered WMT metrics are BERTscore, BLEU, BLEURT-20, COMET, Calibri-COMET22, MEE4, MaTESe, MetricX-23, MetricX-23-b, MetricX-23-c, XCOMET-Ensemble, XCOMET-XL, XCOMET-XXL, XLsim, YiSi-1, chrF, docWMT22CometDA, eBLEU, embed\_llama, f200spBLEU, instructscore, mre-score-labse-regular, prismRef and tokengram\_F. The first seven form the pool, and the remainder are candidates in both WMT populations.

\paragraph{Implementation and replay checks.}
The replay rebuilds each replacement score vector. It computes each mean in the same left-to-right order for full-matrix evaluation and single queries, never by subtracting a member from a rounded total. Observed and endpoint flip calculations therefore agree numerically. For rational budget $\alpha=a/d$, it computes $b_H=\lfloor2kMTa/d\rfloor$ by exact integer division.

Safe retention and query priorities use a label-blind permutation of test-item indices generated by NumPy \texttt{default\_rng(...).permutation} with seed $\text{split seed}+700001$. The priorities remain fixed throughout each path. Random acquisition uses the same NumPy permutation API with seed $\text{split seed}+99173$.  Reindexing carries each output\textquotesingle s original tie priority with it, preserving features and zero-query masks. For each candidate, ridge uses every reference output and its flip count for the current pool and cutoff. Features are the $M$ ECDF-normalized base scores, signed cutoff margin $c_P(i)-t_0$, base-pool score range and $p_i/M$, with no additional standardization. The objective is the sum of squared residuals plus a squared $\ell_2$ penalty of coefficient 1, with an unpenalized intercept. The normal equations are solved with \texttt{numpy.linalg.solve}; the recorded replay uses NumPy 2.0.2. The ridge penalty is fixed at one; available records do not document held-out tuning or identify the original BLAS library and platform versions. Test labels and candidate test values are not fitting inputs. Width chooses candidate index order within its chosen row. Reference priority chooses the maximum predicted upper-cost reduction, resolving candidate ties by increasing index.

An active-row entry query reveals one candidate score; known base scores determine its member-replacement outcomes and flip count. At each stopping or checkpoint state, the implementation stores the mask, retained-count bounds, query count and upper-cost sum. Certificate computation uses revealed scores and endpoint bounds only. A separate audit uses the complete cached matrix to recompute replies, reject duplicate queries and reconstruct recorded intervals and masks. For each recorded stopping target, it checks actual and upper-budget feasibility, $V^-\le V^*\le V^+$ and the corresponding stopping condition. Its access to hidden scores and $V^*$ is for evaluation, not acquisition or certification.

The accounting records $Tn_{\rm ref}$ candidate-reference entries, $Mn_{\rm ref}$ base-reference entries and $MN$ base-test entries. The full-information controls inspect $TN$ candidate-test entries; CRC additionally uses candidate calibration scores. Active query fractions divide by $mT$, with $m$ the number of selected rows having $p_i>0$, whereas panel fractions divide by $NT$. These entry counts are not GPU-call counts: reward models can batch completions, and document-level metrics can share computation across segments.

For each policy, commitment coverage is $|C|/N$ (including committed exclusions) and selected coverage is $|S_0\cap C|/|S_0|$, using the policy's own base list $S_0$. The automatically returned count is $K=|S_0\cap C|$; $K/k$ is its ratio to the nominal count before review. RewardBench precision is the preferred-response count among returned outputs divided by $K$. Reported ratios and paired differences are computed within each split before aggregation. Active inspection is defined as zero when $m=0$ (no eligible work), and this value enters the 40-split mean. Other ratios with empty denominators, such as returned precision, are undefined and omitted from their respective means. Display rounding is applied only after aggregation. Cross-table returned and queried counts use exact integer sums divided by the number of splits, with decimal half-up rounding to one place. The selected-coverage summaries have 40 defined splits per configuration; RewardBench precision comparisons and the WMT24 reference-priority quality tables likewise have 40.

No active set is empty in the original-pool budget experiments or primary WMT24 experiments. In composition sensitivity, thirteen-stream WMT23 with random pool 04, median pooling and $q=.05$ has an empty active set in all 40 splits. In mapping sensitivity, WMT24 ja-zh, quarter reference, midrank ECDF, median, $q=.05$ has one empty-active split (index 33). Its reference-priority positive-gap mean is 24.283\% over all 40 splits, versus 24.905\% conditional on the 39 nonempty splits. The tables use the former convention. 

The original WMT23 score panels contain no quality labels, so their original precision fields remain missing. Appendix~\ref{app:quality} separately aligns official MQM after freezing the masks and paths.

\FloatBarrier
\section{Comparison Methods}
\label{app:baselines}
Full-information controls use exact observed costs and $b_H$; the zero-query control uses endpoint upper costs. Table~\ref{tab:control-objects} summarizes their different decision objects and guarantees. The four original acquisition orders (fixed, random, largest width and reference priority) share the selected-retention solver, information boundary, risk budget and stopping tolerance. Their query counts therefore compare algorithms on one certification task; frontier priority is evaluated separately (Appendix~\ref{app:frontier}). The controls below compare selection objectives or statistical guarantees and have different scoring requirements.
\begin{itemize}
\item Zero-query retention maximizes returned baseline-selected count using upper costs $Tp_i$, with no candidate-test queries.
\item Batch uniform threshold chooses the largest feasible threshold among zero and the observed costs $F_i$. The total cost of all selected and unselected outputs at or below it must not exceed $b_H$; cost ties are included together.
\item Batch cardinality sorts all items by cost, then by label-blind tie priority, and takes the maximum feasible prefix. It maximizes committed-decision count, including exclusions.
\item Batch selected retention maximizes returned baseline-selected count using true costs $F_i$, applying the sort only to selected outputs. Nonselected outputs with $p_i=0$ are also committed, at zero churn cost. Its returned count is $V^*$.
\item Fixed-count retention recomputes the complete top-$k$ list for every replacement. It sums membership flips for this different rule and maximizes retention within its base list under the same nominal risk normalization.
\end{itemize}

\begin{table*}[t]
\centering\small\setlength{\tabcolsep}{3pt}
\begin{tabular}{llll}
\toprule
Control & Decision object & Guarantee or objective & Candidate scores\\
\midrule
Zero-query retention & Baseline-selected outputs & Endpoint-safe retained count & None on test\\
Batch selected retention & Baseline-selected outputs & Optimal count under churn & Full test matrix\\
Batch cardinality & Selected and excluded decisions & Optimal committed count & Full test matrix\\
Batch uniform threshold & Outputs at or below a cost threshold & Feasible tied cost level & Full test matrix\\
Fixed-count retention & Replacement-specific top-$k$ lists & Retention under a different rule & Full test matrix\\
Envelope CRC & Output-cost threshold & Expected churn & Calibration + test\\
Grouped SCoRE & Group acceptance or abstention & Expected group loss & Calibration\\
\bottomrule
\end{tabular}
\caption{Objects and information used by the controls. The last column lists candidate-score access after the shared reference stage; reference-stage costs appear in the common-information tables of Appendix~\ref{app:review-score}. The acquisition policies alone share the same selected-retention target and stopping criterion; the other rows compare objectives or guarantees.}
\label{tab:control-objects}
\end{table*}

\paragraph{Native and matched fixed-count controls.}
The native-tie fixed-count control selects the $k$ largest consensus--item-index keys in the baseline and every replacement. Its matched-tie counterpart uses the same seeded label-blind priority as the retention solver. Both are evaluated with full information on six panels, two pooling rules, four depths and 40 splits: 1,920 paired split instances at $\alpha=.10$.

The matched control orders every baseline and complete-batch replacement by decreasing consensus, then increasing the same seeded label-blind priority used by the retention solver. It refits neither maps nor cutoff and changes no query path. We recompute membership flips, maximum-retention masks, own-baseline coverage, average and replacement-specific churn. Labels are not loaded. Table~\ref{tab:matched-tie-control} summarizes all 48 panel--pooling--depth settings; \nolinkurl{outputs/matched-tie-control/} supplies masks and per-split records.

Relative to native ties, 452 baseline lists and 453 retained sets change; replacement memberships change in 1,869 configurations. Returned count changes in only one configuration, by one output, but average churn changes in 1,435. The sign of fixed-count-minus-frozen returned count is unchanged in all 1,920 configurations. Thus count ordering is robust to this tie correction on these panels, while memberships and risk are not generally invariant. This removes a tie-priority difference, but fixed-count and frozen-cutoff rules still define different baseline tasks. We do not apply the row-local median entry formula to globally coupled top-$k$ membership or claim a new online certificate.

All 3,840 control records (1,920 native and 1,920 frozen) reconcile with the original results. An independent sort and cardinality implementation checks matched masks, costs and priority transport under reindexing on 96 configurations across all panels, depths and pooling rules. All 5,760 records across native, matched and frozen controls meet their integer churn budgets. These are cached-score checks, not evaluator execution or independent population replications.
\begin{table*}[t]
\centering\small
\begin{tabular}{llrrrrrr}
\toprule
Panel & Pooling & $q$ & Frozen Cov. & Matched Cov. & \shortstack{$\Delta K_{\rm tie}$\\Matched--Native} & \shortstack{Changed\\sets/40} & \shortstack{Changed\\signs/40} \\
\midrule
RewardBench 2 & median & 0.05 & 99.6 & 99.4 & 0.025 & 10 & 0 \\
RewardBench 2 & median & 0.10 & 100.0 & 100.0 & 0.000 & 13 & 0 \\
RewardBench 2 & median & 0.20 & 100.0 & 100.0 & 0.000 & 11 & 0 \\
RewardBench 2 & median & 0.30 & 100.0 & 100.0 & 0.000 & 11 & 0 \\
RewardBench 2 & mean & 0.05 & 100.0 & 100.0 & 0.000 & 1 & 0 \\
RewardBench 2 & mean & 0.10 & 100.0 & 100.0 & 0.000 & 1 & 0 \\
RewardBench 2 & mean & 0.20 & 100.0 & 100.0 & 0.000 & 0 & 0 \\
RewardBench 2 & mean & 0.30 & 100.0 & 100.0 & 0.000 & 0 & 0 \\
WMT23, 13 & median & 0.05 & 99.5 & 100.0 & 0.000 & 32 & 0 \\
WMT23, 13 & median & 0.10 & 100.0 & 100.0 & 0.000 & 15 & 0 \\
WMT23, 13 & median & 0.20 & 100.0 & 100.0 & 0.000 & 14 & 0 \\
WMT23, 13 & median & 0.30 & 100.0 & 100.0 & 0.000 & 10 & 0 \\
WMT23, 13 & mean & 0.05 & 99.7 & 100.0 & 0.000 & 1 & 0 \\
WMT23, 13 & mean & 0.10 & 100.0 & 100.0 & 0.000 & 1 & 0 \\
WMT23, 13 & mean & 0.20 & 100.0 & 100.0 & 0.000 & 2 & 0 \\
WMT23, 13 & mean & 0.30 & 100.0 & 100.0 & 0.000 & 1 & 0 \\
WMT23, 12 & median & 0.05 & 99.9 & 100.0 & 0.000 & 19 & 0 \\
WMT23, 12 & median & 0.10 & 100.0 & 100.0 & 0.000 & 13 & 0 \\
WMT23, 12 & median & 0.20 & 100.0 & 100.0 & 0.000 & 10 & 0 \\
WMT23, 12 & median & 0.30 & 100.0 & 100.0 & 0.000 & 13 & 0 \\
WMT23, 12 & mean & 0.05 & 99.1 & 100.0 & 0.000 & 1 & 0 \\
WMT23, 12 & mean & 0.10 & 100.0 & 100.0 & 0.000 & 3 & 0 \\
WMT23, 12 & mean & 0.20 & 100.0 & 100.0 & 0.000 & 0 & 0 \\
WMT23, 12 & mean & 0.30 & 100.0 & 100.0 & 0.000 & 0 & 0 \\
WMT24 en-de & median & 0.05 & 100.0 & 100.0 & 0.000 & 22 & 0 \\
WMT24 en-de & median & 0.10 & 100.0 & 100.0 & 0.000 & 20 & 0 \\
WMT24 en-de & median & 0.20 & 100.0 & 100.0 & 0.000 & 13 & 0 \\
WMT24 en-de & median & 0.30 & 100.0 & 100.0 & 0.000 & 13 & 0 \\
WMT24 en-de & mean & 0.05 & 99.9 & 100.0 & 0.000 & 15 & 0 \\
WMT24 en-de & mean & 0.10 & 100.0 & 100.0 & 0.000 & 5 & 0 \\
WMT24 en-de & mean & 0.20 & 100.0 & 100.0 & 0.000 & 2 & 0 \\
WMT24 en-de & mean & 0.30 & 100.0 & 100.0 & 0.000 & 2 & 0 \\
WMT24 en-es & median & 0.05 & 100.0 & 100.0 & 0.000 & 21 & 0 \\
WMT24 en-es & median & 0.10 & 100.0 & 100.0 & 0.000 & 18 & 0 \\
WMT24 en-es & median & 0.20 & 100.0 & 100.0 & 0.000 & 16 & 0 \\
WMT24 en-es & median & 0.30 & 100.0 & 100.0 & 0.000 & 16 & 0 \\
WMT24 en-es & mean & 0.05 & 99.9 & 100.0 & 0.000 & 19 & 0 \\
WMT24 en-es & mean & 0.10 & 100.0 & 100.0 & 0.000 & 4 & 0 \\
WMT24 en-es & mean & 0.20 & 100.0 & 100.0 & 0.000 & 6 & 0 \\
WMT24 en-es & mean & 0.30 & 100.0 & 100.0 & 0.000 & 3 & 0 \\
WMT24 ja-zh & median & 0.05 & 99.8 & 100.0 & 0.000 & 15 & 0 \\
WMT24 ja-zh & median & 0.10 & 100.0 & 100.0 & 0.000 & 16 & 0 \\
WMT24 ja-zh & median & 0.20 & 100.0 & 100.0 & 0.000 & 20 & 0 \\
WMT24 ja-zh & median & 0.30 & 100.0 & 100.0 & 0.000 & 11 & 0 \\
WMT24 ja-zh & mean & 0.05 & 97.5 & 99.8 & 0.000 & 2 & 0 \\
WMT24 ja-zh & mean & 0.10 & 98.5 & 100.0 & 0.000 & 5 & 0 \\
WMT24 ja-zh & mean & 0.20 & 99.7 & 100.0 & 0.000 & 5 & 0 \\
WMT24 ja-zh & mean & 0.30 & 100.0 & 100.0 & 0.000 & 2 & 0 \\
\bottomrule
\end{tabular}
\caption{Matched label-blind top-$k$ control, $\alpha=.10$, full information. Cov. averages splitwise selected coverage (\%) using each rule\textquotesingle s own baseline. $\Delta K$ is matched-minus-native returned count, averaged over 40 splits. Changed sets/40 counts changed retained sets under this tie change. Changed signs/40 counts changes in the sign of fixed-count-minus-frozen returned count; a zero returned-count difference has sign zero. This compares selection tasks, not online query orders. In this matched control, every baseline and replacement ranking uses decreasing consensus followed by increasing seeded label-blind priority.}
\label{tab:matched-tie-control}
\end{table*}


\begin{table*}[t]
\centering\small
\begin{tabular}{llrrrr}
\toprule
Panel & Policy & Risk & Sel.\ cov.\ (\%) & Returned (outputs) & Precision \\
\midrule
RewardBench 2 & No queries & 0.053 & 80.9 & 95.5 & 0.844 \\
RewardBench 2 & Envelope CRC & 0.079 & 72.2 & 85.6 & 0.868 \\
RewardBench 2 & Batch threshold & 0.081 & 72.9 & 86.4 & 0.868 \\
RewardBench 2 & Batch cardinality & 0.099 & 77.3 & 91.5 & 0.865 \\
RewardBench 2 & Batch retention & 0.082 & 99.6 & 118.4 & 0.823 \\
RewardBench 2 & Fixed-count retention & 0.096 & 99.3 & 114.3 & 0.828 \\
WMT23, 13 streams & No queries & 0.056 & 84.9 & 117.9 & -- \\
WMT23, 13 streams & Envelope CRC & 0.067 & 77.0 & 107.9 & -- \\
WMT23, 13 streams & Batch threshold & 0.085 & 86.0 & 119.8 & -- \\
WMT23, 13 streams & Batch cardinality & 0.096 & 91.1 & 126.5 & -- \\
WMT23, 13 streams & Batch retention & 0.082 & 99.5 & 137.3 & -- \\
WMT23, 13 streams & Fixed-count retention & 0.056 & 100.0 & 122.0 & -- \\
WMT23, 12 streams & No queries & 0.042 & 99.2 & 112.4 & -- \\
WMT23, 12 streams & Envelope CRC & 0.062 & 95.2 & 106.9 & -- \\
WMT23, 12 streams & Batch threshold & 0.071 & 98.6 & 111.5 & -- \\
WMT23, 12 streams & Batch cardinality & 0.072 & 98.9 & 112.1 & -- \\
WMT23, 12 streams & Batch retention & 0.044 & 99.9 & 113.7 & -- \\
WMT23, 12 streams & Fixed-count retention & 0.043 & 100.0 & 112.0 & -- \\
\bottomrule
\end{tabular}
\caption{Risk is the mean of replacement-average churn $R(C)$ over 40 splits; Returned is the mean output count $K$. Median pooling, $q=.05$, $\alpha=.10$: full-information controls and zero-query retention. Selected coverage is $K/|S_0|$ using each policy\textquotesingle s base list. Counts and splitwise ratios are separately averaged over 40 splits. The fixed-count retention row uses native item-index ties; its matched label-blind counterpart is in Table~\ref{tab:matched-tie-control}. RewardBench precision is the preferred fraction among returned outputs; these WMT23 controls omit human-quality labels (see Appendix~\ref{app:quality} for the separate MQM diagnostic). CRC controls expectation; batch controls constrain realized churn. Count means use decimal half-up rounding after aggregating integer records.}
\label{tab:baselines}
\end{table*}


The CRC expectation control uses fully observed candidate scores to compute $d_i=F_i/H$ for every calibration and test item. Its group loss is $L_\lambda(U)=B|U|^{-1}\sum_{i\in U}d_i\one[d_i\le\lambda]$, with $B=N/(2k)$ fixed to the test batch's nominal scale. For calibration groups $U_1,\ldots,U_n$, define
\[
 \hat R_n(\lambda)=\frac1n\sum_{a=1}^{n}L_\lambda(U_a).
\]
The mathematical control chooses the maximum threshold from zero and observed calibration $d_i$ values satisfying
\[
 \frac{n\hat R_n(\lambda)+B\lambda}{n+1}\le\alpha.
\]
Zero is always feasible and includes zero-disagreement items. The following finite-grid argument supplies the expectation guarantee for this threshold convention, using the envelope correction of \citet{yeh2025conformal} and the exchangeability principle of \citet{angelopoulos2024crc}.

\paragraph{Finite-grid expectation guarantee.}
Condition on the reference-fitted maps, cutoff, frozen evaluator family and design quantities $H,N,k$, hence $B=N/(2k)$. Because $F_i$ is an integer, each normalized output cost $d_i=F_i/H$ lies on the fixed grid $\Lambda=\{0,1/H,\ldots,1\}$. The group losses $L_\lambda(U)$ are computed from these costs. Assume the loss vectors $(L_\lambda(U_a))_{\lambda\in\Lambda}$ for the $n$ calibration groups and one future group are conditionally exchangeable. Each group loss is nondecreasing in $\lambda$, is zero at zero, and is at most $B\lambda$. Let $\hat\lambda$ be any searched grid point satisfying the correction in exact arithmetic. Define the auxiliary, permutation-invariant threshold
\[
 \tilde\lambda=\max\left\{\lambda\in\Lambda:
 \frac{\sum_{a=1}^{n+1}L_\lambda(U_a)}{n+1}\le\alpha\right\}.
\]
The calibration correction and $L_{\hat\lambda}(U_{n+1})\le B\hat\lambda$ make $\hat\lambda$ feasible for this auxiliary problem, so $\hat\lambda\le\tilde\lambda$. Monotonicity and conditional exchangeability give
\[
 \begin{aligned}
 \mathbb E[L_{\hat\lambda}(U_{n+1})]
 &\le \mathbb E[L_{\tilde\lambda}(U_{n+1})]\\
 &=\mathbb E\!\left[\frac{1}{n+1}\sum_{a=1}^{n+1}L_{\tilde\lambda}(U_a)\right]
 \le\alpha,
 \end{aligned}
\]
where expectations are conditional on the fixed quantities above. Restricting the search to zero and observed calibration values preserves the expectation guarantee. Every accepted threshold satisfies the correction and is feasible for the auxiliary problem. The argument needs no maximization over all of $\Lambda$. Finiteness avoids a continuity assumption at the inclusive threshold's jumps.

RewardBench groups have four outputs; the two WMT23 panels have thirteen and twelve. With their fixed group counts, $N$ is 2,300, 2,431 and 2,244, respectively; $q$ fixes $k$ before test scores are used. For equal-size test groups sharing $\hat\lambda$, the batch churn is their average group loss. Applying the bound to each future group and using linearity of expectation gives expected batch churn at most $\alpha$. This assumes exchangeability; it does not prove it for the NLP population or bound every realized batch. Unequal-size WMT24 documents do not admit this unweighted reduction. The deterministic query certificate instead enforces the current batch budget under every declared score completion.

\paragraph{Numerical calibration.}
Historical runs evaluate the correction in floating point. An integer-arithmetic audit of all 4,800 calibrations in the original-pool budget grid confirms the exact inequality for every chosen threshold. One RewardBench mean instance (split 6, $q=.10$, $\alpha=.15$) selects $29/49$ instead of the exact maximum $30/49$. Both return 245 selected outputs, but their committed exclusions and churn differ. All 960 primary thresholds match the exact maxima. The tables retain the audited historical runs. The code package also provides an integer implementation enforcing the displayed correction exactly; the expectation argument above requires this exact feasibility condition.



\FloatBarrier
\section{All Depths and Pooling Rules}
Tables~\ref{tab:all-tolerant} and~\ref{tab:all-exact} report acquisition at positive-gap and exact stopping. The complete baseline grids in \nolinkurl{supplement/result_grids.csv} use each control\textquotesingle s own baseline-selected set as the denominator of selected coverage. Because these sets can differ, coverage differences do not measure retention gains on an identical output set. Preference precision is reported separately in Appendix~\ref{app:quality}. Acquisition rankings change with depth. Under mean pooling on the three original panels, reference priority inspects fewer active entries than largest width at $q=.05$. Largest width inspects fewer at $q=.30$. This reversal holds under both stopping rules (Tables~\ref{tab:all-tolerant} and~\ref{tab:all-exact}). Inspection is therefore compared at a fixed pooling rule, depth and stopping target, alongside initial coverage.
\begin{table*}[t]
\centering\small
\begin{tabular}{llrrrrrrrr}
\toprule
 & & & & \multicolumn{4}{c}{Active inspection (\%)} & & \\
Panel & Pooling & $q$ & Z (\%) & Fixed & Random & Width & Ref. & Cov. (\%) & Gap (outputs) \\
\midrule
RewardBench 2 & median & 0.05 & 0.0 & 63.8 & 64.0 & 63.0 & 61.4 & 97.3 & 3.0 \\
RewardBench 2 & median & 0.10 & 2.5 & 49.6 & 49.4 & 47.6 & 45.1 & 97.8 & 5.2 \\
RewardBench 2 & median & 0.20 & 7.5 & 14.4 & 14.2 & 14.5 & 14.1 & 97.9 & 9.5 \\
RewardBench 2 & median & 0.30 & 97.5 & 0.0 & 0.1 & 0.0 & 0.0 & 99.9 & 1.1 \\
RewardBench 2 & mean & 0.05 & 0.0 & 90.9 & 91.5 & 91.7 & 86.5 & 97.5 & 3.0 \\
RewardBench 2 & mean & 0.10 & 0.0 & 87.5 & 88.3 & 87.9 & 83.5 & 97.8 & 5.2 \\
RewardBench 2 & mean & 0.20 & 0.0 & 78.4 & 79.2 & 76.3 & 76.0 & 97.9 & 9.7 \\
RewardBench 2 & mean & 0.30 & 0.0 & 64.7 & 65.5 & 58.8 & 67.0 & 97.9 & 14.3 \\
WMT23, 13 streams & median & 0.05 & 10.0 & 57.9 & 58.4 & 65.0 & 45.1 & 97.7 & 2.9 \\
WMT23, 13 streams & median & 0.10 & 100.0 & 0.0 & 0.0 & 0.0 & 0.0 & 100.0 & 0.1 \\
WMT23, 13 streams & median & 0.20 & 100.0 & 0.0 & 0.0 & 0.0 & 0.0 & 100.0 & 0.0 \\
WMT23, 13 streams & median & 0.30 & 100.0 & 0.0 & 0.0 & 0.0 & 0.0 & 100.0 & 0.0 \\
WMT23, 13 streams & mean & 0.05 & 0.0 & 91.3 & 92.2 & 92.6 & 85.7 & 97.5 & 3.0 \\
WMT23, 13 streams & mean & 0.10 & 0.0 & 88.4 & 89.3 & 89.7 & 83.2 & 97.8 & 5.5 \\
WMT23, 13 streams & mean & 0.20 & 0.0 & 56.4 & 57.3 & 51.3 & 57.4 & 97.9 & 10.1 \\
WMT23, 13 streams & mean & 0.30 & 0.0 & 39.2 & 40.0 & 34.1 & 45.1 & 97.9 & 15.2 \\
WMT23, 12 streams & median & 0.05 & 92.5 & 3.9 & 3.8 & 3.9 & 3.8 & 99.6 & 0.4 \\
WMT23, 12 streams & median & 0.10 & 87.5 & 2.2 & 2.2 & 2.3 & 1.8 & 99.6 & 1.1 \\
WMT23, 12 streams & median & 0.20 & 97.5 & 0.6 & 0.5 & 0.5 & 0.5 & 99.9 & 0.3 \\
WMT23, 12 streams & median & 0.30 & 97.5 & 0.3 & 0.3 & 0.3 & 0.4 & 99.9 & 0.8 \\
WMT23, 12 streams & mean & 0.05 & 0.0 & 89.5 & 90.2 & 89.9 & 85.4 & 97.0 & 2.8 \\
WMT23, 12 streams & mean & 0.10 & 0.0 & 86.0 & 86.8 & 86.5 & 81.3 & 97.8 & 5.0 \\
WMT23, 12 streams & mean & 0.20 & 0.0 & 60.3 & 61.3 & 56.1 & 59.9 & 97.9 & 9.5 \\
WMT23, 12 streams & mean & 0.30 & 0.0 & 51.0 & 52.2 & 45.9 & 54.5 & 97.9 & 14.1 \\
\bottomrule
\end{tabular}
\caption{All depths and pooling rules at $\alpha=.10$, $V^+-V^-\le\lceil.02|S_0|\rceil$. Order columns are mean queried fractions of $mT$ in percent; empty active sets contribute zero. Z is the zero-query stopping percentage over 40 splits. Cov.\ is reference-priority $K/|S_0|$ in percent; Gap is its retained-count gap in outputs. These overlapping splits provide descriptive summaries.}
\label{tab:all-tolerant}
\end{table*}

\begin{table*}[t]
\centering\small
\begin{tabular}{llrrrrrrrr}
\toprule
 & & & & \multicolumn{4}{c}{Active inspection (\%)} & & \\
Panel & Pooling & $q$ & Z (\%) & Fixed & Random & Width & Ref. & Cov. (\%) & Gap (outputs) \\
\midrule
RewardBench 2 & median & 0.05 & 0.0 & 73.8 & 73.9 & 72.0 & 71.1 & 99.6 & 0.0 \\
RewardBench 2 & median & 0.10 & 0.0 & 58.3 & 58.3 & 55.5 & 53.4 & 100.0 & 0.0 \\
RewardBench 2 & median & 0.20 & 2.5 & 23.5 & 23.2 & 23.4 & 22.0 & 100.0 & 0.0 \\
RewardBench 2 & median & 0.30 & 85.0 & 0.8 & 0.8 & 0.8 & 0.8 & 100.0 & 0.0 \\
RewardBench 2 & mean & 0.05 & 0.0 & 93.7 & 93.7 & 93.8 & 90.2 & 100.0 & 0.0 \\
RewardBench 2 & mean & 0.10 & 0.0 & 90.5 & 90.4 & 90.0 & 86.9 & 100.0 & 0.0 \\
RewardBench 2 & mean & 0.20 & 0.0 & 82.0 & 82.0 & 79.5 & 79.5 & 100.0 & 0.0 \\
RewardBench 2 & mean & 0.30 & 0.0 & 69.7 & 69.8 & 63.0 & 71.0 & 100.0 & 0.0 \\
WMT23, 13 streams & median & 0.05 & 10.0 & 68.4 & 67.6 & 71.2 & 54.3 & 99.5 & 0.0 \\
WMT23, 13 streams & median & 0.10 & 97.5 & 0.2 & 0.2 & 0.2 & 0.2 & 100.0 & 0.0 \\
WMT23, 13 streams & median & 0.20 & 100.0 & 0.0 & 0.0 & 0.0 & 0.0 & 100.0 & 0.0 \\
WMT23, 13 streams & median & 0.30 & 100.0 & 0.0 & 0.0 & 0.0 & 0.0 & 100.0 & 0.0 \\
WMT23, 13 streams & mean & 0.05 & 0.0 & 94.1 & 94.0 & 94.6 & 88.8 & 99.7 & 0.0 \\
WMT23, 13 streams & mean & 0.10 & 0.0 & 90.9 & 91.0 & 91.4 & 86.3 & 100.0 & 0.0 \\
WMT23, 13 streams & mean & 0.20 & 0.0 & 62.2 & 62.1 & 55.7 & 62.1 & 100.0 & 0.0 \\
WMT23, 13 streams & mean & 0.30 & 0.0 & 47.4 & 47.2 & 40.2 & 51.6 & 100.0 & 0.0 \\
WMT23, 12 streams & median & 0.05 & 90.0 & 5.7 & 5.5 & 5.5 & 5.3 & 99.9 & 0.0 \\
WMT23, 12 streams & median & 0.10 & 80.0 & 5.2 & 5.1 & 5.9 & 4.8 & 100.0 & 0.0 \\
WMT23, 12 streams & median & 0.20 & 97.5 & 0.9 & 0.8 & 0.8 & 0.8 & 100.0 & 0.0 \\
WMT23, 12 streams & median & 0.30 & 95.0 & 0.9 & 0.9 & 0.8 & 0.9 & 100.0 & 0.0 \\
WMT23, 12 streams & mean & 0.05 & 0.0 & 92.2 & 92.2 & 92.0 & 88.7 & 99.1 & 0.0 \\
WMT23, 12 streams & mean & 0.10 & 0.0 & 88.6 & 88.6 & 88.4 & 84.3 & 100.0 & 0.0 \\
WMT23, 12 streams & mean & 0.20 & 0.0 & 66.0 & 66.0 & 60.6 & 64.8 & 100.0 & 0.0 \\
WMT23, 12 streams & mean & 0.30 & 0.0 & 58.1 & 58.2 & 51.4 & 60.0 & 100.0 & 0.0 \\
\bottomrule
\end{tabular}
\caption{All depths and pooling rules at $\alpha=.10$, $V^+-V^-=0$. Order columns are mean queried fractions of $mT$ in percent; empty active sets contribute zero. Z is the zero-query stopping percentage over 40 splits. Cov.\ is reference-priority $K/|S_0|$ in percent; Gap is its retained-count gap in outputs. These overlapping splits provide descriptive summaries.}
\label{tab:all-exact}
\end{table*}


\section{Post Hoc Quality Diagnostics at Matched Count}
\label{app:quality}
The stability policy optimizes selected retention under a churn budget. It does not optimize preference precision. RewardBench labels identify the preferred completion among each prompt's four outputs. For a returned set of size $K$, we compare its preference precision with the precision of the $K$ outputs having the highest base-consensus scores. The same seeded label-blind secondary key resolves ties. The comparator uses the same consensus and $K\le|S_0|$, so its top-$K$ outputs lie within the baseline cutoff list. 

This post hoc diagnostic leaves frozen fits, acquisition paths, stopping states and protocol unchanged. RewardBench and WMT23 informed development; this statement does not assert label-blind development history. For each split, the comparator matches returned count but is not constrained by the churn budget. Averaging and rounding follow Appendix~\ref{app:protocol}. All RewardBench comparisons have 40 defined splits (Table~\ref{tab:query-quality}).
\begin{table*}[t]
\centering\small
\begin{tabular}{lllr rrr}
\toprule
Pooling & $q$ & Strategy & $K$ & Precision & Top-$K$ & Difference \\
\midrule
median & 0.05 & fixed & 115.7 & 0.823 & 0.828 & -0.0048 \\
median & 0.05 & random & 115.6 & 0.825 & 0.828 & -0.0033 \\
median & 0.05 & largest\_width & 115.6 & 0.827 & 0.828 & -0.0008 \\
median & 0.05 & reference\_priority & 115.6 & 0.827 & 0.828 & -0.0009 \\
median & 0.10 & fixed & 229.1 & 0.702 & 0.706 & -0.0045 \\
median & 0.10 & random & 229.1 & 0.705 & 0.706 & -0.0015 \\
median & 0.10 & largest\_width & 229.1 & 0.705 & 0.706 & -0.0006 \\
median & 0.10 & reference\_priority & 229.1 & 0.701 & 0.706 & -0.0046 \\
median & 0.20 & fixed & 451.4 & 0.586 & 0.588 & -0.0020 \\
median & 0.20 & random & 451.4 & 0.587 & 0.588 & -0.0013 \\
median & 0.20 & largest\_width & 451.4 & 0.585 & 0.588 & -0.0030 \\
median & 0.20 & reference\_priority & 451.4 & 0.586 & 0.588 & -0.0027 \\
median & 0.30 & fixed & 686.0 & 0.511 & 0.511 & -0.0003 \\
median & 0.30 & random & 686.0 & 0.511 & 0.511 & -0.0003 \\
median & 0.30 & largest\_width & 686.0 & 0.511 & 0.511 & -0.0003 \\
median & 0.30 & reference\_priority & 686.0 & 0.511 & 0.511 & -0.0003 \\
mean & 0.05 & fixed & 117.3 & 0.825 & 0.827 & -0.0027 \\
mean & 0.05 & random & 117.3 & 0.827 & 0.827 & 0.0001 \\
mean & 0.05 & largest\_width & 117.3 & 0.829 & 0.827 & 0.0012 \\
mean & 0.05 & reference\_priority & 117.3 & 0.826 & 0.827 & -0.0009 \\
mean & 0.10 & fixed & 231.8 & 0.705 & 0.707 & -0.0019 \\
mean & 0.10 & random & 231.8 & 0.710 & 0.707 & 0.0028 \\
mean & 0.10 & largest\_width & 231.8 & 0.710 & 0.707 & 0.0032 \\
mean & 0.10 & reference\_priority & 231.8 & 0.709 & 0.707 & 0.0014 \\
mean & 0.20 & fixed & 452.8 & 0.591 & 0.593 & -0.0017 \\
mean & 0.20 & random & 452.8 & 0.594 & 0.593 & 0.0004 \\
mean & 0.20 & largest\_width & 452.8 & 0.595 & 0.593 & 0.0023 \\
mean & 0.20 & reference\_priority & 452.8 & 0.588 & 0.593 & -0.0048 \\
mean & 0.30 & fixed & 675.1 & 0.530 & 0.530 & -0.0004 \\
mean & 0.30 & random & 675.1 & 0.530 & 0.530 & -0.0003 \\
mean & 0.30 & largest\_width & 675.1 & 0.533 & 0.530 & 0.0034 \\
mean & 0.30 & reference\_priority & 675.1 & 0.529 & 0.530 & -0.0015 \\
\bottomrule
\end{tabular}
\caption{RewardBench preference diagnostic at matched returned count. All results use $\alpha=.10$ and the primary stopping criterion $V^+-V^-\le\lceil.02|S_0(P)|\rceil$. Within each of 40 grouped splits, the returned count $K$ determines the same-consensus top-$K$ comparator with label-blind seeded tie priorities. Precision and Top-$K$ are preferred-response fractions in the certified set and its matched-count comparator. Difference is certified-set precision minus comparator precision, in proportion units; multiplying by 100 gives percentage points. Entries are splitwise means, and differences are computed before display rounding. The comparator is constructed without a churn-budget constraint. This descriptive analysis was added after the inspection protocol; labels never enter acquisition or stopping. Count means use decimal half-up rounding after aggregating integer records.}
\label{tab:query-quality}
\end{table*}


\subsection{Official WMT23 MQM alignment}
All 173,784 metric-score entries match the frozen 24-metric, 557-segment, thirteen-stream panel using strict equality of parsed Python float values, without a tolerance. This is numerical equality, not raw-file byte identity. The 44 official-file acquisition identities are in \nolinkurl{review-revision/acceptance-review-implementation/wmt23-official-files.json}. We align MQM by system name and original segment identifier, separately from the original score builder. Twelve MT systems have labels on 460 of 557 segments each, giving 5,520 observed labels among 7,241 outputs. The \texttt{synthetic\_ref} output has no MQM label. Higher released MQM scores are better; missing annotations remain missing.

We preserve all original splits, mappings, masks and acquisition paths. The analysis covers both WMT23 populations, both pooling rules, four depths and 40 splits, with all four entry orders at exact and positive-gap stops, plus zero-query and full-information masks: 6,400 records. Labels do not choose the acquisition order, stopping point or protocol. For each returned set, we report total count $K$, annotated count $K_{\rm lab}$ and annotated coverage. We compare the returned annotated subset\textquotesingle s mean MQM with the mean for the $K_{\rm lab}$ highest-consensus outputs among all annotated test outputs, using the original label-blind tie key. This matches observed-label count; the comparator need not satisfy the churn budget. It does not estimate quality on unannotated outputs.

Table~\ref{tab:wmt23-native-quality} gives reference-priority positive-gap results at every depth; complete records retain all orders and stops. At median nominal $q=.05$, the mean splitwise share of returned outputs with MQM annotations is 28.3\% on thirteen streams and 77.3\% on twelve streams. The corresponding returned-minus-comparator mean MQM differences are approximately $-.0020$ and $-.0002$, respectively. Other configurations have differences of either sign. Low coverage, especially when \texttt{synthetic\_ref} is admitted, prevents extending these MQM comparisons to all returned outputs. The 40 overlapping splits provide descriptive averages, not independent samples for a significance test. A separate implementation rebuilds labels and both baseline masks and recomputes every diagnostic record.
\begin{table*}[t]
\centering\small
\begin{tabular}{llrrrrrr}
\toprule
Panel & Pooling & $q$ & $K$ & $K_{\rm lab}$ & Obs. (\%) & MQM & $\Delta$ \\
\midrule
WMT23 (13) & median & 0.05 & 134.8 & 38.5 & 28.3 & -0.3402 & -0.0020 \\
WMT23 (13) & median & 0.10 & 241.3 & 101.0 & 41.1 & -0.4141 & 0.0000 \\
WMT23 (13) & median & 0.20 & 489.1 & 297.8 & 60.7 & -0.7860 & 0.0000 \\
WMT23 (13) & median & 0.30 & 730.9 & 490.4 & 67.0 & -1.3198 & 0.0000 \\
WMT23 (13) & mean & 0.05 & 123.4 & 43.6 & 35.4 & -0.3046 & 0.0092 \\
WMT23 (13) & mean & 0.10 & 241.9 & 112.2 & 45.8 & -0.5335 & 0.0016 \\
WMT23 (13) & mean & 0.20 & 476.4 & 281.6 & 58.8 & -0.8979 & -0.0163 \\
WMT23 (13) & mean & 0.30 & 714.8 & 472.0 & 65.9 & -1.3320 & 0.0031 \\
WMT23 (12) & median & 0.05 & 113.3 & 87.6 & 77.3 & -0.3996 & -0.0002 \\
WMT23 (12) & median & 0.10 & 227.4 & 185.1 & 81.5 & -0.6429 & -0.0013 \\
WMT23 (12) & median & 0.20 & 457.3 & 374.4 & 81.9 & -1.0051 & 0.0013 \\
WMT23 (12) & median & 0.30 & 680.8 & 559.3 & 82.1 & -1.5591 & 0.0013 \\
WMT23 (12) & mean & 0.05 & 113.7 & 88.2 & 77.3 & -0.4973 & 0.0063 \\
WMT23 (12) & mean & 0.10 & 222.2 & 175.2 & 79.1 & -0.6437 & -0.0042 \\
WMT23 (12) & mean & 0.20 & 442.1 & 358.4 & 81.0 & -1.1279 & -0.0090 \\
WMT23 (12) & mean & 0.30 & 666.9 & 542.7 & 81.3 & -1.5970 & -0.0293 \\
\bottomrule
\end{tabular}
\caption{Post hoc WMT23 MQM diagnostic for the original reference-priority positive-gap stop, $\alpha=.10$, full-reference right ECDF. $K$ counts all returned outputs and $K_{\rm lab}$ the labeled ones. Obs. averages splitwise labeled fractions. MQM is higher-is-better on labeled outputs; $\Delta$ subtracts the consensus-ranked comparator with the same labeled count. Each row averages 40 overlapping splits; no population inference or missing-label imputation is made. All four acquisition methods, both stops and the two controls are included in the reproducibility results.}
\label{tab:wmt23-native-quality}
\end{table*}


\FloatBarrier
\section{Safe Retention Frontiers}
Figure~\ref{fig:frontiers} shows the three development panels under both pooling rules at $q=.05$. The single-split trajectory in Table~\ref{tab:query-prefix-example} uses a different depth. It shows how two query orders can certify different retained counts after the same number of observations.
\begin{figure*}[t]
\centering
\includegraphics[width=\textwidth]{query-frontiers.pdf}
\caption{At $q=.05$, $\alpha=.10$, colored curves show safe selected coverage $V^-/|S_0|$ and dotted lines show $V^*/|S_0|$. No optimistic frontier is plotted. The horizontal axis is the allowed fraction of $mT$ active entries; actual queries can be fewer after early exact stopping. Every plotted mask satisfies the risk budget. Curves and dotted lines are arithmetic means of splitwise coverage over 40 splits. Vertical ranges differ across panels.}
\label{fig:frontiers}
\end{figure*}

Table~\ref{tab:query-prefix-example} follows the first seeded split of thirteen-stream WMT23 with mean pooling at $q=.30$, whose aggregate comparison favors largest width. Both orders start with 646 safe outputs and an optimistic count of 715. After 1,122 queries, largest width certifies 694 safe outputs, compared with 674 for reference priority. Largest width reaches the 15-output tolerance after 1,249 queries, whereas reference priority needs 1,715. Both certify the same full-information count of 715 at exact stopping. The batch, risk budget and certificate criterion stay fixed; the trace illustrates the ordering difference without identifying its cause.
\begin{table*}[t]
\centering\small
\begin{tabular}{lrrrrrr}
\toprule
& \multicolumn{3}{c}{Reference priority} & \multicolumn{3}{c}{Largest width} \\
Stage & Queries & $V^-$ & $V^+$ & Queries & $V^-$ & $V^+$ \\
\midrule
Initial & 0 & 646 & 715 & 0 & 646 & 715 \\
10\% allowance & 449 & 652 & 715 & 449 & 662 & 715 \\
25\% allowance & 1122 & 674 & 715 & 1122 & 694 & 715 \\
First positive-gap stop & 1715 & 700 & 715 & 1249 & 700 & 715 \\
Exact stop & 2054 & 715 & 715 & 1558 & 715 & 715 \\
\bottomrule
\end{tabular}
\caption{Query-prefix illustration from the first seeded split of thirteen-stream WMT23, mean pooling, $q=.30$ and $\alpha=.10$. Both orders begin with 646 safe outputs and 4,488 active entries. Allowance rows compare the two strategies after the same number of queries; stopping rows report each strategy's own first stopping state. The allowed gap is 15 outputs. Counts come from the existing traces, not a new evaluation sample.}
\label{tab:query-prefix-example}
\end{table*}



\section{Budget and Composition Evidence}
\label{app:review-composition}
The budget grid contains $3\times2\times4\times40\times5=4{,}800$ risk-budget configurations. The 960 budget-independent panel--pooling--depth--split settings each use four strategies, yielding 3,840 orders. The four original orders (fixed, random, width and reference) depend on base geometry, reference fits, revealed entries and label-blind priorities, but not the risk budget. For each budget, we locate its first positive-gap and exact stopping prefixes on the same budget-independent acquisition path. The saved $\alpha=.10$ records, masks, checkpoints and trace prefixes permit independent replay. Five separate stopping problems do not count as five independent data samples.

For composition, the random-pool RNG seeds are 620271 (RewardBench) and 620272 (WMT). We sample 20 distinct seven-evaluator pools, excluding the original combination but allowing its individual evaluators. Full and twelve-stream WMT use the same evaluator subsets. The stored evaluator indices in \nolinkurl{review-revision/composition-protocol.json} determine the replay; the historical pool-generation API is not established. The metadata-defined WMT pool uses seed 620273: three members from BLEU, chrF, eBLEU, f200spBLEU and tokengram\_F, plus four from the remaining names. This category is based on evaluator identity; the remaining group is not restricted to learned metrics. All 62 added panel--pool units run both aggregators, four depths and 40 splits, yielding 19,840 configurations and 79,360 acquisition paths. Candidate sets are the full remaining evaluator set, and the reference cutoff and ridge are refitted for each pool. For a fixed evaluator and reference-output population, ECDF maps can be reused across pools and budgets. Changing the base pool still requires refitting the consensus cutoff and ridge targets.

The companion grids report reference priority at all budgets and depths; the composition tables give descriptive pool-level distributions. Complete result paths and field meanings for acquisition orders, masks, stopping targets, quality diagnostics, witness decompositions and replacement-specific losses are listed in \nolinkurl{RESULT_DATA_DICTIONARY.md}. The anonymous repository includes the listed summaries and numerical query traces. Original third-party text and full score arrays remain outside it; redistribution clearance is still incomplete. Pool-level quantiles summarize sensitivity across sampled pools, not population confidence intervals.
\subsection{Invariant and budget-admissible zero-query retention}
\label{app:zero-decomposition}
Let $I=\{i\in S_0:p_i=0\}$ and let $C_0$ be the zero-query mask computed with $Tp_i$. Every invariant selected output is retained for free. For each split,
\[
 \frac{|S_0\cap C_0|}{|S_0|}
 =\frac{|I|}{|S_0|}+\frac{|(S_0\cap C_0)\setminus I|}{|S_0|}.
\]
Both terms divide by $|S_0|$. The second counts retained noninvariant outputs whose combined upper cost fits the budget. Zero-query stopping additionally requires the initial retained-count gap to meet the tolerance. Tables~\ref{tab:zero-decomposition} and~\ref{tab:zero-decomposition-mean} report all six panels at four depths under median and mean pooling, respectively, using each panel\textquotesingle s original evaluator pool and averaging splitwise fractions. Removing the \texttt{synthetic\_ref} output stream changes the output population as well as the refitted maps and cutoff. The twelve- and thirteen-stream comparison does not isolate deletion from these protocol changes.
\begin{table*}[t]
\centering\small
\begin{tabular}{lrrrrr}
\toprule
Panel & $q$ & Invariant & Noninv. & Zero & Z \\
\midrule
RewardBench 2 & 0.05 & 46.6 & 34.3 & 80.9 & 0.0 \\
RewardBench 2 & 0.1 & 50.5 & 34.4 & 84.9 & 2.5 \\
RewardBench 2 & 0.2 & 59.6 & 35.1 & 94.6 & 7.5 \\
RewardBench 2 & 0.3 & 66.5 & 33.3 & 99.8 & 97.5 \\
WMT23, 13 streams & 0.05 & 54.4 & 30.6 & 84.9 & 10.0 \\
WMT23, 13 streams & 0.1 & 78.8 & 21.2 & 100.0 & 100.0 \\
WMT23, 13 streams & 0.2 & 77.3 & 22.7 & 100.0 & 100.0 \\
WMT23, 13 streams & 0.3 & 76.4 & 23.6 & 100.0 & 100.0 \\
WMT23, 12 streams & 0.05 & 75.1 & 24.1 & 99.2 & 92.5 \\
WMT23, 12 streams & 0.1 & 71.7 & 27.7 & 99.4 & 87.5 \\
WMT23, 12 streams & 0.2 & 73.5 & 26.3 & 99.9 & 97.5 \\
WMT23, 12 streams & 0.3 & 72.0 & 27.9 & 99.9 & 97.5 \\
WMT24 en-de & 0.05 & 85.3 & 14.2 & 99.5 & 95.0 \\
WMT24 en-de & 0.1 & 77.9 & 21.8 & 99.7 & 87.5 \\
WMT24 en-de & 0.2 & 76.2 & 23.8 & 100.0 & 100.0 \\
WMT24 en-de & 0.3 & 76.2 & 23.8 & 100.0 & 100.0 \\
WMT24 en-es & 0.05 & 83.7 & 16.1 & 99.8 & 97.5 \\
WMT24 en-es & 0.1 & 78.2 & 21.8 & 99.9 & 100.0 \\
WMT24 en-es & 0.2 & 76.6 & 23.4 & 100.0 & 100.0 \\
WMT24 en-es & 0.3 & 78.1 & 21.9 & 100.0 & 100.0 \\
WMT24 ja-zh & 0.05 & 77.8 & 20.3 & 98.1 & 77.5 \\
WMT24 ja-zh & 0.1 & 77.1 & 22.2 & 99.3 & 87.5 \\
WMT24 ja-zh & 0.2 & 75.5 & 24.4 & 99.9 & 97.5 \\
WMT24 ja-zh & 0.3 & 77.9 & 22.1 & 100.0 & 100.0 \\
\bottomrule
\end{tabular}
\caption{Median pooling, $\alpha=.10$, full-reference right ECDF, original evaluator pools. Invariant and Noninv. are selected-output shares retained through $p_i=0$ and budget-admissible $p_i>0$, respectively. Before rounding, their sum is zero-query selected coverage (Zero). Shares divide by $|S_0|$ within each split, then average over 40 splits. Z is the percentage of splits stopping at zero queries when $V^+-V^-\le\lceil.02|S_0|\rceil$. Invariant, Noninv., Zero and Z are percentages, rounded separately; $q$ is the nominal selection fraction.}
\label{tab:zero-decomposition}
\end{table*}

\begin{table*}[t]
\centering\small
\begin{tabular}{lrrrrr}
\toprule
Panel & $q$ & Invariant & Noninv. & Zero & Z \\
\midrule
RewardBench 2 & 0.05 & 0.6 & 19.7 & 20.3 & 0.0 \\
RewardBench 2 & 0.1 & 7.2 & 24.5 & 31.7 & 0.0 \\
RewardBench 2 & 0.2 & 22.2 & 29.5 & 51.7 & 0.0 \\
RewardBench 2 & 0.3 & 37.3 & 33.4 & 70.8 & 0.0 \\
WMT23, 13 streams & 0.05 & 0.0 & 20.3 & 20.3 & 0.0 \\
WMT23, 13 streams & 0.1 & 0.0 & 19.9 & 19.9 & 0.0 \\
WMT23, 13 streams & 0.2 & 49.2 & 28.5 & 77.6 & 0.0 \\
WMT23, 13 streams & 0.3 & 59.0 & 28.7 & 87.7 & 0.0 \\
WMT23, 12 streams & 0.05 & 0.0 & 20.6 & 20.6 & 0.0 \\
WMT23, 12 streams & 0.1 & 3.2 & 21.1 & 24.3 & 0.0 \\
WMT23, 12 streams & 0.2 & 47.9 & 28.0 & 75.8 & 0.0 \\
WMT23, 12 streams & 0.3 & 53.3 & 29.6 & 82.9 & 0.0 \\
WMT24 en-de & 0.05 & 0.0 & 22.0 & 22.0 & 0.0 \\
WMT24 en-de & 0.1 & 29.9 & 24.8 & 54.7 & 0.0 \\
WMT24 en-de & 0.2 & 41.6 & 25.7 & 67.3 & 0.0 \\
WMT24 en-de & 0.3 & 48.4 & 28.0 & 76.4 & 0.0 \\
WMT24 en-es & 0.05 & 0.0 & 21.2 & 21.2 & 0.0 \\
WMT24 en-es & 0.1 & 27.8 & 23.8 & 51.6 & 0.0 \\
WMT24 en-es & 0.2 & 40.7 & 25.6 & 66.3 & 0.0 \\
WMT24 en-es & 0.3 & 51.4 & 27.7 & 79.2 & 0.0 \\
WMT24 ja-zh & 0.05 & 0.0 & 21.4 & 21.4 & 0.0 \\
WMT24 ja-zh & 0.1 & 0.0 & 21.8 & 21.8 & 0.0 \\
WMT24 ja-zh & 0.2 & 11.7 & 26.4 & 38.1 & 0.0 \\
WMT24 ja-zh & 0.3 & 24.1 & 29.2 & 53.3 & 0.0 \\
\bottomrule
\end{tabular}
\caption{Mean pooling, $\alpha=.10$, full-reference right ECDF, original evaluator pools. Invariant and Noninv. are selected-output shares retained through $p_i=0$ and budget-admissible $p_i>0$, respectively. Before rounding, their sum is zero-query selected coverage (Zero). Shares divide by $|S_0|$ within each split, then average over 40 splits. Z is the percentage of splits stopping at zero queries when $V^+-V^-\le\lceil.02|S_0|\rceil$. Invariant, Noninv., Zero and Z are percentages, rounded separately; $q$ is the nominal selection fraction.}
\label{tab:zero-decomposition-mean}
\end{table*}


\subsection{Complete composition and budget results}
Evaluator composition changes inspection requirements (Tables~\ref{tab:review-composition-median}, \ref{tab:review-composition-mean} and~\ref{tab:review-balanced}). For thirteen-stream WMT23 at $q=.05$, the 10th--90th percentile range of pool-level mean reference-priority inspection is 0.0--67.8\% under median pooling and 79.0--84.4\% under mean pooling. These ranges describe the 20 sampled pools. 
\begin{table*}[t]
\centering\small
\begin{tabular}{llrrrrr}
\toprule
Panel & Pooling & $q$ & Zero (\%) & Ref. (\%) & 10--90\% & $\Delta$ Width \\
\midrule
RewardBench 2 & median & 0.05 & 81.6 & 55.8 & 43.7--68.9 & -6.3 \\
RewardBench 2 & median & 0.10 & 87.2 & 39.0 & 20.8--49.3 & -6.7 \\
RewardBench 2 & median & 0.20 & 96.2 & 9.0 & 0.9--22.0 & -1.0 \\
RewardBench 2 & median & 0.30 & 100.0 & 0.0 & 0.0--0.7 & 0.0 \\
WMT23, 13 streams & median & 0.05 & 98.6 & 2.5 & 0.0--67.8 & 0.0 \\
WMT23, 13 streams & median & 0.10 & 99.4 & 1.4 & 0.0--18.5 & 0.0 \\
WMT23, 13 streams & median & 0.20 & 100.0 & 0.0 & 0.0--3.1 & 0.0 \\
WMT23, 13 streams & median & 0.30 & 100.0 & 0.0 & 0.0--1.1 & 0.0 \\
WMT23, 12 streams & median & 0.05 & 98.1 & 6.4 & 2.8--29.3 & -0.4 \\
WMT23, 12 streams & median & 0.10 & 98.6 & 3.6 & 0.6--13.2 & -0.1 \\
WMT23, 12 streams & median & 0.20 & 99.6 & 0.9 & 0.0--6.6 & -0.1 \\
WMT23, 12 streams & median & 0.30 & 99.9 & 0.1 & 0.0--4.6 & 0.0 \\
\bottomrule
\end{tabular}
\caption{Primary stopping uses $V^+-V^-\le\lceil.02|S_0|\rceil$. Median pooling: Random evaluator composition at $\alpha=.10$, 20 pools with 40 splits each. Zero is the median of pool-wise mean zero-query selected coverage. Ref.\ is the median pool-wise mean queried fraction of $mT$ in percent, with descriptive 10th--90th percentiles; empty active sets contribute zero. Within each pool we average splitwise Reference minus Width active-inspection differences. $\Delta$ is the median of these 20 differences in percentage points; positive values mean more Reference queries.}
\label{tab:review-composition-median}
\end{table*}

\begin{table*}[t]
\centering\small
\begin{tabular}{llrrrrr}
\toprule
Panel & Pooling & $q$ & Zero (\%) & Ref. (\%) & 10--90\% & $\Delta$ Width \\
\midrule
RewardBench 2 & mean & 0.05 & 20.9 & 86.9 & 84.1--94.4 & -4.8 \\
RewardBench 2 & mean & 0.10 & 29.9 & 83.8 & 81.3--88.5 & -4.4 \\
RewardBench 2 & mean & 0.20 & 52.2 & 75.4 & 74.6--77.3 & -0.3 \\
RewardBench 2 & mean & 0.30 & 70.9 & 66.2 & 62.6--67.3 & 5.6 \\
WMT23, 13 streams & mean & 0.05 & 20.8 & 81.6 & 79.0--84.4 & -4.3 \\
WMT23, 13 streams & mean & 0.10 & 31.8 & 78.7 & 74.6--81.4 & -4.3 \\
WMT23, 13 streams & mean & 0.20 & 72.9 & 60.8 & 55.0--68.3 & 2.2 \\
WMT23, 13 streams & mean & 0.30 & 82.8 & 51.9 & 48.4--59.8 & 6.6 \\
WMT23, 12 streams & mean & 0.05 & 21.1 & 81.2 & 78.5--87.0 & -3.9 \\
WMT23, 12 streams & mean & 0.10 & 40.9 & 76.3 & 71.3--79.7 & -3.5 \\
WMT23, 12 streams & mean & 0.20 & 69.4 & 63.9 & 61.1--69.3 & 0.1 \\
WMT23, 12 streams & mean & 0.30 & 79.7 & 57.4 & 54.0--63.5 & 5.8 \\
\bottomrule
\end{tabular}
\caption{Primary stopping uses $V^+-V^-\le\lceil.02|S_0|\rceil$. Mean pooling: Random evaluator composition at $\alpha=.10$, 20 pools with 40 splits each. Zero is the median of pool-wise mean zero-query selected coverage. Ref.\ is the median pool-wise mean queried fraction of $mT$ in percent, with descriptive 10th--90th percentiles; empty active sets contribute zero. Within each pool we average splitwise Reference minus Width active-inspection differences. $\Delta$ is the median of these 20 differences in percentage points; positive values mean more Reference queries.}
\label{tab:review-composition-mean}
\end{table*}

\begin{table*}[t]
\centering\small
\begin{tabular}{llrrrrrr}
\toprule
Panel & Pooling & $q$ & Zero & Fixed & Random & Width & Ref. \\
\midrule
WMT23, 13 streams & median & 0.05 & 79.1 & 89.6 & 90.0 & 89.9 & 85.5 \\
WMT23, 13 streams & median & 0.10 & 98.8 & 3.7 & 3.7 & 3.0 & 3.2 \\
WMT23, 13 streams & median & 0.20 & 100.0 & 0.0 & 0.0 & 0.0 & 0.0 \\
WMT23, 13 streams & median & 0.30 & 99.9 & 0.2 & 0.2 & 0.1 & 0.1 \\
WMT23, 13 streams & mean & 0.05 & 20.4 & 85.4 & 86.5 & 85.6 & 81.5 \\
WMT23, 13 streams & mean & 0.10 & 28.1 & 83.3 & 84.3 & 84.1 & 78.5 \\
WMT23, 13 streams & mean & 0.20 & 81.3 & 49.4 & 50.5 & 47.3 & 51.1 \\
WMT23, 13 streams & mean & 0.30 & 85.6 & 44.0 & 44.8 & 41.2 & 48.0 \\
WMT23, 12 streams & median & 0.05 & 97.8 & 10.4 & 10.8 & 10.6 & 10.1 \\
WMT23, 12 streams & median & 0.10 & 99.5 & 1.7 & 1.8 & 1.6 & 1.5 \\
WMT23, 12 streams & median & 0.20 & 99.2 & 2.0 & 2.1 & 1.7 & 1.8 \\
WMT23, 12 streams & median & 0.30 & 99.3 & 1.5 & 1.5 & 1.2 & 1.3 \\
WMT23, 12 streams & mean & 0.05 & 21.2 & 85.7 & 86.7 & 86.3 & 81.7 \\
WMT23, 12 streams & mean & 0.10 & 52.0 & 74.7 & 75.8 & 73.6 & 70.9 \\
WMT23, 12 streams & mean & 0.20 & 74.1 & 61.2 & 62.3 & 59.5 & 59.7 \\
WMT23, 12 streams & mean & 0.30 & 80.0 & 56.7 & 57.9 & 54.4 & 57.3 \\
\bottomrule
\end{tabular}
\caption{Primary stopping uses $V^+-V^-\le\lceil.02|S_0|\rceil$. WMT pool with three surface-overlap metrics and four remaining metrics (not necessarily learned). All depths, $\alpha=.10$, 40 split means. Zero is zero-query $K/|S_0|$ in percent; order columns are queried $mT$ percentages, with zero for empty active sets.}
\label{tab:review-balanced}
\end{table*}

\begin{figure*}[t]
\centering
\includegraphics[width=\textwidth]{review-budgets.pdf}
\caption{Median pooling at $q=.05$, all five budgets. Points average queried fractions of $mT$ active selected entries over 40 splits, stopping at $V^+-V^-\le\lceil.02|S_0|\rceil$ outputs. Empty active sets contribute zero inspection (Appendix~\ref{app:protocol}). Orders are shared across budgets; first stopping states are recorded separately.}
\label{fig:review-budgets}
\end{figure*}
Full five-budget grids for all three original panels and both pooling rules are retained as \texttt{review-budget-*} in \nolinkurl{supplement/result_grids.csv}; the companion JSON preserves captions and units.

\FloatBarrier
\section{WMT24 Admission and Quality Diagnostics}
\label{app:review-quality}
\begin{table*}[t]
\centering\small
\begin{tabular}{llrrrr}
\toprule
 &  & \multicolumn{3}{c}{Selected coverage (\%)} & \\
Language & Pooling & Zero & Opt. & Ref. & Active (\%) \\
\midrule
en-de & median & 99.5 & 100.0 & 99.9 & 1.3 \\
en-de & mean & 22.0 & 99.9 & 97.9 & 83.3 \\
en-es & median & 99.8 & 100.0 & 99.9 & 1.1 \\
en-es & mean & 21.2 & 99.9 & 97.8 & 82.1 \\
ja-zh & median & 98.1 & 99.8 & 99.3 & 11.3 \\
ja-zh & mean & 21.4 & 97.5 & 96.0 & 85.9 \\
\bottomrule
\end{tabular}
\caption{Ref.\ stops at $V^+-V^-\le\lceil .02|S_0|\rceil$; Opt.\ uses full information. WMT24, $q=.05$, $\alpha=.10$: 40 document-grouped split means. Zero, Opt.\ and Ref.: zero-query, full-information and reference-priority coverage. Active: reference inspection ($100Q/(mT)$), zero for empty active sets. Admission uses metadata and finite scores.}
\label{tab:review-external}
\end{table*}

We obtained the official MTME archive from \url{https://data.statmt.org/wmt26/mt-metrics-eval-v2.tgz}, the replacement host identified by the maintainer in \url{https://github.com/google-research/mt-metrics-eval/issues/27}. The prepared source/code packages contain \path{review-revision/acquisition/official-hosting.json} and \path{review-revision/external-admission.json}, identifying the archive and used files. The analysis uses the WMT24 files inside this archive; the wmt26 path is the hosting location.

In these files, paragraph granularity is labeled \texttt{seg}; these units are paragraphs, not individual sentences. Standard references are refB for en-de and refA for en-es/ja-zh. Using system-name metadata, we exclude human references as outputs and retain other MT systems common to the admitted metrics, including benchmark outliers. Every retained paragraph has finite scores for every retained MT system and admitted default-reference metric. Score files are parsed into floating-point arrays, with systems sorted by name and paragraphs in source order. On the common finite paragraphs, deduplication compares SHA-256 signatures of the array byte representations, keeping the alphabetically first metric for each signature. It does not compare raw text files. One duplicate is removed in each English-source panel; no paragraph or MT system is removed on the basis of MQM, churn or query performance.

The en-de, en-es and ja-zh panels contain 171, 171 and 197 source documents, respectively. For split $r\in\{0,\ldots,39\}$, the document-shuffle seed is $924000+r$. Let $D$ be the number of documents. Unique document identifiers are collected in first-appearance order and shuffled with Python \texttt{random.Random(seed).shuffle}, and partitioned using $d_{\rm third}=\lfloor D/3\rfloor$. The first two partitions each contain $d_{\rm third}$ documents and the remainder forms test: 171 documents give 57/57/57 and 197 give 65/65/67. Documents remain intact; output counts depend on document lengths. The second partition is reserved for calibration and is not used by these WMT24 query experiments. Unequal document sizes make test count $N$ and nominal $k$ vary across splits. ECDF maps and cutoff use all outputs in reference documents. Tie and random-query seeds add 700001 and 99173, respectively, to the current document-shuffle seed. Each panel uses $M=7$; en-de and en-es have $T=14$, and ja-zh has $T=15$. Ridge features and the unnormalized squared-error objective with penalty one follow Appendix~\ref{app:protocol}. We use both $\varepsilon=\lceil.02|S_0|\rceil$ and exact stopping. The exact ordered base and candidate names and removed duplicates appear in \nolinkurl{supplement/wmt24_evaluator_manifest.csv}, derived from the admission record. We evaluate all three admitted directions. The pre-results freeze covered the risk objective and panel-inclusion criteria; the later numerical repair is described in the numerical-query-priorities paragraph below. Evaluator-training independence is not claimed.

\paragraph{Numerical query priorities.}
\label{app:wmt24-priority-repair}
For WMT24 entry acquisition, each predicted flip count is converted separately to an integer fixed-point value $e_{ir}=\operatorname{rint}(10^{12}\hat f_{ir})$, with nearest-even rounding. Active rows are ranked in ascending order of $10^{12}\ell_i+\sum_{r\text{ unqueried}}e_{ir}$. Within the chosen row, candidates are ordered by decreasing predicted upper-cost reduction from the rounded per-cell values; candidate index breaks ties. Row ties use the frozen, seeded label-blind item priority defined in Appendix~\ref{app:protocol}. Each value is rounded before summation. 

These fixed-point values affect query order only; observed flip costs, budgets and certificates retain integer arithmetic. RewardBench and WMT23 entry acquisition retain their original priority arithmetic; only WMT24 entry acquisition uses this fixed-point rule. Column acquisition follows the separate rule in Appendix~\ref{app:mapping-column}. The fixed-point rule for WMT24 entry acquisition was introduced after an exact-replay failure and was outside the pre-results freeze. The first fresh replay differed in 22 result fields for en-de, mean pooling, split 37, $q=.10$, reference priority. The two runs used 3,063 and 3,059 queries and both ended with $(V^-,V^+)=(477,487)$; their masks and quality diagnostics differed. The recorded first query-order difference was at query 71. The original prediction arrays were not retained, and the records do not establish a particular floating-point backend as the cause. After introducing fixed-point priorities, all 960 WMT24 configurations were rerun and checked against an independent replay; the unaffected original-panel computations were reused.

We align native MQM by exact system name and original zero-based paragraph index. Annotations are available for 8,279 of 25,948 en-de outputs, 8,722 of 22,954 en-es outputs and 7,840 of 15,884 ja-zh outputs. Values missing in the official release stay missing. Selection uses the complete metric panel; it is not filtered to annotated systems or paragraphs. For a returned list of size $K$, the original matched-$K$ comparator uses the same consensus and label-blind tie priorities. Its observed-label mean and the certified list's observed-label mean may involve different numbers of annotated outputs. For a descriptive diagnostic, we extract the $K_{\rm lab}$ annotated returned outputs without changing the certified mask. We compare their mean with the top $K_{\rm lab}$ consensus outputs among all test outputs having MQM. This comparator matches annotated count and concerns only the observed-label subpopulation. MQM means use only annotated outputs within each list. A difference is defined only when both compared means are defined; splitwise differences are then averaged without imputing missing values. All reference-priority quality comparisons reported here have 40 defined splits. Neither comparator is constructed under the churn-budget constraint. These diagnostics do not estimate quality for unannotated outputs.

Tables~\ref{tab:review-mqm-median} and~\ref{tab:review-mqm-mean} report the median and mean diagnostics separately.
\begin{table*}[t]
\centering\small
\begin{tabular}{llrrrrr}
\toprule
Language & Pooling & $q$ & Obs. (\%) & MQM & $\Delta\mathrm{MQM}_K$ & $\Delta\mathrm{MQM}_{K_{\rm lab}}$ \\
\midrule
en-de & median & 0.05 & 37.2 & -0.0695 & -0.0010 & -0.0010 \\
en-de & median & 0.10 & 36.7 & -0.2688 & 0.0014 & 0.0009 \\
en-de & median & 0.20 & 38.2 & -0.7180 & 0.0000 & 0.0000 \\
en-de & median & 0.30 & 39.0 & -1.1547 & 0.0000 & 0.0000 \\
en-es & median & 0.05 & 30.9 & -0.0668 & 0.0000 & 0.0000 \\
en-es & median & 0.10 & 34.0 & -0.1181 & -0.0001 & -0.0001 \\
en-es & median & 0.20 & 38.3 & -0.1791 & -0.0001 & -0.0001 \\
en-es & median & 0.30 & 40.1 & -0.2218 & 0.0000 & 0.0000 \\
ja-zh & median & 0.05 & 54.8 & -0.1467 & 0.0018 & 0.0021 \\
ja-zh & median & 0.10 & 56.2 & -0.3327 & -0.0003 & 0.0001 \\
ja-zh & median & 0.20 & 59.8 & -0.6149 & -0.0007 & -0.0008 \\
ja-zh & median & 0.30 & 61.7 & -0.9359 & 0.0000 & 0.0000 \\
\bottomrule
\end{tabular}
\caption{Median pooling: native MQM after reference-priority stopping. All results use $\alpha=.10$ and the primary stopping criterion $V^+-V^-\le\lceil.02|S_0(P)|\rceil$. Entries average 40 split values; differences precede rounding. Obs.\ is the percentage of returned outputs with MQM annotations. MQM is higher-is-better. $\Delta\mathrm{MQM}_K$ is the certified set\textquotesingle s observed-label mean MQM minus the matched-$K$ comparator\textquotesingle s mean. $\Delta\mathrm{MQM}_{K_{\rm lab}}$ uses the top-$K_{\rm lab}$ annotated-output comparator, matching the number of annotated returned outputs. The matched-$K$ comparator matches total count; each MQM mean uses the annotated outputs in its own list. Positive differences favor the certified set. Missing labels are not imputed; both diagnostics concern annotated outputs.}
\label{tab:review-mqm-median}
\end{table*}

\begin{table*}[t]
\centering\small
\begin{tabular}{llrrrrr}
\toprule
Language & Pooling & $q$ & Obs. (\%) & MQM & $\Delta\mathrm{MQM}_K$ & $\Delta\mathrm{MQM}_{K_{\rm lab}}$ \\
\midrule
en-de & mean & 0.05 & 36.1 & -0.0636 & 0.0003 & 0.0040 \\
en-de & mean & 0.10 & 36.7 & -0.2389 & -0.0068 & -0.0065 \\
en-de & mean & 0.20 & 37.8 & -0.6264 & -0.0030 & -0.0034 \\
en-de & mean & 0.30 & 39.2 & -1.1102 & -0.0113 & -0.0102 \\
en-es & mean & 0.05 & 29.7 & -0.0577 & -0.0016 & -0.0002 \\
en-es & mean & 0.10 & 34.8 & -0.1010 & 0.0006 & -0.0005 \\
en-es & mean & 0.20 & 38.8 & -0.1609 & 0.0011 & 0.0030 \\
en-es & mean & 0.30 & 40.5 & -0.2081 & -0.0017 & -0.0018 \\
ja-zh & mean & 0.05 & 55.8 & -0.3203 & -0.0015 & -0.0009 \\
ja-zh & mean & 0.10 & 60.7 & -0.4475 & 0.0057 & 0.0053 \\
ja-zh & mean & 0.20 & 64.0 & -0.6797 & 0.0049 & 0.0063 \\
ja-zh & mean & 0.30 & 62.4 & -0.9202 & 0.0012 & 0.0039 \\
\bottomrule
\end{tabular}
\caption{Mean pooling: native MQM after reference-priority stopping. All results use $\alpha=.10$ and the primary stopping criterion $V^+-V^-\le\lceil.02|S_0(P)|\rceil$. Entries average 40 split values; differences precede rounding. Obs.\ is the percentage of returned outputs with MQM annotations. MQM is higher-is-better. $\Delta\mathrm{MQM}_K$ is the certified set\textquotesingle s observed-label mean MQM minus the matched-$K$ comparator\textquotesingle s mean. $\Delta\mathrm{MQM}_{K_{\rm lab}}$ uses the top-$K_{\rm lab}$ annotated-output comparator, matching the number of annotated returned outputs. The matched-$K$ comparator matches total count; each MQM mean uses the annotated outputs in its own list. Positive differences favor the certified set. Missing labels are not imputed; both diagnostics concern annotated outputs.}
\label{tab:review-mqm-mean}
\end{table*}


Correlation pooled across inputs alone does not establish discrimination among outputs for the same input. To illustrate this, official source-only and reference-only sentinel score streams \citep{perrella2024sentinel} are evaluated descriptively on the same observed MQM entries. Pooled Pearson correlation uses all matching entries; within-paragraph pair accuracy uses annotated system pairs with unequal MQM scores and gives one half to metric ties. The probes contain no candidate translation information and tie all systems within a paragraph. These streams are probes; their diagnostic performance does not trigger evaluator removal, cutoff refitting or candidate-family changes. This differs from removing WMT23\textquotesingle s \texttt{synthetic\_ref} output stream, which changes the output axis, not the evaluator axis.
Table~\ref{tab:review-sentinel} reports the two sentinel diagnostics without changing the certified selection protocol.
\begin{table*}[t]
\centering\small
\begin{tabular}{llrrrr}
\toprule
Language & Probe & Pearson & Pairs & Metric ties (\%) & Accuracy \\
\midrule
en-de & Source & 0.257 & 39819 & 100.0 & 0.500 \\
en-de & Reference & 0.245 & 39819 & 100.0 & 0.500 \\
en-es & Source & 0.207 & 21736 & 100.0 & 0.500 \\
en-es & Reference & 0.195 & 21736 & 100.0 & 0.500 \\
ja-zh & Source & 0.244 & 29850 & 100.0 & 0.500 \\
ja-zh & Reference & 0.252 & 29850 & 100.0 & 0.500 \\
\bottomrule
\end{tabular}
\caption{Official missing-information sentinel probes on native annotated WMT24 entries. Pooled Pearson can be positive despite complete within-paragraph ties across MT systems. Pair accuracy assigns one half to metric ties and excludes human-score ties. This descriptive diagnostic changes no candidate family and does not identify invalid real metrics.}
\label{tab:review-sentinel}
\end{table*}


\FloatBarrier
\section{Variation across Individual Replacements}
\label{app:individual-replacements}
\begin{table*}[t]
\centering\small
\begin{tabular}{lrrrr}
\toprule
Panel & $q$ & Avg. & Mean max & Over (\%) \\
\midrule
RewardBench 2 & 0.05 & 0.077 & 0.164 & 17.3 \\
RewardBench 2 & 0.10 & 0.066 & 0.128 & 9.3 \\
RewardBench 2 & 0.20 & 0.047 & 0.088 & 0.3 \\
RewardBench 2 & 0.30 & 0.041 & 0.068 & 0.0 \\
WMT23, 13 streams & 0.05 & 0.078 & 0.234 & 38.0 \\
WMT23, 13 streams & 0.10 & 0.035 & 0.072 & 0.3 \\
WMT23, 13 streams & 0.20 & 0.035 & 0.073 & 0.0 \\
WMT23, 13 streams & 0.30 & 0.033 & 0.069 & 0.1 \\
WMT23, 12 streams & 0.05 & 0.043 & 0.114 & 6.1 \\
WMT23, 12 streams & 0.10 & 0.048 & 0.112 & 5.5 \\
WMT23, 12 streams & 0.20 & 0.042 & 0.085 & 1.6 \\
WMT23, 12 streams & 0.30 & 0.040 & 0.082 & 1.1 \\
WMT24 en-de & 0.05 & 0.022 & 0.063 & 2.4 \\
WMT24 en-de & 0.10 & 0.036 & 0.081 & 1.1 \\
WMT24 en-de & 0.20 & 0.034 & 0.073 & 0.0 \\
WMT24 en-de & 0.30 & 0.032 & 0.065 & 0.0 \\
WMT24 en-es & 0.05 & 0.029 & 0.070 & 2.4 \\
WMT24 en-es & 0.10 & 0.036 & 0.081 & 0.9 \\
WMT24 en-es & 0.20 & 0.036 & 0.075 & 0.6 \\
WMT24 en-es & 0.30 & 0.031 & 0.064 & 0.0 \\
WMT24 ja-zh & 0.05 & 0.046 & 0.106 & 9.3 \\
WMT24 ja-zh & 0.10 & 0.048 & 0.110 & 8.3 \\
WMT24 ja-zh & 0.20 & 0.047 & 0.103 & 4.9 \\
WMT24 ja-zh & 0.30 & 0.041 & 0.088 & 1.5 \\
\bottomrule
\end{tabular}
\caption{Primary stopping uses $V^+-V^-\le\lceil.02|S_0|\rceil$. Median pooling: reference-priority returns at $\alpha=.10$, all depths. Avg.\ averages the splitwise replacement-average churn $R(C)$ over 40 splits. Mean max averages the maximum replacement-specific churn within each split; it is not the maximum across all splits. Over is the mean percentage of replacements exceeding .10. Each split separately satisfies $R(C)\le\alpha$; individual replacements need not.}
\label{tab:review-replacement-median}
\end{table*}

\begin{table*}[t]
\centering\small
\begin{tabular}{lrrrr}
\toprule
Panel & $q$ & Avg. & Mean max & Over (\%) \\
\midrule
RewardBench 2 & 0.05 & 0.068 & 0.193 & 14.6 \\
RewardBench 2 & 0.10 & 0.059 & 0.144 & 12.3 \\
RewardBench 2 & 0.20 & 0.045 & 0.104 & 4.1 \\
RewardBench 2 & 0.30 & 0.036 & 0.073 & 0.0 \\
WMT23, 13 streams & 0.05 & 0.072 & 0.327 & 27.9 \\
WMT23, 13 streams & 0.10 & 0.056 & 0.211 & 14.2 \\
WMT23, 13 streams & 0.20 & 0.031 & 0.093 & 0.3 \\
WMT23, 13 streams & 0.30 & 0.031 & 0.072 & 0.0 \\
WMT23, 12 streams & 0.05 & 0.069 & 0.290 & 20.3 \\
WMT23, 12 streams & 0.10 & 0.049 & 0.185 & 8.8 \\
WMT23, 12 streams & 0.20 & 0.037 & 0.094 & 0.6 \\
WMT23, 12 streams & 0.30 & 0.036 & 0.080 & 0.0 \\
WMT24 en-de & 0.05 & 0.063 & 0.175 & 23.7 \\
WMT24 en-de & 0.10 & 0.032 & 0.075 & 0.3 \\
WMT24 en-de & 0.20 & 0.035 & 0.087 & 1.1 \\
WMT24 en-de & 0.30 & 0.030 & 0.072 & 0.2 \\
WMT24 en-es & 0.05 & 0.056 & 0.177 & 17.0 \\
WMT24 en-es & 0.10 & 0.037 & 0.087 & 1.3 \\
WMT24 en-es & 0.20 & 0.035 & 0.081 & 0.8 \\
WMT24 en-es & 0.30 & 0.029 & 0.068 & 0.2 \\
WMT24 ja-zh & 0.05 & 0.078 & 0.203 & 24.5 \\
WMT24 ja-zh & 0.10 & 0.077 & 0.177 & 23.6 \\
WMT24 ja-zh & 0.20 & 0.071 & 0.161 & 19.8 \\
WMT24 ja-zh & 0.30 & 0.061 & 0.141 & 14.0 \\
\bottomrule
\end{tabular}
\caption{Primary stopping uses $V^+-V^-\le\lceil.02|S_0|\rceil$. Mean pooling: reference-priority returns at $\alpha=.10$, all depths. Avg.\ averages the splitwise replacement-average churn $R(C)$ over 40 splits. Mean max averages the maximum replacement-specific churn within each split; it is not the maximum across all splits. Over is the mean percentage of replacements exceeding .10. Each split separately satisfies $R(C)\le\alpha$; individual replacements need not.}
\label{tab:review-replacement-mean}
\end{table*}

Only replacement-average churn is constrained. For example, consider mean pooling, $q=.05$ and reference priority at primary stopping. For thirteen-stream WMT23 in this configuration, the mean replacement-average churn across splits is $.072$. The mean of the within-split maximum replacement churn is $.327$, and on average 27.9\% of replacements exceed $\alpha=.10$ (Table~\ref{tab:review-replacement-mean}). For each returned mask, we descriptively report the largest replacement-specific churn and the fraction exceeding $\alpha$, averaging each statistic across splits. The mean of split-specific maxima is not the maximum over all splits; neither diagnostic is a certificate constraint.
Table~\ref{tab:review-replacement-median} gives the corresponding odd-median results.
\paragraph{Risk across individual replacements.}
For each of the 15,360 saved masks from six panels, two pooling rules, four depths, 40 splits, four acquisition orders and two stopping targets, we reconstruct all $MT$ member-replacement outcomes. Each individual risk divides committed membership flips by the same $2k$; their average matches $R(C)$. All masks satisfy the average budget $.10$, while 6,985 have at least one replacement above $.10$. Among 960 reference-priority primary masks per pooling rule, this occurs in 308 median and 541 mean cases. Across the 960 masks for each pooling rule, the mean per-mask maximum risk is .0945 for median pooling and .1404 for mean pooling. The corresponding mean replacement-average risks are .0427 and .0496. The largest individual risk across all masks is .50. These observations illustrate the distinction between average and per-replacement constraints, not violations of the stated certificate. The accompanying per-panel, depth, policy and stopping summaries include the 95th percentile and fraction of replacements above budget. Splits overlap, so these are descriptive frequencies. Scripts and all summaries are in \nolinkurl{review-revision/granularity-risk-implementation/}.


\section{A Same-Loss Statistical Comparison}
\label{app:review-score}
SCoRE provides marginal deployment risk (MDR) and selective deployment risk (SDR) control for general bounded losses \citep{bai2026score}. We use its marginal criterion because churn has a fixed batch denominator. A selected-output denominator would control a different quantity. For a prompt/source group $U$ of $g$ outputs, let
\[
 L(U)=\frac{\sum_{i\in U\cap S_0}F_i}{gH}\in[0,1].
\]
The group decision $\psi(U)$ accepts or rejects its baseline-selected outputs jointly. Structurally invariant decisions may additionally be committed at zero churn cost. For $n_{\rm test}$ equal-size groups, $N=g n_{\rm test}$ and the original risk satisfies
\[
 R(C)=\frac{N}{2k}\frac1{n_{\rm test}}\sum_U L(U)\psi(U).
\]
For equal-size groups, the identity converts churn into $B=N/(2k)$ times average group loss. SCoRE-MDR at target $\alpha/B$ therefore bounds expected churn by $\alpha$, assuming conditional whole-group exchangeability after independent reference fitting. This algebraic reduction does not establish the assumption for a fixed benchmark. The method accepts a group\textquotesingle s selected outputs together, unlike output-level retention.

For a group containing selected outputs, its reference-ridge priority is
\[
 \hat v(U)=\frac{\sum_{i\in U\cap S_0}\sum_r\hat f_{ir}}{H|U\cap S_0|}.
\]
When $|U\cap S_0|=0$, we set $\hat v(U)=1$. Only selected outputs enter this prediction; the loss still divides by the full group size $g$. The same candidate-specific ridge model and feature definitions predict calibration and test flip costs. Calibration losses use the fully scored candidate calibration matrix. For calibration groups, set $v_a=\hat v(U_a)$; for the test group, set $v=\hat v(U)$. With $n$ calibration groups, SCoRE's Proposition~4.4 at nominal level $\alpha/B$ and $\gamma=\alpha/B$ accepts a test group with priority $v$ when
\[
 \frac{1+\sum_{a=1}^{n}L(U_a)\one[v_a\le v]}{n+1}\le\frac{\alpha}{B}.
\]
We implement the inequality by exact integer cross-multiplication of group flip sums and the rational budget. Candidate test scores are unavailable until decisions are fixed. Each of the 960 original-pool configurations runs all five budgets. A separate implementation recomputes calibration costs, all 1,518,400 group decisions and test churn; 300 leave-one-out arithmetic checks additionally verify the symmetric construction on finite instances. These finite checks test implementation consistency; they neither prove the published risk theorem nor establish exchangeability of the benchmark groups.

Of 4,800 panel-and-budget records, 70 have realized batch churn above their specified budget. Outputs recur across overlapping splits and budgets, so 70/4,800 is descriptive, not an estimate from independent batches. An expectation guarantee does not require every batch to meet its budget. These exceedances therefore do not contradict the marginal expectation guarantee under the stated exchangeability assumptions. High coverage without candidate test queries likewise does not establish the deterministic completion certificate. Grouped SCoRE requires fully scored base and candidate calibration data; query certification uses reference and base-test information plus the candidate test scores it inspects. These costs are recorded separately.

WMT24 documents contain unequal numbers of outputs. An unweighted document-level MDR objective would average document risks, not reproduce the output-average loss, so we do not apply this equal-size-group reduction there. A weighted adaptation requires a separate derivation.

The complete grouped-SCoRE grids, including every budget-exceedance count, are retained as \texttt{review-score-*} in \nolinkurl{supplement/result_grids.csv}.

SCOPE uses bidirectional judgment probabilities to compute Bidirectional Preference Entropy (BPE), which these scalar score panels do not provide \citep{badshah2026scope}. Transductive full-ranking inference observes the true relative ranks of a calibration subset and assumes tie-free, exchangeable rankings \citep{fermanian2025transductive}. Its rank-coverage guarantees would require a separate derivation for churn and retained-count certification. We do not evaluate such an adaptation.

\paragraph{Setup and amortization.}
Each configuration records $Mn_{\rm ref}$ base-reference entries, $Tn_{\rm ref}$ candidate-reference entries and $MN$ base-test entries. CRC and grouped SCoRE also record fully scored base and candidate calibration entries. The earlier WMT23 correlation check reads all $24\times557\times13=173{,}784$ retained scores, including outputs later assigned to test splits. It is excluded from post-freeze counts, so replay does not establish acquisition savings from the start of panel construction. Reference scoring could be reused across batches only while fitted maps, cutoff protocol and declared family remain applicable. The 40 overlapping, separately refitted splits do not evaluate batches sharing one frozen setup and therefore do not measure deployment amortization.

Removing a candidate changes the certified family. A claim about the original family still needs conservative treatment of its unobserved scores. Endpoint bounds avoid queries while retaining the candidate in the family.

Table~\ref{tab:review-workload} reports the corresponding information-access accounting.
\begin{table*}[t]
\centering\small
\begin{tabular}{llrrrr}
\toprule
Panel & Method & Setup & Cal. & Cand. test & Accounted total \\
\midrule
RewardBench 2 & Reference & 48244 & 0 & 285.0 & 48529.0 \\
RewardBench 2 & Full batch & 48244 & 0 & 16100.0 & 64344.0 \\
RewardBench 2 & CRC & 48244 & 32144 & 16100.0 & 96488.0 \\
RewardBench 2 & SCoRE & 48244 & 32144 & 0.0 & 80388.0 \\
WMT23, 13 streams & Reference & 74737 & 0 & 523.2 & 75260.2 \\
WMT23, 13 streams & Full batch & 74737 & 0 & 41327.0 & 116064.0 \\
WMT23, 13 streams & CRC & 74737 & 57720 & 41327.0 & 173784.0 \\
WMT23, 13 streams & SCoRE & 74737 & 57720 & 0.0 & 132457.0 \\
WMT23, 12 streams & Reference & 68988 & 0 & 43.3 & 69031.3 \\
WMT23, 12 streams & Full batch & 68988 & 0 & 38148.0 & 107136.0 \\
WMT23, 12 streams & CRC & 68988 & 53280 & 38148.0 & 160416.0 \\
WMT23, 12 streams & SCoRE & 68988 & 53280 & 0.0 & 122268.0 \\
\bottomrule
\end{tabular}
\caption{Full batch denotes the full-information benchmark, not an online acquisition policy. At $q=.05$, median pooling, $\alpha=.10$: logical score-entry accounting. Setup includes pool/candidate reference and pool test entries; Cal.\ includes pool/candidate calibration entries; Cand. test excludes base-test scores already counted in Setup. Accounted total sums these three columns after panel construction; it omits preprocessing reads of the full panel. Values average 40 splits and do not measure runtime. Endpoint zero-query certification can omit candidate-reference scoring; no setup-optimality claim is made.}
\label{tab:review-workload}
\end{table*}

\paragraph{Common information accounting.}
Tables~\ref{tab:common-RewardBench2-median}--\ref{tab:common-WMT23-without-synthetic-ref-mean} summarize four methods across three original panels, two pooling rules, four depths and 40 splits: 3,840 method--configuration records at $\alpha=.10$. Candidate reference, calibration and test observations share the denominator $NT$; the records separately list pool-score stages. Setup is charged as recorded, not optimized for deployment. Zero-query retention and reference priority certify the current batch; CRC and SCoRE retain their expectation guarantees and calibration assumptions. SCoRE decides whole groups. A common numerical loss and cost unit do not make these guarantees identical. Panel preprocessing and full-matrix posthoc verification are outside the selection-visible counts.
\begin{table*}[t]
\centering\small
\setlength{\tabcolsep}{4pt}
\begin{tabular}{rlrrrrrrrr}
\toprule
$q$ & Method & $K$ & Cov. & Risk & Over & Ref. & Cal. & Test & Total \\
\midrule
0.05 & Zero & 95.5 & 80.9 & 0.053 & 0 & 0.998 & 0.000 & 0.000 & 0.998 \\
0.05 & Reference & 115.6 & 97.3 & 0.077 & 0 & 0.998 & 0.000 & 0.018 & 1.016 \\
0.05 & CRC & 85.6 & 72.2 & 0.079 & 2 & 0.998 & 0.998 & 1.000 & 2.997 \\
0.05 & SCoRE & 114.1 & 96.4 & 0.078 & 1 & 0.998 & 0.998 & 0.000 & 1.997 \\
0.10 & Zero & 198.2 & 84.9 & 0.048 & 0 & 0.998 & 0.000 & 0.000 & 0.998 \\
0.10 & Reference & 229.1 & 97.8 & 0.066 & 0 & 0.998 & 0.000 & 0.024 & 1.022 \\
0.10 & CRC & 186.1 & 79.5 & 0.079 & 6 & 0.998 & 0.998 & 1.000 & 2.997 \\
0.10 & SCoRE & 234.3 & 100.0 & 0.070 & 0 & 0.998 & 0.998 & 0.000 & 1.997 \\
0.20 & Zero & 435.7 & 94.6 & 0.043 & 0 & 0.998 & 0.000 & 0.000 & 0.998 \\
0.20 & Reference & 451.4 & 97.9 & 0.047 & 0 & 0.998 & 0.000 & 0.012 & 1.010 \\
0.20 & CRC & 436.5 & 94.7 & 0.089 & 9 & 0.998 & 0.998 & 1.000 & 2.997 \\
0.20 & SCoRE & 460.9 & 100.0 & 0.050 & 0 & 0.998 & 0.998 & 0.000 & 1.997 \\
0.30 & Zero & 685.9 & 99.8 & 0.041 & 0 & 0.998 & 0.000 & 0.000 & 0.998 \\
0.30 & Reference & 686.0 & 99.9 & 0.041 & 0 & 0.998 & 0.000 & 0.000 & 0.998 \\
0.30 & CRC & 687.0 & 100.0 & 0.091 & 0 & 0.998 & 0.998 & 1.000 & 2.997 \\
0.30 & SCoRE & 687.1 & 100.0 & 0.041 & 0 & 0.998 & 0.998 & 0.000 & 1.997 \\
\bottomrule
\end{tabular}
\caption{RewardBench 2, median pooling, $\alpha=.10$: common recorded churn and candidate-score accounting. Risk is replacement-average churn, averaged across the 40 splits. $K$ is returned count; Cov.\ is selected coverage in percent; Over counts batches above budget out of 40. Ref., Cal., Test and Total are logical candidate entries divided by the same test $NT$; Total includes the three stages and can exceed one. K, coverage, risk and stage-cost ratios are splitwise means; Over is a count of splits. Reference priority stops at $V^+-V^-\le\lceil.02|S_0|\rceil$. Zero and Reference give deterministic batch feasibility; CRC and SCoRE give expectation guarantees under group exchangeability, so Over does not refute these guarantees. SCoRE decides whole groups. Reference setup is charged as recorded, not minimized. Pool scores and panel-preprocessing reads are separate; full matrices used only for posthoc validation are outside selection-visible information. Full pool-stage counts accompany the records.}
\label{tab:common-RewardBench2-median}
\end{table*}
\begin{table*}[t]
\centering\small
\setlength{\tabcolsep}{4pt}
\begin{tabular}{rlrrrrrrrr}
\toprule
$q$ & Method & $K$ & Cov. & Risk & Over & Ref. & Cal. & Test & Total \\
\midrule
0.05 & Zero & 24.0 & 20.3 & 0.012 & 0 & 0.998 & 0.000 & 0.000 & 0.998 \\
0.05 & Reference & 117.3 & 97.5 & 0.068 & 0 & 0.998 & 0.000 & 0.045 & 1.043 \\
0.05 & CRC & 97.3 & 81.1 & 0.095 & 15 & 0.998 & 0.998 & 1.000 & 2.997 \\
0.05 & SCoRE & 118.7 & 98.8 & 0.072 & 0 & 0.998 & 0.998 & 0.000 & 1.997 \\
0.10 & Zero & 74.9 & 31.7 & 0.005 & 0 & 0.998 & 0.000 & 0.000 & 0.998 \\
0.10 & Reference & 231.8 & 97.8 & 0.059 & 0 & 0.998 & 0.000 & 0.080 & 1.078 \\
0.10 & CRC & 207.3 & 87.4 & 0.093 & 12 & 0.998 & 0.998 & 1.000 & 2.997 \\
0.10 & SCoRE & 237.0 & 100.0 & 0.063 & 0 & 0.998 & 0.998 & 0.000 & 1.997 \\
0.20 & Zero & 238.7 & 51.7 & 0.003 & 0 & 0.998 & 0.000 & 0.000 & 0.998 \\
0.20 & Reference & 452.8 & 97.9 & 0.045 & 0 & 0.998 & 0.000 & 0.119 & 1.117 \\
0.20 & CRC & 454.6 & 98.3 & 0.097 & 15 & 0.998 & 0.998 & 1.000 & 2.997 \\
0.20 & SCoRE & 462.5 & 100.0 & 0.048 & 0 & 0.998 & 0.998 & 0.000 & 1.997 \\
0.30 & Zero & 487.7 & 70.8 & 0.006 & 0 & 0.998 & 0.000 & 0.000 & 0.998 \\
0.30 & Reference & 675.1 & 97.9 & 0.036 & 0 & 0.998 & 0.000 & 0.126 & 1.124 \\
0.30 & CRC & 689.0 & 100.0 & 0.084 & 0 & 0.998 & 0.998 & 1.000 & 2.997 \\
0.30 & SCoRE & 689.3 & 100.0 & 0.039 & 0 & 0.998 & 0.998 & 0.000 & 1.997 \\
\bottomrule
\end{tabular}
\caption{RewardBench 2, mean pooling, $\alpha=.10$: common recorded churn and candidate-score accounting. Risk is replacement-average churn, averaged across the 40 splits. $K$ is returned count; Cov.\ is selected coverage in percent; Over counts batches above budget out of 40. Ref., Cal., Test and Total are logical candidate entries divided by the same test $NT$; Total includes the three stages and can exceed one. K, coverage, risk and stage-cost ratios are splitwise means; Over is a count of splits. Reference priority stops at $V^+-V^-\le\lceil.02|S_0|\rceil$. Zero and Reference give deterministic batch feasibility; CRC and SCoRE give expectation guarantees under group exchangeability, so Over does not refute these guarantees. SCoRE decides whole groups. Reference setup is charged as recorded, not minimized. Pool scores and panel-preprocessing reads are separate; full matrices used only for posthoc validation are outside selection-visible information. Full pool-stage counts accompany the records.}
\label{tab:common-RewardBench2-mean}
\end{table*}
\begin{table*}[t]
\centering\small
\setlength{\tabcolsep}{4pt}
\begin{tabular}{rlrrrrrrrr}
\toprule
$q$ & Method & $K$ & Cov. & Risk & Over & Ref. & Cal. & Test & Total \\
\midrule
0.05 & Zero & 117.9 & 84.9 & 0.056 & 0 & 0.989 & 0.000 & 0.000 & 0.989 \\
0.05 & Reference & 134.8 & 97.7 & 0.078 & 0 & 0.989 & 0.000 & 0.013 & 1.002 \\
0.05 & CRC & 107.9 & 77.0 & 0.067 & 1 & 0.989 & 0.989 & 1.000 & 2.979 \\
0.05 & SCoRE & 114.9 & 82.2 & 0.049 & 0 & 0.989 & 0.989 & 0.000 & 1.979 \\
0.10 & Zero & 241.3 & 100.0 & 0.035 & 0 & 0.989 & 0.000 & 0.000 & 0.989 \\
0.10 & Reference & 241.3 & 100.0 & 0.035 & 0 & 0.989 & 0.000 & 0.000 & 0.989 \\
0.10 & CRC & 239.1 & 99.1 & 0.064 & 0 & 0.989 & 0.989 & 1.000 & 2.979 \\
0.10 & SCoRE & 241.4 & 100.0 & 0.035 & 0 & 0.989 & 0.989 & 0.000 & 1.979 \\
0.20 & Zero & 489.1 & 100.0 & 0.035 & 0 & 0.989 & 0.000 & 0.000 & 0.989 \\
0.20 & Reference & 489.1 & 100.0 & 0.035 & 0 & 0.989 & 0.000 & 0.000 & 0.989 \\
0.20 & CRC & 487.2 & 99.6 & 0.060 & 0 & 0.989 & 0.989 & 1.000 & 2.979 \\
0.20 & SCoRE & 489.1 & 100.0 & 0.035 & 0 & 0.989 & 0.989 & 0.000 & 1.979 \\
0.30 & Zero & 730.9 & 100.0 & 0.033 & 0 & 0.989 & 0.000 & 0.000 & 0.989 \\
0.30 & Reference & 730.9 & 100.0 & 0.033 & 0 & 0.989 & 0.000 & 0.000 & 0.989 \\
0.30 & CRC & 730.3 & 99.9 & 0.066 & 0 & 0.989 & 0.989 & 1.000 & 2.979 \\
0.30 & SCoRE & 730.9 & 100.0 & 0.033 & 0 & 0.989 & 0.989 & 0.000 & 1.979 \\
\bottomrule
\end{tabular}
\caption{WMT23, 13, median pooling, $\alpha=.10$: common recorded churn and candidate-score accounting. Risk is replacement-average churn, averaged across the 40 splits. $K$ is returned count; Cov.\ is selected coverage in percent; Over counts batches above budget out of 40. Ref., Cal., Test and Total are logical candidate entries divided by the same test $NT$; Total includes the three stages and can exceed one. K, coverage, risk and stage-cost ratios are splitwise means; Over is a count of splits. Reference priority stops at $V^+-V^-\le\lceil.02|S_0|\rceil$. Zero and Reference give deterministic batch feasibility; CRC and SCoRE give expectation guarantees under group exchangeability, so Over does not refute these guarantees. SCoRE decides whole groups. Reference setup is charged as recorded, not minimized. Pool scores and panel-preprocessing reads are separate; full matrices used only for posthoc validation are outside selection-visible information. Full pool-stage counts accompany the records.}
\label{tab:common-WMT23-median}
\end{table*}
\begin{table*}[t]
\centering\small
\setlength{\tabcolsep}{4pt}
\begin{tabular}{rlrrrrrrrr}
\toprule
$q$ & Method & $K$ & Cov. & Risk & Over & Ref. & Cal. & Test & Total \\
\midrule
0.05 & Zero & 24.0 & 20.3 & 0.015 & 0 & 0.989 & 0.000 & 0.000 & 0.989 \\
0.05 & Reference & 123.4 & 97.5 & 0.072 & 0 & 0.989 & 0.000 & 0.045 & 1.035 \\
0.05 & CRC & 116.2 & 91.8 & 0.074 & 0 & 0.989 & 0.989 & 1.000 & 2.979 \\
0.05 & SCoRE & 96.5 & 75.7 & 0.051 & 2 & 0.989 & 0.989 & 0.000 & 1.979 \\
0.10 & Zero & 48.0 & 19.9 & 0.012 & 0 & 0.989 & 0.000 & 0.000 & 0.989 \\
0.10 & Reference & 241.9 & 97.8 & 0.056 & 0 & 0.989 & 0.000 & 0.085 & 1.074 \\
0.10 & CRC & 242.6 & 98.1 & 0.080 & 2 & 0.989 & 0.989 & 1.000 & 2.979 \\
0.10 & SCoRE & 246.9 & 99.8 & 0.061 & 0 & 0.989 & 0.989 & 0.000 & 1.979 \\
0.20 & Zero & 377.6 & 77.6 & 0.012 & 0 & 0.989 & 0.000 & 0.000 & 0.989 \\
0.20 & Reference & 476.4 & 97.9 & 0.031 & 0 & 0.989 & 0.000 & 0.059 & 1.048 \\
0.20 & CRC & 484.4 & 99.6 & 0.058 & 0 & 0.989 & 0.989 & 1.000 & 2.979 \\
0.20 & SCoRE & 486.5 & 100.0 & 0.034 & 0 & 0.989 & 0.989 & 0.000 & 1.979 \\
0.30 & Zero & 638.5 & 87.7 & 0.017 & 0 & 0.989 & 0.000 & 0.000 & 0.989 \\
0.30 & Reference & 714.8 & 97.9 & 0.031 & 0 & 0.989 & 0.000 & 0.057 & 1.047 \\
0.30 & CRC & 729.3 & 99.9 & 0.065 & 0 & 0.989 & 0.989 & 1.000 & 2.979 \\
0.30 & SCoRE & 730.0 & 100.0 & 0.034 & 0 & 0.989 & 0.989 & 0.000 & 1.979 \\
\bottomrule
\end{tabular}
\caption{WMT23, 13, mean pooling, $\alpha=.10$: common recorded churn and candidate-score accounting. Risk is replacement-average churn, averaged across the 40 splits. $K$ is returned count; Cov.\ is selected coverage in percent; Over counts batches above budget out of 40. Ref., Cal., Test and Total are logical candidate entries divided by the same test $NT$; Total includes the three stages and can exceed one. K, coverage, risk and stage-cost ratios are splitwise means; Over is a count of splits. Reference priority stops at $V^+-V^-\le\lceil.02|S_0|\rceil$. Zero and Reference give deterministic batch feasibility; CRC and SCoRE give expectation guarantees under group exchangeability, so Over does not refute these guarantees. SCoRE decides whole groups. Reference setup is charged as recorded, not minimized. Pool scores and panel-preprocessing reads are separate; full matrices used only for posthoc validation are outside selection-visible information. Full pool-stage counts accompany the records.}
\label{tab:common-WMT23-mean}
\end{table*}
\begin{table*}[t]
\centering\small
\setlength{\tabcolsep}{4pt}
\begin{tabular}{rlrrrrrrrr}
\toprule
$q$ & Method & $K$ & Cov. & Risk & Over & Ref. & Cal. & Test & Total \\
\midrule
0.05 & Zero & 112.4 & 99.2 & 0.042 & 0 & 0.989 & 0.000 & 0.000 & 0.989 \\
0.05 & Reference & 113.3 & 99.6 & 0.043 & 0 & 0.989 & 0.000 & 0.001 & 0.990 \\
0.05 & CRC & 106.9 & 95.2 & 0.062 & 1 & 0.989 & 0.989 & 1.000 & 2.979 \\
0.05 & SCoRE & 109.6 & 97.2 & 0.037 & 0 & 0.989 & 0.989 & 0.000 & 1.979 \\
0.10 & Zero & 226.6 & 99.4 & 0.047 & 0 & 0.989 & 0.000 & 0.000 & 0.989 \\
0.10 & Reference & 227.4 & 99.6 & 0.048 & 0 & 0.989 & 0.000 & 0.001 & 0.990 \\
0.10 & CRC & 225.3 & 98.8 & 0.075 & 3 & 0.989 & 0.989 & 1.000 & 2.979 \\
0.10 & SCoRE & 228.5 & 100.0 & 0.048 & 0 & 0.989 & 0.989 & 0.000 & 1.979 \\
0.20 & Zero & 456.8 & 99.9 & 0.041 & 0 & 0.989 & 0.000 & 0.000 & 0.989 \\
0.20 & Reference & 457.3 & 99.9 & 0.042 & 0 & 0.989 & 0.000 & 0.000 & 0.990 \\
0.20 & CRC & 456.4 & 99.7 & 0.076 & 1 & 0.989 & 0.989 & 1.000 & 2.979 \\
0.20 & SCoRE & 457.6 & 100.0 & 0.042 & 0 & 0.989 & 0.989 & 0.000 & 1.979 \\
0.30 & Zero & 680.3 & 99.9 & 0.040 & 0 & 0.989 & 0.000 & 0.000 & 0.989 \\
0.30 & Reference & 680.8 & 99.9 & 0.040 & 0 & 0.989 & 0.000 & 0.000 & 0.990 \\
0.30 & CRC & 681.3 & 100.0 & 0.079 & 0 & 0.989 & 0.989 & 1.000 & 2.979 \\
0.30 & SCoRE & 681.6 & 100.0 & 0.040 & 0 & 0.989 & 0.989 & 0.000 & 1.979 \\
\bottomrule
\end{tabular}
\caption{WMT23, 12, median pooling, $\alpha=.10$: common recorded churn and candidate-score accounting. Risk is replacement-average churn, averaged across the 40 splits. $K$ is returned count; Cov.\ is selected coverage in percent; Over counts batches above budget out of 40. Ref., Cal., Test and Total are logical candidate entries divided by the same test $NT$; Total includes the three stages and can exceed one. K, coverage, risk and stage-cost ratios are splitwise means; Over is a count of splits. Reference priority stops at $V^+-V^-\le\lceil.02|S_0|\rceil$. Zero and Reference give deterministic batch feasibility; CRC and SCoRE give expectation guarantees under group exchangeability, so Over does not refute these guarantees. SCoRE decides whole groups. Reference setup is charged as recorded, not minimized. Pool scores and panel-preprocessing reads are separate; full matrices used only for posthoc validation are outside selection-visible information. Full pool-stage counts accompany the records.}
\label{tab:common-WMT23-without-synthetic-ref-median}
\end{table*}
\begin{table*}[t]
\centering\small
\setlength{\tabcolsep}{4pt}
\begin{tabular}{rlrrrrrrrr}
\toprule
$q$ & Method & $K$ & Cov. & Risk & Over & Ref. & Cal. & Test & Total \\
\midrule
0.05 & Zero & 22.0 & 20.6 & 0.015 & 0 & 0.989 & 0.000 & 0.000 & 0.989 \\
0.05 & Reference & 113.7 & 97.0 & 0.069 & 0 & 0.989 & 0.000 & 0.046 & 1.035 \\
0.05 & CRC & 103.8 & 89.0 & 0.077 & 4 & 0.989 & 0.989 & 1.000 & 2.979 \\
0.05 & SCoRE & 89.2 & 77.6 & 0.048 & 1 & 0.989 & 0.989 & 0.000 & 1.979 \\
0.10 & Zero & 55.6 & 24.3 & 0.010 & 0 & 0.989 & 0.000 & 0.000 & 0.989 \\
0.10 & Reference & 222.2 & 97.8 & 0.049 & 0 & 0.989 & 0.000 & 0.080 & 1.069 \\
0.10 & CRC & 222.6 & 98.2 & 0.076 & 2 & 0.989 & 0.989 & 1.000 & 2.979 \\
0.10 & SCoRE & 226.3 & 99.5 & 0.054 & 0 & 0.989 & 0.989 & 0.000 & 1.979 \\
0.20 & Zero & 342.0 & 75.8 & 0.011 & 0 & 0.989 & 0.000 & 0.000 & 0.989 \\
0.20 & Reference & 442.1 & 97.9 & 0.037 & 0 & 0.989 & 0.000 & 0.063 & 1.053 \\
0.20 & CRC & 450.3 & 99.7 & 0.075 & 0 & 0.989 & 0.989 & 1.000 & 2.979 \\
0.20 & SCoRE & 451.6 & 100.0 & 0.040 & 0 & 0.989 & 0.989 & 0.000 & 1.979 \\
0.30 & Zero & 563.3 & 82.9 & 0.015 & 0 & 0.989 & 0.000 & 0.000 & 0.989 \\
0.30 & Reference & 666.9 & 97.9 & 0.036 & 0 & 0.989 & 0.000 & 0.078 & 1.067 \\
0.30 & CRC & 680.3 & 99.9 & 0.076 & 0 & 0.989 & 0.989 & 1.000 & 2.979 \\
0.30 & SCoRE & 681.0 & 100.0 & 0.039 & 0 & 0.989 & 0.989 & 0.000 & 1.979 \\
\bottomrule
\end{tabular}
\caption{WMT23, 12, mean pooling, $\alpha=.10$: common recorded churn and candidate-score accounting. Risk is replacement-average churn, averaged across the 40 splits. $K$ is returned count; Cov.\ is selected coverage in percent; Over counts batches above budget out of 40. Ref., Cal., Test and Total are logical candidate entries divided by the same test $NT$; Total includes the three stages and can exceed one. K, coverage, risk and stage-cost ratios are splitwise means; Over is a count of splits. Reference priority stops at $V^+-V^-\le\lceil.02|S_0|\rceil$. Zero and Reference give deterministic batch feasibility; CRC and SCoRE give expectation guarantees under group exchangeability, so Over does not refute these guarantees. SCoRE decides whole groups. Reference setup is charged as recorded, not minimized. Pool scores and panel-preprocessing reads are separate; full matrices used only for posthoc validation are outside selection-visible information. Full pool-stage counts accompany the records.}
\label{tab:common-WMT23-without-synthetic-ref-mean}
\end{table*}


\section{Reference Mapping and Query Granularity}
\label{app:mapping-column}
\subsection{Reference-only mapping sensitivity}
We retain the original 40 group or document splits, all test outputs, both pooling rules, four selection depths and all four acquisition methods. Within each reference split, nested prefixes of the original shuffled group order retain all, half or one quarter of the reference groups. Fractional group counts are rounded down, with at least one group retained. The calibration and test assignments remain unchanged. For each evaluator, the right empirical CDF is the original protocol. For a raw score $s$, the midrank variant is $(n_{<s}+n_{\le s})/(2n_{\rm ref})$, where both counts and $n_{\rm ref}$ use the current reference subset. The reference subset supplies both mapping and cutoff fitting, and each cutoff remains fixed during certification. Ridge feature definitions, regularization, priority rules, integer risk arithmetic and stopping criteria are unchanged. Fitted predictions and query order can change with the reference subset and mapping. WMT24 entry acquisition retains the integer-tick comparisons and ties specified in Appendix~\ref{app:review-quality}.

The six mapping protocols run at $\alpha=.10$ on RewardBench, both WMT23 variants and the three WMT24 directions, yielding 11,520 configurations and 46,080 entry-acquisition paths. These are reference-protocol ablations on existing data, not new independent samples. Table~\ref{tab:mapping-summary} summarizes nominal $q=.05$; the supplement's \texttt{mapping-all-*} grids cover every depth. Each range spans the six protocol-specific means over 40 splits, not a confidence interval. Machine-readable summaries also retain each mapping and reference size separately, including exact stopping, setup costs and the agreement/opposition witness decomposition.

Mapping affects baseline selection and, under median pooling, its one-vote boundary. In thirteen-stream WMT23 with median pooling at nominal $q=.05$, changing full-reference right ECDF to midrank reduces the mean fraction of baseline-selected outputs exactly at the cutoff from 28.8\% to 7.3\%; mean splitwise base-list Jaccard similarity is 0.772. This diagnostic was computed after execution without selecting a preferred mapping. The same change reduces the mean minimum certificate size for exact stopping from 415.9 to 5.1 entries. Conversely, WMT24 ja-zh under right ECDF needs 11.3\% of active entries with full reference groups and 34.0\% with a quarter. Its zero-query selected coverage falls from 98.1\% to 92.9\%, while $|S_0|/k$ rises from 1.27 to 1.66. These are different selection protocols; smaller inspection counts do not by themselves identify a preferable returned set. The original right ECDF remains the primary protocol; midrank is a sensitivity variant. A deployment must fix its mapping and reference-size rule before interpreting candidate-test outcomes. We pooled all 5,760 median protocol configurations: six reference protocols, six panels, four depths and 40 splits. The zero-query gap $(|S_0|-V^-_0)/|S_0|$ had Spearman correlation .997 with the normalized exact minimum $Q_{\rm exact}/|S_0|$. Within each panel, the correlation was at least .999. The active fraction $m/|S_0|$ gave .569 and the cutoff-tie fraction .219. The gap uses only reference fits and visible base scores. It is therefore available before candidate scoring as a diagnostic of the required evidence, without determining the evidence in advance. The two quantities share the active safe rows, so part of this agreement is structural. These analyses were conducted on previously inspected panels and were not prospectively preregistered.

\begin{table*}[t]
\centering\small
\begin{tabular}{llrrrrr}
\toprule
Panel & Pooling & $|S_0|/k$ & \shortstack{Zero coverage\\(\%)} & Active (\%) & $Q_{\rm exact}$ & Jaccard \\
\midrule
RewardBench 2 & median & 0.99--1.03 & 80.9--82.6 & 53.5--61.4 & 178.9--196.6 & 0.84--1.00 \\
RewardBench 2 & mean & 1.01--1.05 & 20.3--21.5 & 85.3--86.6 & -- & 0.90--1.00 \\
WMT23, 13 streams & median & 1.04--1.15 & 84.9--99.8 & 0.3--45.1 & 5.1--415.9 & 0.72--1.00 \\
WMT23, 13 streams & mean & 1.02--1.12 & 20.3--21.3 & 84.6--87.5 & -- & 0.78--1.00 \\
WMT23, 12 streams & median & 1.01--1.20 & 95.5--99.2 & 3.8--24.1 & 41.1--220.6 & 0.66--1.00 \\
WMT23, 12 streams & mean & 1.04--1.14 & 20.4--22.1 & 83.3--85.4 & -- & 0.73--1.00 \\
WMT24 en-de & median & 1.09--1.37 & 96.7--99.5 & 1.3--13.1 & 37.1--408.2 & 0.65--1.00 \\
WMT24 en-de & mean & 1.08--1.35 & 22.0--23.6 & 80.8--83.9 & -- & 0.64--1.00 \\
WMT24 en-es & median & 1.11--1.35 & 96.9--99.8 & 1.1--14.4 & 18.1--343.9 & 0.71--1.00 \\
WMT24 en-es & mean & 1.11--1.29 & 20.9--22.0 & 81.6--83.4 & -- & 0.71--1.00 \\
WMT24 ja-zh & median & 1.25--1.66 & 92.9--98.1 & 10.8--34.0 & 194.1--1143.5 & 0.61--1.00 \\
WMT24 ja-zh & mean & 1.20--1.48 & 21.4--25.4 & 82.2--85.9 & -- & 0.43--1.00 \\
\bottomrule
\end{tabular}
\caption{Mapping sensitivity at $q=.05$, $\alpha=.10$. Ranges span six protocol-specific 40-split means (full/half/quarter reference, right/midrank ECDF), keeping test groups fixed. Zero is zero-query selected coverage in percent. Active is the mean percentage of active entries queried by reference priority, $100Q/(mT)$, at the stop $V^+-V^-\le\lceil.02|S_0|\rceil$; it is zero when $m=0$. $Q_{\rm exact}$ is the 40-split mean entry witness under median pooling, an exact target. Jaccard compares base selected sets with full-reference right ECDF. Active and witness columns use different stopping targets and do not form an online-to-lower-bound ratio. Ranges are not confidence intervals.}
\label{tab:mapping-summary}
\end{table*}



\subsection{Reference-cutoff bootstrap}
\label{app:cutoff-bootstrap}
A paired group bootstrap compares cutoff-only sensitivity with joint map-and-cutoff sensitivity. For each of six panels and splits 0--7, we draw 50 reference samples with replacement at the original prompt/source/document unit. At $q=.05$, both pooling rules use two arms: retaining original right ECDF maps and resampling the reference-consensus cutoff, or refitting maps and cutoff on the same sampled groups. Test groups and all candidates stay fixed. The generator is NumPy \texttt{default\_rng}, with seed $620275+1000h+r$, where $h$ is the zero-based panel index and $r$ the split index. The panel order is RewardBench~2, thirteen-stream WMT23, twelve-stream WMT23, WMT24 en-de, en-es and ja-zh. Reference groups retain their first-occurrence order; each draw samples the original number of groups with replacement, using the same sampled groups for both arms. This yields 9,600 protocol instances. We recompute initial bounds, the full-information count and, for median pooling, matching exact and primary minima; no labels, reference-priority refit or online acquisition enter this analysis.

For all 9,600 records, saved accounting checks verify unique configuration keys, positive sample counts, integer budgets, cutoff and Jaccard ranges, ordered capacity bounds, gap identities, active-plus-invariant selected counts and the zero-query stopping condition. Mean-pooling records carry no median-only witness value. A separate bisect-based ECDF and explicit replacement replay reconstructs 24 fixed full-test cases (split zero, draw zero, every panel, arm and pooling rule), not every bootstrap record. Table~\ref{tab:cutoff-bootstrap} summarizes selection Jaccard, zero-query primary stops and median exact minima. On thirteen-stream WMT23 median, cutoff-only and map-plus-cutoff Jaccard medians are .997 and .959; zero-query primary-stop fractions are 21.0\% and 25.5\%. We also compute the sample cutoff standard deviation within each 50-draw split and average these deviations over eight splits. On thirteen-stream WMT23 median, the cutoff-only and map-plus-cutoff averages are .0115 and .0078 on the normalized score scale; all 24 cells accompany the records. A draw changes the fixed selection protocol, so a smaller certificate concerns a changed task. The saved percentiles describe these draws, not population confidence intervals or independent replications.
\begin{table*}[t]
\centering\small
\begin{tabular}{lrrrrrrr}
\toprule
Panel & Pool & $J_t$ & $J_{\rm map+t}$ & \shortstack{Zero stops\\$t$ (\%)} & \shortstack{Zero stops\\map+$t$ (\%)} & $Q_{0,t}$ & $Q_{0,\rm map+t}$ \\
\midrule
RewardBench 2 & median & 0.951 & 0.909 & 0.0 & 0.0 & 153.0 & 153.0 \\
RewardBench 2 & mean & 0.933 & 0.933 & 0.0 & 0.0 & -- & -- \\
WMT23, 13 streams & median & 0.997 & 0.959 & 21.0 & 25.5 & 312.0 & 244.0 \\
WMT23, 13 streams & mean & 0.892 & 0.890 & 0.0 & 0.0 & -- & -- \\
WMT23, 12 streams & median & 0.847 & 0.842 & 82.3 & 80.8 & 0.0 & 0.0 \\
WMT23, 12 streams & mean & 0.862 & 0.855 & 0.0 & 0.0 & -- & -- \\
WMT24 en-de & median & 0.855 & 0.856 & 92.0 & 90.0 & 0.0 & 0.0 \\
WMT24 en-de & mean & 0.853 & 0.846 & 1.0 & 1.5 & -- & -- \\
WMT24 en-es & median & 0.851 & 0.854 & 98.0 & 95.3 & 0.0 & 0.0 \\
WMT24 en-es & mean & 0.877 & 0.877 & 0.5 & 0.0 & -- & -- \\
WMT24 ja-zh & median & 0.841 & 0.831 & 77.5 & 74.5 & 0.0 & 0.0 \\
WMT24 ja-zh & mean & 0.829 & 0.836 & 0.0 & 0.0 & -- & -- \\
\bottomrule
\end{tabular}
\caption{Reference-group bootstrap at $q=.05$, $\alpha=.10$: 50 paired draws in each of eight splits, 400 instances per cell. $t$ resamples the reference-consensus cutoff while retaining original maps; map+$t$ refits ECDFs and cutoff on the same resampled prompt/source/document groups. Test groups and the evaluator family stay fixed. $J$ is median baseline-selection Jaccard against the original split. Zero reports the fraction meeting the primary count-gap stop without queries; $Q_0$ is the median matching exact entry minimum, reported only for median pooling. Every draw defines a changed fixed protocol. Percentiles in the saved records are descriptive bootstrap summaries, not population confidence intervals. No label-based tuning or new online acquisition is performed.}
\label{tab:cutoff-bootstrap}
\end{table*}


\subsection{Whole-evaluator query controls}
A column oracle reveals all $N$ scores from one candidate evaluator for the current test batch. It is charged $N$ logical score entries, including invariant and nonselected outputs. The same certificate solver updates the active selected rows with the newly observed flips. Fixed column order follows candidate index, random column order uses the original query seed plus 571921, and reference column order sorts candidates by their summed predicted upper-cost reduction across active selected rows. Let $e_{ir}=\operatorname{round}_{\rm even}(10^{12}\hat f_{ir})$ be prediction ticks. Columns sort by decreasing $\sum_{i\in S_0:p_i>0}(10^{12}p_i-e_{ir})$, with candidate-index ties. This static column priority uses reference and visible base information only, not hidden candidate test scores. It aggregates the entry method\textquotesingle s per-cell predictions at the column level.

The original full-reference right-ECDF configuration runs all three column orders at every depth and pooling rule on 40 splits, yielding 5,760 paths. The positive-gap and exact stops are recorded separately. Table~\ref{tab:column-summary} uses the positive-gap stop, and the supplement's \texttt{column-all-*} grids give every depth. Entry and column protocols share the risk budget and stopping criterion, although their positive-gap returned sets can differ. All-column percentages divide by $NT$; dividing them by $mT$ would mix scores outside the active selected rows with active-entry inspection. The entry witness $Q_{\rm exact}$ is not an unscaled bound on column count. Under the same odd-median completion model and exact stopping criterion, any successful certificate using $c$ columns satisfies $Q_{\rm exact}\le mc\le Nc$; the entry witness therefore lower-bounds its score-entry cost. This exact bound does not apply directly to the positive-gap results. Appendix~\ref{app:column-witness} gives the converted bound $\lceil Q_{\rm exact}/m\rceil$ for $m>0$ and exact column witnesses; $m=0$ needs no queries.

At $q=.05$, $\alpha=.10$ and primary stopping, RewardBench reference-priority entry paths access all seven candidate evaluators on average under both pooling rules. Reference column paths access 5.0 under median and 6.7 under mean. The queried-score fractions are reported in Table~\ref{tab:column-summary}. These results separate revealed entries from the number of candidate test columns involved. Reference and base setup costs remain recorded for each protocol. Cached-score replay establishes neither model-loading savings nor end-to-end inference savings.

\begin{table*}[t]
\centering\small
\begin{tabular}{llrrrrrrr}
\toprule
Panel & Pooling & Entry (\%) & Touched & Fixed (\%) & Random (\%) & Ref. (\%) & Cols. & $T$ \\
\midrule
RewardBench 2 & median & 1.77 & 7.0 & 69.3 & 71.4 & 71.8 & 5.0 & 7 \\
RewardBench 2 & mean & 4.52 & 7.0 & 99.6 & 99.6 & 96.1 & 6.7 & 7 \\
WMT23, 13 streams & median & 1.27 & 15.1 & 67.5 & 62.9 & 40.1 & 6.8 & 17 \\
WMT23, 13 streams & mean & 4.53 & 17.0 & 95.4 & 94.7 & 90.1 & 15.3 & 17 \\
WMT23, 12 streams & median & 0.11 & 1.3 & 4.1 & 4.3 & 4.1 & 0.7 & 17 \\
WMT23, 12 streams & mean & 4.61 & 17.0 & 92.5 & 92.6 & 89.9 & 15.3 & 17 \\
WMT24 en-de & median & 0.03 & 0.7 & 2.0 & 2.3 & 1.3 & 0.2 & 14 \\
WMT24 en-de & mean & 4.69 & 14.0 & 92.1 & 92.1 & 89.6 & 12.6 & 14 \\
WMT24 en-es & median & 0.03 & 0.4 & 1.3 & 1.4 & 1.4 & 0.2 & 14 \\
WMT24 en-es & mean & 4.76 & 14.0 & 90.5 & 90.5 & 88.6 & 12.4 & 14 \\
WMT24 ja-zh & median & 0.36 & 3.3 & 12.7 & 12.3 & 12.3 & 1.9 & 15 \\
WMT24 ja-zh & mean & 5.46 & 15.0 & 93.0 & 93.5 & 92.8 & 13.9 & 15 \\
\bottomrule
\end{tabular}
\caption{Whole-evaluator queries at $q=.05$, $\alpha=.10$, full-reference right ECDF, stopping when $V^+-V^-\le\lceil.02|S_0|\rceil$. Entry is reference-priority entry inspection; Touched counts candidate columns accessed by entry queries; it is not a scoring-call count. Fixed, Random and Ref.\ are column orders, charged $N$ scores per evaluator. Percentages divide by $NT$; Cols.\ counts reference-order evaluators. $T$ is the candidate count; the replacement family has $H=MT$ members. Counts and ratios average 40 splits separately. No inference time is measured.}
\label{tab:column-summary}
\end{table*}



\subsection{Scope of the completion and query model}
\label{app:model-scope}
Observed score correlations may guide query order without changing the endpoint certificate. Using relations among missing scores to tighten bounds requires guarantees of these relations for the true scores. Learned low-rank approximations or concentrated posteriors alone provide no such guarantee. Point predictions do not remove any consistent completion. Both joint endpoint vectors remain attainable, so predictions cannot tighten the frontiers in this model. A margin ruling out every single-member flip is already captured by $p_i=0$; base-pool variance alone does not restrict missing candidate scores. Fixed replacement weights are treated below; the odd-median unit-count formulas assume uniform weights and Boolean entry outcomes $y_{ir}$. Mean pooling can have several flip-count levels and different $p_i$ across rows; Theorem~\ref{thm:joint} handles their shared evidence without the median binary decomposition. The same structure holds when replacing member $j$ substitutes the candidate score into a consensus nondecreasing in each score. A fixed-weight mean in which the candidate takes the replaced member's weight is one case. Each flip indicator is then monotone in the candidate score, and scores zero and one attain all upper and lower row costs at once. Proposition~\ref{prop:stop}, Theorem~\ref{thm:joint} and the comparison chain therefore apply with the induced caps $p_i$; the odd-median closed forms need equal votes and do not. Renormalized candidate weights fall outside this argument, since they can make the zero-flip endpoint unattainable. Simultaneous replacements change median boundary geometry; continuous mixtures also require a specified risk measure. These are different certification models.

\paragraph{Fixed replacement weights.}
The weights apply to replacement scenarios, not to scores within the pooling rule. For specified $w_{jr}\ge0$ with $\sum_{j,r}w_{jr}=1$, define the candidate contribution
\[
 g_{ir}(x)=\sum_j w_{jr}\one[\sigma_P(i)\ne\sigma_{P^{j\to r}(x)}(i)].
\]
The weighted risk is $R_w(C)=\sum_{i\in C}\sum_r g_{ir}(s_r(i))/(2k)$. For selected rows, let $a_{ir}=g_{ir}(0)$. The observed sum is a valid lower bound, and adding $\sum_{r\in U_i}a_{ir}$ gives its upper bound; unselected rows use the opposite endpoint. Nonnegative weights preserve joint endpoint attainment and the sorted-prefix retained-count argument. The budget here is the real-valued $2k\alpha$, without an integer floor unless contributions are first rescaled to a common integer unit. Nonuniform weights can give candidate-specific endpoint caps; Theorem~\ref{thm:joint} and the median closed forms are stated for the original uniform family. For irrational weights, exclusion compares the sum of the smallest lower costs with the real budget by strict inequality. For rational weights, rescaling to integer contributions gives an integer budget. Lemma~\ref{lem:order-threshold} then holds with thresholds up to the largest lower cost, so the exclusion condition and the entry-level integer program of Appendix~\ref{app:compact-joint} carry over. Entry $(i,r)$ raises the lower cost by $g_{ir}(s_r(i))$ and lowers the upper cost by $a_{ir}-g_{ir}(s_r(i))$. This program defines the weighted $Q_{\rm joint}$; Theorem~\ref{thm:joint} gives its row-state form only for uniform weights. The chain $Q_{\rm joint}\le Q_{\cup,\rm sep}\le Q_{\rm split}\le2Q_{\rm joint}$ also holds. Observations only tighten both bounds, since $0\le g_{ir}\le a_{ir}$. A union of a safety witness and an exclusion witness for one count is therefore a joint certificate. This gives the first inequality, and deduplication gives the second. A joint witness proves both tasks for its count, which gives the third. Uniform weights belong to the weighted class, so the factor-two families of Appendix~\ref{app:deduplicated} keep the factor two tight. We tested seeded nonuniform weights on the 960 original mean instances. Each instance draws integer candidate weights $w_r\in\{1,2,3\}$ and sets $w_{jr}=w_r/(MW)$ with $W=\sum_rw_r$. Scaling by $MW$ gives integer contributions $w_rf_{ir}$ and budget $\lfloor2k\alpha MW\rfloor$. Full retention stayed feasible in 894 instances, so exclusion was automatic there. For the other 66, each joint, safety or exclusion program had a 60-second limit, and each union program 120 seconds. At exact stopping, 25 comparisons closed, and all 25 showed a strict joint gain. The mean saving was 22.2 entries, at most 1.30\% of the union. At primary stopping, 34 comparisons closed and six showed a gain. Over these six gains the mean saving was 61.2 entries, with at most 2.44\%. The remaining 41 exact and 32 primary comparisons stayed open. The closed weighted gains reached at most 2.44\% of the union, against 2.61\% under uniform weights. Records are in \nolinkurl{outputs/scope-extensions/weighted_records.jsonl}.

Restricting only unobserved completions to a discrete set retaining both joint endpoint assignments leaves these coordinate bounds and capacities unchanged, with observed scores, base scores and cutoff fixed. An arbitrary grid restriction gives no evidence of the true missing scores lying in the set. Reference-sample minima and maxima alone do not establish hard support for unseen test scores.

A known hard constraint $\mathcal K$ satisfied by the true matrix restricts completions to $\Omega_{\mathcal K}\subseteq\Omega$. Every original safe certificate remains valid, but endpoint bounds can cease to be tight and the minimum sufficient evidence can change. An observed correlation alone does not establish $\mathcal K$. Original minimum-query lower bounds need not remain lower bounds for this restricted model. If a learned set contains the true matrix only with a stated probability, certification within this set inherits the containment event; marginal coverage for individual entries does not establish simultaneous validity under adaptive acquisition. Such a change requires its own information and guarantee specification.

\paragraph{Two tractable hard constraints.}
\label{app:constrained}
Two constraint families keep the endpoint structure. Both act on candidate test scores of the active selected rows and define their geometry from visible base scores. Their validity also requires constants $D_r$ or $L_r$ holding for the true candidate scores, and the Lipschitz envelopes use queried candidate scores. The first caps each candidate's deviation from base consensus: $|s_r(i)-c_P(i)|\le D_r$. The second is evaluator-wise Lipschitz continuity in base-score space: $|s_r(i)-s_r(i')|\le L_rd(i,i')$, with $d(i,i')=\max_j|s_j(i)-s_j(i')|$.
\begin{proposition}[Attainable envelopes under deviation and Lipschitz constraints]
\label{prop:constrained}
Suppose the true matrix satisfies one of these constraints with known $D_r$ or $L_r$. Let $O_r$ hold the queried active selected rows of candidate $r$; $L_r=\infty$ imposes no constraint. Under the deviation constraint, set $\lambda_r(i)=\max(0,c_P(i)-D_r)$ and $\mu_r(i)=\min(1,c_P(i)+D_r)$. Under the Lipschitz constraint, set
\[
 \lambda_r(i)=\max\Big(0,\max_{i'\in O_r}\big(s_r(i')-L_rd(i,i')\big)\Big),
\]
and define $\mu_r(i)$ by the corresponding minimum capped at one; an empty maximum is zero and an empty minimum is one. Every consistent completion satisfies $\lambda_r\le s_r\le\mu_r$ on unqueried entries. Setting all unqueried entries to $\lambda$, or all to $\mu$, gives consistent completions. Hence row bounds $u_i=\sum_{r:i\in O_r}f_{ir}+\sum_{r:i\notin O_r}f_i(\lambda_r(i))$ and $\ell_i=\sum_{r:i\in O_r}f_{ir}+\sum_{r:i\notin O_r}f_i(\mu_r(i))$ are both attained, and Proposition~\ref{prop:stop} holds with these bounds.
\end{proposition}
\begin{proof}
The deviation sets form a box, so its corners are consistent. For the Lipschitz case, each map $i\mapsto s_r(i')-L_rd(i,i')$ is $L_r$-Lipschitz under the pseudometric $d$, and so is a maximum of such maps and the constant zero. For $i\in O_r$, the constraint gives $s_r(i')-L_rd(i,i')\le s_r(i)$, with equality at $i'=i$, so $\lambda_r$ matches every observation. The upper envelope is symmetric. Columns are constrained separately, so the envelopes combine across candidates. Selected-row flip counts are nonincreasing in the candidate score, so $\lambda$ attains every upper row cost and $\mu$ every lower row cost. The argument of Proposition~\ref{prop:stop} then applies unchanged.
\end{proof}
A Lipschitz query narrows other rows in its column, so one entry can tighten several rows. The closed-form minima of Section~\ref{sec:median} assume no such spillover and do not apply to this model. Deviation caps add no spillover. For odd medians, an entry whose range lies on one side of the cutoff has a fixed outcome. Every other entry keeps both outcomes. The residual witnesses of Appendix~\ref{app:residual} then apply with fixed entries as history. With the oracle caps of the experiment below, 0.015\% of active entries in 960 median configurations (six panels, four depths, 40 splits) at $\alpha=.10$ had fixed outcomes. The exact minimum fell in 27 configurations, 26 of them on thirteen-stream WMT23. The summed minimum fell by 0.48\% in total and by at most 8.2\% in any one. At primary stopping the total fell by 0.71\%. Records are in \nolinkurl{outputs/scope-extensions/median_deviation_minima.jsonl}.

We evaluate these constraints with the smallest batch-valid constants $D_r$ and $L_r$, computed from the hidden active-row scores. This oracle analysis measures the tightening available within each constraint family on the current batches. The reference-priority query order is held fixed, so only the certificate bounds and stopping points change. Table~\ref{tab:constrained} reports nominal $q=.05$ and $.10$ on all six panels, both pooling rules and 40 splits. Narrower ranges can only tighten the bounds, so neither constraint increased a query count. Deviation caps raised the zero-query safe count in 500 of 960 configurations: 480 with mean pooling and 20 with median pooling. At $q=.05$ on RewardBench, mean safe retention rose from 24.0 to 54.0 outputs. Primary-stop queries fell in 387 configurations, by 1.38\% in total. Lipschitz bounds never raised a zero-query count, because each bound needs an observed entry in its column. They reduced primary-stop queries in 453 configurations, by 1.51\% in total. The largest effect was on thirteen-stream WMT23 with median pooling at $q=.05$, where mean queries fell from 523.2 to 252.1. Oracle Lipschitz constants were finite for 76.5\% of candidates. The other candidates gave different scores to outputs with identical base scores. Reference-estimated constants held for every candidate of the test batch in 231 configurations for deviation caps and in 439 for Lipschitz bounds. Reference data alone therefore did not certify these constraints on most batches.
\begin{table*}[t]
\centering\small
\begin{tabular}{@{}llrrrrrrrr@{}}
\toprule
& & \multicolumn{3}{c}{Zero-query safe count} & \multicolumn{3}{c}{Primary-stop queries} & \multicolumn{2}{c}{Ref.\ constants valid (\%)}\\
\cmidrule(lr){3-5}\cmidrule(lr){6-8}\cmidrule(l){9-10}
Panel & Pool & None & Dev. & Lip. & None & Dev. & Lip. & Dev. & Lip.\\
\midrule
\multicolumn{10}{@{}l}{$q=.05$}\\
RewardBench~2 & median & 95.5 & 95.5 & 95.5 & 285.0 & 285.0 & 284.9 & 96.4 & 98.9\\
RewardBench~2 & mean & 24.0 & 54.0 & 24.0 & 727.9 & 723.2 & 723.9 & 93.2 & 88.6\\
WMT23 (13) & median & 117.9 & 118.4 & 117.9 & 523.2 & 518.3 & 252.1 & 87.5 & 92.4\\
WMT23 (13) & mean & 24.0 & 47.7 & 24.0 & 1873.7 & 1847.7 & 1851.7 & 90.9 & 88.4\\
WMT23 (12) & median & 112.4 & 112.4 & 112.4 & 43.3 & 43.3 & 42.9 & 91.3 & 96.8\\
WMT23 (12) & mean & 22.0 & 49.0 & 22.0 & 1757.7 & 1706.3 & 1733.4 & 88.2 & 77.4\\
WMT24 en-de & median & 461.9 & 462.0 & 461.9 & 44.1 & 44.1 & 26.1 & 92.1 & 98.4\\
WMT24 en-de & mean & 84.4 & 139.0 & 84.4 & 5619.0 & 5529.0 & 5526.4 & 73.6 & 86.4\\
WMT24 en-es & median & 410.5 & 410.5 & 410.5 & 33.2 & 33.2 & 33.2 & 93.8 & 97.9\\
WMT24 en-es & mean & 74.5 & 146.3 & 74.5 & 4962.1 & 4853.0 & 4858.2 & 78.4 & 82.1\\
WMT24 ja-zh & median & 333.7 & 333.7 & 333.7 & 276.4 & 276.4 & 275.9 & 91.8 & 95.7\\
WMT24 ja-zh & mean & 54.0 & 63.1 & 54.0 & 4402.5 & 4399.7 & 4395.3 & 77.3 & 84.8\\
\multicolumn{10}{@{}l}{$q=.10$}\\
RewardBench~2 & median & 198.2 & 198.2 & 198.2 & 378.4 & 378.4 & 378.4 & 89.3 & 96.1\\
RewardBench~2 & mean & 74.9 & 128.4 & 74.9 & 1289.6 & 1278.7 & 1283.0 & 90.0 & 83.9\\
WMT23 (13) & median & 241.3 & 241.3 & 241.3 & 0.0 & 0.0 & 0.0 & 77.5 & 97.6\\
WMT23 (13) & mean & 48.0 & 117.6 & 48.0 & 3518.8 & 3501.0 & 3500.3 & 77.9 & 76.2\\
WMT23 (12) & median & 226.6 & 226.6 & 226.6 & 29.8 & 29.8 & 29.7 & 84.9 & 93.1\\
WMT23 (12) & mean & 55.6 & 137.7 & 55.6 & 3047.4 & 3004.4 & 3030.1 & 76.8 & 70.0\\
WMT24 en-de & median & 921.3 & 921.3 & 921.3 & 25.7 & 25.7 & 25.3 & 80.9 & 98.6\\
WMT24 en-de & mean & 485.6 & 616.3 & 485.6 & 6538.0 & 6378.6 & 6473.0 & 74.3 & 89.6\\
WMT24 en-es & median & 809.4 & 809.4 & 809.4 & 0.0 & 0.0 & 0.0 & 85.9 & 98.6\\
WMT24 en-es & mean & 406.0 & 546.2 & 406.0 & 6038.1 & 5891.2 & 5955.6 & 80.4 & 84.5\\
WMT24 ja-zh & median & 624.4 & 624.4 & 624.4 & 159.1 & 159.1 & 158.9 & 86.2 & 98.5\\
WMT24 ja-zh & mean & 108.5 & 175.8 & 108.5 & 7834.2 & 7817.8 & 7823.6 & 69.5 & 84.2\\
\bottomrule
\end{tabular}
\caption{Certification under oracle-valid deviation caps (Dev.) and Lipschitz bounds in base-score space (Lip.), $\alpha=.10$, means over 40 splits. Queries follow the unchanged reference-priority path. The last two columns give the share of candidates whose reference-estimated constant also holds on the test batch; these candidate-level shares differ from the configuration-level counts in the text.}
\label{tab:constrained}
\end{table*}


\paragraph{Task-preserving transformations.}
For odd-sized median pooling, apply a common strictly increasing bijection $h:[0,1]\to[0,1]$ to every score and replace the cutoff by $h(t_0)$. Median order statistics commute with $h$, and cutoff membership is unchanged for every replacement. Transporting each unobserved score by $h$ gives a bijection of completion matrices, so flip costs and certificate minima are unchanged. Separate evaluator-wise ECDF refits are not such a common transformation. For mean pooling, the corresponding property holds for a common positive affine map $h(x)=ax+b$, $a>0$, with the cutoff and score-completion range transported together. A nonlinear common map need not commute with averaging. These transformations give sufficient conditions for task equivalence. Outside them, equivalence must be checked, not assumed; the ECDF-refitting experiments above change baseline selected sets.

\paragraph{Evidence requirements under changed protocols.}
Table~\ref{tab:protocol-evidence} connects each changed median protocol to its initial count gap, invariant and budget-admissible zero-query retention, and matching exact entry minimum. We reaggregate 5,760 median configurations (six reference protocols, six panels, four depths and 40 splits) without new query runs. A lower minimum under a different reference map is evidence about a different selected set, not a same-instance acquisition improvement.
\begin{table*}[t]
\centering\small
\setlength{\tabcolsep}{4pt}
\begin{tabular}{lllrrrrrr}
\toprule
Panel & Ref. & ECDF & $|S_0|$ & Initial gap & Inv. & Noninv. & $Q_0$ & Jaccard \\
\midrule
RewardBench 2 & 1 & Right & 118.9 & 23.4 & 55.5 & 40.0 & 196.6 & 1.000 \\
RewardBench 2 & 1 & Mid & 118.9 & 23.3 & 55.7 & 40.0 & 189.1 & 0.994 \\
RewardBench 2 & 1/2 & Right & 114.3 & 20.8 & 53.5 & 40.0 & 178.9 & 0.896 \\
RewardBench 2 & 1/2 & Mid & 114.5 & 21.1 & 53.5 & 40.0 & 187.5 & 0.897 \\
RewardBench 2 & 1/4 & Right & 114.0 & 21.0 & 53.0 & 40.0 & 180.2 & 0.836 \\
RewardBench 2 & 1/4 & Mid & 114.2 & 21.0 & 53.3 & 40.0 & 179.9 & 0.836 \\
WMT23, 13 & 1 & Right & 137.9 & 20.0 & 76.7 & 41.2 & 415.9 & 1.000 \\
WMT23, 13 & 1 & Mid & 126.4 & 0.4 & 93.6 & 32.5 & 5.1 & 0.772 \\
WMT23, 13 & 1/2 & Right & 139.3 & 14.8 & 83.7 & 40.9 & 326.7 & 0.906 \\
WMT23, 13 & 1/2 & Mid & 128.9 & 2.1 & 92.0 & 34.8 & 33.0 & 0.739 \\
WMT23, 13 & 1/4 & Right & 140.9 & 12.6 & 88.3 & 40.0 & 239.4 & 0.823 \\
WMT23, 13 & 1/4 & Mid & 133.5 & 3.4 & 95.1 & 35.0 & 71.8 & 0.724 \\
WMT23, 12 & 1 & Right & 114.0 & 1.6 & 85.5 & 26.9 & 41.1 & 1.000 \\
WMT23, 12 & 1 & Mid & 117.6 & 3.8 & 84.9 & 29.0 & 63.2 & 0.762 \\
WMT23, 12 & 1/2 & Right & 113.5 & 2.7 & 85.9 & 24.9 & 60.7 & 0.821 \\
WMT23, 12 & 1/2 & Mid & 116.0 & 3.9 & 84.5 & 27.6 & 64.5 & 0.718 \\
WMT23, 12 & 1/4 & Right & 133.9 & 8.6 & 96.8 & 28.5 & 220.6 & 0.724 \\
WMT23, 12 & 1/4 & Mid & 131.3 & 5.9 & 95.2 & 30.2 & 98.1 & 0.659 \\
WMT24 en-de & 1 & Right & 464.6 & 2.7 & 400.3 & 61.6 & 37.1 & 1.000 \\
WMT24 en-de & 1 & Mid & 466.1 & 2.9 & 401.4 & 61.8 & 40.2 & 0.992 \\
WMT24 en-de & 1/2 & Right & 506.7 & 13.0 & 407.5 & 86.2 & 245.5 & 0.774 \\
WMT24 en-de & 1/2 & Mid & 510.2 & 12.5 & 414.3 & 83.3 & 238.9 & 0.775 \\
WMT24 en-de & 1/4 & Right & 581.5 & 21.0 & 451.9 & 108.7 & 408.2 & 0.666 \\
WMT24 en-de & 1/4 & Mid & 578.0 & 16.2 & 457.2 & 104.7 & 341.6 & 0.651 \\
WMT24 en-es & 1 & Right & 411.8 & 1.4 & 347.7 & 62.7 & 19.0 & 1.000 \\
WMT24 en-es & 1 & Mid & 413.1 & 1.3 & 350.2 & 61.6 & 18.1 & 0.997 \\
WMT24 en-es & 1/2 & Right & 441.2 & 9.0 & 353.0 & 79.2 & 124.2 & 0.785 \\
WMT24 en-es & 1/2 & Mid & 440.9 & 9.1 & 356.0 & 75.9 & 125.3 & 0.778 \\
WMT24 en-es & 1/4 & Right & 493.0 & 18.2 & 385.3 & 89.5 & 334.8 & 0.714 \\
WMT24 en-es & 1/4 & Mid & 499.4 & 18.8 & 388.9 & 91.6 & 343.9 & 0.710 \\
WMT24 ja-zh & 1 & Right & 345.5 & 11.8 & 266.0 & 67.7 & 271.6 & 1.000 \\
WMT24 ja-zh & 1 & Mid & 338.5 & 13.1 & 257.6 & 67.8 & 194.1 & 0.789 \\
WMT24 ja-zh & 1/2 & Right & 396.8 & 25.8 & 302.1 & 68.9 & 625.0 & 0.778 \\
WMT24 ja-zh & 1/2 & Mid & 364.3 & 21.7 & 275.3 & 67.3 & 348.7 & 0.707 \\
WMT24 ja-zh & 1/4 & Right & 448.4 & 52.8 & 325.9 & 69.7 & 1143.5 & 0.661 \\
WMT24 ja-zh & 1/4 & Mid & 398.1 & 36.7 & 291.7 & 69.7 & 606.0 & 0.614 \\
\bottomrule
\end{tabular}
\caption{Median protocol sensitivity at $q=.05$, $\alpha=.10$, means over 40 fixed test splits. Ref.\ is the fraction of reference groups. Initial gap is $|S_0|-V(u^{(0)})$; Inv.\ and Noninv.\ are output counts decomposing zero-query retention. $Q_0$ counts score entries in the matching exact minimum for each changed protocol. Jaccard compares its baseline selected set with full-reference right ECDF on the same split. These joint descriptions connect changed selection tasks to evidence requirements; they are not same-instance acquisition interventions or an endorsement of a mapping.}
\label{tab:protocol-evidence}
\end{table*}

\FloatBarrier

\section{Risk Scaling and Selection Objects}
\label{app:risk-interpretation}
For a batch-fixed positive denominator $d$, define $R_d(C)=\sum_{i\in C}F_i/(dH)$. Then $R_d=(2k/d)R(C)$, and using $\alpha_d=2k\alpha/d$ preserves the integer budget and every feasible mask. Choosing $d=|S_0|$ requires a nonempty base list; choosing $d=N$ scales by all $NH$ possible membership flips. Keeping the same numerical $\alpha$ instead changes the allowed flip count. For example, with $k=10$, $|S_0|=30$ and $\alpha=.10$, the original replacement-average allowance is two flips. Base-list normalization preserves it at $\alpha_d=1/15$; using $.10$ would allow three. A denominator equal to the returned count depends on the mask and has no such fixed rescaling.

For two $k$-element subsets of an $N$-output batch, the largest symmetric difference is $2\min(k,N-k)$. The nominal $2k$ scale is a general upper bound, attained when $2k\le N$. Frozen-cutoff lists can have other sizes, so it is not a universal bound on their churn. An allowance of $\rho$ average membership changes corresponds to $\alpha=\rho/(2k)$. Requiring each replacement to satisfy its own budget instead assigns outputs a vector of replacement costs. The longest feasible prefix of one sorted scalar cost no longer solves the retention problem (Appendix~\ref{app:per-replacement}).

The all-mask replacement-risk diagnostics are reported together in Appendix~\ref{app:individual-replacements}.

Fixed-count and frozen-cutoff stability can differ even without ties. Consider consensus scores $(.90,.80,.10)$ with $k=2$ and cutoff $.75$. If a replacement changes the third score to $.85$, frozen-cutoff selection adds this output and preserves the first two memberships. Top-two selection instead adds the third and excludes the second. This example compares selection objects under the same score change; it is not a comparison of query algorithms. 

The grouped SCoRE reduction in Appendix~\ref{app:review-score} targets the same nonnegative loss with a fixed batch denominator. Further pruning an accepted set cannot increase this loss, so a subsequent certification step limited to pruning the accepted set preserves the marginal expectation bound under the original calibration assumptions. Selected-risk guarantees have a changing denominator and need a separate argument. Re-certifying each future batch already bounds its realized churn within the declared family; calibration does not extend this family or establish semantic correctness.

For any certified mask $C$, a subsequent content policy can restrict commitments to $C'\subseteq C$. Nonnegative $F_i$ and the fixed denominator give $R(C')\le R(C)\le\alpha$. The pruned returned count $K'=|S_0(P)\cap C'|$ has certified regret at most $V^+-K'$ relative to the original optimum. If $C=C^-$ is returned at gap tolerance $\varepsilon$, then $K=V^-$ and this bound is at most $\varepsilon+(K-K')$. The realized regret $V^*-K'$ can be evaluated when $V^*$ is known, through exact certification or full-information evaluation. Improved preference precision still requires label evidence.

\subsection{Per-replacement churn caps}
\label{app:per-replacement}
\begin{table*}[t]
\centering\small
\begin{tabular}{@{}llrrrrrrr@{}}
\toprule
Panel & Pool & $|S_0|$ & $V^*$ & $V^*_{\max}$ & $V^-_0$ & $V^-_{0,\max}$ & Avg.\ (\%) & Per-repl.\ (\%)\\
\midrule
RewardBench~2 & median & 118.9 & 118.4 & 101.7 & 95.5 & 93.1 & 61.4 & 74.0\\
RewardBench~2 & mean & 120.2 & 120.2 & 97.0 & 24.0 & 23.7 & 86.5 & 99.5\\
WMT23 (13) & median & 137.9 & 137.3 & 101.7 & 117.9 & 100.8 & 45.1 & 76.0\\
WMT23 (13) & mean & 126.7 & 126.1 & 68.0 & 24.0 & 24.0 & 85.7 & 99.5\\
WMT23 (12) & median & 114.0 & 113.7 & 108.2 & 112.4 & 106.2 & 3.8 & 42.3\\
WMT23 (12) & mean & 117.4 & 116.1 & 67.8 & 22.0 & 22.0 & 85.4 & 99.3\\
WMT24 en-de & median & 464.6 & 464.6 & 457.1 & 461.9 & 452.8 & 1.3 & 11.9\\
WMT24 en-de & mean & 461.3 & 461.0 & 389.0 & 84.4 & 84.4 & 83.3 & 97.0\\
WMT24 en-es & median & 411.8 & 411.8 & 406.6 & 410.5 & 404.2 & 1.1 & 12.4\\
WMT24 en-es & mean & 413.5 & 413.0 & 350.9 & 74.5 & 74.5 & 82.1 & 96.0\\
WMT24 ja-zh & median & 345.5 & 343.6 & 329.7 & 333.7 & 319.5 & 11.3 & 42.3\\
WMT24 ja-zh & mean & 324.1 & 310.5 & 251.1 & 54.0 & 54.0 & 85.9 & 95.3\\
\bottomrule
\end{tabular}
\caption{Average-budget and per-replacement targets at $q=.05$, $\alpha=.10$, means over 40 splits. $V^*$ and $V^*_{\max}$ are full-information retained counts; $V^-_0$ and $V^-_{0,\max}$ are zero-query safe counts. The last two columns give active-entry inspection at the primary stop on the same reference-priority path.}
\label{tab:per-replacement}
\end{table*}

The average budget bounds expected churn when the executed replacement is drawn uniformly from the declared family. An operator who executes one replacement from the declared family may require every scenario to meet the budget. For row $i$, write $\phi_{i,jr}=\one[\sigma_P(i)\ne\sigma_{P^{j\to r}}(i)]$ and $\beta=\lfloor2k\alpha\rfloor$. The per-replacement target keeps a mask $C$ with $\sum_{i\in C}\phi_{i,jr}\le\beta$ for all $H$ scenarios and maximizes the returned count. Any such mask also meets the average budget, because $\sum_{j,r}\sum_{i\in C}\phi_{i,jr}\le H\beta\le b_H$. The per-replacement full-information optimum $V^*_{\max}$ therefore never exceeds $V^*$. The same target arises from adversarial weights. The weighted risk $R_w(C)$ is linear in the weights. Its maximum over normalized nonnegative weights therefore sits at a single scenario and equals the largest per-replacement churn. Certifying against every weighting is therefore certifying per-replacement caps.

\begin{proposition}[Per-replacement retention is NP-hard]
\label{prop:per-replacement-hard}
For every fixed odd $M$, under odd-median or mean pooling with the candidate family part of the input, deciding from full information whether some mask retains at least $K$ outputs under per-replacement caps is NP-complete. The average-budget problem is solved by sorting.
\end{proposition}
\begin{proof}
Membership in NP is immediate. Reduce from set packing: given sets $A_1,\ldots,A_n\subseteq\{1,\ldots,T\}$ and an integer $K$, decide whether $K$ sets are pairwise disjoint. Use $n$ outputs, cutoff $t_0=1/2$ and, for $M=2h+1$, base scores $.9$ on members $1,\ldots,h+1$ and $.1$ on the others. Every row is a selected boundary row with the same voting members. Give candidate $r$ score $.1$ on row $i$ when $r\in A_i$ and $.9$ otherwise. Replacing a voting member $j$ by candidate $r$ flips row $i$ exactly when $r\in A_i$; replacing a nonvoting member flips nothing. The same scores give the same flips under mean pooling, since $(h+.9)/M\ge1/2>(h+.1)/M$. Choose $\alpha=1/(2k)$, so $\beta=1$. A mask meets every cap exactly when no two retained rows share an element, so the largest retained count equals the largest packing. With the average budget, sorting the scalar costs $F_i$ gives the optimum.
\end{proof}

Partial observations keep the endpoint structure. For each base member $j$, the indicator $\phi_{i,jr}$ is a monotone function of the candidate score, nonincreasing on selected rows. Let $\kappa_{ij}$ be its value at score zero. An unqueried entry $(i,r)$ contributes between zero and $\kappa_{ij}$ to scenario $(j,r)$; a queried entry contributes its observed value. Score zero on every unqueried entry attains all upper scenario loads at once, and score one attains all lower loads. Let $V^-_{\max}$ and $V^+_{\max}$ maximize the returned count under the upper and lower loads. Then $V^-_{\max}\le V^*_{\max}\le V^+_{\max}$, the mask attaining $V^-_{\max}$ meets every cap for every consistent completion, and both bounds are attained. Computing them is NP-hard by Proposition~\ref{prop:per-replacement-hard}; we solve them by integer programming. The joint-evidence characterization of Theorem~\ref{thm:joint} does not extend directly, because exclusion no longer reduces to one sorted threshold.

\paragraph{Empirical comparison.}
We use the six core panels, both pooling rules, four depths, 40 splits and $\alpha=.10$, giving 1,920 instances. Reference priority does not depend on the budget, so one query path serves both targets. For each path we locate the first primary and exact stops of each target by bisection; both gaps are monotone in the path prefix. HiGHS solves each capacity program with a 120-second limit; every program closed within it. Table~\ref{tab:per-replacement} reports nominal $q=.05$. Across the 1,920 instances, the per-replacement optimum was lower in 899. The primary stop needed more queries in 1,401 instances, fewer in 31 and the same number in 488. Under mean pooling at $q=.05$, full-information retention fell from 126.1 to 68.0 outputs on thirteen-stream WMT23 and from 120.2 to 97.0 on RewardBench~2. Primary-stop inspection rose to 95.3--99.5\% of active entries. Under median pooling, it rose from 1.1--61.4\% to 11.9--76.0\%. Records are in \nolinkurl{outputs/scope-extensions/}.

\FloatBarrier
\section{Exact Information Costs Across Query Granularities}
\label{app:granularity}
The objective remains certifying churn-feasible retention within a count tolerance. Changing the query unit, prior observations or pooling rule changes the evidence problem; a pooling change can also change the selected set. In this appendix, finite enumerations check stated instances and implementations; the general claims rest on the proofs, and constructed instances are evaluated apart from the NLP panels.
\begin{table*}[t]
\centering\small
\begin{tabular}{lrrrrrr}
\toprule
Panel & $q$ & Safety & Exclusion & Offline & Fixed & Ref. \\
\midrule
RewardBench 2 & 0.05 & 158.8 & 37.9 & 196.6 & 338.8 & 327.6 \\
RewardBench 2 & 0.10 & 249.2 & 0.0 & 249.2 & 487.3 & 446.9 \\
WMT23, 13 streams & 0.05 & 319.5 & 96.5 & 415.9 & 774.0 & 622.6 \\
WMT23, 13 streams & 0.10 & 1.3 & 0.0 & 1.3 & 2.6 & 3.4 \\
WMT23, 12 streams & 0.05 & 20.5 & 20.6 & 41.1 & 59.0 & 56.0 \\
WMT23, 12 streams & 0.10 & 31.1 & 0.0 & 31.1 & 81.0 & 74.6 \\
\bottomrule
\end{tabular}
\caption{Count means use decimal half-up rounding after aggregating integer records. Exact median stopping, $\alpha=.10$: mean entry counts over 40 splits. Safety counts agreement entries; Exclusion counts opposition entries. They sum to the offline minimum before separate rounding. For every instance, Fixed and Ref.\ online counts are no smaller than this minimum.}
\label{tab:certificates}
\end{table*}

\subsection{Acquisition when the initial gap requires queries}
\label{app:nontrivial}
Before any candidate-test query, all lower costs are zero, so $V^+_0=|S_0|$. The zero-query mask has count $V^-_0$. We stratify configurations by whether $|S_0|-V^-_0>\varepsilon$, using the same initial state and tolerance for every acquisition policy. This condition, not high initial coverage alone, determines whether queries are needed. For each panel, depth, pooling rule and stopping target, we report the subset size and conditional query counts for all four strategies. Exact and positive-gap targets can give different subsets. When the subset is empty, its means are undefined.

At $q=.05$, $\alpha=.10$ and the primary two-percent stop, thirteen-stream WMT23 median requires queries on 36/40 splits; reference priority inspects 50.1\% of active entries on these splits, versus 45.1\% across all 40. Twelve-stream WMT23 requires queries on 3/40, with conditional inspection 51.0\% versus 3.8\% overall. RewardBench median requires queries on all 40, so its conditional and overall inspection both equal 61.4\%. Mean pooling requires queries on every split of these original panels. Tables~\ref{tab:nontrivial-median-tolerant}--\ref{tab:nontrivial-mean-exact} cover all six panels and four depths, including empty subsets. These summaries describe the existing acquisition records; they do not tune protocols or provide independent-population inference.
\begin{table*}[t]
\centering\small


\begin{tabular}{llrrrrrrrr}
\toprule
Panel & $q$ & $n$ & Fixed & Random & Width & Ref. & Ref.\ active & $K$ & Gap \\
\midrule
RewardBench 2 & 0.05 & 40 & 296.3 & 296.8 & 292.2 & 285.0 & 61.4 & 115.6 & 3.0 \\
RewardBench 2 & 0.10 & 39 & 426.6 & 425.4 & 409.0 & 388.1 & 46.2 & 230.1 & 5.3 \\
RewardBench 2 & 0.20 & 37 & 210.8 & 207.7 & 211.2 & 204.8 & 15.2 & 453.8 & 9.8 \\
RewardBench 2 & 0.30 & 1 & 27.0 & 38.0 & 35.0 & 33.0 & 1.8 & 735.0 & 15.0 \\
WMT23, 13 & 0.05 & 36 & 738.3 & 746.0 & 817.2 & 581.4 & 50.1 & 132.5 & 3.2 \\
WMT23, 13 & 0.10 & 0 & -- & -- & -- & -- & -- & -- & -- \\
WMT23, 13 & 0.20 & 0 & -- & -- & -- & -- & -- & -- & -- \\
WMT23, 13 & 0.30 & 0 & -- & -- & -- & -- & -- & -- & -- \\
WMT23, 12 & 0.05 & 3 & 588.0 & 584.7 & 585.7 & 577.7 & 51.0 & 167.0 & 4.0 \\
WMT23, 12 & 0.10 & 5 & 283.2 & 282.6 & 303.0 & 238.2 & 14.8 & 300.6 & 6.6 \\
WMT23, 12 & 0.20 & 1 & 712.0 & 658.0 & 605.0 & 641.0 & 20.1 & 551.0 & 12.0 \\
WMT23, 12 & 0.30 & 1 & 637.0 & 588.0 & 563.0 & 683.0 & 14.8 & 821.0 & 17.0 \\
WMT24 en-de & 0.05 & 2 & 1066.0 & 1094.5 & 1247.0 & 881.0 & 25.8 & 592.0 & 12.5 \\
WMT24 en-de & 0.10 & 5 & 256.4 & 263.2 & 313.8 & 205.2 & 3.6 & 1563.2 & 32.6 \\
WMT24 en-de & 0.20 & 0 & -- & -- & -- & -- & -- & -- & -- \\
WMT24 en-de & 0.30 & 0 & -- & -- & -- & -- & -- & -- & -- \\
WMT24 en-es & 0.05 & 1 & 1484.0 & 1439.0 & 1293.0 & 1329.0 & 45.4 & 733.0 & 15.0 \\
WMT24 en-es & 0.10 & 0 & -- & -- & -- & -- & -- & -- & -- \\
WMT24 en-es & 0.20 & 0 & -- & -- & -- & -- & -- & -- & -- \\
WMT24 en-es & 0.30 & 0 & -- & -- & -- & -- & -- & -- & -- \\
WMT24 ja-zh & 0.05 & 9 & 1278.1 & 1275.8 & 1288.1 & 1228.4 & 50.2 & 558.2 & 12.1 \\
WMT24 ja-zh & 0.10 & 5 & 1374.0 & 1352.6 & 1466.2 & 1273.0 & 33.7 & 1001.2 & 21.0 \\
WMT24 ja-zh & 0.20 & 1 & 336.0 & 329.0 & 595.0 & 359.0 & 9.7 & 765.0 & 16.0 \\
WMT24 ja-zh & 0.30 & 0 & -- & -- & -- & -- & -- & -- & -- \\
\bottomrule
\end{tabular}
\caption{Median entry acquisition with $\varepsilon=\lceil .02|S_0|\rceil$ and $\alpha=.10$. Each panel--depth setting has 40 splits. The shared subset has initial gap greater than $\varepsilon$; $n$ gives its size. Query counts are conditional means for fixed, random, largest-width and reference-priority acquisition. Ref.\ active is the mean $100Q/(mT)$. For reference priority, $K=V^-$ is returned count and Gap is $V^+-V^-$; both average the same initially non-stopping splits. Empty subsets are marked --. Integer count means use decimal half-up rounding; ratios are averaged separately. These overlapping splits give descriptive results.}
\label{tab:nontrivial-median-tolerant}
\end{table*}

\begin{table*}[t]
\centering\small


\begin{tabular}{llrrrrrrrr}
\toprule
Panel & $q$ & $n$ & Fixed & Random & Width & Ref. & Ref.\ active & $K$ & Gap \\
\midrule
RewardBench 2 & 0.05 & 40 & 338.8 & 339.2 & 331.6 & 327.6 & 71.1 & 118.4 & 0.0 \\
RewardBench 2 & 0.10 & 40 & 487.3 & 487.1 & 464.3 & 446.9 & 53.4 & 234.3 & 0.0 \\
RewardBench 2 & 0.20 & 39 & 321.6 & 316.9 & 320.5 & 300.5 & 22.6 & 461.8 & 0.0 \\
RewardBench 2 & 0.30 & 6 & 98.5 & 92.5 & 95.0 & 89.7 & 5.1 & 730.5 & 0.0 \\
WMT23, 13 & 0.05 & 36 & 859.9 & 851.0 & 891.1 & 691.8 & 60.3 & 135.3 & 0.0 \\
WMT23, 13 & 0.10 & 1 & 105.0 & 129.0 & 125.0 & 137.0 & 9.2 & 368.0 & 0.0 \\
WMT23, 13 & 0.20 & 0 & -- & -- & -- & -- & -- & -- & -- \\
WMT23, 13 & 0.30 & 0 & -- & -- & -- & -- & -- & -- & -- \\
WMT23, 12 & 0.05 & 4 & 590.3 & 579.3 & 578.0 & 560.3 & 53.1 & 165.3 & 0.0 \\
WMT23, 12 & 0.10 & 8 & 404.8 & 399.6 & 454.5 & 373.1 & 23.9 & 286.9 & 0.0 \\
WMT23, 12 & 0.20 & 1 & 1191.0 & 1074.0 & 985.0 & 1017.0 & 31.8 & 563.0 & 0.0 \\
WMT23, 12 & 0.30 & 2 & 829.5 & 806.0 & 738.0 & 776.0 & 17.2 & 839.5 & 0.0 \\
WMT24 en-de & 0.05 & 3 & 969.7 & 1009.3 & 1048.3 & 802.3 & 24.3 & 639.7 & 0.0 \\
WMT24 en-de & 0.10 & 5 & 1245.0 & 1259.0 & 1389.2 & 1011.6 & 18.6 & 1595.8 & 0.0 \\
WMT24 en-de & 0.20 & 0 & -- & -- & -- & -- & -- & -- & -- \\
WMT24 en-de & 0.30 & 0 & -- & -- & -- & -- & -- & -- & -- \\
WMT24 en-es & 0.05 & 2 & 1066.0 & 1014.0 & 950.0 & 927.5 & 31.9 & 716.0 & 0.0 \\
WMT24 en-es & 0.10 & 3 & 339.3 & 350.0 & 456.0 & 298.3 & 8.7 & 890.3 & 0.0 \\
WMT24 en-es & 0.20 & 1 & 440.0 & 429.0 & 724.0 & 333.0 & 6.3 & 1423.0 & 0.0 \\
WMT24 en-es & 0.30 & 0 & -- & -- & -- & -- & -- & -- & -- \\
WMT24 ja-zh & 0.05 & 13 & 1194.5 & 1203.3 & 1207.6 & 1144.0 & 49.9 & 528.4 & 0.0 \\
WMT24 ja-zh & 0.10 & 9 & 1412.2 & 1409.3 & 1515.4 & 1290.3 & 36.7 & 929.8 & 0.0 \\
WMT24 ja-zh & 0.20 & 2 & 664.5 & 668.5 & 873.0 & 688.5 & 17.8 & 1243.0 & 0.0 \\
WMT24 ja-zh & 0.30 & 0 & -- & -- & -- & -- & -- & -- & -- \\
\bottomrule
\end{tabular}
\caption{Median entry acquisition with $\varepsilon=0$ and $\alpha=.10$. Each panel--depth setting has 40 splits. The shared subset has initial gap greater than $\varepsilon$; $n$ gives its size. Query counts are conditional means for fixed, random, largest-width and reference-priority acquisition. Ref.\ active is the mean $100Q/(mT)$. For reference priority, $K=V^-$ is returned count and Gap is $V^+-V^-$; both average the same initially non-stopping splits. Empty subsets are marked --. Integer count means use decimal half-up rounding; ratios are averaged separately. These overlapping splits give descriptive results.}
\label{tab:nontrivial-median-exact}
\end{table*}

\begin{table*}[t]
\centering\small


\begin{tabular}{llrrrrrrrr}
\toprule
Panel & $q$ & $n$ & Fixed & Random & Width & Ref. & Ref.\ active & $K$ & Gap \\
\midrule
RewardBench 2 & 0.05 & 40 & 763.6 & 768.5 & 770.6 & 727.9 & 86.5 & 117.3 & 3.0 \\
RewardBench 2 & 0.10 & 40 & 1349.4 & 1362.4 & 1356.2 & 1289.6 & 83.5 & 231.8 & 5.2 \\
RewardBench 2 & 0.20 & 40 & 1978.4 & 1998.9 & 1925.8 & 1919.6 & 76.0 & 452.8 & 9.7 \\
RewardBench 2 & 0.30 & 40 & 1957.3 & 1982.9 & 1781.6 & 2028.7 & 67.0 & 675.1 & 14.3 \\
WMT23, 13 & 0.05 & 40 & 1984.8 & 2004.0 & 2013.2 & 1873.7 & 85.7 & 123.4 & 3.0 \\
WMT23, 13 & 0.10 & 40 & 3729.6 & 3768.0 & 3782.5 & 3518.8 & 83.2 & 241.9 & 5.5 \\
WMT23, 13 & 0.20 & 40 & 2381.8 & 2423.1 & 2167.6 & 2419.2 & 57.4 & 476.4 & 10.1 \\
WMT23, 13 & 0.30 & 40 & 2067.3 & 2111.6 & 1800.6 & 2363.8 & 45.1 & 714.8 & 15.2 \\
WMT23, 12 & 0.05 & 40 & 1827.4 & 1843.5 & 1841.1 & 1757.7 & 85.4 & 113.7 & 2.8 \\
WMT23, 12 & 0.10 & 40 & 3219.9 & 3246.8 & 3239.9 & 3047.4 & 81.3 & 222.2 & 5.0 \\
WMT23, 12 & 0.20 & 40 & 2439.8 & 2481.9 & 2273.8 & 2413.8 & 59.9 & 442.1 & 9.5 \\
WMT23, 12 & 0.30 & 40 & 2797.3 & 2863.3 & 2520.6 & 2979.8 & 54.5 & 666.9 & 14.1 \\
WMT24 en-de & 0.05 & 40 & 5848.3 & 5904.3 & 5920.0 & 5619.0 & 83.3 & 451.6 & 9.7 \\
WMT24 en-de & 0.10 & 40 & 6809.8 & 6896.9 & 6785.0 & 6538.0 & 69.7 & 900.5 & 19.0 \\
WMT24 en-de & 0.20 & 40 & 10060.3 & 10220.6 & 9911.6 & 9591.7 & 63.0 & 1743.6 & 36.0 \\
WMT24 en-de & 0.30 & 40 & 11020.1 & 11247.0 & 10569.2 & 11135.4 & 58.0 & 2591.7 & 53.4 \\
WMT24 en-es & 0.05 & 40 & 5173.8 & 5222.9 & 5228.2 & 4962.1 & 82.1 & 404.5 & 8.8 \\
WMT24 en-es & 0.10 & 40 & 6335.6 & 6414.1 & 6272.6 & 6038.1 & 72.4 & 789.1 & 16.6 \\
WMT24 en-es & 0.20 & 40 & 8977.9 & 9121.1 & 8836.9 & 8509.9 & 64.4 & 1535.9 & 31.9 \\
WMT24 en-es & 0.30 & 40 & 8750.2 & 8914.4 & 8309.6 & 8924.9 & 55.6 & 2301.2 & 47.5 \\
WMT24 ja-zh & 0.05 & 40 & 4560.0 & 4599.3 & 4590.0 & 4402.5 & 85.9 & 306.7 & 6.9 \\
WMT24 ja-zh & 0.10 & 40 & 8142.3 & 8189.5 & 8182.0 & 7834.2 & 86.0 & 562.3 & 12.1 \\
WMT24 ja-zh & 0.20 & 40 & 13303.9 & 13385.2 & 13133.2 & 12544.0 & 82.5 & 1103.9 & 23.1 \\
WMT24 ja-zh & 0.30 & 40 & 15658.6 & 15797.3 & 15180.5 & 14436.1 & 76.1 & 1633.2 & 33.9 \\
\bottomrule
\end{tabular}
\caption{Mean entry acquisition with $\varepsilon=\lceil .02|S_0|\rceil$ and $\alpha=.10$. Each panel--depth setting has 40 splits. The shared subset has initial gap greater than $\varepsilon$; $n$ gives its size. Query counts are conditional means for fixed, random, largest-width and reference-priority acquisition. Ref.\ active is the mean $100Q/(mT)$. For reference priority, $K=V^-$ is returned count and Gap is $V^+-V^-$; both average the same initially non-stopping splits. Empty subsets are marked --. Integer count means use decimal half-up rounding; ratios are averaged separately. These overlapping splits give descriptive results.}
\label{tab:nontrivial-mean-tolerant}
\end{table*}

\begin{table*}[t]
\centering\small


\begin{tabular}{llrrrrrrrr}
\toprule
Panel & $q$ & $n$ & Fixed & Random & Width & Ref. & Ref.\ active & $K$ & Gap \\
\midrule
RewardBench 2 & 0.05 & 40 & 786.6 & 786.5 & 787.8 & 759.2 & 90.2 & 120.2 & 0.0 \\
RewardBench 2 & 0.10 & 40 & 1395.2 & 1394.8 & 1387.1 & 1341.4 & 86.9 & 237.0 & 0.0 \\
RewardBench 2 & 0.20 & 40 & 2069.8 & 2068.5 & 2006.4 & 2006.6 & 79.5 & 462.5 & 0.0 \\
RewardBench 2 & 0.30 & 40 & 2111.0 & 2111.6 & 1906.5 & 2148.8 & 71.0 & 689.3 & 0.0 \\
WMT23, 13 & 0.05 & 40 & 2042.8 & 2041.7 & 2053.1 & 1940.4 & 88.8 & 126.1 & 0.0 \\
WMT23, 13 & 0.10 & 40 & 3833.5 & 3838.8 & 3856.5 & 3647.6 & 86.3 & 247.3 & 0.0 \\
WMT23, 13 & 0.20 & 40 & 2622.2 & 2621.9 & 2352.5 & 2615.6 & 62.1 & 486.5 & 0.0 \\
WMT23, 13 & 0.30 & 40 & 2486.0 & 2473.2 & 2108.1 & 2685.8 & 51.6 & 730.0 & 0.0 \\
WMT23, 12 & 0.05 & 40 & 1879.9 & 1880.6 & 1879.4 & 1823.0 & 88.7 & 116.1 & 0.0 \\
WMT23, 12 & 0.10 & 40 & 3315.4 & 3314.0 & 3311.7 & 3160.6 & 84.3 & 227.2 & 0.0 \\
WMT23, 12 & 0.20 & 40 & 2668.5 & 2666.6 & 2447.8 & 2611.1 & 64.8 & 451.6 & 0.0 \\
WMT23, 12 & 0.30 & 40 & 3180.4 & 3181.9 & 2816.3 & 3274.4 & 60.0 & 681.0 & 0.0 \\
WMT24 en-de & 0.05 & 40 & 5999.6 & 6004.5 & 6017.4 & 5809.8 & 86.2 & 461.0 & 0.0 \\
WMT24 en-de & 0.10 & 40 & 7121.7 & 7122.1 & 7007.2 & 6847.4 & 73.1 & 919.4 & 0.0 \\
WMT24 en-de & 0.20 & 40 & 10716.2 & 10718.8 & 10388.9 & 10201.6 & 67.2 & 1779.6 & 0.0 \\
WMT24 en-de & 0.30 & 40 & 12079.3 & 12092.0 & 11327.7 & 11985.3 & 62.5 & 2645.1 & 0.0 \\
WMT24 en-es & 0.05 & 40 & 5310.6 & 5312.2 & 5316.6 & 5126.9 & 84.9 & 413.0 & 0.0 \\
WMT24 en-es & 0.10 & 40 & 6611.1 & 6612.9 & 6471.5 & 6320.6 & 75.9 & 805.7 & 0.0 \\
WMT24 en-es & 0.20 & 40 & 9562.9 & 9563.1 & 9263.8 & 9054.4 & 68.5 & 1567.7 & 0.0 \\
WMT24 en-es & 0.30 & 40 & 9699.3 & 9684.2 & 9008.4 & 9703.0 & 60.5 & 2348.7 & 0.0 \\
WMT24 ja-zh & 0.05 & 40 & 4674.0 & 4678.6 & 4672.9 & 4554.2 & 89.1 & 310.5 & 0.0 \\
WMT24 ja-zh & 0.10 & 40 & 8339.0 & 8340.0 & 8332.5 & 8093.6 & 89.1 & 570.1 & 0.0 \\
WMT24 ja-zh & 0.20 & 40 & 13731.9 & 13735.6 & 13549.6 & 13115.0 & 86.4 & 1123.6 & 0.0 \\
WMT24 ja-zh & 0.30 & 40 & 16379.6 & 16383.9 & 15846.2 & 15245.9 & 80.3 & 1667.1 & 0.0 \\
\bottomrule
\end{tabular}
\caption{Mean entry acquisition with $\varepsilon=0$ and $\alpha=.10$. Each panel--depth setting has 40 splits. The shared subset has initial gap greater than $\varepsilon$; $n$ gives its size. Query counts are conditional means for fixed, random, largest-width and reference-priority acquisition. Ref.\ active is the mean $100Q/(mT)$. For reference priority, $K=V^-$ is returned count and Gap is $V^+-V^-$; both average the same initially non-stopping splits. Empty subsets are marked --. Integer count means use decimal half-up rounding; ratios are averaged separately. These overlapping splits give descriptive results.}
\label{tab:nontrivial-mean-exact}
\end{table*}

\FloatBarrier

\subsection{Multiplicative entry discovery ratios}
An offline constructor selects certificate entries using the complete score matrix. We define the multiplicative discovery ratio as $Q_{\rm online}/Q_{\rm exact}$ for exact median entry certification, conditional on $Q_{\rm exact}>0$. Splits with $Q_{\rm exact}=0$ have undefined ratios and are reported separately; all acquisition policies stop without queries on these splits. Ratios are computed per split before aggregation. Table~\ref{tab:entry-overhead-summary} reports the median, interpolated 90th percentile and maximum ratios for reference priority at $q=.05$. The supplement's \texttt{entry-overhead-all-*} grids include every original depth and order. Across six panels, four depths, six reference protocols, four acquisition orders and 40 splits, this exact median analysis gives 23,040 records and 576 panel--depth--protocol--order summary cells. Overlapping splits give descriptive distributions, not population confidence intervals or a competitive guarantee.

RewardBench~2 has a mean exact-certificate size of 196.625 entries; reference priority queries 327.575 entries on average. Their ratio, about 1.666, differs from the mean per-instance ratio 1.895 and the median per-instance ratio 1.968. WMT24 en-de and en-es have only three and two positive witnesses at this depth, limiting interpretation of their conditional ratios. Positive-gap, mean and whole-column paths have different witness targets and are excluded from this entry-ratio denominator.
\begin{table*}[t]
\centering\small


\begin{tabular}{lrrrrr}
\toprule
Panel & $Q_{\rm exact}>0$ & $Q_{\rm exact}=0$ & Median & P90 & Max \\
\midrule
RewardBench 2 & 40 & 0 & 1.968 & 2.195 & 2.878 \\
WMT23, 13 streams & 36 & 4 & 1.662 & 2.156 & 2.554 \\
WMT23, 12 streams & 4 & 36 & 2.058 & 2.519 & 2.672 \\
WMT24 en-de & 3 & 37 & 1.824 & 2.110 & 2.182 \\
WMT24 en-es & 2 & 38 & 2.172 & 2.390 & 2.445 \\
WMT24 ja-zh & 13 & 27 & 1.873 & 2.533 & 3.061 \\
\bottomrule
\end{tabular}
\caption{Exact median entry-query discovery ratios at $q=.05$, $\alpha=.10$, full-reference right ECDF and reference priority. Ratios are $Q_{\rm online}/Q_{\rm exact}$ on positive-witness splits; zero-witness splits are separate and use no online queries. $Q_{\rm exact}>0$ and $Q_{\rm exact}=0$ report split counts; each panel has 40 splits. P90 uses linear interpolation; overlapping splits supply descriptive distributions, not population intervals.}
\label{tab:entry-overhead-summary}
\end{table*}



\subsection{Odd-median column certificate characterization}
\label{app:column-witness}
A column query reveals one candidate's scores for all $N$ outputs and costs $N$ score entries. Selected interior rows have zero churn cost; their scores remain included in the $N$ entries charged for each column query. Let $Y\in\{0,1\}^{m\times T}$ contain the boundary rows\textquotesingle\ opposition bits. Under bit budget $b$, define
\[
 V_B(w)=\max\Big\{|C|:C\subseteq\mathcal B,\ \sum_{i\in C}w_i\le b\Big\},
\]
and true optimum $L=V_B(z)$. For a queried column subset $A$, write $o_i(A)=\sum_{r\in A}y_{ir}$. Every consistent completion has row cost between $o_i(A)$ and $o_i(A)+T-|A|$. The minimum number of column queries for exact certification is
\begin{align*}
 K_{\rm col}&=\min_{A\subseteq\{1,\ldots,T\}}|A|\\
 &\text{subject to}\\
 V_B(o(A)+T-|A|)&=V_B(o(A))=L.
\end{align*}
Both capacities here concern boundary rows, with the same bit budget. Adding the known interior count gives the full selected-retention optimum.

\begin{proposition}[Exact column condition]
\label{prop:column-condition}
Order the observed counts as $o_{(1)}(A)\le\cdots\le o_{(m)}(A)$. A subset $A$ certifies exact optimality if and only if
\begin{align*}
 \sum_{i=1}^{L}o_{(i)}(A)+L(T-|A|)&\le b,\\
 \sum_{i=1}^{L+1}o_{(i)}(A)&>b\quad\text{if }L<m.
\end{align*}
The first condition uses an empty sum when $L=0$; the second is omitted when $L=m$. For $m>0$,
\[
 \left\lceil\frac{Q_{\rm exact}}m\right\rceil
 \le K_{\rm col}\le\min(T,Q_{\rm exact}).
\]
If $m=0$ or $Q_{\rm exact}=0$, $K_{\rm col}=0$.
\end{proposition}
\begin{proof}
The upper and lower costs share their order because the unobserved-column offset is common to every row. Their cheapest prefixes therefore have the displayed costs. The true optimum lies between their capacities. Since upper-cost capacity cannot exceed the true optimum $L$, it equals $L$ exactly when the cheapest $L$-row upper-cost prefix is feasible. The capacity under lower costs is at most $L$ exactly when the cost sum of its cheapest $(L+1)$-row prefix exceeds the budget. This proves necessity and sufficiency, including both endpoints.

A $c$-column witness reveals at most $mc$ relevant boundary entries. Keeping only the boundary entries yields an entry certificate, so $Q_{\rm exact}\le mc$. Conversely, reveal every column touched by a minimum entry witness. Extra observations cannot decrease safe capacity or increase optimistic capacity, so they cannot enlarge the gap; at most $Q_{\rm exact}$ columns are touched. Revealing all $T$ columns always suffices. A zero entry witness certifies the same problem before any column query; with no boundary rows optimality is already known.
\end{proof}

For these candidate counts, we compute exact column minima by enumerating subsets and testing Proposition~\ref{prop:column-condition}. Encode each row's opposition support as a $T$-bit mask and merge equal supports with their multiplicities. For each column subset, counting opposition bits in its intersection with each row mask gives a histogram of observed costs. Histogram prefixes test both conditions without enumerating returned masks. With $D_{\rm sup}$ distinct row supports, worst-case time is $O(mT+2^T(D_{\rm sup}+T))$ and space is $O(m+2^T+D_{\rm sup}+T)$ beyond the input. This bound treats a $T$-bit mask intersection and population count as unit-cost operations, as when masks fit a machine word with a population-count primitive. Multiword implementations must include their bit-operation costs; the original seven-member-pool panels considered here have $T\le17$. The entry lower bound and feasible online column counts restrict the subset sizes requiring enumeration. Every reported optimum is additionally checked using raw-row sorted costs and a separate enumeration of every $(K_{\rm col}-1)$-subset. Revealing more columns cannot invalidate a certificate, so failure of every subset at this size rules out every smaller subset. When $K_{\rm col}=0$ there is no smaller level.

\paragraph{Evidence shared across rows.}
Take $m=T=2$, $b=0$ and row costs $(1,1)$. The matrices
\[
 \begin{pmatrix}1&0\\1&0\end{pmatrix}
 \quad\text{and}\quad
 \begin{pmatrix}1&0\\0&1\end{pmatrix}
\]
share $L=0$ and $Q_{\rm exact}=2$, but have minimum column counts one and two. They are realizable with seven-member pools having four scores above a $.5$ cutoff and three below. For these selected median boundary rows, candidate--output scores below the cutoff implement opposition bits; scores at or above it implement agreement bits. Row sums determine entry certificate size; column certificate size also depends on whether the needed evidence occurs in the same columns. More generally, for any $m\ge T\ge1$, give each row one opposition and set $b=0$. Shared support on one candidate gives $K_{\rm col}=1$; assigning the opposition cyclically across all $T$ candidates gives $K_{\rm col}=T$. Both have $Q_{\rm exact}=m$ and identical row costs. They attain the lower and upper bounds of Proposition~\ref{prop:column-condition}, respectively.

When $b=0$ and every boundary row has a true opposition, $L=0$ and the exact column condition requires at least one observed opposition in every row. Associate each candidate column with the rows where it has an opposition. The minimum column certificate is a minimum set cover of the boundary rows. Covering formulations and queries revealing coefficients across constraints already occur in uncertainty exploration \citep{megow2023set,schloter2025knapsack}. This specialization identifies the coupling in the present median model; it does not supply an online approximation or competitive guarantee.

\subsection{Minimum residual entry evidence for odd medians}
\label{app:residual}
An existing observation can supply agreement or opposition evidence before further entries are queried. This matters when a column is bought once and the remaining evidence is acquired cell by cell. We condition the minimum entry witness on arbitrary revealed cells, keeping the original odd-median completion model and bit budget $b$.

Let $\mathcal O$ be the already observed cells, with agreement counts $a_i$, opposition counts $o_i$ and upper costs $u_i=T-a_i$. The true row costs are $z_i$, and $L^*=V_B(z)$. Interior outputs add the same known count to both count bounds and cancel from the gap. For $0\le s\le L^*$, define
\[
 G_s(\mathcal O)=\max\left\{0,
 \min_{\substack{|C|=s\\\sum_{i\in C}z_i\le b}}
 \sum_{i\in C}u_i-b\right\}.
\]
The minimum ranges over boundary rows; for $s=0$ its sum is zero. It counts the fewest new agreements needed to make an $s$-row set safe. In this offline minimization, true costs restrict the sets to the actually feasible ones. Existing agreements determine the additional evidence needed for each such set. A dynamic program indexed by retained-row count and total true cost computes the minimum. Let $D_j(s,c)$ be the minimum upper-cost sum for $s$ boundary rows among the first $j$ whose total true cost is exactly $c$. Initialize $D_0(0,0)=0$ and all other states to $+\infty$. Then
\[
 \begin{aligned} D_j(s,c)&=\min\{D_{j-1}(s,c),\\
 &\quad D_{j-1}(s-1,c-z_j)+u_j\}.\end{aligned}
\]
Infeasible indices are omitted. The resulting minimum is
\[
 G_s(\mathcal O)=\max\Big\{0,\min_{0\le c\le b}D_m(s,c)-b\Big\}.
\]
The cost dimension can be capped at $\min(b,mT)$. Writing $B_{\rm DP}=\min(b,mT)+1$, a dense implementation takes $O(m(L^*+1)B_{\rm DP})$ time; the supplied implementation stores only reachable states. This is an offline computation using true costs $z_i$, not an acquisition policy.

For an exclusion target $L^*<p\le m$, put
\[
 v_i(t)=\max\{o_i,\min(z_i,t)\}.
\]
If the $p$ smallest $o_i$ already sum above $b$, set $E_p(\mathcal O)=0$. Otherwise, over integer levels $0\le t\le T$, define
\begin{align*}
 \lambda_p&=\min\left\{t:
 \sum_{j=1}^{p}v_{(j)}(t)>b\right\},\\
 E_p(\mathcal O)&=\sum_i\big(v_i(\lambda_p)-o_i\big)\\
 &\quad-\left(\sum_{j=1}^{p}v_{(j)}(\lambda_p)-b-1\right).
\end{align*}
At each level $t$, the values $v_i(t)$ are sorted anew; this order need not match the true-cost order. For $p>m$, set $E_p(\mathcal O)=0$. Thus existing opposition evidence is retained, not charged again.

\begin{proposition}[Minimum residual entry witness]
\label{prop:residual}
For integer $\varepsilon\ge0$, the minimum additional entry count certifying a gap at most $\varepsilon$ is
\[
 Q_\varepsilon(\mathcal O)=
 \min_{s=\max(0,L^*-\varepsilon)}^{L^*}
 \left[G_s(\mathcal O)+E_{s+\varepsilon+1}(\mathcal O)\right].
\]
With no initial observations this reduces to Corollary~\ref{cor:gap-witness}.
\end{proposition}
\begin{proof}
For a truly feasible $s$-row set $C$, each additional agreement reduces its upper-cost sum by one. At least $\max(0,\sum_{i\in C}u_i-b)$ additional agreements are necessary. The available agreements total $\sum_{i\in C}(u_i-z_i)$, enough to attain this number because $\sum_{i\in C}z_i\le b$. Minimizing over $C$ gives $G_s$.

For exclusion, write $B=b+1$ and assume the observed prefix does not already exclude $p$ rows. The critical level exists since $p>L^*$. Moreover $\lambda_p\le z_{(p)}$. Rowwise, $v_i(t)\ge\min(z_i,t)$, so their ordered values satisfy the same inequality. At $t=z_{(p)}$, the $p$ smallest truncated true costs sum to $\sum_{j=1}^p z_{(j)}\ge B$; hence the $p$ smallest $v_i(t)$ also sum to at least $B$.

Consider any sufficient final opposition counts $w_i\in[o_i,z_i]$, and let $t=w_{(p)}$. For $w_i\le t$, we have $v_i(t)\ge w_i$ because $o_i\le w_i\le z_i$; for $w_i>t$, $v_i(t)\ge t$. Thus $v_i(t)\ge\min(w_i,t)$ rowwise. Since $t=w_{(p)}$, truncating at $t$ preserves the sum of the $p$ smallest counts. Their sum cannot decrease when $\min(w,t)$ is replaced by $v(t)$, so it remains at least $B$, implying $t\ge\lambda_p$. Every row with $o_i>t$ lies outside this prefix. Consequently the total new opposition count is at least
\[
 B+\sum_{j=1}^{m-p}\big[\max(o,t)\big]_{[j]}-\sum_i o_i,
\]
where brackets order values from largest to smallest. This expression is nondecreasing in $t$. At $t=\lambda_p$ it equals $E_p$: because $\lambda_p\le z_{(p)}$, the largest $m-p$ values of $v(\lambda_p)$ equal the largest $m-p$ values of $\max(o,\lambda_p)$.

To attain the bound, start with $v(\lambda_p)$ and remove new opposition counts one at a time among the current $p$ smallest row counts until the sum of these counts is $B$, reordering after each removal. A removable count has $w_i>o_i$ and $w_i\le w_{(p)}$; reducing it decreases the prefix sum by exactly one. Such a count exists whenever the sum remains above $B$. Otherwise the cheapest $p$ counts are unchanged observations, and all positive increments lie above their threshold. Every row with a positive increment has final count at most $\lambda_p$. Since all such rows lie above the prefix threshold, the threshold is below $\lambda_p$. The vector $v(\lambda_p-1)$ would then have a prefix sum at least as large, contradicting minimality of $\lambda_p$. The construction removes exactly the excess subtracted in $E_p$.

Any successful gap certificate has safe count $s\in[\max(0,L^*-\varepsilon),L^*]$. If $p=s+\varepsilon+1\le m$, it must prove no $p$-row set feasible under any consistent completion. If $p>m$, no exclusion evidence is needed. New agreement and opposition cells are disjoint, so their lower bounds add. Conversely, the two constructions can be united: agreements change upper costs and oppositions change lower costs. They certify the required gap and attain the displayed minimum.
\end{proof}

\subsection{Joint entry evidence beyond Boolean costs}
\label{app:joint-evidence}
Mean pooling changes how observations serve a certificate. For a selected row with base-score sum $S_i$, replacement of member $j$ flips membership when $x<Mt_0-(S_i-s_j(i))$. Thus $f_i(x)$ is an integer staircase, although scores are continuous. The two endpoint costs remain attainable, as required by Proposition~\ref{prop:stop}; intermediate query outcomes can tighten both bounds, raising the lower cost and lowering the upper cost.

Consider two selected rows with identical base scores $(.1,.7,.9)$, cutoff $.5$ and one candidate scoring $.6$ on each row. Take $M=3$, $N=2$, $q=.5$, $k=1$ and $\alpha=1/6$, so $b_H=1$. Each row has $p_i=2$ and true cost one. No queries give $(V^-,V^+)=(0,2)$, one query gives $(1,2)$, and two give $(1,1)$. Safety requires one query and exclusion requires two, but their union requires two, not three. This example shows the evidence overlap absent from odd-median boundary rows.

\paragraph{Attainable row states.}
For initial observations $\mathcal O$, let $U_i$ be the unobserved candidates on active row $i$. Define
\[
 \mathcal D_i=\left\{\left(|J|,\sum_{r\in J}f_{ir}\right):J\subseteq U_i\right\}.
\]
Initialize the row-state set with $(0,0)$. For each remaining candidate, retain existing states and add $(a+1,v+f_{ir})$ from each $(a,v)$. Store one witnessing subset for each pair. There are at most $(T+1)(Tp_i+1)$ pairs. The offline constructor uses the true contributions to build these states; the verifier subsequently sees only the chosen observations.

\begin{lemma}[Excluding a cardinality by one threshold]
\label{lem:order-threshold}
For nonnegative integer lower costs $\ell$ and $1\le p\le m$,
\[
 \sum_{j=1}^{p}\ell_{(j)}=
 \max_{t\in\{0,\ldots,TM\}}\left[pt-\sum_i(t-\ell_i)_+\right],
\]
provided every $\ell_i\le TM$.
\end{lemma}
\begin{proof}
For any $p$-row set $A$, the bracket is at most $\sum_{i\in A}\ell_i$: bounding each $(t-\ell_i)_+$ below by $t-\ell_i$ on $A$ proves this inequality. At $t=\ell_{(p)}$, all values below $t$ belong to the cheapest $p$ rows, while any remaining rows in this prefix equal $t$. The bracket therefore equals their sum. Integer costs make exceeding $b_H$ equivalent to reaching $b_H+1$.
\end{proof}

\paragraph{Proof of Theorem~\ref{thm:joint}.}
Any successful observation set has active safe count $s=V^--m_0$ in $[\max(0,L^*-\varepsilon),L^*]$, where $m_0$ is the invariant selected count. Its queried cells define row states, and a safe prefix supplies $C$. If $p=s+\varepsilon+1\le m$, optimistic active count $V^+-m_0$ is at most $s+\varepsilon$, so the cheapest $p$ lower costs exceed $b_H$. Lemma~\ref{lem:order-threshold} supplies the required threshold. Thus every successful certificate appears in the minimization.

Conversely, witnessing row subsets produce the stated bounds. The safe-set constraint gives at least $s$ active retained outputs. The threshold condition gives optimistic active count at most $s+\varepsilon$; for $p>m$, this bound is automatic. Hence the count gap is at most $\varepsilon$. Each additional cell belongs to exactly one row, so its cost is counted once even when it improves both bounds. Minimizing constructs a certificate of exactly the minimum size. The proof uses joint endpoint attainability, not attainability of every integer between the endpoints.

\paragraph{Row-state dynamic program.}
For fixed $s,t$, put $D=pt-b_H-1$ and discard negative $D$. For each processed row prefix, a state $(n,u,d)$ stores the minimum additional query count. Here $n$ counts safe rows, $u$ is their upper-cost sum, and $d$ is the accumulated threshold deficit. Initialize $(0,0,0)$ with cost zero. For an attainable row state $(a,v)$ and a safe-row indicator $y\in\{0,1\}$, update
\[
 \begin{aligned}
 n'&=n+y,\\
 u'&=u+y(u_i^0-ap_i+v),\\
 d'&=d+(t-\ell_i^0-v)_+,
 \end{aligned}
\]
adding query cost $a$. Discard states with more than $s$ safe rows, upper cost above $b_H$ or deficit above $D$; read the minimum at $n=s$ and reconstruct its row subsets and safe set. When $p>m$, omit the deficit coordinate and threshold. Minimizing over the permitted $s$ and thresholds gives Theorem~\ref{thm:joint}, with safety and exclusion coupled through the same row states.

Writing $R=\max_i|\mathcal D_i|$, the dense computation for fixed $s,t$ takes $O(mR(s+1)(b_H+1)(D+1))$ time and $O((s+1)(b_H+1)(D+1))$ working memory, apart from witness storage. The implementation stores reachable states. In the single-member family, $p_i\le M$ and row costs are at most $TM$. If $b_H\ge\sum_iTp_i$, zero queries suffice; otherwise $b_H<mTM$ and $D\le mTM$. Together with $R\le(T+1)(TM+1)$, these bounds make the algorithm polynomial in the explicit counts $m,T,M$, although the dense state space can be large. For generic integer contributions with binary-encoded large caps, the same budget-indexed algorithm is pseudo-polynomial. The odd-median formulas provide a much smaller specialization.


\paragraph{Complete search cost.}
Row states can be cached once in $O(\sum_i T|\mathcal D_i|)$ time and $O(\sum_i|\mathcal D_i|)$ space, apart from stored subsets. There are at most $\min(\varepsilon,L^*)+1$ safe targets and $TM+1$ thresholds per target. Multiplying the fixed-$s,t$ bound by these counts bounds total search time; the working state grid is reused. Thus the explicit single-member model is tractable in polynomial time even when its dense implementation is impractical. Binary-encoded generic costs need not satisfy these polynomial bounds. An incumbent query count $q_0$ permits discarding partial states already costing at least $q_0$, since remaining query costs are nonnegative. This pruning seeks strict improvement and preserves a previously constructed witness when no improvement exists.

\begin{algorithm}[t]
\caption{Offline minimum joint witness by row states}
\label{alg:offline-joint}
\begin{algorithmic}[1]
\REQUIRE True $m\times T$ contributions $0\le f_{ir}\le p_i\le M$, base-pool size $M$, consistent observations, integer budget $b_H\ge0$ and integer tolerance $\varepsilon\ge0$, under Theorem~\ref{thm:joint}.
\STATE Construct and cache each $\mathcal D_i$; compute $L^*$ from true row costs.
\STATE Initialize an incumbent with all remaining active entries.
\FOR{$s=\max(0,L^*-\varepsilon),\ldots,L^*$}
\STATE Set $p=s+\varepsilon+1$.
\FOR{each integer $t\in\{0,\ldots,TM\}$ with $pt\ge b_H+1$ (or once without $t$ if $p>m$)}
\STATE Initialize $(n,u,d)=(0,0,0)$ with cost zero; omit $d$ when $p>m$.
\STATE Apply the row transitions above; keep the cheapest witness for each reachable state.
\STATE Compare feasible terminal states ($n=s$, $u\le b_H$ and, if $p\le m$, $d\le pt-b_H-1$) against the incumbent.
\ENDFOR
\ENDFOR
\RETURN The minimum query count and its reconstructed entries.
\end{algorithmic}
\end{algorithm}

\FloatBarrier
\paragraph{Even-sized median pools.}
\label{app:even-median}
The separation between Boolean and shared evidence concerns replacement outcomes, not only the name of the pooling rule.
\begin{corollary}[Joint evidence for even medians]
\label{cor:even-median}
Let $M$ be even and define median consensus as the average of the two central scores. Under the same fixed cutoff, uniformly weighted single-member replacements and completion model without cross-entry constraints, Theorem~\ref{thm:joint} characterizes the minimum additional entry certificate size. Its row-state dynamic program and outcome-count integer formulation remain valid.
\end{corollary}
\begin{proof}
Median consensus is coordinatewise nondecreasing. For a base-selected row, replacing any member by $1$ preserves selection, while replacing it by $0$ maximizes the exclusion indicator for this replacement. Thus $f_i(1)=0$, $f_i(0)=p_i$, and $f_i(x)$ is a nonincreasing integer-valued function bounded by $p_i\le M$. The Cartesian product of the entry domains makes both endpoint cost vectors jointly attainable. Querying $(i,r)$ raises the lower cost by $f_{ir}$ and lowers the upper cost by $p_i-f_{ir}$. These are precisely the bounds and reachable row states used in Theorem~\ref{thm:joint}; its safety, threshold-exclusion and attainability arguments therefore apply unchanged.
\end{proof}
For example, take two rows with base scores $(.4,.8)$, cutoff $.5$ and one candidate scoring $.4$ on each row. With $k=1$ and $\alpha=1/4$, the flip budget is one. Both rows have endpoint cap two and true cost one. The successive count bounds are $(0,2)$, $(1,2)$ and $(1,1)$ after zero, one and two queries. Safety needs one observation and exclusion needs two, while their union needs two. Even medians therefore need not have the disjoint evidence underlying the odd-median closed forms.

\paragraph{Solver implementation and resources.}
The integer programs use Python 3.9.6, NumPy 2.0.2 and SciPy 1.13.1 on macOS 15.6 arm64, through \texttt{scipy.optimize.milp} on an Apple M4 Pro. The solver runtime banner reports HiGHS 1.2.0. The original SciPy calls explicitly set the remaining time limit, zero relative MIP gap and presolve. Later direct-HiGHS calls also set one thread and random seed zero. These call families do not fully record effective feasibility or integrality tolerances for the historical targets; present-day defaults cannot recover them. The explicit settings for the new secondary analysis are reported in Appendix~\ref{app:deduplicated}; they do not reconstruct omitted historical settings. A numerical lower bound $L_B$ is converted to $\max(0,\lceil L_B-10^{-6}\rceil)$; reconstructed witnesses are checked in integer arithmetic. This combines numerical optimality bounds with exact feasibility checks, not rational proof certification. Worker processes set BLAS/OpenMP thread limits to one. The environment and attempt ledger report machine details, time limits and retained failures; process memory measurements are not isolated solver peaks.

\paragraph{A compact outcome-count formulation.}
\label{app:compact-joint}
Invariant selected rows cancel from the count gap. For unit-cost entry witnesses, unqueried entries within the same active row are interchangeable when their true flip contributions are equal. For active row $i$, let $d_{iv}$ count remaining candidates with contribution $v\in\{0,\ldots,p_i\}$. Querying $a_{iv}$ of them gives the same row bounds regardless of which candidates are chosen. This groups integer outcomes without discretizing scores.

For a fixed safe target $s$, write $p=s+\varepsilon+1$. Use integer counts $0\le a_{iv}\le d_{iv}$, binary safe-row indicators $y_i$, and continuous products $z_{iv}$ constrained by
\begin{align*}
 0\le z_{iv}&\le a_{iv},\\
 z_{iv}&\le d_{iv}y_i,\\
 z_{iv}&\ge a_{iv}-d_{iv}(1-y_i).
\end{align*}
These inequalities make $z_{iv}=a_{iv}y_i$ exactly. Minimize $\sum_{i,v}a_{iv}$ subject to
\begin{align*}
 \sum_i y_i&=s,\\
 \sum_i u_i^0y_i-\sum_{i,v}(p_i-v)z_{iv}&\le b_H.
\end{align*}
If $p\le m$, also use an integer threshold $0\le t\le TM$ and nonnegative continuous deficits $w_i$ satisfying
\begin{align*}
 w_i&\ge t-\ell_i^0-\sum_v v a_{iv},\\
 pt-\sum_i w_i&\ge b_H+1.
\end{align*}
For $p>m$, omit the threshold and deficit constraints.

\begin{proposition}[Outcome-count equivalence]
\label{prop:compact-joint}
Minimizing this program over the same safe targets as Theorem~\ref{thm:joint} gives the same minimum additional entry count. With $G=\sum_i|\{v:d_{iv}>0\}|\le m\min(T,M+1)$, each target has $O(G+m)$ variables and constraints.
\end{proposition}
\begin{proof}
Every entry witness determines counts $a_{iv}$ and a safe set of $s$ rows. The product constraints encode its safe upper-cost sum. Lemma~\ref{lem:order-threshold} supplies a threshold and deficits satisfying the exclusion constraint. Conversely, choose any $a_{iv}$ distinct remaining entries from each outcome group. They realize lower costs $\ell_i^0+\sum_v v a_{iv}$ and upper costs $u_i^0-\sum_v(p_i-v)a_{iv}$. The indicated $s$ rows are safe, and the deficits exclude every $p$-row set. Each query is counted once. Both mappings preserve objective value, proving equivalence. There are at most $M+1$ outcome groups per row and at most $T$ nonempty groups; the variable and constraint bounds follow. These size bounds do not bound integer-program search time.
\end{proof}
The ungrouped formulation uses one binary variable per unqueried entry, treating each entry as a separate group of capacity one. Both formulations use the true matrix offline. Outcome grouping is specific to unit-cost entries; whole-column actions still depend on candidate identity.

\paragraph{Offline threshold screening.}
For exact stopping, $s=L^*$ and $p=L^*+1$. Every partial lower cost satisfies $\ell_i\le F_i$, so a threshold is impossible if
\[
 pt-\sum_i(t-F_i)_+<b_H+1.
\]
Discarding these thresholds and solving the compact program at every remaining integer $t\in\{0,\ldots,TM\}$ preserves the global minimum. If $p>m$, no exclusion threshold is needed. This decomposition uses true costs offline; it does not expose them to the online policy or claim a global optimum from a partially solved threshold.

\paragraph{Separate witnesses as a comparison.}
For each safe target $s$, independently minimize safety queries and exclusion queries for $p=s+\varepsilon+1$, then minimize their sum over the same target range. Call this additive upper bound $Q_{\rm split}$. It may count shared cells twice and is not the size of an optimally deduplicated union. The union of these separately constructed witnesses is feasible because both bounds improve monotonically under additional observations. Therefore $Q_{\rm joint}\le Q_{\rm split}$. Equality holds for odd-median boundary evidence under Proposition~\ref{prop:residual}; intermediate contributions under mean or even-median pooling can make the inequality strict.

\subsection{Joint-certificate validation and NLP subbatches}
\label{app:joint-validation}
Independent query-subset enumeration checks 6,316 capped-integer cases, 768 scaled median cases and 400 realizable mean cases with exact rational scores. Histories, endpoint outcomes, intermediate flips, ties, budgets and count tolerances are included. The 768 scaled odd-median cases use empty histories and match Corollary~\ref{cor:gap-witness}. The generator \nolinkurl{scripts/check_joint_certificates.py} specifies the enumerated domains and the seeded random instances; subset enumeration is exhaustive within each generated instance, not over every possible input.

For each of six evaluated panels, both pooling rules, four depths and forty splits, take the first four baseline-selected outputs in the original label-blind tie order, or all of them if fewer exist. Reference maps, the cutoff and complete candidate family stay fixed. Set $N$ to this subbatch size, reset $k=\max(1,\operatorname{round}(qN))$ and use $\alpha=.10$. The original $q$ still determines the original reference cutoff; this construction defines a separate batch, not a subset certificate for the original budget. Invariant rows are included and require no queries. All 1,920 realized subbatches have four outputs; the smaller-batch fallback is not exercised by these records (Table~\ref{tab:joint-small}).

We compute minima at $\varepsilon\in\{0,1\}$ from empty history and from a fixed quarter-history: reveal the first $\lfloor mT/4\rfloor$ eligible cells in seeded row then candidate-index order. These 1,920 batch configurations give 7,680 joint minima and matching $Q_{\rm split}$ comparisons. A separate binary integer optimization enumerates all safe and excluded output subsets, instead of using the threshold identity or row states. All 7,680 small-batch optimization checks report zero optimizer optimality gap; their witnesses are checked in integer arithmetic. This solver gap is distinct from the retained-count gap. Both independent formulations return the same minimum additional joint query count on all 7,680 tasks, with no unresolved instances. Complete records include reconstructed query sets and safe masks.

Under exact mean certification, $Q_{\rm joint}<Q_{\rm split}$ in 186/960 empty-history instances; their differences total 306 entries, or .319 per instance. With fixed quarter-history, 105 instances have a total difference of 157 additional entries. At $\varepsilon=1$, these counts are 74 and 23, respectively. All 3,840 median history/tolerance comparisons have zero difference. These reductions from an additive upper bound are modest and instance-dependent. The sum of separate minima remains valid but need not equal the minimum joint certificate size.

\begin{table*}[t]
\centering\small
\begin{tabular}{lrrrrrrr}
\toprule
Panel & $q$ & $Q_{\rm joint}$ & $Q_{\rm split}$ & \shortstack{$Q_{\rm joint}<Q_{\rm split}$\\splits/40} & $Q_{\rm joint}(\mathcal O)$ & $Q_{\rm split}(\mathcal O)$ & \shortstack{$Q_{\rm joint}<Q_{\rm split}$\\splits/40} \\
\midrule
RewardBench 2 & 0.05 & 19.63 & 19.95 & 13 & 13.85 & 13.95 & 4 \\
RewardBench 2 & 0.10 & 19.05 & 19.30 & 6 & 13.30 & 13.45 & 3 \\
RewardBench 2 & 0.20 & 13.95 & 14.48 & 11 & 10.73 & 11.05 & 8 \\
RewardBench 2 & 0.30 & 12.93 & 13.08 & 5 & 9.55 & 9.70 & 5 \\
WMT23, 13 streams & 0.05 & 39.18 & 39.30 & 4 & 25.48 & 25.58 & 4 \\
WMT23, 13 streams & 0.10 & 49.33 & 49.65 & 8 & 35.93 & 36.05 & 3 \\
WMT23, 13 streams & 0.20 & 21.05 & 21.18 & 5 & 15.25 & 15.33 & 3 \\
WMT23, 13 streams & 0.30 & 15.90 & 16.40 & 8 & 10.88 & 11.10 & 3 \\
WMT23, 12 streams & 0.05 & 43.25 & 43.25 & 0 & 29.53 & 29.53 & 0 \\
WMT23, 12 streams & 0.10 & 46.38 & 46.43 & 2 & 33.08 & 33.08 & 0 \\
WMT23, 12 streams & 0.20 & 19.30 & 19.63 & 5 & 13.83 & 14.00 & 3 \\
WMT23, 12 streams & 0.30 & 19.08 & 19.60 & 13 & 13.18 & 13.35 & 7 \\
WMT24 en-de & 0.05 & 37.85 & 38.08 & 7 & 27.35 & 27.50 & 5 \\
WMT24 en-de & 0.10 & 32.55 & 32.73 & 5 & 23.48 & 23.50 & 1 \\
WMT24 en-de & 0.20 & 21.33 & 21.73 & 12 & 14.73 & 14.98 & 8 \\
WMT24 en-de & 0.30 & 18.83 & 19.10 & 6 & 13.08 & 13.25 & 5 \\
WMT24 en-es & 0.05 & 40.68 & 40.95 & 10 & 28.98 & 29.13 & 5 \\
WMT24 en-es & 0.10 & 29.98 & 30.05 & 3 & 21.88 & 21.90 & 1 \\
WMT24 en-es & 0.20 & 23.03 & 23.23 & 6 & 16.23 & 16.38 & 4 \\
WMT24 en-es & 0.30 & 18.55 & 18.70 & 3 & 13.20 & 13.30 & 2 \\
WMT24 ja-zh & 0.05 & 40.13 & 40.70 & 13 & 28.40 & 28.75 & 8 \\
WMT24 ja-zh & 0.10 & 40.25 & 40.83 & 12 & 29.65 & 29.83 & 6 \\
WMT24 ja-zh & 0.20 & 36.83 & 37.35 & 11 & 28.50 & 28.85 & 7 \\
WMT24 ja-zh & 0.30 & 28.58 & 29.55 & 18 & 20.78 & 21.20 & 10 \\
\bottomrule
\end{tabular}
\caption{Exact mean entry certificates on the predefined subbatches of at most four baseline-selected outputs. Each row averages 40 splits, retaining original maps, cutoff and all candidates while resetting $N$, $k$ and $b_H$ for this subbatch. The first group starts empty; the second reports minimum additional entries after observing the first quarter of eligible cells in seeded row and candidate-index order. $Q_{\rm split}$ adds independently minimized safety and exclusion sizes over the same safe-target range. The $Q_{\rm joint}<Q_{\rm split}$ columns count splits below the additive bound; they do not compare against the best deduplicated union. An independent subset formulation gives matching numerical optima; reconstructed witnesses pass integer feasibility checks. These are not minima for the original full test batches.}
\label{tab:joint-small}
\end{table*}


The four original complete acquisition policies are also run on these same batches from empty history through exact stopping, retaining their original reference-fitted predictions and tie conventions. Their first stops at $\varepsilon=0$ and one give 15,360 records. Table~\ref{tab:joint-discovery} reports mean-pooling discovery ratios at $q=.05$, averaged over splits with positive matching subbatch minima. Full-panel mean acquisition is not divided by these small-batch minima. These analyses use previously exposed panels and provide descriptive evidence.
\begin{table*}[t]
\centering\small
\begin{tabular}{lrrrrrrr}
\toprule
Panel & $\varepsilon$ & Mean $Q_{\rm joint}$ & $n_+$ & Fixed & Random & Width & Reference \\
\midrule
RewardBench 2 & 0 & 19.63 & 40 & 1.405 & 1.447 & 1.456 & 1.288 \\
RewardBench 2 & 1 & 12.73 & 40 & 1.741 & 2.029 & 2.180 & 1.710 \\
WMT23, 13 streams & 0 & 39.18 & 40 & 1.795 & 1.844 & 1.836 & 1.774 \\
WMT23, 13 streams & 1 & 23.18 & 40 & 2.666 & 3.179 & 3.312 & 3.061 \\
WMT23, 12 streams & 0 & 43.25 & 40 & 1.587 & 1.658 & 1.642 & 1.406 \\
WMT23, 12 streams & 1 & 26.75 & 40 & 2.090 & 2.431 & 2.481 & 2.080 \\
WMT24 en-de & 0 & 37.85 & 40 & 1.562 & 1.591 & 1.573 & 1.474 \\
WMT24 en-de & 1 & 24.48 & 40 & 1.956 & 2.362 & 2.452 & 2.130 \\
WMT24 en-es & 0 & 40.68 & 40 & 1.384 & 1.421 & 1.422 & 1.324 \\
WMT24 en-es & 1 & 26.83 & 40 & 1.689 & 2.127 & 2.215 & 1.858 \\
WMT24 ja-zh & 0 & 40.13 & 40 & 1.563 & 1.563 & 1.561 & 1.361 \\
WMT24 ja-zh & 1 & 25.68 & 40 & 1.954 & 2.156 & 2.217 & 1.818 \\
\bottomrule
\end{tabular}
\caption{Mean-pooling discovery ratios from empty history on the same predefined subbatches at $q=.05$, $\alpha=.10$. The four policies use their original full definitions, including WMT24 priorities. Policy columns average $Q_{\rm online}/Q_{\rm joint}$ over the $n_+$ splits with positive matching minima; zero minima have undefined ratios. Mean $Q_{\rm joint}$ counts entry observations; policy columns are dimensionless mean splitwise ratios. Counts and certificate gaps concern this separate small-batch task, not the full-panel acquisition records.}
\label{tab:joint-discovery}
\end{table*}


\paragraph{Even-median validation.}
For pools of size two and four, exact rational scores on the quarter-step grid generate 3,690 two-row cases with two candidates, cutoff ties, endpoint and intermediate outcomes, empty and partial histories, and exact and positive-gap targets. A separate reconstruction enumerates all 38,887 eligible query subsets across these cases and agrees with the saved joint minima. These finite checks complement Corollary~\ref{cor:even-median}. The scripts \texttt{check\_even\_median.py} and \texttt{verify\_rational\_cases.py} in \path{review-revision/}\allowbreak\path{even-median-resource-20261007/} specify case generation and independent exhaustive reconstruction, respectively.

We also fix even pools of size six and eight on all six NLP panels, using the first $M$ evaluators in the existing order and all remaining evaluators as candidates. Six panels, two even pool sizes, 40 splits, four depths and two stopping targets give 3,840 full-batch tasks. Each task refits its consensus cutoff on the same reference outputs, retains right-ECDF normalization and $\alpha=.10$, and uses either the primary gap or exact stopping. Integer ECDF rank numerators make comparisons of the two central scores exact, including cutoff ties. Per-task records retain lower and feasible upper query bounds after a nominal 60-second outcome-MILP budget. Table~\ref{tab:even-median-panels} averages these bounds over 160 targets per row. These exposed-panel, offline results neither evaluate a new acquisition policy nor compare $M=6,7,8$ as the same selection task.
\begin{table*}[t]
\centering\small
\begin{tabular}{lrrlrrrr}
\toprule
Panel & $M$ & $T$ & Stop & Closed & Zero & Mean LB & Mean UB \\
\midrule
RewardBench & 6 & 8 & primary & 160 & 0 & 355.8 & 355.8 \\
RewardBench & 6 & 8 & exact & 160 & 0 & 434.2 & 434.2 \\
RewardBench & 8 & 6 & primary & 160 & 62 & 55.8 & 55.8 \\
RewardBench & 8 & 6 & exact & 160 & 47 & 78.5 & 78.5 \\
WMT23 (13 streams) & 6 & 18 & primary & 160 & 143 & 9.2 & 9.2 \\
WMT23 (13 streams) & 6 & 18 & exact & 160 & 123 & 23.0 & 23.0 \\
WMT23 (13 streams) & 8 & 16 & primary & 160 & 159 & 0.019 & 0.019 \\
WMT23 (13 streams) & 8 & 16 & exact & 160 & 155 & 1.6 & 1.6 \\
WMT23 (12 streams) & 6 & 18 & primary & 160 & 141 & 20.5 & 20.5 \\
WMT23 (12 streams) & 6 & 18 & exact & 160 & 126 & 41.7 & 41.7 \\
WMT23 (12 streams) & 8 & 16 & primary & 160 & 156 & 3.5 & 3.5 \\
WMT23 (12 streams) & 8 & 16 & exact & 160 & 150 & 8.0 & 8.0 \\
WMT24 en--de & 6 & 15 & primary & 160 & 115 & 298.8 & 298.8 \\
WMT24 en--de & 6 & 15 & exact & 160 & 92 & 522.3 & 522.3 \\
WMT24 en--de & 8 & 13 & primary & 160 & 157 & 4.0 & 4.0 \\
WMT24 en--de & 8 & 13 & exact & 160 & 151 & 17.4 & 17.4 \\
WMT24 en--es & 6 & 15 & primary & 160 & 113 & 241.4 & 241.4 \\
WMT24 en--es & 6 & 15 & exact & 160 & 86 & 420.6 & 420.6 \\
WMT24 en--es & 8 & 13 & primary & 160 & 158 & 4.3 & 4.3 \\
WMT24 en--es & 8 & 13 & exact & 160 & 155 & 8.5 & 8.5 \\
WMT24 ja--zh & 6 & 16 & primary & 160 & 93 & 447.0 & 447.0 \\
WMT24 ja--zh & 6 & 16 & exact & 160 & 82 & 628.4 & 628.4 \\
WMT24 ja--zh & 8 & 14 & primary & 160 & 147 & 48.5 & 48.5 \\
WMT24 ja--zh & 8 & 14 & exact & 160 & 140 & 71.2 & 71.2 \\
\bottomrule
\end{tabular}
\caption{Even-median joint evidence on complete test batches, averaging all four depths and 40 splits (160 targets) per row. Closed counts numerically closed minima; Zero counts zero-query certificates. Each target uses $\alpha=.10$ and either $\varepsilon=\lceil .02|S_0|\rceil$ (primary) or zero (exact). Mean lower and upper bounds include every target. The positive mean below 0.1 is shown to three decimals; an unresolved target is not replaced by a claimed minimum. Pools of size six and eight use the first members in the frozen evaluator order. Changing $M$ also changes the candidate family and cutoff, so these are separate certification tasks.}
\label{tab:even-median-panels}
\end{table*}


\subsection{Resource scaling and full-test mean evidence}
\label{app:joint-scale}

\paragraph{Matched-target resource comparison.}
Table~\ref{tab:controlled-resources} compares row-state DP, entry MILP and outcome-count MILP on the same 1,920 original full-mean targets. Methods run sequentially in fresh single-thread processes on the same workstation, once per target, with their order rotated deterministically across sorted targets. The nominal 60-second algorithm budget includes row-state or model construction; imports and input loading precede timing. Construction time is deducted before the MILP call receives the remaining budget. No external wall-clock deadline terminates the process. All 111 attempts exceeding 60 seconds remain unresolved; measured elapsed times, including the maximum 66.91 seconds, are retained. DP additionally permits at most 100,000 states and 10 million transitions. We record sampled whole-process RSS, final query bounds, model dimensions and status for all 5,760 attempts, including unresolved cases. The all-entry witness supplies the initial feasible bound; time to the first improving solver incumbent was not measured. Previously obtained numerical minima check every reported interval and new numerical closure in this resource study.

An interrupted attempt has no recoverable elapsed-time measurement and is excluded from completed-attempt timing summaries. The run ledger preserves the interruption and subsequent continuation. Pauses between attempts permitted separate validation and compilation; the timings do not establish performance on an otherwise idle machine. Complete task-level records include $m,T,G$, model sizes and resource-limit reasons.
\begin{table*}[t]
\centering\small
\begin{tabular}{llrrrrrr}
\toprule
Method & Stop & Closed & Mean s & Median s & Max s & RSS MiB & Gap \\
\midrule
Row DP & primary & 4 & 8.82 & 9.04 & 12.34 & 138.6 & 7371.6 \\
Row DP & exact & 82 & 6.83 & 7.20 & 10.57 & 137.7 & 7279.6 \\
Entry MILP & primary & 920 & 4.09 & 0.46 & 62.50 & 1116.9 & 169.2 \\
Entry MILP & exact & 933 & 2.26 & 0.07 & 66.91 & 1081.8 & 0.7 \\
Outcome MILP & primary & 931 & 2.21 & 0.05 & 60.12 & 618.9 & 132.7 \\
Outcome MILP & exact & 945 & 1.49 & 0.02 & 60.34 & 420.4 & 0.1 \\
\bottomrule
\end{tabular}
\caption{Sequential single-thread measurements on the same 1,920 full mean targets, with 960 targets per stopping criterion. RSS is the maximum sampled whole-process memory across attempts in the row; Gap averages final feasible query counts minus numerical lower bounds over all 960 attempts in each row, each retaining a feasible incumbent. Each attempt has a nominal 60-second algorithm budget, not an external wall-clock cutoff; DP also retains its 100,000-state and 10-million-transition caps. All 5,760 saved attempts with resource measurements enter timing and memory summaries, including attempts with unresolved targets; the interrupted attempt without recoverable timing is excluded. RSS is sampled whole-process memory, not solver-only allocation.}
\label{tab:controlled-resources}
\end{table*}


\paragraph{Resource tails and instance structure.}
Table~\ref{tab:resource-strata} reaggregates the same outcome-count attempts by $G$. Within $2000\le G<5000$, 205/218 primary attempts close, and the 95th percentile reaches 60.03 seconds. The three primary instances with $G\ge5000$ all close within 0.41 seconds. Larger group counts do not consistently require longer solution times on this grid. Panel, active count, budget and exclusion structure vary together; these are descriptive strata, not controlled interventions. Saved summaries include $m,T,G$, maximum formulation sizes, status counts and timing tails, including every unresolved attempt.
\begin{table*}[t]
\centering\small
\begin{tabular}{llrrrrrr}
\toprule
$G$ & Stop & Closed / attempts & $m$ range & Max vars & Median s & 95th s & Max s \\
\midrule
0--499 & exact & 173/174 & 53--276 & 1227 & 0.008 & 0.07 & 60.01 \\
0--499 & primary & 173/174 & 53--276 & 1233 & 0.014 & 0.10 & 60.01 \\
500--1999 & exact & 557/565 & 170--921 & 5181 & 0.018 & 1.63 & 60.12 \\
500--1999 & primary & 550/565 & 170--921 & 5183 & 0.042 & 1.80 & 60.11 \\
2000--4999 & exact & 212/218 & 706--1860 & 11116 & 0.082 & 15.32 & 60.34 \\
2000--4999 & primary & 205/218 & 706--1860 & 11116 & 0.179 & 60.03 & 60.12 \\
5000+ & exact & 3/3 & 1907--2025 & 12381 & 0.170 & 0.17 & 0.17 \\
5000+ & primary & 3/3 & 1907--2025 & 12381 & 0.360 & 0.40 & 0.40 \\
\bottomrule
\end{tabular}
\caption{Outcome-count MILP resource strata from the existing 1,920 full-mean attempts in Table~\ref{tab:controlled-resources}. $G$ counts row--outcome groups and $m$ active rows. All attempts, including unresolved ones, enter the time summaries; quantiles use linear interpolation. Closure uses the original nominal 60-second algorithm budget. These are posthoc descriptive strata: $G$, panel, active count and budget co-vary, so differences are not controlled scaling effects. Machine-readable summaries also stratify by $T$.}
\label{tab:resource-strata}
\end{table*}




The cached row-state implementation constructs each row table once, removes invariant query actions and prunes states whose query count cannot improve a feasible full-remaining-entry incumbent. We fixed six panels, splits 0--7, $q\in\{.05,.30\}$ and selected sizes 4, 8, 16, 32, 64 and full before examining these new runs. Empty and fixed quarter-histories with exact and primary gaps give 2,304 cases. All saved incumbents satisfy their integer certificate conditions. Independent subset MILP checks 396 predefined cases: every case of size at most four and split-zero empty exact eight-row cases. Larger completed DP minima are not all independently subset-enumerated. Table~\ref{tab:joint-dp-resources} reports completed pairs and resource-limited attempts, including the limits.
\begin{table*}[t]
\centering\small
\begin{tabular}{lrrrr}
\toprule
Selected size & Both closed & Median (s) & 90th pct. (s) & 90th MiB \\
\midrule
4 & 384/384 & 0.010 & 0.276 & 207.5 \\
8 & 318/384 & 0.273 & 4.157 & 195.5 \\
16 & 202/384 & 3.287 & 5.976 & 174.4 \\
32 & 60/384 & 4.990 & 7.538 & 155.5 \\
64 & 62/384 & 5.733 & 8.384 & 158.0 \\
full & 16/384 & 6.587 & 12.915 & 187.5 \\
\bottomrule
\end{tabular}
\caption{Cached row-state DP resource study: six panels, splits 0--7, $q\in\{.05,.30\}$, empty/quarter histories and exact/primary gaps give 384 instances per size. Subbatches reset $N,k,b_H$; full preserves the original test task. Both closed counts completed joint and separate minima. Limits are 100,000 reachable states, ten million transitions and 60 seconds shared by the two optimizations. Resource-limited rows retain valid witnesses without claiming minima. Times include unresolved attempts; RSS is sampled every 10 ms and includes Python, inputs and caches. Six concurrent workers share a workstation; these are not controlled speedup measurements.}
\label{tab:joint-dp-resources}
\end{table*}


For the matching 576 empty-history exact tasks, entry and outcome-count MILPs each numerically close 573 joint/separate pairs within the initial 60-second budget (Table~\ref{tab:joint-compact-resources}). The two formulations agree on every pair both numerically close; they also agree with all 264 completed joint DP minima; 263 of these tasks also have a completed separate optimum and enter the paired summary. Both additionally match all 7,680 earlier small-batch minima. Append-only 600-second retries closed two more pairs, and exhaustive fixed-threshold decomposition closed the final joint objective. The merged records thus contain 576 closed pairs; Table~\ref{tab:joint-compact-resources} keeps the initial-run outcomes. The last 64-row ja-zh task has joint minimum 555 and separate minimum 569. Its variable-threshold run and a diagnostic 600-second retry had retained bounds 554--555. Offline screening left exactly $t=8,\ldots,16$; each of the nine fixed-threshold programs attained matching numerical primal and dual objectives; the minimum across their nine optimal values was 555. The resource-limited records are retained.
\begin{table*}[t]
\centering\small
\begin{tabular}{lrrrrr}
\toprule
Size & Row DP & Entry MILP & Outcome MILP & Outcome vars & Median s \\
\midrule
4 & 96/96 & 96/96 & 96/96 & 45 & 0.017 \\
8 & 81/96 & 96/96 & 96/96 & 76 & 0.033 \\
16 & 48/96 & 96/96 & 96/96 & 132 & 0.059 \\
32 & 14/96 & 95/96 & 95/96 & 281 & 0.107 \\
64 & 16/96 & 95/96 & 95/96 & 497 & 0.154 \\
full & 8/96 & 95/96 & 95/96 & 11116 & 0.128 \\
\bottomrule
\end{tabular}
\caption{Matching empty-history exact mean tasks at $\alpha=.10$: 96 cases per size. Counts require both joint and separate optima within the initial resource budget. Each MILP backend receives 60 seconds shared between these optimizations; DP additionally has state/transition limits. Outcome MILP groups equal integer flip contributions (Proposition~\ref{prop:compact-joint}); its median time includes unresolved attempts. Outcome MILP numerically closes 573/576 joint-and-separate pairs in the initial run. Append-only retries numerically close two more pairs; exhaustive feasible-threshold decomposition numerically closes the last joint minimum, completing all 576 pairs. Initial censoring remains in this table. Outcome vars counts the largest outcome-count formulation represented in each row, not a bound on solve time.}
\label{tab:joint-compact-resources}
\end{table*}


The 96 full instances preserve original test $N$, nominal $k$, maps, cutoff, candidate family and budget. Their active-row counts range from 53 to 1,860; the largest outcome formulation has 11,116 variables and 13,448 constraints. Twelve additional split-zero full-instance probes sample process RSS every 10 ms: peaks range from 88.1 to 115.4 MiB, including Python, inputs, solver and resident caches. Concurrent load and sampled allocation preclude a controlled backend speedup or exact solver-memory claim.

To vary dimensions beyond the cached panels, we fix $N\in\{2048,8192\}$, $M\in\{7,15\}$ and $T\in\{17,65\}$ before running the complete outcome-count joint model. Base and candidate scores are independently drawn from $U(0,1)$ in a declared PCG64 order, with seed $731191+N+100M+T$. Mean pooling uses a construction-fixed cutoff $.5$, $q=.30$, $\alpha=.10$, empty history and exact stopping. Every selected row remains in the problem. Each target runs in its own process, with one HiGHS thread, zero relative MIP gap and a nominal 60-second budget including model work. Protocol and version records accompany \nolinkurl{review-revision/discovery-bound-implementation/scale/}.

Four of eight targets numerically close; the other four retain their bounds in Table~\ref{tab:controlled-resource-grid}. Independent explicit replacements reconstruct all costs and validate every feasible witness. The largest model has 39,097 variables and 52,154 constraints; measured process peaks span 156.0--635.1 MiB. At the largest grid point, the true optimum retains every selected output, making exclusion automatic; this instance closes quickly. Thus dimensions alone do not determine difficulty. These constructed-score, exact-stop probes measure computational scope, not NLP generalization, evaluator inference time or a uniform scalability guarantee.

\begin{table*}[t]
\centering\small
\begin{tabular}{rrrrrrrrr}
\toprule
$N$ & $M$ & $T$ & $m$ & $G$ & Joint bounds & Seconds & Peak MiB & Closed\\
\midrule
2048 & 7 & 17 & 790 & 3182 & 12778--12778 & 24.24 & 158.2 & yes\\
2048 & 7 & 65 & 759 & 3883 & 46602--46673 & 60.18 & 305.8 & no\\
2048 & 15 & 17 & 634 & 3483 & 10648--10648 & 21.08 & 226.6 & yes\\
2048 & 15 & 65 & 619 & 5122 & 39767--39767 & 22.92 & 311.6 & yes\\
8192 & 7 & 17 & 3178 & 13219 & 0--54026 & 61.56 & 309.5 & no\\
8192 & 7 & 65 & 3142 & 16079 & 192553--194598 & 63.20 & 635.1 & no\\
8192 & 15 & 17 & 2529 & 13747 & 42795--42830 & 61.67 & 595.2 & no\\
8192 & 15 & 65 & 2419 & 19294 & 132536--132536 & 2.01 & 156.0 & yes\\
\bottomrule
\end{tabular}
\caption{Controlled mean-score resource grid: all eight predefined exact joint targets. Bounds are a numerical lower bound and an independently checked feasible entry count; zero can be an uninformative lower bound. Seconds include model construction, optimization and witness validation, which can exceed the nominal 60-second budget. Peak process memory includes inputs and Python. One instance per point and concurrent background work preclude a calibrated runtime scaling law.}
\label{tab:controlled-resource-grid}
\end{table*}

Returning to the 96 full-batch NLP instances, all matching exact mean minima are positive. The four original exact-stop policies give 384 same-instance discovery counts, all at least the matching minimum (Table~\ref{tab:joint-full-discovery}). In 4/96 instances, $Q_{\rm joint}$ is below the additive $Q_{\rm split}$, by 2,645 of 435,876 separately counted entries (0.607\%). Appendix~\ref{app:deduplicated} separates double counting from improvements beyond optimal deduplication. A compact exact formulation therefore extends the evidence analysis to full tasks; these records do not establish broadly large gains from sharing or an optimal online policy. These results are descriptive: they use eight overlapping splits of panels already examined during method development.
\begin{table*}[t]
\centering\small
\begin{tabular}{lrrrrrr}
\toprule
Panel & $q$ & Mean $Q_{\rm joint}$ & Fixed & Random & Width & Reference \\
\midrule
RewardBench 2 & 0.05 & 668.4 & 1.141 & 1.141 & 1.145 & 1.095 \\
RewardBench 2 & 0.30 & 1560.6 & 1.336 & 1.334 & 1.204 & 1.359 \\
WMT23, 13 streams & 0.05 & 1689.9 & 1.174 & 1.173 & 1.176 & 1.102 \\
WMT23, 13 streams & 0.30 & 1851.8 & 1.504 & 1.516 & 1.289 & 1.631 \\
WMT23, 12 streams & 0.05 & 1735.3 & 1.125 & 1.129 & 1.128 & 1.082 \\
WMT23, 12 streams & 0.30 & 2370.9 & 1.476 & 1.476 & 1.295 & 1.524 \\
WMT24 en-de & 0.05 & 4698.9 & 1.132 & 1.132 & 1.133 & 1.065 \\
WMT24 en-de & 0.30 & 9857.6 & 1.295 & 1.295 & 1.212 & 1.283 \\
WMT24 en-es & 0.05 & 3920.0 & 1.116 & 1.117 & 1.117 & 1.060 \\
WMT24 en-es & 0.30 & 7553.6 & 1.322 & 1.322 & 1.222 & 1.345 \\
WMT24 ja-zh & 0.05 & 4289.8 & 1.142 & 1.150 & 1.147 & 1.090 \\
WMT24 ja-zh & 0.30 & 13957.3 & 1.271 & 1.271 & 1.219 & 1.174 \\
\bottomrule
\end{tabular}
\caption{Full-test mean entry minima and online discovery from empty history at exact stopping, $\alpha=.10$, splits 0--7. Every row averages eight original full tasks with positive minima; ratio columns average the splitwise $Q_{\rm online}/Q_{\rm joint}$, not a ratio of means. All four original query paths are unchanged. Across 96 minima and 384 paths, every online count is at least its matching minimum. Mean $Q_{\rm joint}$ counts entry observations; policy columns are dimensionless. Joint minima are below the additive sum in 4/96 tasks, by 2,645 of 435,876 separately counted entries (0.607\%). The best deduplicated comparison is in Table~\ref{tab:deduplicated-summary}. These denominators are not used for positive-tolerance stops.}
\label{tab:joint-full-discovery}
\end{table*}



\FloatBarrier
\subsection{Joint evidence beyond deduplicating separate optima}
\label{app:deduplicated}
For each safe target $s$ in Theorem~\ref{thm:joint}, let $\mathcal A_s$ be the set of all minimum-size additional safety witnesses, and let $\mathcal B_{s+\varepsilon+1}$ be the set of all minimum-size additional exclusion witnesses. Exclusion is automatic beyond the active-row count, with the empty witness. Define
\[
 Q_{\cup,\rm sep}=\min_s\min_{A\in\mathcal A_s,\,B\in\mathcal B_{s+\varepsilon+1}}|A\cup B|.
\]
The history is fixed and not charged again. This comparison ranges over all tied component optima, not one pair returned by a solver.

\begin{proposition}[Deduplicated separate evidence]
\label{prop:deduplicated}
For unit-cost entry certification under the same instance, history and tolerance,
\[
 Q_{\rm joint}\le Q_{\cup,\rm sep}\le Q_{\rm split}\le 2Q_{\rm joint}.
\]
Within each active row, group unqueried entries with equal true contribution. If two witnesses choose $a_g,b_g$ entries from a group, their minimum possible union size is $\max(a_g,b_g)$. Minimizing $\sum_g\max(a_g,b_g)$ at separately optimal component cardinalities exactly characterizes $Q_{\cup,\rm sep}$.
\end{proposition}
\begin{proof}
The union preserves safety and exclusion by monotonicity, proving the first inequality. Choosing the target and witnesses attaining the additive bound gives the second. For the last inequality, take a minimum joint witness $X$ and its safe target. The same $X$ is feasible for each component, so each component minimum is at most $|X|$. Their sum at this target is at most $2|X|$, and minimizing over targets cannot increase it. This also covers $Q_{\rm joint}=0$, automatic exclusion and fixed prior observations. Within a contribution group, identities do not affect either row bound. Nested choices attain union size $\max(a_g,b_g)$; no union can contain fewer than either component. Summing and optimizing over targets proves the characterization. This exchange argument uses unit-cost entry access; column identities need not be interchangeable.
\end{proof}


\paragraph{What a query history can establish about joint advantage.}
Instance-wise savings concern the realized hidden scores. A stronger question asks whether a history alone establishes strict savings for every consistent completion. Here each $Q$ measures the cost of additional reveal actions after this history; past observations are free. Both component minima use the same action family.

\begin{proposition}[No uniform residual advantage from a history]
\label{prop:residual-no-uniform}
Fix any consistent history $\mathcal O$ under the completion model of Section~\ref{sec:partial} with no cross-entry constraints, an integer count tolerance $\varepsilon\ge0$, and either entry or whole-column queries with fixed nonnegative additive action costs. There exists a completion consistent with $\mathcal O$ for which
\[
 Q_{\rm joint}(\mathcal O)=Q_{\cup,\rm sep}(\mathcal O)
 =Q_{\rm split}(\mathcal O).
\]
Consequently, no history can certify a strictly positive residual joint advantage uniformly over all its consistent completions.
\end{proposition}
\begin{proof}
Complete every unobserved selected-row score at its zero-flip endpoint. This is admissible because no cross-entry constraints exclude this simultaneous assignment. The completed row costs equal their current lower costs $\ell^0$. Let $L_0=V(\ell^0)-m_0$ be the optimal active count in this completion. Every additional query reveals zero contributions, so the optimistic active count remains $L_0$. A gap certificate must therefore establish safety for at least $s_0=\max(0,L_0-\varepsilon)$ active rows. Conversely, any such safety witness is a joint certificate: the history already excludes every count above $L_0$, and $s_0+\varepsilon\ge L_0$. The minimum safety witness for $s_0$ attains the joint minimum. Pairing it with the empty additional exclusion witness gives a separate union and additive cost of the same value. The inequalities in Proposition~\ref{prop:deduplicated}, or the fixed-action version below, give equality throughout. This also covers an already satisfied gap and zero-cost actions.
\end{proof}

The proposition concerns residual minima, not the original empty-history minima of an instance eventually revealed by queries. It does not rule out predicting realized savings or designing efficient acquisition. A uniformly valid positive residual test requires an additional constraint excluding the zero-contribution completion. A statistical prediction and a deterministic certificate then have different roles.

\paragraph{Discovery at fixed pool size and risk budget.}
The separation between a short witness and finding it persists with the numerical parameters of the primary study. It concerns exact stopping, not the primary positive tolerance.
\begin{proposition}[Exact discovery with fixed parameters]
\label{prop:fixed-discovery}
For every integer $g\ge0$, there is a family of odd-median instances with $M=T=7$, $q=1/20$, $\alpha=1/10$, $k=20g+17$, $N=20k$ and $|S_0|=k$, each with a minimum exact entry certificate of size one. For every deterministic policy correct under the full completion model, some instance in this family requires at least $49g+42$ queries. For every randomized always-valid policy, some fixed instance in the family has expected query cost at least $(49g+43)/2$. Thus no batch-size-independent competitive factor holds for exact discovery even with these parameters fixed.
\end{proposition}
\begin{proof}
Set $m=7g+6$. The $m$ selected boundary rows have four base scores $.75$ and three $.25$. A further $k-m$ selected rows have all base scores $.75$, and the remaining $N-k$ rows have all base scores $.25$. At cutoff $t_0=3/5$, only the boundary rows can change membership. Candidate scores are $.25$ everywhere except one active candidate--output entry with score $.75$. Each active opposition contributes four flips; the unique agreement contributes zero.

These normalized values and cutoff can be realized by the reference procedure. Use 20 reference outputs indexed $a=0,\ldots,19$. For evaluator $j=0,\ldots,13$, set its raw reference scores to $1+((a+3j)\bmod20)$. Each evaluator's right ECDF maps raw scores 5 and 15 to $.25$ and $.75$. The first seven reference columns have maximum median consensus $3/5$. Since $k_{\rm ref}=1$, the fitted cutoff is $3/5$. The reference panel is fixed across all instances and supplies no information about the agreement's location.

The flip budget and opposition budget are
\[
 \begin{aligned}
 b_H&=\lfloor2k\alpha MT\rfloor=196g+166,\\
 b&=\lfloor b_H/4\rfloor=49g+41=mT-1.
 \end{aligned}
\]
All $k$ selected outputs are truly feasible. Initially, however, the upper costs allow only $k-1$ and the lower costs allow $k$. Observing the unique agreement reduces the total upper cost from $4mT$ to $4(mT-1)\le b_H$, certifying all $k$. Hence the exact offline minimum is one.

Before this entry is observed, the all-opposition completion remains consistent and permits only $k-1$ selected outputs, while a completion with an agreement at any unseen active entry permits $k$. These two completions prevent exact stopping. The policy must be correct on the full completion set; it has no guarantee of exactly one agreement in the hidden matrix. Place the agreement last on a deterministic policy's all-opposition path to force $mT$ active queries. For a uniformly chosen agreement location, any fixed random seed induces an order with expected discovery position at least $(mT+1)/2$. Averaging first over locations and then over seeds gives a fixed location with at least this expected cost. Queries outside the active entries can only add cost.
\end{proof}

The initial gap is one. Therefore every $\varepsilon\ge1$, including the primary $\lceil .02|S_0|\rceil$ criterion, allows stopping without queries on this family. The lower bound neither explains the measured positive-tolerance overhead nor establishes how often hard instances occur in NLP. It shows why a short exact offline certificate alone does not imply an efficient exact online policy. The reference construction, explicit replacement checks for all 273 agreement positions at $g=0,1,2$, and exact-arithmetic checks of 1,003 values of $g$ are provided in \nolinkurl{review-revision/discovery-bound-implementation/}; finite checks supplement the proof.

\paragraph{Discovery at the primary positive tolerance.}
A different construction keeps the primary gap nontrivial as the batch grows and proves Proposition~\ref{prop:positive-discovery}.
\begin{proof}
Take $M=T=7$, $\alpha=1/10$, $|S_0|=m$, $\varepsilon=e=\lceil m/50\rceil$, $s=m-e-1$, $k=5s+4$ and $N=20k$, so $q=k/N=1/20$. All instances use the same reference-fitted maps and cutoff. Use 20 reference outputs and raw scores $1+((a+3j)\bmod20)$ for output $a=0,\ldots,19$ and evaluator $j=0,\ldots,13$. Each right ECDF maps raw rank $v$ to $v/20$. Since $k_{\rm ref}=1$, the fitted mean cutoff is the largest reference consensus, $t_0=83/140$. Give every selected test row the base scores
\[
 (1,3/20,3/10,9/20,3/5,3/4,9/10),
\]
whose mean equals $t_0$. Every base score is positive, so a candidate score zero causes seven flips, whereas score one causes none. Set all other base rows to zero; they remain unselected even after replacing one base score by one, since $1/7<t_0$. These test scores are realized by the same reference ECDFs, using raw ranks or zero.

All candidate scores on selected rows are zero except two entries in one hidden row, in fixed candidate columns, whose scores are one. Thus every selected row has cap seven and upper cost 49 initially. The hidden row has true cost 35; every other selected row has true cost 49. The common reference data reveal neither the hidden row nor its test scores. The budget is
\[
 b_H=\lfloor 2k\alpha MT\rfloor=49s+39.
\]
The initial safe and optimistic counts are $s$ and $m$, with gap $e+1$. Revealing both exceptional entries lowers one upper cost to 35. The safe count becomes $s+1$, the optimistic count remains $m$, and the gap becomes $e$. No single entry suffices. A zero-flip reply lowers one upper cost only to 42, leaving the safe count $s$ because $42>39$; a seven-flip reply leaves all upper costs unchanged. In either case the total lower cost is at most seven and all $m$ rows remain optimistically feasible. Hence the minimum is exactly two.

Consider any history of at most $m-1$ queries with only seven-flip replies on selected rows. Its lower-cost sum is at most $7(m-1)\le b_H$, while its upper costs remain 49, so the gap stays $e+1$. Here $e\le m/2$ for $m\ge2$, which verifies the displayed budget inequality. At least one selected row has not been queried. Placing the two exceptional entries in this row gives a family member consistent with the history. A deterministic policy therefore cannot finish before $m$ queries in the worst case. Queries on invariant rows cannot shorten this history.

For randomized policies, choose the hidden row uniformly among the $m$ selected rows, independently of the policy randomness. After $t<m$ queries, at most $t$ rows have been touched along the all-seven-flip path. With probability at least $1-t/m$, the hidden row is untouched and the policy cannot stop. If $Q$ is its query count, then
\[
 \begin{aligned}
 \mathbb E Q&\ge\sum_{t=0}^{m-1}\Pr(Q>t)\\
 &\ge\sum_{t=0}^{m-1}(1-t/m)=(m+1)/2.
 \end{aligned}
\]
Averaging over rows gives a fixed instance with at least this expected cost. Policies must be valid on the full completion set; the exceptional-row pattern is not supplied as a promise. Since $N=20(5s+4)$ grows linearly with $m$, the cost ratio to the two-entry minimum is unbounded in batch size.
\end{proof}

The realized baseline-selected fraction is $|S_0|/N=m/N$. It establishes a worst-case discovery obstruction, not the frequency or cause of overhead on NLP panels. Exact score reconstruction checks the reference cutoff and replacement outcomes. Independent enumeration checks all 16,384 and 2,097,152 entry subsets for $m=2,3$; further finite checks cover 2,002 parameter values and 63,466 adversarial histories. The check script is \nolinkurl{review-revision/stable-batch-implementation/check_positive_discovery.py}.

\paragraph{Fixed query actions.}
For a fixed family of reveal actions with nonnegative additive costs, define each component optimum using the same actions, history and safe-count target. Define the best separate union by minimizing its action cost over all tied component optima and admissible targets. Then
\[
 Q_{\rm joint}\le Q_{\cup,\rm sep}\le Q_{\rm split}\le2Q_{\rm joint}.
\]
A union preserves both component certificates and costs no more than their separately charged sum. A minimum joint action set is feasible for each component at its attained safe target, so each component minimum costs at most the joint set. These two observations prove the inequalities. The bound applies to fixed-price column actions; the within-row exchange formula for entries does not apply to columns. History-dependent batching and loading costs require a different cost model.

Strict column advantage is possible even though the entry tight family does not establish column tightness. Take mean pooling with $M=T=3$, $N=2$, $k=1$, cutoff $1/2$ and flip budget five, corresponding to $\alpha=5/18$. The two base-score rows are $(2/5,3/5,4/5)$ and $(1/10,7/10,9/10)$. Candidate-score rows $(2/5,2/5,2/5)$ and $(2/5,3/5,2/5)$ give caps $(3,2)$ and flip contributions
\[
 \begin{pmatrix}1&1&1\\2&1&2\end{pmatrix}.
\]
Exactly one row can be retained. Column two alone is the unique one-column safety witness. Columns one and three form the unique minimum exclusion witness: their lower costs sum to six, while either other pair sums to five. That pair also certifies row one safe, with upper cost five. Thus $Q_{\rm joint}=2$ and $Q_{\cup,\rm sep}=Q_{\rm split}=3$ columns. This realizes a ratio of $3/2$, without establishing a tight factor-two column separation.


\begin{corollary}[Intermediate contributions bound joint advantage]
\label{cor:overlap-budget}
Under Theorem~\ref{thm:joint}, let $\mathcal X^*$ be the set of all minimum additional joint witnesses for the fixed history and tolerance. For $X\in\mathcal X^*$, define
\[
 \begin{aligned}
 I(X)&=\{(i,r)\in X:0<f_{ir}<p_i\},\\
 d_*&=\min_{X\in\mathcal X^*}|I(X)|.
 \end{aligned}
\]
Then
\[
 Q_{\cup,\rm sep}\le Q_{\rm split}\le Q_{\rm joint}+d_*.
\]
In particular, an optimal witness with no intermediate contributions implies $Q_{\cup,\rm sep}=Q_{\rm split}=Q_{\rm joint}$. This includes endpoint-only unobserved contributions, without requiring equal row caps.
\end{corollary}
\begin{proof}
Take any $X\in\mathcal X^*$ and its certified safe target $s$. Retain only entries with $p_i-f_{ir}>0$ to form $A_X$. Removing the other entries leaves every upper cost unchanged, so $A_X$ still certifies safety at $s$. Retain only entries with $f_{ir}>0$ to form $B_X$. Every lower cost is unchanged, so $B_X$ still certifies exclusion at $s+\varepsilon+1$. When exclusion is automatic, the empty witness also suffices. Previously observed entries are held fixed throughout.

The component minima at this target are at most $|A_X|$ and $|B_X|$. Their sum is at most
\[
 |A_X\cup B_X|+|A_X\cap B_X|
 \le |X|+|I(X)|.
\]
Minimizing over safe targets and then over $X\in\mathcal X^*$ proves the bound. The equality conclusion follows from Proposition~\ref{prop:deduplicated}.
\end{proof}
Strict joint advantage requires $d_*>0$, but $d_*>0$ does not imply strict advantage. Corollary~\ref{cor:overlap-budget} applies to arbitrary row counts, prior observations and integer count-gap tolerances, but does not characterize all instances with strict joint gain. For a particular optimal witness $X$, $|I(X)|$ is a valid, possibly looser upper bound on the gain; obtaining $d_*$ requires minimizing over all tied joint optima.

Direct query-subset enumeration checks 14,626 empty-history instances with caps one or two and at most three rows and two candidates, plus 1,200 seeded instances with arbitrary histories, zero caps and up to four rows. Every count tolerance from zero to the row count is checked. Separately optimal witnesses include every tie.

\paragraph{Minimizing intermediate evidence.}
We compute $d_*$ on the existing 1,920 full mean targets by fixing the previously numerically closed cardinality $Q_{\rm joint}$ and minimizing $\sum_{i,v}\one\{0<v<p_i\}a_{iv}$ in the outcome-count formulation, subject to $\sum_{i,v}a_{iv}=Q_{\rm joint}$. Here $a_{iv}$ is the number of queried entries from outcome group $(i,v)$. Every admissible safe-count target remains eligible; within-row groups preserve all contribution multiplicities. This is a secondary objective over the full joint-optimal set, not a count from one saved witness. When the true optimal active count is $m$, exclusion is automatic and the smallest admissible safe count suffices by monotonicity.

Exact full retention has a simpler integer solution. Let $j=Q_{\rm joint}$ and let $E(a)$ and $I(a)$ sum the $a$ largest reductions $p_i-f_{ir}$ among unobserved endpoint and intermediate cells, respectively. Let $D=\sum_i u_i^{(0)}-b_H$. Then
\[
 d_*=\min\{d:E(j-d)+I(d)\ge D\},
\]
where the two class sizes restrict admissible $d$. Exact retention of every active row makes the sum of upper costs the sole remaining constraint. For a fixed $d$, the two sorted prefixes maximize its reduction among all $j$-entry sets with $d$ intermediate cells; therefore the first feasible $d$ is minimal. Selecting these prefixes reconstructs a witness. This argument fixes the previously numerically closed cardinality $j$ and minimizes intermediate entries among $j$-entry witnesses; it does not prove $j$ minimal among certificate cardinalities.

For the remaining cases we retain all numerical bounds and integer-check each reconstructed query set. Three append-only rounds use nominal wall-time budgets of 5, 60 and 120 seconds per unresolved target. The remaining budget is computed before each target model is assembled; model assembly and solver termination can make elapsed time exceed this budget. Maximum recorded elapsed times are 8.06, 61.82 and 121.20 seconds, respectively. The secondary runs explicitly request zero relative gap and $10^{-7}$ primal, dual and MIP feasibility tolerances. Independent query-subset enumeration checks 240 seeded cases with histories, varying caps and tolerances, plus 80 cases restricted to nonzero secondary objectives and non-full retention. Both entry and grouped formulations match all 320 cases (640 comparisons). Full-grid checks validate each saved witness; numerical lower bounds are not rational proof certificates.

Table~\ref{tab:overlap-diagnostic} reports \dStarClosed{} closed secondary objectives and \dStarUnclosed{} remaining intervals. Across both stops, \dStarZero{} targets have $d_*=0$, compared with 457 zero counts in the original saved witnesses. All have no strict gain. At least \dStarPositiveNoGain{} no-gain targets have $d_*>0$. Although \dStarUnclosed{} objective values retain intervals, every target is classified as zero or positive. Every known union gain respects the bound, and \dStarTighter{} new witnesses improve the original intermediate-count upper bound. Minimizing over ties certifies more zero values, but these do not add empirical equality coverage beyond the simpler full-retention condition: all 1,024 zero cases lie among its 1,806 targets (Table~\ref{tab:explanatory-overlap}). A positive value remains insufficient for advantage. The 89 strict gains, including 63 on ja-zh, retain their original values. These overlapping-split summaries neither estimate population frequencies nor identify a causal effect.
\begin{table*}[t]
\centering\small
\begin{tabular}{lrrrr}
\toprule
Panel & $d_*=0$ & $d_*>0$ & Unresolved sign & Closed objective \\
\midrule
RewardBench 2 & 13 / 4 & 147 / 156 & 0 / 0 & 160 / 160 \\
WMT23, 13 & 132 / 108 & 28 / 52 & 0 / 0 & 159 / 160 \\
WMT23, 12 & 98 / 65 & 62 / 95 & 0 / 0 & 159 / 160 \\
WMT24 en-de & 147 / 121 & 13 / 39 & 0 / 0 & 160 / 160 \\
WMT24 en-es & 152 / 132 & 8 / 28 & 0 / 0 & 160 / 160 \\
WMT24 ja-zh & 28 / 24 & 132 / 136 & 0 / 0 & 139 / 142 \\
\bottomrule
\end{tabular}
\caption{Secondary minimization of intermediate entries over all minimum joint witnesses in the full mean grid. Cells give primary / exact counts out of 160 targets (four depths and 40 overlapping splits) per panel, using full-reference right ECDF and $\alpha=.10$. Primary uses $\varepsilon=\lceil.02|S_0|\rceil$; exact uses zero. The first three columns classify targets using integer-checked witnesses and bounds from Corollary~\ref{cor:overlap-budget} or numerical optimization. Solver-based classifications depend on the validity of the numerical bounds. Closed objective counts equality of lower and upper bounds, conditional on the previously numerically closed joint minimum. The remaining objectives retain intervals; a positive sign does not imply an exact objective value.}
\label{tab:overlap-diagnostic}
\end{table*}


\paragraph{Matched multirow counterexample.}
Intermediate counts also fail to determine the gain when basic instance summaries are fixed. Consider the two contribution matrices
\[
 F^{A}=\begin{pmatrix}0&0&0\\0&1&3\\1&2&2\end{pmatrix},\qquad
 F^{B}=\begin{pmatrix}0&0&0\\0&2&2\\1&1&3\end{pmatrix}.
\]
Both have three rows, three candidates, caps $p_i=3$, row totals $(0,4,5)$, empty history, budget $b_H=8$ and exact stopping. Their optimal retained count is two; their safety and exclusion minima are four and five, respectively. Both have $Q_{\rm joint}=7$ and $d_*=4$, with four intermediate cells in the complete matrix. Yet $Q_{\cup,\rm sep}$ is eight for $F^A$ and seven for $F^B$. Thus identical row totals, component minima, joint minimum and intermediate counts do not determine the best separate union. The difference lies in how minimum safety and exclusion witnesses can overlap.

These are realizable mean scores. Use base scores $(.50,.51,.52)$ on each row and $t_0=.5$. Candidate scores $1,.485,.475,0$ realize contributions $0,1,2,3$, respectively. Set $M=T=N=3$, $q=1/3$, $k=1$ and $\alpha=4/9$. Exact rational replacement calculations recover both matrices. Independent enumeration of all $2^9$ query subsets per matrix verifies every minimum and all component ties. The search uses three rows, three candidates and common cap three. It enumerates nondecreasing row triples from $\{0,1,2,3\}$, then unordered multisets of three such rows, and every integer budget with true retained count one or two. This removes within-row entry permutations and row permutations from the entry-query search, yielding 16,720 instances. The pair refutes sufficiency of these summaries, not a population hypothesis or every multirow criterion.

\paragraph{Deduplication and joint optimization on NLP tasks.} The two-row example in Appendix~\ref{app:joint-evidence} has $Q_{\rm split}=3$ but $Q_{\cup,\rm sep}=Q_{\rm joint}=2$, so its apparent saving is entirely double counting. We therefore separate $Q_{\rm split}-Q_{\cup,\rm sep}$ from $Q_{\cup,\rm sep}-Q_{\rm joint}$ on the existing 3,840 mean small-batch tasks and 96 complete exact tasks, retaining their original histories, budgets and targets. In 3,544 tasks, equality follows from the previously numerically closed joint and additive optima together with the theoretical bound sandwich. For the 392 remaining tasks, an outcome-count program fixes both component cardinalities before minimizing their union. A separate binary-entry model enumerates row subsets on all 388 small tasks; entry-variable threshold programs check the four full tasks. An exhaustive query-subset implementation also checks 160 fixed-seed numerical cases with at most nine entries, including histories and positive tolerances.

Table~\ref{tab:deduplicated-full} identifies the four full-batch cases with gains beyond optimal deduplication.
\begin{table*}[t]
\centering\small
\begin{tabular}{llrrrrr}
\toprule
Task & History & $\varepsilon$ & $n$ & $Q_{\rm joint}<Q_{\cup,\rm sep}$ & Double count & Joint reduction \\
\midrule
Full & empty & 0 & 96 & 4 & 2573 & 72 \\
Small & empty & 0 & 960 & 84 & 190 & 116 \\
Small & empty & 1 & 960 & 43 & 37 & 71 \\
Small & quarter & 0 & 960 & 38 & 106 & 51 \\
Small & quarter & 1 & 960 & 9 & 17 & 18 \\
\bottomrule
\end{tabular}
\caption{Mean entry evidence: separating deduplication from joint optimization. The strict-gain column counts instances. Double count sums $Q_{\rm split}-Q_{\cup,\rm sep}$; Joint reduction sums $Q_{\cup,\rm sep}-Q_{\rm joint}$ in additional entry observations. Each small row covers six panels, four depths and 40 splits, with the original reset subbatch budget. Full tasks preserve the original test batch and budget at two depths and eight splits per panel. All component ties are optimized; equal joint/additive cases use the proved bound sandwich. Histories are not charged again.}
\label{tab:deduplicated-summary}
\end{table*}
\begin{table*}[t]
\centering\small
\begin{tabular}{lrrrrr}
\toprule
Panel & Split & $q$ & $Q_{\rm joint}$ & $Q_{\cup,\rm sep}$ & $Q_{\rm split}$ \\
\midrule
WMT23, 12 streams & 6 & 0.05 & 2273 & 2293 & 2521 \\
WMT23, 12 streams & 7 & 0.05 & 2347 & 2361 & 2602 \\
WMT24 ja-zh & 1 & 0.05 & 5149 & 5167 & 6023 \\
WMT24 ja-zh & 4 & 0.05 & 10103 & 10123 & 11371 \\
\bottomrule
\end{tabular}
\caption{The four full mean instances with joint minima below optimally deduplicated separate witnesses. All use empty history, exact stopping and $\alpha=.10$. The other 92 full tasks have equality throughout. Both outcome-count and binary-entry programs numerically close these union optima; original time-limited attempts remain in the records. All costs count entry observations, not scorer calls or execution time.}
\label{tab:deduplicated-full}
\end{table*}

\paragraph{Small gains beyond deduplication.}
Across the 96 full-batch tasks, $Q_{\rm split}-Q_{\rm joint}$ totals 2,645 entries. Optimal deduplication accounts for 2,573, and joint optimization for the remaining 72 across four instances. These 72 entries equal 0.0166\% of the aggregate $Q_{\cup,\rm sep}$ of 433,303. Thus joint selection can improve on the best union of separately optimal proofs, but the full-batch improvements here are small. The comparison concerns offline evidence, not a gain achieved by an online policy. Small-task results remain a separate target with reset batch size and budget. Saved records retain query sets, component minima, all-target checks and solver states. One primary outcome run reported infeasibility despite an independently feasible witness; disabling presolve closed the same model. One full entry check reached its initial time limit and was resolved by an append-only retry. Failed and unresolved primary attempts are preserved.


\begin{proposition}[Strict gain beyond optimal deduplication]
\label{prop:strict-joint}
For every even $T\ge2$, a realizable $M=T+1$ mean-pooling instance with two selected outputs, empty history and exact stopping satisfies
\[
 Q_{\rm joint}=T+1,\qquad Q_{\cup,\rm sep}=Q_{\rm split}=2T.
\]
Thus the factor-two upper bound in Proposition~\ref{prop:deduplicated} is asymptotically tight even after optimal deduplication. It concerns offline unit-cost entry certificates, not online query competitiveness.
\end{proposition}
\begin{proof}
Use endpoint caps $(T+1,T)$, flip rows $(1,\ldots,1)$ and $(1,0,\ldots,0)$, and budget $b_H=T$. Their true costs are $T$ and one, so the optimal retained count is one. After $a$ queries, the first row's upper cost is $T(T+1)-aT$; certifying it requires all $T$ entries. On the second row, querying $a_2$ entries gives upper cost $T(T-a_2)+e$, where $e\in\{0,1\}$ records whether its positive entry was queried. Minimum safety therefore uses its $T-1$ zero entries. Excluding the two-row set requires observed total cost at least $T+1$, hence all $T+1$ positive entries across both rows. These separately minimum witnesses are disjoint, giving union and additive counts $2T$.

A joint witness queries the first row completely and the positive entry of the second row. It retains the first row and excludes retaining both, using $T+1$ entries. Every exact certificate must exclude retaining both and therefore reveal these $T+1$ positive entries. The joint minimum follows.

For a score realization, set cutoff $1/2$ and define the first base row by $a_j=1/2+j/(4M^2)$ for $j=0,\ldots,M-1$. Set the second row to $b_0=0$ and $b_{j+1}=M/(2(M-1))+j/(4M^2)$ for $j=0,\ldots,M-2$. Both means exceed the cutoff. Replacing base member $j$ flips the first row\textquotesingle s membership precisely when the candidate score is below $\tau_j=M/2-\sum_v a_v+a_j$; replace $a$ by $b$ for the second row. The first row has $M$ positive distinct thresholds; the second has $M-1$. In each row, the midpoint of the largest two positive thresholds realizes exactly one flip. Use this score in every first-row candidate and in one second-row candidate; set all other second-row candidate scores to one. All scores lie in $[0,1]$. With $N=2$, $q=1/2$, $k=1$ and $\alpha=1/(2M)$, the budget is $T$.
\end{proof}
Table~\ref{tab:joint-example} instantiates the construction at $T=6$. The separately shortest safety proof retains row~2; joint optimization retains row~1 using evidence already needed for exclusion. Exact rational score checks and independent query-subset enumeration verify the construction; NLP reductions above remain measured on their own fixed panels.



\paragraph{Separation with a fixed evaluator pool.}
The preceding family increases the pool size with the candidate count. The same mechanism operates with a fixed pool; the integer budget determines how much evidence the separate minima can share.
\begin{proposition}[Fixed-pool separation]
\label{prop:fixed-pool-separation}
Fix integers $M\ge3$, $T\ge1$ and flip budget $b=b_H$ with $T\le b<(M-1)T$. Write $b=r(M-1)+t$ with $0\le t<M-1$. There is a realizable mean-pooling instance with two selected outputs, $k=1$, $\alpha=b/(2MT)$, empty history and exact stopping satisfying
\[
 Q_{\rm joint}=T+1,\qquad Q_{\rm split}=2T+1-r,
\]
\[
 Q_{\cup,\rm sep}=\begin{cases}
 T+1,&r=0,\\
 2T+1-r-\one[t>0],&r\ge1.
 \end{cases}
\]
Strict joint advantage holds precisely when $r\ge1$ and $T-r>\one[t>0]$ within this family.
\end{proposition}
\begin{proof}
Use caps $(M,M-1)$. Give the first row $T$ positive integer contributions at most $M-1$, summing to $b$; such a row exists because $T\le b<(M-1)T$. Give the second row one contribution of one and $T-1$ zeros. The true costs are $b$ and one, so the optimum retains one row. Excluding both requires all $T+1$ positive entries, since their total is $b+1$. This unique minimum exclusion witness also certifies the first row, proving $Q_{\rm joint}=T+1$.

Certifying the first row requires all its entries: every unqueried entry has strictly positive upper-cost slack. With $a$ queries on the second row, its upper cost is $(T-a)(M-1)+e$, where $e$ is one if its positive entry was queried and zero otherwise. If $r\ge1$, minimum safety uses $a=T-r<T$ queries on this row. With $e=0$ the upper cost is $b-t$, whereas $e=1$ is feasible precisely when $t>0$. A minimum safety witness can therefore share at most $\one[t>0]$ entries with the unique minimum exclusion witness, and this overlap is attainable. This proves both the additive and best-union formulas, including all ties. If $r=0$, both rows require $T$ queries for safety. Choosing the first row makes its safety witness a subset of the minimum exclusion witness, giving union size $T+1$.

For a score realization, use cutoff $1/2$, first-row base scores $a_j=1/2+j/(4M^2)$ for $j=0,\ldots,M-1$, and second-row scores $b_0=1/(16M^2)$ and $b_{j+1}=M/(2(M-1))+j/(4M^2)$ for $j=0,\ldots,M-2$. Both base means exceed the cutoff. Their replacement thresholds $M/2-\sum_v a_v+a_j$ (and analogously for the second row) are distinct; respectively $M$ and $M-1$ lie in $(0,1)$. Candidate scores strictly between consecutive positive thresholds realize each needed flip count. A score strictly above the largest threshold realizes zero. Taking $N=2$, $q=1/2$ and the stated $\alpha$ gives integer budget $b$.
\end{proof}

The same instance is valid for any $\alpha$ with $\lfloor2MT\alpha\rfloor=b$. Fix $M$ and $1/(2M)\le\alpha<(M-1)/(2M)$, and take this rounded budget as $T$ grows. Then
\[
 \frac{Q_{\cup,\rm sep}}{Q_{\rm joint}}
 \longrightarrow 2-\frac{2M\alpha}{M-1}.
\]
At $M=7$ and $\alpha=.10$, the limit is $53/30$ without a divisibility restriction on $T$. The original specialization $T=30n$, $b=42n$ gives costs $30n+1$ and $53n+1$. These are attained ratios in this family, not tight upper bounds for every fixed-pool instance. The exact formula also identifies equality cases: at $r=0$ the first row supplies a minimum safety witness, and at $r=T-1$, $t>0$ a one-entry safety witness lies in the exclusion witness.

At gap tolerance one, either safety for one row or exclusion of both suffices. Their minima are $T-r$ and $T+1$, respectively, so joint and best separate-union costs both become $T-r$. Because $|S_0|=2$ here, the primary rule gives $\varepsilon=\lceil.02|S_0|\rceil=1$, which removes the strict joint advantage.

\paragraph{Local score perturbations.}
The construction of Proposition~\ref{prop:fixed-pool-separation} does not require cutoff ties. Choose each positive-count candidate at the midpoint between its adjacent thresholds, and each zero-count candidate midway between the largest threshold and one. Consider the distances of the constructed base and candidate scores from zero and one, together with the absolute cutoff margins of base means, observed replacement means and replacement means at candidate endpoints zero and one. Let $\delta$ be half the smallest of these distances and margins. All these distances are positive. Perturbing each score by less than $\delta$, with cutoff and budget fixed, changes each mean by less than $\delta$. Base memberships, endpoint caps and observed flip contributions therefore stay unchanged, preserving both certificate minima. This is local robustness of the normalized-score construction, not robustness to refitting ECDF maps.

Exact rational checks cover 595 constructed instances, including every integer budget $T\le b<(M-1)T$ for $M=3,\ldots,9$ and $T=1,\ldots,6$. Literal enumeration examines 882,928 query subsets across 589 cases. An independent outcome-count enumeration covers 545,800 states across all 595 cases, optimizing over all tied component minima. Both agree with the exact and gap-one formulas; 223 cases have no strict joint advantage. At $M=7$, $\alpha=.10$ and $T=7,17$, the rounded budgets are 9 and 23. Joint costs are 8 and 18, versus best separate-union costs 13 and 31 and additive costs 14 and 32. At $T=30,60,90$, joint costs are 31, 61 and 91, versus separate-union costs 54, 107 and 160. The constructed score margins permit $\delta=1/5488$ in these three cases.

\paragraph{When separate minima are incompatible.}
Changing the retained set can make a longer safety proof share more exclusion evidence. We give a prefix-sum criterion for any number of active rows when the true optimum retains all but one, including a class with tied minimum exclusion witnesses. Consider $m\ge2$ selected active rows, empty history, exact stopping and integer budget $b=b_H$, with
\[
 \sum_{i=1}^{m-1}F_{(i)}\le b<\sum_{i=1}^m F_i.
\]
Thus $L^*=m-1$. For each $(m-1)$-row set $C$, define reductions $d_{ir}=p_i-f_{ir}$ and required reduction $D_C=T\sum_{i\in C}p_i-b$. Let $a_C$ be the length of the shortest descending reduction prefix in $C$ reaching $D_C$: zero if $D_C\le0$, and infinity if no prefix suffices. Set $a=\min_C a_C$.

Order all $mT$ contributions as $c_1\ge\cdots\ge c_{mT}$, put $C_j=\sum_{v=1}^jc_v$, and let $e=\min\{j:C_j>b\}$. For the unique and tied cases specified below, we represent the family of minimum exclusion witnesses as
\[
 \mathcal E=\{H\cup B':B'\subseteq B,\ |B'|=\kappa\}.
\]
If $e=mT$, the only minimum exclusion witness is the full entry set: take it as $H$, with $B=\varnothing$ and $\kappa=0$. For $e<mT$, the top-$e$ set is unique when $C_e-c_e+c_{e+1}\le b$; use the same convention. In the remaining case, suppose $C_e=b+1$; take $H$ to contain entries with contribution greater than $c_e$, $B$ the entries equal to $c_e$, and $\kappa=e-|H|$. All other entries form $Z$. For $X\in\{H,B,Z\}$, let $P_C^X(t)$ sum the $t$ largest reductions among entries of $X$ in rows $C$, with $P_C^X(0)=0$ and value $-\infty$ when fewer than $t$ entries exist.

\begin{proposition}[Near-full compatibility with exclusion ties]
\label{prop:compatibility}
For a row set $C$, write $h_C=\sum_{(i,r)\in H:i\in C}d_{ir}$ and $b_C=|B\cap(C\times[T])|$. Suppose
\[
 h_C+P_C^B(\min\{\kappa,b_C\})\ge D_C
\]
for some $(m-1)$-row set $C$. Let $\mathcal F_C$ be the set of all nonnegative integer triples $(u,v,w)$ satisfying
\[
 \begin{aligned}
 u+v+w&=a,\\
 P_C^H(u)+P_C^B(v)+P_C^Z(w)&\ge D_C.
 \end{aligned}
\]
Define
\[
 \delta=\min_{C:a_C=a}\ \min_{(u,v,w)\in\mathcal F_C}
       \bigl[w+(v-\kappa)_+\bigr].
\]
Then
\[
 Q_{\rm joint}=e,\qquad Q_{\cup,\rm sep}=e+\delta.
\]
Hence strict joint advantage is equivalent to $\delta>0$ within this class. Sorting and prefix sums compute the tied case in $O(m^2T\log(mT)+ma^2)$ time; the unique case reduces to $O(m^2T\log(mT))$.
\end{proposition}
\begin{proof}
Exclusion requires total observed contribution greater than $b$, so every exclusion witness has at least $e$ entries. If $e=mT$, there is only one size-$e$ set. If $e<mT$ and the unique-case inequality holds, every other size-$e$ set has sum at most $C_e-c_e+c_{e+1}\le b$. In the remaining tied case with $C_e=b+1$, integer contributions force every size-$e$ exclusion witness to attain this maximum sum. It must contain all entries above $c_e$ and choose exactly $\kappa$ entries tied at $c_e$. Thus $\mathcal E$ is the complete family in either case.

For fixed $C$, the mandatory entries in rows $C$ contribute their reductions, while up to $\kappa$ boundary entries in $C$ can contribute the largest available reductions. Remaining boundary slots can be filled without reducing this sum. The assumed inequality therefore says some $E\in\mathcal E$ is safe for $m-1$ rows. It attains the exclusion lower bound, proving $Q_{\rm joint}=e$.

Every minimum safety witness has size $a$ and lies in a set $C$ with $a_C=a$; otherwise an entry outside a set made safe could be removed. Suppose such a witness takes $u,v,w$ entries from $H,B,Z$. Its best possible reduction is the sum of the three prefixes, attained by taking these entries. A size-$e$ exclusion witness contains all $u$ mandatory entries and can share at most $\min(v,\kappa)$ chosen boundary entries. It shares no $Z$ entries. Its union with the safety witness therefore has minimum size $e+w+(v-\kappa)_+$. Minimizing over the feasible counts and row sets proves the formula. For each of $m$ row sets, sorting at most $mT$ reductions and scanning $u,v$ costs $O(mT\log(mT)+a^2)$. If $B$ is empty, $v=0$ and $w=a-u$, leaving only $u$ to scan.
\end{proof}
A singleton minimum-exclusion family gives the unique-exclusion case. In the construction of Proposition~\ref{prop:strict-joint}, $e=T+1$, $a=T-1$ and $\delta=T-1$. The tied case optimizes over all minimum exclusion witnesses, not arbitrarily breaking a tie. It still requires a minimum exclusion witness also certifying safety. It does not cover arbitrary exclusion families, nonzero tolerance or larger retained-count deficits.


\paragraph{A realizable multirow family.}
For any $m,T\ge2$, use $M=T$, caps $p_i=T$, contribution rows $(0,\ldots,0,1)$ and $m-1$ copies of $(1,\ldots,1)$, and budget $b=(m-1)T$. All but one row can be retained. Exclusion requires every positive entry, so its unique minimum $E$ has size $(m-1)T+1$. It certifies the $m-1$ all-one rows. For any $(m-1)$-row set, the required upper-cost reduction is $D=(m-1)T(T-1)$. A query reduces the upper cost by $T-1$, or by $T$ at one of the $T-1$ zero entries. With $a$ queries, reduction is at most $a(T-1)+(T-1)$. Reaching $D$ requires $a\ge(m-1)T-1$, with equality only if all $T-1$ zeros are queried. This bound is attained by taking these zeros and $(m-2)T$ positive entries within a set containing the first row and $m-2$ other rows. Thus
\[
 Q_{\rm joint}=(m-1)T+1,\qquad Q_{\cup,\rm sep}=mT.
\]
Here $\delta=T-1$: allowing a longer safety proof avoids all zero entries and reduces the union.

For a mean-score realization, give every row base scores $a_j=1/2+j/(4T^2)$, $j=0,\ldots,T-1$, and cutoff $1/2$. All $T$ replacement thresholds $T/2-\sum_v a_v+a_j$ are positive and distinct. A candidate midway between the largest two thresholds realizes contribution one; a candidate score of one realizes zero. With $N=m$, $q=1/m$, $k=1$ and $\alpha=(m-1)/(2T)$, the budget is exactly $(m-1)T$. Exact rational arithmetic checks all 25 members with $m,T\in\{2,\ldots,6\}$; independent enumeration checks the 12 members with $mT\le12$, totaling 19,600 query subsets. Larger checked members use the explicit witnesses and criterion, not exhaustive subset enumeration.

\paragraph{Finite checks and empirical scope.}
Independent subset enumeration checks 16,149 canonical near-full instances: $(m,T)=(3,2)$ with caps in $\{1,2,3\}$, $(4,2)$ with caps in $\{1,2\}$, and $(3,3)$ with every cap equal to three. Within each grid we enumerate row contribution multisets, row multisets and every integer budget satisfying $L^*=m-1$. Across 5,688,256 query subsets, the unique-exclusion formulas agree in all 2,220 applicable instances: 615 have strict gain and 1,605 have equality. Another 2,424 strict-gain instances fall outside this unique-exclusion criterion. These finite checks test the formulas, not a general proof or NLP prevalence.

The earlier two-row check remains a specialization: all 14,674 instances with $T\in\{1,2,3\}$, caps in $\{1,2,3\}$ and budgets satisfying $F_1,F_2\le b<F_1+F_2$ and $b<Tp_1,Tp_2$ span 913,084 query subsets in total. Within this search, the criterion applies to 2,812 instances spanning 171,772 subsets; all agree with enumeration (742 gains and 2,070 equalities). Returning to the 16,149-instance multirow grid, allowing the tied maximum-sum family adds 1,203 applicable instances: 237 gains and 966 equalities. The formulas agree with freshly enumerated minima in all 3,423 applicable instances, spanning 1,061,120 query subsets. The other instances are outside the sufficient conditions; they are not classified by this test. Among 960 exact NLP targets, five have active optimum $m-1$. All five satisfy the tied-family condition, but none of their minimum exclusion witnesses can certify safety. The extended criterion therefore still has zero coverage on this NLP grid. Positive-tolerance targets are outside its scope.


\paragraph{Evidence position and entry discovery.}
\label{app:entry-placement}
Use the two-row score construction of Proposition~\ref{prop:strict-joint}, with caps $(T+1,T)$, and vary only the position $j\in\{1,\ldots,T\}$ of the second-row positive entry. Every row retains its contribution multiset, so $Q_{\rm joint}=T+1$ and $Q_{\cup,\rm sep}=2T$ throughout. A fixed order querying the first row and then the second, each in candidate-index order, stops exactly after $T+j$ queries: the first row is safe after $T$ observations, and exclusion becomes certified when the second-row positive entry is seen. Reversing the row order takes $2T$ queries.

We enumerate every position for $T\in\{2,4,6,8,10,12\}$, using these two fixed orders and largest width (row-one ties first, candidate index ascending). All 42 score instances use exact rational realizations and exact stopping from empty history. The 126 paths are checked at every reply and first stop; independent subset enumeration verifies the joint minima for $T\le6$. At $T=6$, the first fixed order takes 7--12 queries, the reverse order always 12, and width 11--12, while the joint minimum remains seven. These controlled positions separate evidence size from discoverability under entry access. They do not establish how often this mechanism explains the NLP results or evaluate learned reference priorities. Complete scores, positions and paths accompany the code.


\subsection{Full-batch mean evidence: numerical minima, sensitivity and discovery cost}
\label{app:full-mean-grid}
We retain the six core panels, four depths and 40 splits, full-reference right ECDF, complete candidate families and $\alpha=.10$. For each of the 960 instances, we optimize both $\varepsilon=0$ and $\varepsilon=\lceil.02|S_0|\rceil$, starting from empty candidate-test history. Batch size, cutoff, selected set and budget are unchanged. Outcome-count optimization ranges over every admissible safe target. Equivalent fixed-threshold formulations cover every level not excluded by true costs; continuous relaxations screen levels whose lower bound reaches an existing feasible witness. Strict-improvement searches retain this witness when no smaller one exists. All levels contribute to the global lower bound.

The historical optimization records report numerical closure for all 1,920 joint targets. Each reconstructed integer-feasible witness matches the recorded numerical global lower bound under the rounding convention of Appendix~\ref{app:joint-evidence}. This agreement assumes valid solver bounds and is not an exact minimality certificate. Historical solver feasibility and integrality settings were not fully recorded. For the hard cases left by the initial runs, complete integer safe-target partitions strengthen the fixed-threshold bounds; every branch contributes to the global lower bound. Each witness is reconstructed as actual score entries and checked against both capacities. The original 96 exact targets agree with the earlier full-test study. Comparisons include all four original query orders and the three frontier variants, each using the same instance and stopping target. Tables~\ref{tab:full-mean-primary} and~\ref{tab:full-mean-exact} report all depths and the number of positive minima in each ratio denominator.

For separately optimal safety and exclusion witnesses, the union program minimizes over tied component optima and safe targets. Equal joint and additive optima close the union by a bound sandwich. A feasible union of separately optimal components also closes the comparison when its size reaches the numerically closed joint minimum; other safe targets then need not be solved. More generally, a feasible union of separately optimal components certifies the global minimum when every admissible safe target has a lower bound at least as large as the cardinality of this union. This does not require solving every safe-target subproblem to optimality. Exhaustive integer-threshold partitions, with LP bounds and integer witnesses, resolve the remaining comparisons. A minimum safety component contains no query with zero upper-cost reduction; a minimum exclusion component contains no query with zero lower-cost increase. The \fullMeanUnionProven{} closed union comparisons contain \fullMeanStrictGain{} strict joint improvements, totaling \fullMeanJointSaving{} entries. They comprise 32 of 960 primary and 57 of 960 exact targets. Mean reductions over all targets are 5.34 and 1.19 entries, and both medians are zero. These gains concern offline evidence, not savings achieved by the acquisition policies. Numerical optimality uses solver bounds, not exact rational proof certificates.
\paragraph{Independent numerical-bound audit.}
\label{app:solver-audit}
The independent replay found one invalid numerical lower bound, in a union subproblem. An audit continues all 73 attempts left unresolved or failed by a 202-attempt replay with explicit settings. It allows 180 seconds per objective, and the attempts involve 40 distinct targets. Of these attempts, 21 reach numerical closure, 51 remain unresolved and one returns an inconsistent numerical bound, which is rejected. Integer checks validate the returned witnesses in the 72 non-error attempts. In WMT24 ja-zh, split 16, $q=.10$, at the primary stop, a union subproblem for active safe count 913 reports an optimum of 14,463. The returned component witnesses instead have a feasible union of 14,462 entries, certifying safety for 913 and optimistic count at most 933. This contradicts the subproblem's reported numerical lower bound, which we discard. Disabling presolve yields matching numerical primal and dual values of 14,462 for this same safe target. This invalidates the reported bound for the subproblem but neither refutes nor independently verifies the historical all-target value of 14,423. The implementation retains such witnesses as feasible upper bounds and rejects their conflicting optimality status. Historical merged numerical comparisons and exact minimality proofs remain distinct. Table~\ref{tab:evidence-levels} separates these evidence levels.

\begin{table}[t]\centering\small
\begin{tabular}{@{}p{1.9cm}p{2.6cm}p{2.3cm}@{}}\toprule
Evidence & Establishes & Does not establish\\\midrule
Integer witness check & The revealed entries satisfy the certificate constraints & No smaller witness exists\\
Numerical bound closure & Minimality if the solver bounds are valid & Exact-arithmetic minimality\\
Presolve-off re-solve & The same value under a second solver configuration & Validity on other data\\
Enumeration and proofs & Minimality on covered instances or under theorem conditions & Results on uncovered instances\\\bottomrule
\end{tabular}
\caption{Evidence levels behind reported minima. Integer checks cover every reported witness. Numerical closure covers the historical 1,920-target grid and the 48 SummEval unions. Presolve-off re-solves cover the solver-sensitivity targets and all 48 SummEval unions. Enumeration covers the stated small instances.}
\label{tab:evidence-levels}
\end{table}

\paragraph{Solver sensitivity.}
We drew four targets from each panel, depth and stopping cell, giving 192 of the 1,920 targets (seed 20261012). Each joint objective was re-solved from empty history with the outcome-grouped program, HiGHS presolve on or off, and nominal limits of 10 or 60 seconds. With a 10-second limit, 155 targets closed under each presolve setting. With 60 seconds, 179 closed under each setting. Every closed value equaled the recorded minimum, and no returned witness fell below it. Median solve times for closed targets were 0.32--0.68 seconds; the 90th percentile at 60 seconds was 15.5 seconds with presolve and 20.2 seconds without it. With presolve, the 13 targets still open at 60 seconds came from WMT24 at $q=.30$ and from WMT24 ja-zh. In seven of them the returned witness already equaled the recorded minimum, and only the lower bound stayed open. Solve time had Spearman correlation .592 with the number of program variables, which ranged from 388 to 11,586. Records are in \nolinkurl{outputs/scope-extensions/solver_sensitivity.jsonl}. During exhaustive small-instance checks of the new solvers, a presolve-enabled HiGHS run reported a suboptimal union value as optimal. We therefore reran the per-replacement and weighted analyses without presolve; every value closed in both versions agreed.

\begin{table*}[t]
\centering\small
\setlength{\tabcolsep}{4pt}
\begin{tabular}{lrrrrrrrr}
\toprule
Panel & $q$ & Mean $Q$ & Closed & $n_+$ & Ref. & Width & Frontier & Gain/n \\
\midrule
RewardBench 2 & 0.05 & 663.2 & 40 & 40 & 1.098 & 1.165 & 1.134 & 0/40 \\
RewardBench 2 & 0.10 & 1157.8 & 40 & 40 & 1.114 & 1.173 & 1.139 & 0/40 \\
RewardBench 2 & 0.20 & 1605.5 & 40 & 40 & 1.197 & 1.200 & 1.147 & 0/40 \\
RewardBench 2 & 0.30 & 1441.1 & 40 & 40 & 1.410 & 1.237 & 1.192 & 0/40 \\
WMT23, 13 & 0.05 & 1709.4 & 40 & 40 & 1.105 & 1.197 & 1.132 & 2/40 \\
WMT23, 13 & 0.10 & 3289.3 & 40 & 40 & 1.070 & 1.153 & 1.108 & 0/40 \\
WMT23, 13 & 0.20 & 1712.0 & 40 & 40 & 1.427 & 1.267 & 1.208 & 0/40 \\
WMT23, 13 & 0.30 & 1337.7 & 40 & 40 & 1.976 & 1.350 & 1.297 & 0/40 \\
WMT23, 12 & 0.05 & 1613.3 & 40 & 40 & 1.090 & 1.150 & 1.117 & 5/40 \\
WMT23, 12 & 0.10 & 2834.3 & 40 & 40 & 1.073 & 1.141 & 1.106 & 0/40 \\
WMT23, 12 & 0.20 & 1769.1 & 40 & 40 & 1.381 & 1.285 & 1.246 & 0/40 \\
WMT23, 12 & 0.30 & 1880.7 & 40 & 40 & 1.618 & 1.345 & 1.302 & 0/40 \\
WMT24 en-de & 0.05 & 5144.4 & 40 & 40 & 1.097 & 1.170 & 1.133 & 0/40 \\
WMT24 en-de & 0.10 & 5886.1 & 40 & 40 & 1.120 & 1.157 & 1.123 & 0/40 \\
WMT24 en-de & 0.20 & 8277.0 & 40 & 40 & 1.181 & 1.208 & 1.183 & 0/40 \\
WMT24 en-de & 0.30 & 8383.2 & 40 & 40 & 1.343 & 1.264 & 1.223 & 0/40 \\
WMT24 en-es & 0.05 & 4641.5 & 40 & 40 & 1.067 & 1.132 & 1.102 & 1/40 \\
WMT24 en-es & 0.10 & 5401.0 & 40 & 40 & 1.124 & 1.165 & 1.127 & 0/40 \\
WMT24 en-es & 0.20 & 7306.3 & 40 & 40 & 1.169 & 1.215 & 1.191 & 0/40 \\
WMT24 en-es & 0.30 & 6525.0 & 40 & 40 & 1.387 & 1.279 & 1.241 & 0/40 \\
WMT24 ja-zh & 0.05 & 4115.2 & 40 & 40 & 1.087 & 1.156 & 1.142 & 12/40 \\
WMT24 ja-zh & 0.10 & 7347.5 & 40 & 40 & 1.080 & 1.150 & 1.126 & 11/40 \\
WMT24 ja-zh & 0.20 & 11132.9 & 40 & 40 & 1.131 & 1.196 & 1.178 & 1/40 \\
WMT24 ja-zh & 0.30 & 12272.4 & 40 & 40 & 1.179 & 1.242 & 1.206 & 0/40 \\
\bottomrule
\end{tabular}
\caption{Full-batch mean evidence from empty history, $\alpha=.10$, $\varepsilon=\lceil .02|S_0|\rceil$. All four depths and 40 original splits are included. $Q$ averages the matching minimum joint entry count over all 40 splits; all targets are numerically closed. Closed counts minima whose feasible value matches the numerical global lower bound; $n_+$ counts positive matching minima. Ref., Width and Frontier average $Q_{\rm online}/Q$ only on these $n_+$ matching instances, with zero minima excluded. Gain/n counts $Q_{\rm joint}<Q_{\cup,\rm sep}$ among n closed union comparisons. Historical resource-limited attempts are retained separately. Candidate-test score entries are the cost unit.}
\label{tab:full-mean-primary}
\end{table*}

\begin{table*}[t]
\centering\small
\setlength{\tabcolsep}{4pt}
\begin{tabular}{lrrrrrrrr}
\toprule
Panel & $q$ & Mean $Q$ & Closed & $n_+$ & Ref. & Width & Frontier & Gain/n \\
\midrule
RewardBench 2 & 0.05 & 696.4 & 40 & 40 & 1.092 & 1.137 & 1.110 & 0/40 \\
RewardBench 2 & 0.10 & 1212.0 & 40 & 40 & 1.107 & 1.147 & 1.119 & 0/40 \\
RewardBench 2 & 0.20 & 1699.8 & 40 & 40 & 1.182 & 1.181 & 1.123 & 0/40 \\
RewardBench 2 & 0.30 & 1579.6 & 40 & 40 & 1.362 & 1.207 & 1.152 & 0/40 \\
WMT23, 13 & 0.05 & 1789.9 & 40 & 40 & 1.097 & 1.173 & 1.110 & 5/40 \\
WMT23, 13 & 0.10 & 3387.7 & 40 & 40 & 1.077 & 1.142 & 1.099 & 0/40 \\
WMT23, 13 & 0.20 & 1911.4 & 40 & 40 & 1.378 & 1.231 & 1.181 & 0/40 \\
WMT23, 13 & 0.30 & 1631.7 & 40 & 40 & 1.721 & 1.296 & 1.257 & 0/40 \\
WMT23, 12 & 0.05 & 1694.4 & 40 & 40 & 1.081 & 1.125 & 1.093 & 8/40 \\
WMT23, 12 & 0.10 & 2930.0 & 40 & 40 & 1.077 & 1.129 & 1.097 & 0/40 \\
WMT23, 12 & 0.20 & 1959.6 & 40 & 40 & 1.344 & 1.249 & 1.214 & 0/40 \\
WMT23, 12 & 0.30 & 2174.8 & 40 & 40 & 1.531 & 1.299 & 1.255 & 0/40 \\
WMT24 en-de & 0.05 & 5370.8 & 40 & 40 & 1.086 & 1.140 & 1.106 & 3/40 \\
WMT24 en-de & 0.10 & 6180.4 & 40 & 40 & 1.116 & 1.137 & 1.113 & 0/40 \\
WMT24 en-de & 0.20 & 8856.5 & 40 & 40 & 1.171 & 1.183 & 1.165 & 0/40 \\
WMT24 en-de & 0.30 & 9272.1 & 40 & 40 & 1.304 & 1.225 & 1.189 & 0/40 \\
WMT24 en-es & 0.05 & 4785.9 & 40 & 40 & 1.069 & 1.117 & 1.090 & 2/40 \\
WMT24 en-es & 0.10 & 5655.5 & 40 & 40 & 1.122 & 1.147 & 1.120 & 0/40 \\
WMT24 en-es & 0.20 & 7825.5 & 40 & 40 & 1.160 & 1.190 & 1.171 & 0/40 \\
WMT24 en-es & 0.30 & 7295.5 & 40 & 40 & 1.345 & 1.239 & 1.205 & 0/40 \\
WMT24 ja-zh & 0.05 & 4298.4 & 40 & 40 & 1.079 & 1.130 & 1.119 & 16/40 \\
WMT24 ja-zh & 0.10 & 7693.5 & 40 & 40 & 1.069 & 1.122 & 1.104 & 16/40 \\
WMT24 ja-zh & 0.20 & 11954.3 & 40 & 40 & 1.114 & 1.161 & 1.148 & 7/40 \\
WMT24 ja-zh & 0.30 & 13096.4 & 40 & 40 & 1.168 & 1.217 & 1.177 & 0/40 \\
\bottomrule
\end{tabular}
\caption{Full-batch mean evidence from empty history, $\alpha=.10$, $\varepsilon=0$. All four depths and 40 original splits are included. $Q$ averages the matching minimum joint entry count over all 40 splits; all targets are numerically closed. Closed counts minima whose feasible value matches the numerical global lower bound; $n_+$ counts positive matching minima. Ref., Width and Frontier average $Q_{\rm online}/Q$ only on these $n_+$ matching instances, with zero minima excluded. Gain/n counts $Q_{\rm joint}<Q_{\cup,\rm sep}$ among n closed union comparisons. Historical resource-limited attempts are retained separately. Candidate-test score entries are the cost unit.}
\label{tab:full-mean-exact}
\end{table*}


\paragraph{Where joint gains occur.}
\label{app:joint-mechanism}
Strict reductions beyond optimal deduplication occur in 32 primary and 57 exact targets out of 960 each; both gain medians are zero. Among targets with strict gains, mean savings are 160.34 and 19.96 entries, respectively. WMT24 ja-zh accounts for 24 of the 32 primary improvements and 39 of the 57 exact improvements; RewardBench has none at either stop. Table~\ref{tab:gain-distribution} reports relative and absolute magnitudes. These differences describe the recorded panels and do not identify a causal effect of language or evaluator composition.
\begin{table*}[t]
\centering\small
\begin{tabular}{lrrrrr}
\toprule
Stop & Strict gains / targets & Median & Mean & Mean if positive & Max. relative reduction \\
\midrule
primary & 32 / 960 & 0 & 5.34 & 160.34 & 2.61\% \\
exact & 57 / 960 & 0 & 1.19 & 19.96 & 0.87\% \\
\bottomrule
\end{tabular}
\caption{Distribution of $Q_{\cup,\rm sep}-Q_{\rm joint}$ across all 960 targets per stopping setting, with the same panels, depths, splits, empty histories and $\alpha=.10$ as Table~\ref{tab:cost-contrasts}. Median and mean are across all targets; the conditional mean includes only strict improvements. Absolute savings count score entries. Relative reduction divides by $Q_{\cup,\rm sep}>0$, and the last column reports its maximum, not its average.}
\label{tab:gain-distribution}
\end{table*}


Full retention is feasible in 903 of the 960 instances, hence in 1,806 targets across the two stops. This makes exclusion automatic at either tolerance. Let $V_A$ denote capacity on the active selected rows, equivalently total capacity after appending zero invariant-row costs minus $m_0$. If $L^*=m$, every history satisfies $V_A(\ell)=m$. The count-gap condition is equivalent to $V_A(u)\ge\max(0,m-\varepsilon)$, so only safety evidence is needed and $Q_{\rm joint}=Q_{\cup,\rm sep}=Q_{\rm split}$. Full-information feasibility does not imply zero-query certifiability.

Table~\ref{tab:explanatory-overlap} cross-classifies this condition with the minimized intermediate-entry count. All 1,024 targets with $d_*=0$ are inside the full-retention region, adding zero no-gain coverage beyond this simpler condition. The remaining 25 no-gain targets are non-full primary targets with $d_*>0$. Replaying their saved optimal witnesses shows $V^-+\varepsilon\ge |S_0|$ in every case, so no count-exclusion evidence is needed. More generally, if a minimum joint witness satisfies this inequality, it is a minimum safety witness for its attained count: a shorter safety witness would also be a joint certificate, contradicting minimality. Pairing it with the empty exclusion witness gives $Q_{\rm joint}=Q_{\cup,\rm sep}=Q_{\rm split}$. This is sufficient, not necessary in general. For example, two rows with $T=p_1=p_2=1$, true contributions $(1,1)$, $b=1$ and $\varepsilon=0$ have zero-query safety for one row and require both entries to exclude two rows. Joint and separate costs equal two, although automatic exclusion is impossible. In this grid, all 1,831 no-gain targets have an optimal witness satisfying $V^-+\varepsilon\ge|S_0|$; none of the 89 strict-gain targets has such a witness. The classification uses saved optimal witnesses and is not a pre-query predictor or an independent confirmation sample.

\paragraph{Diagnosing gains before candidate scoring.}
The full-information condition $V^*+\varepsilon<|S_0|$ marks targets requiring exclusion evidence. It flags exactly the 89 strict-gain targets on this grid, but it needs hidden costs. We therefore tested signals available before any candidate test score: reference-ridge predicted row costs, visible base scores, the budget and the tolerance. The ratio of total predicted active cost to $b_H$ ranked strict-gain targets above no-gain targets. Its AUC was .971 at primary stopping and .959 at exact stopping. The zero-query normalized gap $(|S_0|-V^-_0)/|S_0|$ gave .964 and .930. The fraction of active rows with $p_i\ge2$ gave .786 and .747. Replacing true costs by predictions in this condition flagged 210 of 1,920 targets, which contained 75 of the 89 gains. Gains concentrated in WMT24 ja-zh, so these AUC values partly reflected panel differences. Per-target features are in \nolinkurl{outputs/scope-extensions/gain_predictor_records.csv}.

\paragraph{Using the diagnostic during acquisition.}
A switching policy uses joint frontier when the predicted exclusion-required condition flags a target and safety-only frontier otherwise. On the 960 primary targets it averaged 5,209.9 queries, against 5,210.6 for joint and 5,207.6 for safety-only frontier. Switching on the true exclusion-required condition gave 5,206.9. At exact stopping, the four means differed by at most 2.5 entries. The sharing decision moved online cost by a few entries, whereas joint and safety-only frontier exceeded $Q_{\rm joint}$ by 733.66 and 730.67 entries on average. Discovery overhead was thus far larger than the certificate-size reductions from joint optimization on these panels.
\begin{table*}[t]
\centering\small
\begin{tabular}{llrrrr}
\toprule
Stop & Full retention & Targets & $d_*=0$ & Strict gain & Automatic exclusion\\
\midrule
Exact & Yes & 903 & 454 & 0 & 903\\
Exact & No & 57 & 0 & 57 & 0\\
Primary & Yes & 903 & 570 & 0 & 903\\
Primary & No & 57 & 0 & 32 & 25\\
\bottomrule
\end{tabular}
\caption{No-gain explanations on all 1,920 mean targets at $\alpha=.10$. Full retention concerns the true complete scores. Automatic exclusion counts targets with an integer-checked optimal witness satisfying $V^-+\varepsilon\ge |S_0|$; it does not mean zero-query stopping. All 1,831 no-gain targets have such a witness, including 25 non-full primary targets. All 1,024 zero-intermediate optima remain in the full-retention region. These are descriptive properties of saved optimal witnesses, not pre-query predictors or independent trials.}
\label{tab:explanatory-overlap}
\end{table*}


To test sensitivity within these 25 exposed non-full, no-gain cases, we fix their matrices, maps, cutoff and candidate family and multiply the original integer budget by $.75,1,1.25$ (rounding down) and the original tolerance by $0,1,2$. Of 225 targets, 159 joint/union comparisons close numerically: 42 show strict gains and 117 equality; 66 remain unresolved. Fifty targets coincide with previously closed tasks. At the original budget, removing tolerance changes all 25 no-gain cases to strict gains, whereas increasing the budget by 25\% yields equality for all three tested tolerances. This is a deterministic change of certification task on outcome-selected matrices, not evidence of a population effect or a general monotonicity law for joint advantage. Initial nominal limits were two seconds for each complete joint or additive-separate objective and four seconds for union optimization. Retries gave each complete joint or additive-separate objective 15 seconds; union retries allowed ten seconds per solver call within a 30-second total budget. Selected additive-separate rechecks used 60 seconds. The current summary closes 159 of the 225 comparisons and leaves 66 unresolved; an earlier archived audit recorded 151 closed and 74 unresolved. Later retries allowed 60 seconds for the joint objective and, after its closure, 15 seconds per union solver call within a 30-second total budget. These call limits do not allocate an independent budget to every safety or exclusion component. These are nominal limits, with measured times and unresolved bounds recorded. Current bounds are checked against preserved solver calls.

To examine the role of intermediate contributions, we measure the fraction of active entries satisfying $0<f_{ir}<p_i$. We retain all 1,920 targets and stratify by panel, depth and stopping criterion. Thirteen of the 48 cells contain both strict-gain and no-gain instances. The strict-gain subgroup has a larger mean intermediate fraction in eight cells and a smaller one in five. Thus the presence of intermediate contributions permits shared evidence, but its prevalence alone does not explain the observed gains. Split-level costs, initial gaps, exclusion slack and contribution fractions accompany these post hoc descriptive summaries. Overlapping splits do not supply independent population evidence.



\paragraph{Separating evidence costs.}
For each identical instance and stopping target, we compare $Q_{\rm split}-Q_{\cup,\rm sep}$, $Q_{\cup,\rm sep}-Q_{\rm joint}$ and $Q_{\rm online}-Q_{\rm joint}$. The first removes duplicate counting from separately optimal proofs; the second measures the benefit of choosing their evidence jointly. The last measures online discovery above the joint minimum. The first two sum to $Q_{\rm split}-Q_{\rm joint}$; discovery overhead is measured separately. Table~\ref{tab:cost-contrasts} uses the same complete set for all three means: 960 primary and 960 exact targets. All depths and panels remain in the analysis.

The complete records include 1,920 targets and seven policies, with all joint minima and union comparisons numerically closed. Across all 960 primary targets, the mean reference excess is 809.58 entries; the exact-target mean is 796.66. Historical resource-limited attempts remain separate from these final numerically closed values. Ratios use positive numerically closed minima, and zero minima have a separate field. These overlapping-split summaries describe the recorded panels, not a population effect or an online competitive bound. Files in \nolinkurl{outputs/evidence-cost-contrasts/} include all policy-level records and bounds.
\begin{table*}[t]
\centering\small
\begin{tabular}{llrrrr}
\toprule
Panel & Stop & Closed / tasks & Dedup. & Joint & Discovery \\
\midrule
All panels & primary & 960 / 960 & 18.69 & 5.34 & 809.58 \\
All panels & exact & 960 / 960 & 81.79 & 1.19 & 796.66 \\
RewardBench 2 & primary & 160 / 160 & 0.00 & 0.00 & 274.59 \\
RewardBench 2 & exact & 160 / 160 & 0.00 & 0.00 & 267.05 \\
WMT23, 13 & primary & 160 / 160 & 0.60 & 0.75 & 531.77 \\
WMT23, 13 & exact & 160 / 160 & 5.01 & 0.51 & 542.18 \\
WMT23, 12 & primary & 160 / 160 & 3.33 & 1.70 & 525.33 \\
WMT23, 12 & exact & 160 / 160 & 12.35 & 0.78 & 527.55 \\
WMT24-en-de & primary & 160 / 160 & 0.00 & 0.00 & 1298.33 \\
WMT24-en-de & exact & 160 / 160 & 7.86 & 0.23 & 1291.07 \\
WMT24-en-es & primary & 160 / 160 & 1.34 & 0.67 & 1140.28 \\
WMT24-en-es & exact & 160 / 160 & 7.95 & 0.17 & 1160.61 \\
WMT24-ja-zh & primary & 160 / 160 & 106.86 & 28.95 & 1087.16 \\
WMT24-ja-zh & exact & 160 / 160 & 457.57 & 5.43 & 991.53 \\
\bottomrule
\end{tabular}
\caption{Complete-batch mean cost contrasts at $\alpha=.10$, all four depths. On the same closed-union tasks, columns average $Q_{\rm split}-Q_{\cup,\rm sep}$ (Dedup.), $Q_{\cup,\rm sep}-Q_{\rm joint}$ (Joint), and $Q_{\rm ref}-Q_{\rm joint}$ (Discovery), in score entries. Dedup. and Joint sum to $Q_{\rm split}-Q_{\rm joint}$; Discovery is measured separately. Primary uses $\varepsilon=\lceil.02|S_0|\rceil$; exact uses zero. Empty histories, maps, budgets and stopping targets match. Machine-readable records retain all targets, final lower and upper bounds, zero-minimum indicators and all seven policies. Resource-limited attempts remain in separate records.}
\label{tab:cost-contrasts}
\end{table*}


\paragraph{Matched inspection costs across pooling rules.}
Table~\ref{tab:matched-entry-costs} puts the median gap witness and mean joint minimum on the same entry-cost scale as reference acquisition. Within each pooling rule, each difference uses the same panel, split, depth, maps, cutoff, budget and stopping target. The supplement reports all four common policies, separate zero-query and nonzero-query summaries, and empirical quantiles. Undefined ratios at zero minima remain missing. A zero median excess is compatible with positive mean query counts: for WMT24 en-es median primary stopping, 159 of 160 instances stop without queries, while one requires inspection. These descriptive summaries reuse the saved paths; they add no scorer or optimizer runs.
\begin{table*}[t]
\centering\small
\begin{tabular}{lllrrrrr}
\toprule
Panel & Pool & Stop & $\overline{Q}_{\min}$ & $\overline{Q}_{\rm ref}$ & $\overline{\Delta Q}$ & Med. $\Delta Q$ & Zero / $n$ \\
\midrule
RewardBench 2 & median & primary & 119.6 & 213.4 & 93.8 & 88.0 & 43 / 160 \\
RewardBench 2 & median & exact & 157.5 & 270.2 & 112.7 & 111.0 & 35 / 160 \\
RewardBench 2 & mean & primary & 1,216.9 & 1,491.5 & 274.6 & 216.5 & 0 / 160 \\
RewardBench 2 & mean & exact & 1,297.0 & 1,564.0 & 267.1 & 208.5 & 0 / 160 \\
WMT23, 13 & median & primary & 83.1 & 130.8 & 47.7 & 0.0 & 124 / 160 \\
WMT23, 13 & median & exact & 104.3 & 156.5 & 52.2 & 0.0 & 123 / 160 \\
WMT23, 13 & mean & primary & 2,012.1 & 2,543.9 & 531.8 & 437.5 & 0 / 160 \\
WMT23, 13 & mean & exact & 2,180.2 & 2,722.3 & 542.2 & 475.0 & 0 / 160 \\
WMT23, 12 & median & primary & 15.8 & 26.6 & 10.7 & 0.0 & 150 / 160 \\
WMT23, 12 & median & exact & 26.4 & 48.7 & 22.3 & 0.0 & 145 / 160 \\
WMT23, 12 & mean & primary & 2,024.3 & 2,549.7 & 525.3 & 431.0 & 0 / 160 \\
WMT23, 12 & mean & exact & 2,189.7 & 2,717.2 & 527.5 & 435.5 & 0 / 160 \\
WMT24 en-de & median & primary & 10.6 & 17.4 & 6.8 & 0.0 & 153 / 160 \\
WMT24 en-de & median & exact & 27.3 & 46.7 & 19.4 & 0.0 & 152 / 160 \\
WMT24 en-de & mean & primary & 6,922.7 & 8,221.0 & 1,298.3 & 930.0 & 0 / 160 \\
WMT24 en-de & mean & exact & 7,420.0 & 8,711.0 & 1,291.1 & 909.0 & 0 / 160 \\
WMT24 en-es & median & primary & 3.4 & 8.3 & 4.9 & 0.0 & 159 / 160 \\
WMT24 en-es & median & exact & 8.9 & 19.3 & 10.4 & 0.0 & 154 / 160 \\
WMT24 en-es & mean & primary & 5,968.4 & 7,108.7 & 1,140.3 & 853.5 & 0 / 160 \\
WMT24 en-es & mean & exact & 6,390.6 & 7,551.2 & 1,160.6 & 873.0 & 0 / 160 \\
WMT24 ja-zh & median & primary & 53.2 & 111.1 & 57.9 & 0.0 & 145 / 160 \\
WMT24 ja-zh & median & exact & 98.4 & 174.1 & 75.7 & 0.0 & 136 / 160 \\
WMT24 ja-zh & mean & primary & 8,717.0 & 9,804.2 & 1,087.2 & 897.5 & 0 / 160 \\
WMT24 ja-zh & mean & exact & 9,260.6 & 10,252.2 & 991.5 & 714.5 & 0 / 160 \\
\bottomrule
\end{tabular}
\caption{Matched entry costs for reference priority at $\alpha=.10$, full-reference right ECDF and all four depths. $Q_{\min}$ is the median gap witness or mean joint minimum for the same instance and stopping target; $\Delta Q=Q_{\rm ref}-Q_{\min}$. Primary uses $\varepsilon=\lceil.02|S_0|\rceil$ and exact uses zero. Means include zero-query instances; the last column counts them. The supplement reports per-depth and nonzero-query summaries, quantiles and all four common policies. Ratios exclude zero minima. Overlapping splits support descriptive comparisons.}
\label{tab:matched-entry-costs}
\end{table*}


\paragraph{Initial geometry and discovery overhead.}
We stratify the same saved paths by the initial gap $(V^+_0-V^-_0)/|S_0|$, computed before any candidate-test query. The strata are descriptive and were not used to select a query policy. Writing $g_0=(V^+_0-V^-_0)/|S_0|$, the 960 primary mean targets fall into $0<g_0\le.10$, $.10<g_0\le.50$ and $g_0>.50$, with 10, 466 and 484 targets, respectively. Their mean reference query-to-minimum ratios are 2.711, 1.337 and 1.096; mean excess counts are 752.9, 1,206.7 and 428.4 entries. Thus a smaller initial gap does not identify a uniformly smaller multiplicative overhead on these panels. The strata mix panels and depths and do not establish a causal effect or a deployable prediction rule. The supplement retains every policy, both stops and zero-query cases separately.


\paragraph{Complete budget and tolerance grid.}
We additionally evaluate every original mean instance, without selecting on observed gain. Multiplying the integer budget by $.75,1,1.25$ and the primary tolerance by $0,1,2$ gives 8,640 tasks on the same 960 matrices. Full retention is feasible in 786, 903 and 941 instances at the three budgets, respectively (Table~\ref{tab:budget-geometry}). These instances cannot exhibit strict joint advantage: their optimistic count remains the full selected count, so safety alone certifies the required tolerance. The remaining 174, 57 and 19 instances are eligible for, but need not exhibit, an advantage. Changing tolerance does not change the full-information optimum. For all 8,640 tasks we also enumerate whole-column query subsets and optimize joint and separate-union evidence over all admissible targets and tied component optima. Every joint minimum equals its best separate-union minimum. Independent sorted-row capacity calculations verify the witnesses and rule out all subsets with one fewer column; monotonicity then rules out every smaller subset. Each reported component is also minimum for its target, and its union attains the joint lower bound. The 1,920 original targets agree with the matched column study below. The negative result persists across this complete budget/tolerance grid, but the realizable column example above rules out universal equality. This reuses the same matrices; it supplies neither new entry minima nor independent NLP confirmation.
\begin{table}[t]
\centering\small
\begin{tabular}{rrrrrr}
\toprule
$c$ & Full & Non-full & $0$ & $\varepsilon_0$ & $2\varepsilon_0$ \\
\midrule
0.75 & 786 & 174 & 0 & 0 & 0 \\
1.00 & 903 & 57 & 0 & 0 & 2 \\
1.25 & 941 & 19 & 2 & 3 & 6 \\
\bottomrule
\end{tabular}
\caption{Budget geometry on all 960 original mean instances. Budgets are $\lfloor cb_0\rfloor$ for the displayed multiplier $c$; maps, cutoff and candidates stay fixed. Full and Non-full partition the instances by full-information retention. The last three columns count zero-query gap certificates at each tolerance, with $\varepsilon_0=\lceil.02|S_0|\rceil$. All 8,640 corresponding column tasks have equal joint and best separate-union minima. This grid supplies no newly optimized entry minima.}
\label{tab:budget-geometry}
\end{table}


\paragraph{Previously unused WMT22 panels.}
\label{app:wmt22-confirmation}
Before extracting numeric scores, we froze a separate WMT22 protocol covering all three primary metric-task directions: en-de, en-ru and zh-en \citep{freitag2022metrics}. The cached archive had already been acquired. We admitted complete, finite, system-aligned reference-A segment-score streams, excluded reference outputs and removed exactly duplicate numeric streams. Admission and deduplication read the complete aligned panel before splitting; these preprocessing reads are outside online acquisition counts. After splitting, candidate test scores are available to acquisition only through queries. Alphabetical stream order supplies seven base evaluators and all remaining candidates. The resulting panels contain 21, 20 and 21 metric streams, evaluated on 15, 16 and 18 translation systems, respectively. This tests a previously unevaluated, older benchmark year within the same research workflow, not an independent-team replication or an exclusion of benchmark contamination.

Each panel has one fixed, document-disjoint reference/calibration/test partition. With panel index $h=0,1,2$ in the order above, distinct document IDs are collected in first-occurrence order and shuffled by Python \texttt{random.Random} with seed $20261008+1000h$. For $D$ documents, the first $d=\lfloor D/3\rfloor$ form reference, the next $d$ calibration, and the remainder test: 60/60/61 documents for en-de and en-ru, and 85/85/86 for zh-en. The full test batches contain 9,675, 12,208 and 11,088 outputs. Score directions use Pearson correlation with BLEURT-20 on reference outputs only. Negative correlations reverse the score sign; zero correlations or undefined correlations caused by a constant vector retain the original sign. We then fit right ECDFs and the reference-consensus cutoff, fixing both for each instance. Calibration outputs are reserved and unused by these deterministic certificates. Four depths, $q\in\{.05,.10,.20,.30\}$, two pooling rules and two stops give 48 targets at $\alpha=.10$. Acquisition uses cached scores and no quality labels. The four original policies run throughout; mean pooling additionally uses all three frontier variants.

All 12 mean instances permit full retention. At both stopping targets, exclusion is therefore automatic: joint, additive and best separate-union minima coincide in all 24 comparisons. The matching 24 column comparisons also have no strict gain. Table~\ref{tab:wmt22-confirmation} reports every mean instance and the primary discovery counts. These results are consistent with the full-retention side of the automatic-exclusion prediction; they do not independently reproduce the strict joint advantage seen on part of the original grid. They also show the observed prevalence of sharing gains need not rise with more panels alone. One partition per panel does not support a population-level significance claim.

Independent code reconstructs reference mappings and replacement contributions, replays all 132 acquisition paths (888,799 revealed entries), and checks each first stopping state. All 24 mean entry comparisons are numerically closed; the 12 exact entry minima additionally match an independent integer reduction calculation. Direct enumeration checks 28,704 column predecessors, with monotonicity excluding smaller column sets. Positive-tolerance entry lower bounds still rely on numerical optimization. Protocol, admission manifest, code and all 48 summary rows are indexed under \nolinkurl{review-revision/editorial-evidence-followthrough/confirmation/}. New raw score blocks and full traces remain local pending redistribution clearance; the package contains their hashes, protocol, audit receipts and summary, not the raw blocks. End-to-end reproduction additionally requires the source archive and reconstruction commands listed in \nolinkurl{ACQUISITION_INDEX.md}; the package alone does not supply the raw inputs.
\begin{table*}[t]
\centering\small
\begin{tabular}{llrrrrrr}
\toprule
Panel & $q$ & $|S_0|$ & $Q_{\rm joint}$ & Reference & Width & Frontier & Exact $Q_{\rm joint}$ \\
\midrule
en-de & 0.05 & 385 & 3923 & 4032 & 4241 & 4206 & 4035 \\
en-de & 0.1 & 869 & 8657 & 9123 & 9778 & 9698 & 8951 \\
en-de & 0.2 & 1643 & 8202 & 10151 & 9896 & 9668 & 8741 \\
en-de & 0.3 & 2354 & 7197 & 10534 & 9307 & 9292 & 7981 \\
en-ru & 0.05 & 476 & 4472 & 4645 & 5026 & 4873 & 4602 \\
en-ru & 0.1 & 1084 & 10632 & 10934 & 11951 & 11664 & 10918 \\
en-ru & 0.2 & 2309 & 12054 & 13993 & 14242 & 13843 & 12767 \\
en-ru & 0.3 & 3315 & 8746 & 13980 & 11313 & 11121 & 9762 \\
zh-en & 0.05 & 663 & 7535 & 7837 & 8363 & 8258 & 7745 \\
zh-en & 0.1 & 1189 & 12676 & 13268 & 14209 & 13836 & 13068 \\
zh-en & 0.2 & 2183 & 9902 & 12099 & 12061 & 11599 & 10620 \\
zh-en & 0.3 & 3114 & 6231 & 10254 & 8075 & 7760 & 7260 \\
\bottomrule
\end{tabular}
\caption{WMT22 additional-panel mean-pooling check, one document partition per panel. Query counts use entries. The joint, reference, width and frontier query columns use $V^+-V^-\le\lceil.02|S_0|\rceil$; the final query column uses exact stopping. The selected count $|S_0|$ is an instance property. All results have $\alpha=.10$. Every instance permits full retention and has equal joint and optimal separate-union minima at both stops. Frontier is the joint variant; all other policies and median results remain in the 48-row supplement.}
\label{tab:wmt22-confirmation}
\end{table*}


\paragraph{Partitioned optimization.}
Let $m$ be the active-row count and $L$ its full-information capacity. A feasible safe target lies in $s\in[\max(0,L-\varepsilon),L]$. For $s+\varepsilon<m$, the exclusion constraint is represented by an integer threshold, and every level in the formulation's declared domain is considered. For $s+\varepsilon\ge m$, exclusion is automatic; this branch uses the smallest admissible $s$, since increasing the safety target cannot lower its cost. Each branch retains the original constraints and fixes $s$ and, where required, the exclusion threshold. The admissible branch family covers the original feasible set. An LP relaxation supplies a numerical lower bound; unresolved branches can then receive integer optimization. A branch with no usable bound retains the nonnegative lower bound zero. Solver-declared infeasibility and time limits remain distinct statuses.

For each branch, we keep the largest numerical lower bound among its valid attempts. After removing solver-declared infeasible branches, the global numerical lower bound is the minimum of these branch bounds over an exhaustive branch family; any checked feasible witness gives an upper bound. Validity of this lower bound depends on both the branch bounds and the excluded-branch statuses. Integer witness reconstruction alone does not validate these statuses. This rule is applied at both the threshold and safe-target levels. Integer objective bounds use $\lceil d-10^{-6}\rceil$ for the reported numerical bound $d$; $10^{-6}$ is a rounding convention, not a certified error bound on the solver. Closure requires this global numerical lower bound to reach the checked incumbent. The implementation retains branch domains, attempts and coverage checks in \nolinkurl{external-verification-20261006/solver/}; \nolinkurl{PARTITION_AND_ARTIFACT_INDEX.md} in the code package lists the exact files and replay commands.

\paragraph{Implementation checks.}
Saved results include entry witnesses, global numerical bounds and historical attempts, including failures, time limits, aborts and retries. Final closed-target comparisons are separate from these attempts. Reconstructed witnesses are checked against both capacities; solver-bound closure does not constitute independent human proof review. The code package indexes the composed results and saved runs in \path{review-revision/}\allowbreak\path{editorial-evidence-}\allowbreak\path{implementation/}\allowbreak\path{closing-index.csv}, including the remaining numerical-certification boundary.

\paragraph{Cached-score pipeline costs.}
\label{app:cached-timing}
We time the complete original entry policies at $q=.05$, $\alpha=.10$ and exact stopping on all six panels with mean and seven-member median pooling. The fixed timing grid uses splits 0--7, all test outputs and candidate evaluators, and three serial repetitions, giving 384 panel/pooling/split/policy configurations, each repeated three times (1,152 executions). Panel and policy orders rotate across repetitions. Every query count matches the corresponding original path; each reply and final mask is checked. We report the mean of eight splitwise median times, including zero-query stops.

Measurements use Python 3.9, an Apple M4 Pro and single-thread numerical kernels on a shared research machine. Timed components include reference ECDF mapping, cutoff and base-bound setup, ridge fitting for reference priority, and acquisition through exact stopping. Capacity kernels, mask-prefix construction, cached lookup and flip calculation are also timed separately; remaining acquisition includes policy selection, updates and bookkeeping. Instrumentation overhead remains in the total. Imports, kernel compilation and raw-file preload are outside the timed interval; candidate scores are cached, so no evaluator forward pass is measured. Shared-host timings describe this execution environment, not dedicated-hardware latency.

Table~\ref{tab:cached-pipeline-timing} gives complete pipeline times. Reference priority averages 19.9--146.5 ms across the twelve panel/pooling cells. Fewer queries need not imply lower cached runtime: fitting and maintaining a priority can outweigh the saved lookups. The supplement reports every repeat and component summary. The separate 5,760-attempt offline resource study measures certificate optimization, not these online paths. With a nominal 60-second algorithm budget, outcome MILP numerically closes 1,876 of 1,920 targets (931 primary and 945 exact), entry MILP 1,853 and row DP 86. DP also has the 100,000-state and 10-million-transition caps. These resource-limited runs are separate from the final numerical closures in Table~\ref{tab:full-mean-main}.
\begin{table*}[t]
\centering\small
\begin{tabular}{llrrrr}
\toprule
Panel & Pool & Fixed & Random & Width & Reference \\
\midrule
RewardBench & Median & 12.7 & 12.9 & 12.6 & 28.7 \\
RewardBench & Mean & 14.8 & 14.4 & 16.0 & 19.9 \\
WMT23 (13) & Median & 23.5 & 22.8 & 25.0 & 58.3 \\
WMT23 (13) & Mean & 36.1 & 33.0 & 36.0 & 49.1 \\
WMT23 (12) & Median & 15.7 & 15.6 & 15.5 & 46.3 \\
WMT23 (12) & Mean & 33.0 & 34.6 & 35.4 & 48.4 \\
WMT24 en-de & Median & 54.4 & 54.5 & 54.5 & 142.4 \\
WMT24 en-de & Mean & 99.1 & 96.2 & 107.5 & 146.5 \\
WMT24 en-es & Median & 48.7 & 48.8 & 48.8 & 130.4 \\
WMT24 en-es & Mean & 88.2 & 86.4 & 96.0 & 122.2 \\
WMT24 ja-zh & Median & 33.4 & 33.7 & 33.6 & 82.3 \\
WMT24 ja-zh & Mean & 79.6 & 76.3 & 84.7 & 113.6 \\
\bottomrule
\end{tabular}
\caption{Instrumented cached-score pipeline time in milliseconds at $q=.05$, $\alpha=.10$ and exact stopping. Each cell averages eight splitwise medians of three serial runs. Full batches and original policies are used. Times include reference mapping, base setup and acquisition; only reference priority requires the timed ridge fit. Raw-file preload, imports, compilation and evaluator inference are excluded. Component times are supplied separately.}
\label{tab:cached-pipeline-timing}
\end{table*}


\subsection{Mixed entry and column prices}
\label{app:mixed-prices}
For odd-sized median pools with no initial candidate-test observations, fix entry price $\tau>0$ and column prices $c_r>0$. Let $\mathcal O_A=\mathcal B\times A$ denote the boundary evidence revealed by columns $A$. The full column still reveals all $N$ scores; $c_r$ prices this entire action. With additive prices independent of query order, the minimum mixed cost is
\begin{equation}
 C_\varepsilon^{\rm mix}=
 \min_{A\subseteq[T]}\left\{\sum_{r\in A}c_r+
 \tau Q_\varepsilon(\mathcal O_A)\right\}.
 \label{eq:mixed-cost}
\end{equation}
Under these additive, order-independent prices, remove entry actions already covered by a purchased column and place the column actions first. The observations are unchanged and the cost cannot increase. Proposition~\ref{prop:residual} lower-bounds the remaining entry cost and constructs a witness attaining it. Pure entry and pure column access are included. The formula uses the true matrix to choose evidence offline; it supplies neither an online policy nor a batching-dependent execution model. For mean pooling, replace the median residual minimum by the joint residual minimum from Theorem~\ref{thm:joint}, using active-row column observations; the same column-first argument applies.

For example, take $b=0$, $\varepsilon=0$ and opposition rows $(1,0,0)$ and $(0,0,1)$. Both rows must be excluded. With column prices $(1,3,7)$ and entry price two, pure entry access costs four and pure column access costs eight. Querying the first column and the second row's last entry costs three and certifies the same exact count. The first column is cheaper than its needed entries bought one by one, and entry access avoids paying for the expensive last column. This illustrates why remaining cell evidence matters to mixed access.

Our residual implementation constructs the witnessing cells. A separate action-subset enumeration checks all 8,192 pairs of binary truth matrices and observation masks with $(m,T)=(3,2)$ or $(2,3)$, budgets $b\in\{0,\ldots,mT\}$ and tolerances $\varepsilon\in\{0,\ldots,m\}$, giving 200,704 comparisons. It also checks 800 larger random cases and 4,032 mixed-action optima with three fixed price profiles. Empty initial states match the existing formula in 3,136 cases.

\subsection{Exact NLP column witnesses}
Exact column witnesses use the six core panels with median pooling, full-reference right ECDF and $\alpha=.10$. Four depths and 40 splits give 960 instances. Candidate counts are 7 for RewardBench, 17 for each WMT23 panel, 14 for WMT24 en-de and en-es, and 15 for ja-zh; every family is complete. We replay all 2,880 original exact column stops and compare them with the independently checked witnesses. Table~\ref{tab:column-witness-summary} summarizes nominal $q=.05$; the supplement's \texttt{column-witness-all} grid retains every depth and order. Positive-gap column costs in Appendix~\ref{app:mapping-column} remain a separate target.

At RewardBench nominal $q=.05$, the minimum whole-column witness averages \granRbColumns{} of seven columns and \granRbMinimum\% of candidate-test scores; reference order uses \granRbOnlineColumns{} columns at exact stopping. For thirteen-stream WMT23 at the same depth and exact median target, the minimum whole-column witness reveals \granWmtMinimum\% of $NT$ scores. For these RewardBench and thirteen-stream WMT23 configurations, minimum whole-column certificates still reveal a large fraction of candidate test scores even when evidence is chosen with full knowledge of the matrix. Online counts additionally include the cost of discovering suitable columns.
\begin{table*}[t]
\centering\small


\begin{tabular}{lrrrrrr}
\toprule
Panel & $K_{\rm col}>0$ & $Q_{\rm exact}$ & $K_{\rm col}$ & Min. (\%) & Ref. cols & Ratio \\
\midrule
RewardBench 2 & 40 & 196.6 & 4.93 & 70.4 & 5.63 & 1.183 \\
WMT23, 13 streams & 36 & 415.9 & 7.75 & 45.6 & 8.35 & 1.000 \\
WMT23, 12 streams & 4 & 41.1 & 0.83 & 4.9 & 0.95 & 1.100 \\
WMT24 en-de & 3 & 37.1 & 0.25 & 1.8 & 0.25 & 1.000 \\
WMT24 en-es & 2 & 19.0 & 0.23 & 1.6 & 0.28 & 1.125 \\
WMT24 ja-zh & 13 & 271.6 & 2.45 & 16.3 & 2.63 & 1.000 \\
\bottomrule
\end{tabular}
\caption{Exact median column witnesses at $q=.05$, $\alpha=.10$, full-reference right ECDF. $Q_{\rm exact}$ is the entry witness; $K_{\rm col}$ is the minimum whole-column witness. Counts and Min.\ ($100K_{\rm col}/T$) average all 40 splits; Ratio is the median of per-split $K_{\rm ref}/K_{\rm col}$ over $K_{\rm col}>0$ splits, where $K_{\rm ref}$ is the exact reference-order column count. The $K_{\rm col}>0$ column reports split count. Whole-column score cost is $NK_{\rm col}$. These exact costs have a different target from the positive-gap costs in Table~\ref{tab:column-summary}.}
\label{tab:column-witness-summary}
\end{table*}


\subsection{Whole-column evidence for mean and even medians}
\label{app:mean-even-columns}
We extend the same fixed selection protocols to whole-column access for mean pooling with seven base evaluators and even medians with six or eight. Six panels, three pooling configurations, four depths, 40 splits and two stopping targets give $6\times3\times4\times40\times2=5{,}760$ targets. All candidate columns and selected outputs remain in each instance. A column reveals all $N$ scores and costs $N$ logical entries, including scores outside the selected active rows.

For a column subset $A$, the observed row cost is $\ell_i(A)=\sum_{r\in A}f_{ir}$ and its upper bound is $u_i(A)=\ell_i(A)+(T-|A|)p_i$. We minimize $|A|$ subject to $V(\ell(A))-V(u(A))\le\varepsilon$. Complete cardinality enumeration starts at the valid bound $\lceil Q_{\rm joint}/m\rceil$ when $m>0$, because one column reveals $m$ active entries, and ends at the best feasible policy count. Zero active rows require no queries. Integer witnesses and a separate predecessor-subset calculation check the minima.

The fixed and seeded-random column orders retain their original definitions. Reference column priority sorts columns by the sum of reference-predicted upper-cost reductions over active selected rows, quantized at $10^{-12}$ as in the original column protocol. It uses the same ridge features and penalty, with the actual mean or even-median pool and cutoff. All three orders are evaluated at their first valid stopping states; full scores are used only to evaluate observations and offline minima.

Table~\ref{tab:mean-even-columns} shows nominal $q=.05$ results, including zero-query cases. Exact mean column minima reveal 89.1--99.6\% of candidate-test scores across panels; reference order reveals 90.5--99.6\%. For RewardBench, the corresponding exact minima for six- and eight-member medians are 85.0\% and 50.0\%. These pooling changes also change the selected set and replacement family. Every policy comparison uses the minimum for its own unchanged instance. Complete depth, policy and zero-minimum summaries are included in the supplement.
\begin{table*}[t]
\centering\small
\begin{tabular}{llrrrr}
\toprule
Panel & Pool & Min., primary & Min., exact & Ref., exact & Zero, exact \\
\midrule
RewardBench & Mean & 96.1 & 99.6 & 99.6 & 0/40 \\
RewardBench & Median 6 & 79.1 & 85.0 & 91.2 & 0/40 \\
RewardBench & Median 8 & 42.1 & 50.0 & 60.4 & 1/40 \\
WMT23 (13) & Mean & 90.0 & 92.1 & 92.1 & 0/40 \\
WMT23 (13) & Median 6 & 3.3 & 7.1 & 12.5 & 20/40 \\
WMT23 (13) & Median 8 & 0.2 & 0.6 & 0.8 & 37/40 \\
WMT23 (12) & Mean & 89.6 & 91.5 & 91.5 & 0/40 \\
WMT23 (12) & Median 6 & 6.9 & 10.6 & 15.4 & 27/40 \\
WMT23 (12) & Median 8 & 1.6 & 3.8 & 4.8 & 34/40 \\
WMT24 en-de & Mean & 88.9 & 90.0 & 91.1 & 0/40 \\
WMT24 en-de & Median 6 & 3.3 & 4.7 & 5.3 & 34/40 \\
WMT24 en-de & Median 8 & 0.0 & 0.6 & 0.8 & 38/40 \\
WMT24 en-es & Mean & 87.5 & 89.1 & 90.5 & 0/40 \\
WMT24 en-es & Median 6 & 5.5 & 8.7 & 10.0 & 29/40 \\
WMT24 en-es & Median 8 & 1.0 & 1.5 & 1.5 & 38/40 \\
WMT24 ja-zh & Mean & 92.2 & 94.2 & 94.3 & 0/40 \\
WMT24 ja-zh & Median 6 & 50.6 & 56.7 & 59.4 & 6/40 \\
WMT24 ja-zh & Median 8 & 8.4 & 13.2 & 14.6 & 30/40 \\
\bottomrule
\end{tabular}
\caption{Whole-column certification at $q=.05$ and $\alpha=.10$. Percentages use all $NT$ candidate-test scores, averaging all 40 splits. Min. is the offline minimum; Ref. is the reference column order. Primary tolerance is $\lceil.02|S_0|\rceil$; exact tolerance is zero. Median 6 and 8 average their two central scores. Zero reports exact zero-query splits. Full depth and policy results accompany the supplement.}
\label{tab:mean-even-columns}
\end{table*}


\paragraph{Joint evidence versus separate column optima.}
For all 1,920 mean targets, we compare the joint column minimum with the least expensive union over all separately optimal safety and exclusion column sets. The 1,806 full-retention targets have automatic exclusion. For the remaining 114 targets, enumeration of every column subset gives both capacities; subset envelopes retain the largest safe count and smallest optimistic count for each contained subset cardinality. Minimizing over admissible safe counts and all tied component minima gives the best union. Complete subset enumeration determines component minima. An independent sorted-row calculation additionally verifies the capacities of 8,540 predecessor subsets. Each resulting union attains the matching joint minimum, proving equality on every target, including all 89 strict entry-gain targets. A column costs $N$ entries, so the equality also holds in this cost unit. Table~\ref{tab:evidence-regimes} reports both query models on identical tasks.
\begin{table}[t]
\centering\small
\begin{tabular}{llrrr}
\toprule
Stop & Retention & Targets & Entry & Column \\
\midrule
Primary & Full & 903 & 0 & 0 \\
Primary & Non-full & 57 & 32 & 0 \\
Exact & Full & 903 & 0 & 0 \\
Exact & Non-full & 57 & 57 & 0 \\
\bottomrule
\end{tabular}
\caption{Strict joint gains over the best union of separately optimal witnesses on the original mean-pooling grid. Full means the full-information optimum retains every baseline-selected output. Entry and Column count targets with strict gain using their respective query units. Each stop uses the same 960 instances at $\alpha=.10$; the 40 splits overlap.}
\label{tab:evidence-regimes}
\end{table}


The fixed-action bound and realizable $3/2$ separation above show column gains to be possible in the model. On the original 114 non-full-retention targets, every column is necessary. A row-range condition explains this full-revelation requirement for 101 of them using row contribution ranges. Independent finite checks compare the subset-envelope implementation with literal witness-pair enumeration on 400 seeded cases. A further 1,200 seeded histories, each under entry and column actions with positive integer prices, check the zero-contribution residual equality in Proposition~\ref{prop:residual-no-uniform}. These computations support implementation consistency; the proposition rests on its completion argument.


\paragraph{When every column is necessary.}
Restrict all cost vectors to baseline-selected rows, retain capacity $V$ and budget $b_H$, and start with empty history. For $T\ge1$, define the gap when column $r$ alone is missing by
\[
 \delta_r=V(F-f_{\cdot r})-V(F+p-f_{\cdot r}).
\]
Also let $\underline f_i=\min_r f_{ir}$, $\overline f_i=\max_r f_{ir}$ and
\[
 \gamma_{\rm col}=V(F-\underline f)-V(F+p-\overline f).
\]
Both quantities use the full score matrix to explain offline evidence requirements.
\begin{proposition}[Full-column necessity]
\label{prop:all-columns}
Under the completion model of Section~\ref{sec:partial}, the joint minimum equals $T$ unit-cost column queries if and only if $\min_r\delta_r>\varepsilon$. Whenever this holds, the best separate-union minimum also equals $T$. Moreover, $\gamma_{\rm col}\le\min_r\delta_r$, so $\gamma_{\rm col}>\varepsilon$ is a sufficient condition remaining true under every within-row permutation of candidate scores, with base scores, cutoff, budget and tolerance fixed.
\end{proposition}
\begin{proof}
If column $r$ alone is unrevealed, the row bounds are $F_i-f_{ir}$ and $F_i+p_i-f_{ir}$, giving gap $\delta_r$. If some $\delta_r\le\varepsilon$, this omission certifies the target using $T-1$ columns. Otherwise every one-column omission fails; any proper subset is contained in such an omission and has at least its gap by revealing monotonicity. Full revelation always succeeds, proving the equivalence. A best separate union is a joint certificate and uses at most the $T$ available columns, proving equality of the two minima in this case.

Since $\underline f_i\le f_{ir}\le\overline f_i$ and capacity is nonincreasing in each cost,
\[
 \begin{aligned}
 &V(F-f_{\cdot r})-V(F+p-f_{\cdot r})\\
 &\qquad\ge V(F-\underline f)-V(F+p-\overline f).
 \end{aligned}
\]
Within-row permutations preserve $F,\underline f,\overline f,p$ and hence the sufficient condition.
\end{proof}

The range condition holds for 56 of 57 exact and 45 of 57 primary non-full-retention targets in the six-panel mean grid. The omission criterion establishes full revelation for the remaining 13. Independent enumeration checks the range condition on 66,339 finite integer cases; independent omission calculations check the equivalence on all 570 layout/target comparisons below, with full revelation in 566. These checks supplement the proof. The equivalence is a direct consequence of revealing monotonicity: it characterizes full revelation, not all instances of equality or strict advantage between joint and separate-union certificates.

\paragraph{Changing column layout at fixed row costs.}
We retain all 57 non-full-retention mean instances and both stops. Besides the original layout, we sort each row's candidate contributions increasingly; rotate these sorted rows by their row index modulo $T$; or independently permute each row using PCG64 seeds 104729 and 130363. Each permutation transports the corresponding normalized candidate scores, so the contributions remain realizable under the same base scores and cutoff. These synthetic columns need not represent a single deployed evaluator. Row multisets, $F_i$, $p_i$, $V^*$, budgets and tolerances stay fixed. The rowwise bijection preserves all entry witnesses and their minima; it need not preserve a column witness.

All 570 layout/target comparisons have equal joint and best separate-union minima. Every column is necessary in 566 comparisons. The four exceptions use the aligned layout and primary stopping, where both minima drop by one column and exclusion is automatic. All 114 original-layout targets require every column. Independent sorted-row calculations check 9,236 predecessor subsets and each joint witness. This post hoc intervention shows persistence of the observed column equality under these changes of evidence location; it does not establish universal equality or a new NLP task result. Protocols, layout records and condition coverage are in \nolinkurl{review-revision/granularity-risk-implementation/}.

\subsection{Heterogeneous column-price sensitivity}
\label{app:column-prices}
Whole-evaluator certificate feasibility is unchanged when columns have different prices. For exact stopping, minimize $\sum_{r\in A}c_r$ under Proposition~\ref{prop:column-condition}; for positive tolerance, require the two capacities to differ by at most the same $\varepsilon$ as the original replay.

We retain the six original median panels, four depths, 40 splits, complete candidate families, full-reference right ECDF and $\alpha=.10$. In frozen candidate-index order $r=0,\ldots,T-1$, the three price vectors are $(2T-1)\mathbf1$, $(T+2r)_r$ and its reverse. They share total price $T(2T-1)$; reporting cost divided by $2T-1$ makes the mean column price one. These are declared additive action costs, not hardware measurements. Existing column orders and their stopping states remain fixed.

Full subset optimization gives 5,760 profile/target records. A separate raw-row enumeration checks every cheaper subset; all 960 uniform-price exact minima match the original column witnesses. Tables~\ref{tab:column-cost-exact} and~\ref{tab:column-cost-tolerant} summarize all panel--depth settings across 40 splits. Ties in total witness cost are resolved by minimizing column count. In 30 of the 3,840 nonuniform-price records, minimizing total query cost needs more columns than minimizing column count alone. This holds even when ties among minimum-cost witnesses go to the fewest columns. Enumerating all same-price subsets in these 30 records confirms this secondary minimum. For example, RewardBench split 17 at $q=.10$ and positive tolerance has a five-column ascending-price optimum with raw price 55, or $55/13$ after normalization, whereas its minimum column count is four. Thus minimizing column count and minimizing additive query cost can select different evidence, even under the same certification target. Existing-order priced cost ratios compare paid action costs against the minimum under the same prices and stopping target; they do not measure execution time.
\begin{table*}[t]
\centering\small
\begin{tabular}{lrrrrrrrrr}
\toprule
Panel & $q$ & $C_0$ & $C_\uparrow$ & $C_\downarrow$ & $R_0$ & $R_\uparrow$ & $R_\downarrow$ & $n_+$ & More \\
\midrule
RewardBench & 0.05 & 4.925 & 4.513 & 4.648 & 1.162 & 1.264 & 1.332 & 40 & 0/0 \\
RewardBench & 0.10 & 4.050 & 3.408 & 3.498 & 1.098 & 1.241 & 1.408 & 40 & 2/1 \\
RewardBench & 0.20 & 1.825 & 1.248 & 1.254 & 1.128 & 1.620 & 1.705 & 39 & 0/1 \\
RewardBench & 0.30 & 0.150 & 0.081 & 0.081 & 1.000 & 2.048 & 1.667 & 6 & 0/0 \\
WMT23 (13) & 0.05 & 7.750 & 7.841 & 6.942 & 1.095 & 1.145 & 1.227 & 36 & 2/0 \\
WMT23 (13) & 0.10 & 0.050 & 0.027 & 0.027 & 1.000 & 2.000 & 1.667 & 1 & 0/0 \\
WMT23 (13) & 0.20 & 0.000 & 0.000 & 0.000 & -- & -- & -- & 0 & 0/0 \\
WMT23 (13) & 0.30 & 0.000 & 0.000 & 0.000 & -- & -- & -- & 0 & 0/0 \\
WMT23 (12) & 0.05 & 0.825 & 0.764 & 0.762 & 1.300 & 1.345 & 1.646 & 4 & 1/1 \\
WMT23 (12) & 0.10 & 0.750 & 0.621 & 0.528 & 1.418 & 2.063 & 1.826 & 8 & 2/1 \\
WMT23 (12) & 0.20 & 0.150 & 0.100 & 0.117 & 1.000 & 1.561 & 1.234 & 1 & 0/0 \\
WMT23 (12) & 0.30 & 0.150 & 0.091 & 0.107 & 1.000 & 1.298 & 1.919 & 2 & 0/1 \\
WMT24 en-de & 0.05 & 0.250 & 0.207 & 0.207 & 1.000 & 1.606 & 1.269 & 3 & 0/0 \\
WMT24 en-de & 0.10 & 0.375 & 0.274 & 0.265 & 1.100 & 1.555 & 1.476 & 5 & 1/1 \\
WMT24 en-de & 0.20 & 0.000 & 0.000 & 0.000 & -- & -- & -- & 0 & 0/0 \\
WMT24 en-de & 0.30 & 0.000 & 0.000 & 0.000 & -- & -- & -- & 0 & 0/0 \\
WMT24 en-es & 0.05 & 0.225 & 0.185 & 0.204 & 1.125 & 1.677 & 1.628 & 2 & 0/0 \\
WMT24 en-es & 0.10 & 0.100 & 0.074 & 0.072 & 1.333 & 2.012 & 1.586 & 3 & 0/0 \\
WMT24 en-es & 0.20 & 0.025 & 0.017 & 0.013 & 1.000 & 1.000 & 2.571 & 1 & 0/0 \\
WMT24 en-es & 0.30 & 0.000 & 0.000 & 0.000 & -- & -- & -- & 0 & 0/0 \\
WMT24 ja-zh & 0.05 & 2.450 & 2.305 & 2.241 & 1.101 & 1.436 & 1.337 & 13 & 0/0 \\
WMT24 ja-zh & 0.10 & 1.250 & 1.126 & 1.012 & 1.266 & 1.726 & 1.711 & 9 & 2/0 \\
WMT24 ja-zh & 0.20 & 0.100 & 0.072 & 0.067 & 1.500 & 1.915 & 2.602 & 2 & 0/0 \\
WMT24 ja-zh & 0.30 & 0.000 & 0.000 & 0.000 & -- & -- & -- & 0 & 0/0 \\
\bottomrule
\end{tabular}
\caption{Exact-stop column-cost sensitivity under the original median protocol and $\alpha=.10$. Uniform, ascending and descending prices have the same mean, normalized to one. The three $C$ columns average minimum costs over all 40 splits. The $R$ columns average reference-order cost divided by the matching optimum over the $n_+$ splits with positive minimum cost. All three price profiles have the same $n_+$. More reports ascending/descending split counts where even the fewest-column minimum-cost witness uses more columns than the minimum feasible count. Ties in total witness cost minimize column count. Undefined ratios remain missing. Prices are declared action-cost models, not measured scoring times or fees.}
\label{tab:column-cost-exact}
\end{table*}

\begin{table*}[t]
\centering\small
\begin{tabular}{lrrrrrrrrr}
\toprule
Panel & $q$ & $C_0$ & $C_\uparrow$ & $C_\downarrow$ & $R_0$ & $R_\uparrow$ & $R_\downarrow$ & $n_+$ & More \\
\midrule
RewardBench & 0.05 & 4.300 & 3.785 & 3.888 & 1.210 & 1.419 & 1.457 & 40 & 0/0 \\
RewardBench & 0.10 & 3.475 & 2.792 & 2.848 & 1.109 & 1.268 & 1.539 & 39 & 1/0 \\
RewardBench & 0.20 & 1.375 & 0.856 & 0.844 & 1.054 & 1.646 & 1.832 & 37 & 0/0 \\
RewardBench & 0.30 & 0.025 & 0.013 & 0.013 & 1.000 & 1.857 & 1.857 & 1 & 0/0 \\
WMT23 (13) & 0.05 & 6.350 & 6.431 & 5.561 & 1.098 & 1.155 & 1.303 & 36 & 3/0 \\
WMT23 (13) & 0.10 & 0.000 & 0.000 & 0.000 & -- & -- & -- & 0 & 0/0 \\
WMT23 (13) & 0.20 & 0.000 & 0.000 & 0.000 & -- & -- & -- & 0 & 0/0 \\
WMT23 (13) & 0.30 & 0.000 & 0.000 & 0.000 & -- & -- & -- & 0 & 0/0 \\
WMT23 (12) & 0.05 & 0.625 & 0.558 & 0.557 & 1.262 & 1.586 & 1.795 & 3 & 0/0 \\
WMT23 (12) & 0.10 & 0.350 & 0.254 & 0.224 & 1.300 & 2.035 & 2.254 & 5 & 1/0 \\
WMT23 (12) & 0.20 & 0.100 & 0.061 & 0.065 & 1.000 & 1.775 & 1.419 & 1 & 0/0 \\
WMT23 (12) & 0.30 & 0.050 & 0.036 & 0.043 & 1.500 & 1.646 & 2.088 & 1 & 0/1 \\
WMT24 en-de & 0.05 & 0.175 & 0.148 & 0.143 & 1.000 & 1.160 & 1.321 & 2 & 1/0 \\
WMT24 en-de & 0.10 & 0.150 & 0.093 & 0.089 & 1.000 & 1.605 & 1.991 & 5 & 0/0 \\
WMT24 en-de & 0.20 & 0.000 & 0.000 & 0.000 & -- & -- & -- & 0 & 0/0 \\
WMT24 en-de & 0.30 & 0.000 & 0.000 & 0.000 & -- & -- & -- & 0 & 0/0 \\
WMT24 en-es & 0.05 & 0.125 & 0.109 & 0.130 & 1.600 & 1.712 & 1.643 & 1 & 1/0 \\
WMT24 en-es & 0.10 & 0.000 & 0.000 & 0.000 & -- & -- & -- & 0 & 0/0 \\
WMT24 en-es & 0.20 & 0.000 & 0.000 & 0.000 & -- & -- & -- & 0 & 0/0 \\
WMT24 en-es & 0.30 & 0.000 & 0.000 & 0.000 & -- & -- & -- & 0 & 0/0 \\
WMT24 ja-zh & 0.05 & 1.700 & 1.564 & 1.514 & 1.213 & 1.387 & 1.686 & 9 & 2/0 \\
WMT24 ja-zh & 0.10 & 0.575 & 0.519 & 0.460 & 1.223 & 1.487 & 1.473 & 5 & 3/1 \\
WMT24 ja-zh & 0.20 & 0.025 & 0.016 & 0.013 & 3.000 & 4.789 & 5.533 & 1 & 0/0 \\
WMT24 ja-zh & 0.30 & 0.000 & 0.000 & 0.000 & -- & -- & -- & 0 & 0/0 \\
\bottomrule
\end{tabular}
\caption{Column-cost sensitivity at $V^+-V^-\le\lceil.02|S_0|\rceil$ under the original median protocol and $\alpha=.10$. Uniform, ascending and descending prices have the same mean, normalized to one. The three $C$ columns average minimum costs over all 40 splits. The $R$ columns average reference-order cost divided by the matching optimum over the $n_+$ splits with positive minimum cost. All three price profiles have the same $n_+$. More reports ascending/descending split counts where even the fewest-column minimum-cost witness uses more columns than the minimum feasible count. Ties in total witness cost minimize column count. Undefined ratios remain missing. Prices are declared action-cost models, not measured scoring times or fees.}
\label{tab:column-cost-tolerant}
\end{table*}


\subsection{Controlled score panels and implementation checks}
\label{app:synthetic-sharing}
We generate seven-member median pools at cutoff $.5$, with $m\in\{32,128\}$ selected boundary rows, 16 selected interior rows, 16 invariant unselected rows and 16 unselected boundary rows. Base low and high values are drawn from $U(.1,.3)$ and $U(.65,.95)$. Selected boundary, selected interior, invariant unselected and unselected boundary rows have respectively four, five, two and three high votes. For designated selected boundary rows, a first-member high score is set to $.5$, preserving exactly four scores at or above the cutoff. Every selected boundary row has candidate opposition and agreement scores $.25$ and $.75$, with designated agreements tied at $.5$. The other three row types have candidate scores from $U(.1,.95)$. The complete tie indices, random-number order and separate reference seeds are specified by \texttt{make\_panel} in \nolinkurl{scripts/run\_synthetic\_certificates.py}. 

With $T\in\{7,12\}$, each selected boundary row has $d\in\{1,3\}$ opposition bits. The layouts keep row costs fixed and vary opposition locations. Shared support assigns the same $d$ candidate columns to every row; cyclic support rotates a seeded candidate order across rows; independent support draws a uniform $d$-subset per row. 

Eight seeds and bit budgets $0,\lfloor mT/5\rfloor,\lfloor mT/2\rfloor$ give 192 panels and 576 instances. The nominal count is $k=m$, consistent with $q=m/(m+48)$ in the general definition, and rational $\alpha=2b/(7mT)$ gives flip budget exactly $4b$. Reference panels use separate seeds with the same layout; the cutoff is fixed by construction.

For each layout, we run all four entry-acquisition and all three column-acquisition policies from the NLP study. They keep the same ridge features, regularization, query interfaces and both stopping rules. The 192 panels yield 768 entry-acquisition orders, each reused across three budgets, and 1,728 budget-specific column paths. Two stopping targets give $768\times3\times2+1{,}728\times2=8{,}064$ records. Acquisition cannot access unqueried test scores. Complete matrices are used for evaluation and offline witness construction. A separate replay implementation rebuilds explicit single-member flips and every first stopping state. Complete generated base, candidate and reference matrices, paths, results and witness subsets are saved. No third-party score or quality data are required.

Table~\ref{tab:synthetic-column-witnesses} reports the complete controlled comparison, including zero-budget cases. It shows how evidence positions affect column certificate size under odd-median pooling, complementing the entry/column comparison in Section~\ref{sec:query-granularity}. At $m=128,T=12,d=1,b=307$, the common minimum entry certificate size is 1,229; shared support requires ten columns and the other layouts eleven. The layouts preserve each row's contribution multiset and the minimum entry-certificate size. Their column minima differ because informative entries occupy different candidate columns. The experiment does not assess NLP quality or inference time.
\begin{table*}[t]
\centering\small


\begin{tabular}{rrrrrrrr}
\toprule
$m$ & $T$ & $d$ & $b$ & $Q_{\rm exact}$ & Shared & Cyclic & Independent \\
\midrule
32 & 7 & 1 & 0 & 32 & 1.00 & 7.00 & 7.00 \\
32 & 7 & 1 & 44 & 180 & 6.00 & 7.00 & 7.00 \\
32 & 7 & 1 & 112 & 112 & 4.00 & 5.00 & 4.00 \\
32 & 7 & 3 & 0 & 32 & 1.00 & 3.00 & 3.63 \\
32 & 7 & 3 & 44 & 150 & 7.00 & 7.00 & 7.00 \\
32 & 7 & 3 & 112 & 112 & 4.00 & 7.00 & 6.13 \\
32 & 12 & 1 & 0 & 32 & 1.00 & 12.00 & 11.50 \\
32 & 12 & 1 & 76 & 308 & 10.00 & 11.00 & 11.00 \\
32 & 12 & 1 & 192 & 192 & 6.00 & 7.00 & 7.00 \\
32 & 12 & 3 & 0 & 32 & 1.00 & 4.00 & 5.00 \\
32 & 12 & 3 & 76 & 319 & 12.00 & 12.00 & 12.00 \\
32 & 12 & 3 & 192 & 192 & 6.00 & 8.00 & 8.00 \\
128 & 7 & 1 & 0 & 128 & 1.00 & 7.00 & 7.00 \\
128 & 7 & 1 & 179 & 717 & 6.00 & 7.00 & 7.00 \\
128 & 7 & 1 & 448 & 448 & 4.00 & 5.00 & 4.13 \\
128 & 7 & 3 & 0 & 128 & 1.00 & 3.00 & 4.25 \\
128 & 7 & 3 & 179 & 618 & 7.00 & 7.00 & 7.00 \\
128 & 7 & 3 & 448 & 448 & 4.00 & 7.00 & 6.25 \\
128 & 12 & 1 & 0 & 128 & 1.00 & 12.00 & 12.00 \\
128 & 12 & 1 & 307 & 1229 & 10.00 & 11.00 & 11.00 \\
128 & 12 & 1 & 768 & 768 & 6.00 & 7.00 & 7.00 \\
128 & 12 & 3 & 0 & 128 & 1.00 & 4.00 & 7.25 \\
128 & 12 & 3 & 307 & 1300 & 12.00 & 12.00 & 12.00 \\
128 & 12 & 3 & 768 & 768 & 6.00 & 8.00 & 8.00 \\
\bottomrule
\end{tabular}
\caption{Controlled exact-stop column-evidence sharing for seven-member median pools. Each row represents eight seeds in each of three layouts with identical $m,T,d,b$ and minimum entry-witness size. For each seed, we verify a certificate of size $K_{\rm col}$ and rule out every smaller column subset. Shared, Cyclic and Independent report mean minimum column counts across eight seeds; row costs and the entry minimum are identical across layouts. All 576 instances use the complete entry and column acquisition solvers. The synthetic inputs test query geometry, not NLP quality.}
\label{tab:synthetic-column-witnesses}
\end{table*}


\section{Retention Tolerance and Required Evidence}
\label{app:tolerance}
Let $\eta$ denote the nominal count-tolerance fraction. We vary $\eta\in\{0,.005,.01,.02,.05,.10\}$, stopping when the integer gap is at most $\varepsilon_\eta=\lceil\eta|S_0|\rceil$. The original $.02$ remains the main configuration. The six core panels retain the original seven-member pool, full-reference right ECDF, $\alpha=.10$, both pooling rules, four depths, four acquisition methods and forty splits. Acquisition choices do not depend on $\varepsilon$, so the saved paths through exact stopping determine every new first stop. This yields 7,680 paths, 46,080 stopping records and 1,152 configuration means. WMT24 retains its fixed-point priority ties. These are post hoc analyses on exposed panels.

Every cached reply is checked against the corresponding full score matrix. A separate sorted-cost checker reconstructs each returned mask and budget, both count bounds and the full-information optimum. For each stopping state reached after at least one query, the immediately preceding state must have a gap greater than the tolerance. The original $.02$ and exact stopping masks and positions match the saved results. All new masks meet the same churn budget, and realized retained-count regret never exceeds $\varepsilon_\eta$.

Nominal tolerance is not the proportion of incorrect returned outputs. It permits a retained-count gap without relaxing $\alpha$. For $|S_0|>0$, integer rounding gives
\[
 \eta\le\frac{\varepsilon_\eta}{|S_0|}<\eta+\frac1{|S_0|}.
\]
Different nominal tolerances can therefore have identical stopping targets on small lists. Table~\ref{tab:tolerance-witness} reports the effective proportion. Inspection fractions, undefined ratios and split aggregation follow Appendix~\ref{app:protocol}.

The supplement's \texttt{tolerance-all-*} grids show all depths and acquisition orders. At nominal $q=.05$, reference priority with median pooling on RewardBench inspects 71.1\%, 61.4\% and 33.6\% of active entries at nominal tolerances 0, .02 and .10, with selected coverage 99.6\%, 97.3\% and 89.8\%. Thirteen-stream WMT23 has corresponding active-inspection fractions 54.3\%, 45.1\% and 20.9\%. The lower query counts come from relaxed retained-count tolerances; they are not savings for an unchanged exact-certification target.

With mean pooling, reference priority still queries a large fraction of active entries across the tolerance grid. At nominal $q=.05$, reference-priority active inspection across the six panels ranges from 84.9--90.2\% at exact stopping and from 72.4--76.2\% at nominal tolerance .10. On thirteen-stream WMT23 with mean pooling at $q=.30$, reference priority requires more queries than width at all six tolerances. Table~\ref{tab:tolerance-order-check} makes the ranking change explicit: width has the lowest mean at tolerance $.02$, and fixed order at $.10$, among these four policies.
\begin{table}[t]
\centering\small
\begin{tabular}{rrrrr}
\toprule
$\eta$ & Fixed & Random & Width & Ref. \\
\midrule
.02 & 2067.3 & 2111.6 & 1800.6 & 2363.8 \\
.10 & 567.4 & 667.3 & 660.9 & 940.7 \\
\bottomrule
\end{tabular}
\caption{Mean queries on thirteen-stream WMT23, mean pooling, $q=.30$, $\alpha=.10$: the same 40 splits and fixed paths, with stopping gap $\lceil\eta|S_0|\rceil$. Counts are reaggregated from saved records and rounded half up. Width has the lowest mean at $.02$, fixed order at $.10$.}
\label{tab:tolerance-order-check}
\end{table}

\begin{table*}[t]
\centering\small\setlength{\tabcolsep}{4pt}
\begin{tabular}{lrrrrrrr}
\toprule
Panel & $\eta$ & $\varepsilon/|S_0|$ (\%) & $Q_\varepsilon$ & Online queries & Ratio & Positive $Q_\varepsilon$ & Zero-query (\%) \\
\midrule
RewardBench 2 & 0 & 0.0 & 196.6 & 327.6 & 1.895 & 40 & 0.0 \\
RewardBench 2 & 0.005 & 0.9 & 182.6 & 313.6 & 1.924 & 40 & 0.0 \\
RewardBench 2 & 0.01 & 1.6 & 169.2 & 300.3 & 1.951 & 40 & 0.0 \\
RewardBench 2 & 0.02 & 2.5 & 155.0 & 285.0 & 1.992 & 40 & 0.0 \\
RewardBench 2 & 0.05 & 5.4 & 124.0 & 233.7 & 1.950 & 39 & 2.5 \\
RewardBench 2 & 0.1 & 10.4 & 77.9 & 160.2 & 2.110 & 36 & 10.0 \\
\addlinespace
WMT23, thirteen streams & 0 & 0.0 & 415.9 & 622.6 & 1.679 & 36 & 10.0 \\
WMT23, thirteen streams & 0.005 & 0.8 & 400.2 & 590.0 & 1.660 & 36 & 10.0 \\
WMT23, thirteen streams & 0.01 & 1.5 & 349.7 & 559.2 & 1.718 & 36 & 10.0 \\
WMT23, thirteen streams & 0.02 & 2.4 & 332.3 & 523.2 & 1.698 & 36 & 10.0 \\
WMT23, thirteen streams & 0.05 & 5.4 & 234.6 & 409.1 & 1.795 & 34 & 15.0 \\
WMT23, thirteen streams & 0.1 & 10.3 & 137.4 & 252.2 & 1.834 & 26 & 35.0 \\
\addlinespace
WMT23, twelve streams & 0 & 0.0 & 41.1 & 56.0 & 1.952 & 4 & 90.0 \\
WMT23, twelve streams & 0.005 & 1.0 & 38.8 & 52.0 & 1.957 & 4 & 90.0 \\
WMT23, twelve streams & 0.01 & 1.4 & 37.1 & 49.1 & 2.018 & 4 & 90.0 \\
WMT23, twelve streams & 0.02 & 2.4 & 34.2 & 43.3 & 1.954 & 3 & 92.5 \\
WMT23, twelve streams & 0.05 & 5.4 & 28.3 & 33.2 & 1.714 & 2 & 95.0 \\
WMT23, twelve streams & 0.1 & 10.4 & 6.3 & 18.3 & 2.901 & 1 & 97.5 \\
\addlinespace
WMT24 en-de & 0 & 0.0 & 37.1 & 60.2 & 1.775 & 3 & 92.5 \\
WMT24 en-de & 0.005 & 0.6 & 33.5 & 54.2 & 1.586 & 2 & 95.0 \\
WMT24 en-de & 0.01 & 1.1 & 31.4 & 51.0 & 1.597 & 2 & 95.0 \\
WMT24 en-de & 0.02 & 2.2 & 27.2 & 44.1 & 1.604 & 2 & 95.0 \\
WMT24 en-de & 0.05 & 5.1 & 14.3 & 22.6 & 1.642 & 2 & 95.0 \\
WMT24 en-de & 0.1 & 10.1 & 2.6 & 4.6 & 1.779 & 1 & 97.5 \\
\addlinespace
WMT24 en-es & 0 & 0.0 & 19.0 & 46.4 & 2.172 & 2 & 95.0 \\
WMT24 en-es & 0.005 & 0.6 & 17.4 & 42.2 & 2.430 & 1 & 97.5 \\
WMT24 en-es & 0.01 & 1.2 & 16.0 & 38.4 & 2.401 & 1 & 97.5 \\
WMT24 en-es & 0.02 & 2.1 & 13.5 & 33.2 & 2.457 & 1 & 97.5 \\
WMT24 en-es & 0.05 & 5.1 & 5.5 & 11.6 & 2.123 & 1 & 97.5 \\
WMT24 en-es & 0.1 & 10.1 & 0.0 & 0.0 & -- & 0 & 100.0 \\
\addlinespace
WMT24 ja-zh & 0 & 0.0 & 271.6 & 371.8 & 1.887 & 13 & 67.5 \\
WMT24 ja-zh & 0.005 & 0.7 & 254.5 & 344.0 & 1.964 & 13 & 67.5 \\
WMT24 ja-zh & 0.01 & 1.2 & 242.9 & 320.3 & 1.939 & 11 & 72.5 \\
WMT24 ja-zh & 0.02 & 2.2 & 150.6 & 276.4 & 2.383 & 9 & 77.5 \\
WMT24 ja-zh & 0.05 & 5.2 & 105.4 & 161.2 & 1.915 & 5 & 87.5 \\
WMT24 ja-zh & 0.1 & 10.1 & 24.0 & 49.6 & 2.046 & 2 & 95.0 \\
\addlinespace
\bottomrule
\end{tabular}
\caption{Tolerance-specific median entry witnesses at $q=.05$, $\alpha=.10$, with reference-priority acquisition. Witness and online counts average all 40 splits. Ratio averages splitwise $Q_{\rm online}/Q_\varepsilon$ over positive witnesses, whose count is reported separately. Zero witnesses have undefined ratios. Effective tolerance includes integer rounding. These descriptive ratios are not online competitive guarantees.}
\label{tab:tolerance-witness}
\end{table*}
\FloatBarrier


For median paths, Corollary~\ref{cor:gap-witness} supplies the matching offline entry witness $Q_{\varepsilon_\eta}$. At nominal $q=.05$, its RewardBench mean falls from 196.6 entries at exact stopping to 155.0 at $.02$ and \toleranceRBten{} at $.10$. The corresponding thirteen-stream WMT23 means fall from 415.9 to 332.3 and \toleranceWMTten{} entries. Table~\ref{tab:tolerance-witness} separates intrinsic requirements from multiplicative entry discovery ratios. For each split with $Q_{\varepsilon_\eta}>0$, we compute $Q_{\rm online}/Q_{\varepsilon_\eta}$ and then average these ratios across included splits, reporting their number. A ratio of average query counts is a different statistic. A positive-gap query count below $Q_0$ does not violate a lower bound, because the matching minimum is $Q_{\varepsilon_\eta}$ for target gap $\varepsilon_\eta$. These median closed forms do not characterize mean-pooling or whole-evaluator minima. Theorem~\ref{thm:joint} handles joint mean entries; Appendix~\ref{app:full-mean-grid} reports numerically closed full-batch mean minima for the primary and exact targets. Other tolerances have no matching full-grid mean-entry optimization here.

The accompanying records also retain actual regret, certificate gap, zero-query stops, panel fractions, exact-path query savings and each stop's mask. Relative query savings, normalized by the exact-stopping count, are undefined when this count is zero. No tolerance or acquisition method is selected retrospectively to replace the main protocol.


\section{Choosing the Selection and Information Protocol}
\label{app:practical}
Fix the reference mapping, consensus rule, cutoff and replacement family before certification. Changing the mapping or evaluator pool can change the baseline selected set and stability question. When it does, fewer queries under the new protocol do not establish an improvement for the original list. The existing mapping ablations support sensitivity analysis, not a universal midrank recommendation.

Choose the query unit supported by the scoring interface: entries target individual outputs, whereas columns reveal all $N$ scores. For an operational comparison, include reference fitting, base scores, preprocessing, calibration when used and certificate computation. The accounting boundary is specified in Appendix~\ref{app:protocol}; deployment claims require measured execution costs.

Choose the churn budget from the number of membership changes $\rho$ tolerated per replacement, so $\alpha=\rho/(2k)$. The average target lets individual replacements exceed $\rho$ (Appendix~\ref{app:individual-replacements}); the per-replacement target enforces $\rho$ for each one (Appendix~\ref{app:per-replacement}). Choose retention tolerance in integer outputs. If at most $r$ outputs of retained-count shortfall are acceptable, set $\varepsilon=r$. To permit shortfall of at most a fraction $\eta$ of $|S_0|$, use $\varepsilon=\lfloor\eta|S_0|\rfloor$. Our primary stop uses $\lceil.02|S_0|\rceil$ instead, which can exceed 2\% on small lists. Exact stopping certifies the optimal retained count without necessarily revealing all scores; a positive gap bounds retained-count shortfall. A positive tolerance changes the stopping target without relaxing the churn budget $\alpha$. Report invariant retention separately from noninvariant outputs admitted within the budget. Restricting the commitment mask for a subsequent content policy preserves the fixed-denominator budget but can reduce the returned count (Appendix~\ref{app:risk-interpretation}).

If a content or calibrated acceptance rule fixes an eligible set $A_{\rm cal}$ visible before certification, apply the capacity calculation to $S_0\cap A_{\rm cal}$. The same monotone-cost argument gives $V^-_A\le V^*_A\le V^+_A$ and a gap guarantee relative to this restricted full-information optimum. It does not preserve near-optimality over the original $S_0$. Pruning a calibrated mask preserves an expected bound only for the same fixed-denominator, nonnegative loss under its calibration assumptions. A churn guarantee supplies no semantic-quality bound. Report post-certification quality audits separately. If audit labels select a stopping point, statistical quality claims must account for this choice. Under the fixed selection protocol, deterministic churn and count-gap guarantees hold at any chosen state satisfying the certificate conditions of Proposition~\ref{prop:stop}.

Fixed additive column prices can change the cheapest evidence without changing the certification target (Appendix~\ref{app:column-prices}). Real batching, model loading and shared computation can instead make prices depend on the action history. Deployment claims require a verified scorer revision and interface, measured setup and inference costs, and the same target and retained-count accounting. The current priced-witness study does not replace these measurements.


\section{Prediction and Certificate Discovery}
\label{app:acquisition-diagnostics}
\subsection{Prediction accuracy and the stopping objective}
Reference priority uses predicted flips to choose queries, while stopping depends on retention capacities under lower and upper row costs. Low prediction error alone does not identify the observations certifying safety or excluding larger sets. For an active selected row, a queried candidate reduces upper cost by $p_i-f_{ir}$ and raises lower cost by $f_{ir}$. For odd-sized median pools, agreements and oppositions contribute to different sides of the certificate. Ranking rows by predicted total cost does not identify useful evidence across candidates or which rows determine the capacity bounds.

We analyze the six original NLP panels, both pooling rules, four depths and forty splits with full-reference right ECDF and $\alpha=.10$. Four original entry orders give 7,680 paths. Evaluating each at $\eta\in\{0,.02,.10\}$ and $\varepsilon_\eta=\lceil\eta|S_0|\rceil$ gives 23,040 stopping states and 5,760 matched reference-versus-width comparisons. Paths, mappings, cutoffs and stopping rules remain unchanged. Full test matrices are used only for post hoc diagnostics, following the information boundary and aggregation conventions in Appendix~\ref{app:protocol}.

\subsection{Diagnostic definitions and checks}
Let $\hat f_{ir}$ be the clipped ridge prediction and $\mathcal A=\{i\in S_0:p_i>0\}$. Entry prediction error is
\[
 E_{\rm entry}=\frac{1}{mT}\sum_{i\in\mathcal A,r}\frac{|\hat f_{ir}-f_{ir}|}{p_i}.
\]
It is undefined when $m=0$. Initial row-total rank correlation is the Spearman correlation of $\sum_r\hat f_{ir}$ and $F_i$, using average ranks for ties over $\mathcal A$. Constant vectors or fewer than two active rows give an undefined correlation; tables report the number of defined splits.

For reference first-stop count $Q>0$, evaluate error at $t=Q-1$, before the first query meeting the stopping criterion. For $Q=0$, use the initial state $t=0$. Equivalently, $t=\max(0,Q-1)$. Let $U_i(t)$ contain unqueried candidates and $\mathcal J(t)$ the active rows with at least one such candidate. Estimated total cost is $\hat F_i(t)=\ell_i(t)+\sum_{r\in U_i(t)}\hat f_{ir}$. Remaining-row error is
\[
 E_{\rm rem}(t)=\frac{1}{|\mathcal J(t)|}\sum_{i\in\mathcal J(t)}\frac{|\hat F_i(t)-F_i|}{Tp_i}.
\]
It is undefined if no row remains. Its row set depends on the path and target. The error concerns row-total estimates on rows with unqueried entries, normalized by the full row scale $Tp_i$. Because the rows and prefix depend on acquisition and stopping, this diagnostic is not an independent measure of initial prediction accuracy or a causal explanation of query cost.

Paired differences use $\Delta Q=Q_{\rm reference}-Q_{\rm width}$ for the same panel, split, depth, pooling rule and tolerance. Negative values favor reference. For median pooling, excess entry queries are $Q_{\rm reference}-Q_{\varepsilon_\eta}$, relative to the offline witness for the same target. The original full-panel mean ratios use the matching exact subset in Appendix~\ref{app:joint-scale}; the expanded grid covers both targets in Appendix~\ref{app:full-mean-grid}. The small-batch minima in Appendix~\ref{app:joint-validation} remain separate denominators. Records also include fixed-order comparisons, strict wins against all three other methods, remaining-row correlations and within-configuration prediction-error/query-difference correlations. Undefined values remain missing. Forty overlapping splits give descriptive summaries, not independent replications for significance tests.

We check all 22,317,820 cached replies against full matrices and reconstruct the retained-count bounds at every saved stopping state. Reference tie rules are checked at 5,693 diagnostic prefixes, including WMT24 fixed-point priorities. Across 1,920 reference paths, we deduplicate the first-stop positions and their predecessors for three tolerances, retaining positions before the exact stop where a next query exists: 904 median and 4,789 mean prefixes. A separate implementation recomputes ridge predictions and average-rank correlations on 48 configurations, matches all 23,040 original stopping counts and checks 2,016 summary values.


\subsection{Acquisition towards certificate frontiers}
\label{app:frontier}
Frontier priority reuses the reference-only ridge predictions $\hat f_{ir}$, clipped to $[0,p_i]$. The policy operates on active selected rows. Let $V_A$ be capacity restricted to these rows; equivalently, extend their cost vector with zero invariant costs and subtract $m_0$ from $V$. Define the priority only at a query state with $v^-<v^+$, where $v^-=V_A(u)$ and $v^+=V_A(\ell)$ are active counts. Let $\hat F_i=\ell_i+\sum_{r\in U_i}\hat f_{ir}$. Sort $\hat F_i$ with the original label-blind tie order, and set
\[
 s=\min\{v^+,\max\{v^-+1,V_A(\hat F)\}\}.
\]
The first $s$ predicted rows form $A_s$. Because $s>v^-$, no $s$-row set is upper-cost feasible. Their safety deficit is therefore $D_s=\sum_{i\in A_s}u_i-b_H>0$. The priority targets exclusion of $s+1$ active rows, independently of stopping tolerance. When $p=s+1\le m$, set $t=\lceil\hat F_{(p)}\rceil$ and define the exclusion deficit
\[
 D_e=\sum_i(t-\ell_i)_+-pt+b_H+1.
\]
For an unqueried entry, its priority is
\[
 \frac{\one[i\in A_s](p_i-\hat f_{ir})}{D_s}
 +\frac{\min\{\hat f_{ir},(t-\ell_i)_+\}}{D_e},
\]
where the exclusion term is omitted if $p>m$ or $D_e\le0$. The two terms estimate fractional progress towards feasible safety and exclusion proofs, including when the integer count gap stays unchanged. They are point-prediction heuristics, not expected query-optimality claims. Safety-only retains the first term; exclusion-only retains the second. If all priorities are zero, use largest width with the same tie order. Predictions and priorities use nearest-even $10^{-12}$ ticks; ties follow the seeded label-blind row priority, then candidate index. Before acquisition, active rows are sorted by the original tie permutation generated with seed $s_{\rm panel}+\mathrm{split}+700001$. Estimated costs add prediction ticks to observed integer costs converted to the same $10^{-12}$ scale. Acquisition is independent of the stopping tolerance, so one path supplies first primary and exact stops. Each selection recomputes the predicted ordering and scores, taking $O(m\log m+mT)$ time and $O(mT)$ working space. Only the query callback reveals a test contribution. All policies retain the same deterministic interval and stopping guarantees.

\paragraph{Matched acquisition and ablations.}
On all 960 mean instances, joint, safety-only and exclusion-only frontier paths each continue through exact stopping and record the first primary stop. No quality label or unqueried test contribution enters acquisition.

At the primary stop, joint frontier averages \frontierQueries{} entries, compared with \frontierRefQueries{} for reference and \frontierWidthQueries{} for width. It beats reference on only \frontierWins{} of 960 instances, with \frontierLosses{} losses and three ties. The lower mean reflects larger savings on its wins: wins average 420.0 entries and losses 174.0. At $q=.30$, joint frontier wins all 40 splits on every panel except WMT24 ja-zh. At $q=.05$, it loses on at least 34 of 40 splits on every panel. Safety-only averages \frontierSafetyQueries{}, slightly below joint, while exclusion-only averages \frontierExclusionQueries{}. The aggregate comparison thus shows no overall advantage from adding the exclusion term. Tables~\ref{tab:frontier-primary} and~\ref{tab:frontier-exact} give all depths, including unfavorable results. Unlike the original budget-independent orders, frontier uses the current budget in its deficits; cross-budget paths cannot be reused.

\paragraph{When exclusion is required.}
To test the exclusion term where exclusion is required, Table~\ref{tab:frontier-strata} stratifies the saved comparisons by the full-information optimum $V^*$, baseline count $|S_0|$ and tolerance. If $V^*+\varepsilon<|S_0|$, even the largest feasible safe set cannot meet the gap target while the optimistic count remains $|S_0|$, so exclusion evidence is necessary. If $V^*<|S_0|\le V^*+\varepsilon$, sufficient safety evidence can meet the target without lowering this count. The third stratum permits full retention. These conditions are fixed within each matched comparison. They use hidden full-information costs for post hoc diagnosis, not for acquisition, and differ from whether the heuristic's exclusion term is active at a particular history.

At primary stopping, joint frontier saves an average of 20.28 entries over safety-only on the 32 targets requiring exclusion, with 25 wins, six losses and one tie. It uses 3.72 more entries on average on the 903 full-retention targets and 6.40 more on the remaining 25. At exact stopping, it saves 20.21 entries on average on the 57 targets requiring exclusion, but loses on ten. Within this post hoc stratum, the exclusion term thus lowers mean query cost without improving every instance. The stratum uses hidden costs and supplies no online switching rule. These reused, overlapping-split comparisons establish no independent-task confirmation or stepwise causal mechanism.
\begin{table*}[t]
\centering\small
\begin{tabular}{llrrrr}
\toprule
Stop & Instance condition & Targets & Mean $\Delta Q$ & Median $\Delta Q$ & Wins/losses/ties\\
\midrule
Primary & $V^*+\varepsilon<|S_0|$ & 32 & -20.28 & -11 & 25/6/1 \\
Primary & $V^*=|S_0|$ & 903 & +3.72 & +0 & 2/65/836 \\
Primary & $V^*<|S_0|\le V^*+\varepsilon$ & 25 & +6.40 & +3 & 10/14/1 \\
Exact & $V^*+\varepsilon<|S_0|$ & 57 & -20.21 & -7 & 40/10/7 \\
Exact & $V^*=|S_0|$ & 903 & +2.63 & +0 & 5/55/843 \\
\bottomrule
\end{tabular}
\caption{Matched frontier ablation by instance condition on the original mean panels, at $\alpha=.10$ and all four depths. $\Delta Q$ is the joint-frontier query count minus the safety-only count, so negative values and wins favor joint frontier. Both counts use the same instance and stopping target. Primary uses $\varepsilon=\lceil.02|S_0|\rceil$; exact uses zero. The non-full, tolerance-covered stratum is empty at exact stopping. Conditions use full-information $V^*$ for post hoc diagnosis only. These descriptive comparisons reuse saved paths and do not treat overlapping splits as independent samples.}
\label{tab:frontier-strata}
\end{table*}

\paragraph{Tolerance-aware exclusion target.} Frontier priority targets exclusion of $s+1$ rows, whereas Theorem~\ref{thm:joint} requires excluding $s+\varepsilon+1$. A variant targets $s+\varepsilon+1$ and keeps every other rule; with $\varepsilon=0$ it coincides with the original policy, so only primary stopping changes. On the 960 primary targets it averaged 5,209.7 queries, against 5,210.6 for joint and 5,207.6 for safety-only frontier. On the 32 targets requiring exclusion it averaged 8,175.6 queries, against 8,163.3 for joint and 8,183.6 for safety-only frontier. Against joint frontier it had 10 wins, 19 losses and three ties. Aligning the heuristic target with the stopping tolerance therefore does not explain the online results. Records are in \nolinkurl{outputs/scope-extensions/frontier_eps_primary.jsonl}.
\paragraph{Closed-form guidance for odd medians.}
\label{app:median-guided}
For odd medians, Theorem~\ref{thm:safe} and Corollary~\ref{cor:gap-witness} fix the shape of a minimum witness. The witness holds agreements inside a safe target set and oppositions up to a common cap. A closed-form-guided policy fills this shape with reference-ridge opposition propensities $\hat\theta_{ir}=\hat f_{ir}/(h+1)$, clipped to $[0,1]$ and used as probabilities. Write $v^-=V^--m_0$ for the active safe count. At each state, let $\hat z_i$ add the observed oppositions of row $i$ to the predicted opposition probabilities of its unqueried entries. The policy sets $\hat L$ to the capacity of $\hat z$ under budget $b$, with fractional costs, clipped to $[v^-,v^+]$. For each $s$ from $\max(v^-,\hat L-\varepsilon)$ to $\max(v^-,\hat L)$, the target set holds the $s$ rows with the smallest $\hat z_i$, with label-blind tie order. The safety deficit $D_s$ is the excess of the set's upper opposition counts over $b$. Its safety cost divides $D_s$ by the mean agreement probability of the set's unqueried entries. When $p=s+\varepsilon+1\le m$, the cap $t$ is the smallest level at which the $p$ smallest values $\min(\hat z_i,t)$ sum to at least $b+1$. Its exclusion cost counts the entries needed to raise each row's observed oppositions to $t$. It divides this count by the mean opposition probability of the eligible entries and is zero once the observed exclusion deficit $D_e$ vanishes. The policy keeps the $s$ with the smallest summed cost. It then scores each unqueried entry by agreement probability over $D_s$ inside the target set. Rows with fewer than $t$ observed oppositions add opposition probability over $D_e$. Ties go to the label-blind row order and then candidate index; if every score is zero, the policy queries the row of largest width. The policy uses the stopping tolerance, so primary and exact stops use separate paths. On all 960 median instances at $\alpha=.10$, it averaged 87.5 entries at primary stopping. Reference priority averaged 84.6, width 98.1 and the matching minimum $Q_\varepsilon$ 47.6. It beat reference on 41 instances and lost on 142. At $q=.05$ it averaged 189.7 entries versus 200.9 for reference, but it lost on average at every larger depth. At exact stopping it also lost to reference on average, with 124.6 versus 119.3 entries. Knowing the witness shape did not translate into fewer queries with these predictions. For exact stopping, Proposition~\ref{prop:fixed-discovery} rules out a batch-size-independent competitive factor for any odd-median policy, even at the study's fixed parameters. Records are in \nolinkurl{outputs/scope-extensions/median_guided_records.jsonl}.
\paragraph{Implementation checks.}
For the three mean-pooling frontier variants above, the 2,880 paths contain 5,760 stopping states. Replay checks 15,838,138 replies, every stopping mask and the failure to stop at the preceding state. A separate priority calculation checks 14,832 choices over 144 paths spanning all panels and depths.
\begin{table*}[t]
\centering\small
\setlength{\tabcolsep}{4pt}
\begin{tabular}{lrrrrrrrr}
\toprule
Panel & $q$ & Reference & Width & Frontier & S-only & E-only & W/L Ref. & W/L Width \\
\midrule
RewardBench 2 & 0.05 & 727.9 & 770.6 & 752.6 & 743.0 & 770.5 & 0/39 & 36/3 \\
RewardBench 2 & 0.10 & 1289.6 & 1356.2 & 1317.6 & 1316.1 & 1357.0 & 3/37 & 40/0 \\
RewardBench 2 & 0.20 & 1919.6 & 1925.8 & 1839.4 & 1839.4 & 1925.8 & 40/0 & 40/0 \\
RewardBench 2 & 0.30 & 2028.7 & 1781.6 & 1716.9 & 1716.9 & 1781.6 & 40/0 & 40/0 \\
WMT23, 13 & 0.05 & 1873.7 & 2013.2 & 1921.8 & 1904.6 & 2012.0 & 5/34 & 38/0 \\
WMT23, 13 & 0.10 & 3518.8 & 3782.5 & 3637.3 & 3637.3 & 3782.5 & 0/39 & 40/0 \\
WMT23, 13 & 0.20 & 2419.2 & 2167.6 & 2063.8 & 2063.8 & 2167.6 & 39/1 & 40/0 \\
WMT23, 13 & 0.30 & 2363.8 & 1800.6 & 1732.9 & 1732.9 & 1800.6 & 40/0 & 37/3 \\
WMT23, 12 & 0.05 & 1757.7 & 1841.1 & 1800.2 & 1787.5 & 1840.1 & 4/36 & 36/4 \\
WMT23, 12 & 0.10 & 3047.4 & 3239.9 & 3138.9 & 3138.2 & 3239.9 & 3/37 & 40/0 \\
WMT23, 12 & 0.20 & 2413.8 & 2273.8 & 2203.0 & 2203.0 & 2273.8 & 40/0 & 38/2 \\
WMT23, 12 & 0.30 & 2979.8 & 2520.6 & 2437.4 & 2437.4 & 2520.6 & 40/0 & 38/2 \\
WMT24 en-de & 0.05 & 5619.0 & 5920.0 & 5815.1 & 5805.8 & 5920.0 & 3/37 & 37/3 \\
WMT24 en-de & 0.10 & 6538.0 & 6785.0 & 6625.3 & 6625.2 & 6785.3 & 17/23 & 38/2 \\
WMT24 en-de & 0.20 & 9591.7 & 9911.6 & 9730.2 & 9730.2 & 9911.6 & 12/28 & 40/0 \\
WMT24 en-de & 0.30 & 11135.4 & 10569.2 & 10208.6 & 10208.6 & 10569.2 & 40/0 & 39/1 \\
WMT24 en-es & 0.05 & 4962.1 & 5228.2 & 5130.8 & 5109.5 & 5227.8 & 2/38 & 36/4 \\
WMT24 en-es & 0.10 & 6038.1 & 6272.6 & 6092.9 & 6092.9 & 6272.6 & 13/27 & 40/0 \\
WMT24 en-es & 0.20 & 8509.9 & 8836.9 & 8639.3 & 8639.3 & 8836.9 & 7/33 & 38/2 \\
WMT24 en-es & 0.30 & 8924.9 & 8309.6 & 8037.0 & 8037.0 & 8309.6 & 40/0 & 35/5 \\
WMT24 ja-zh & 0.05 & 4402.5 & 4590.0 & 4533.1 & 4531.0 & 4552.7 & 1/39 & 28/10 \\
WMT24 ja-zh & 0.10 & 7834.2 & 8182.0 & 8068.1 & 8069.8 & 8157.2 & 1/39 & 36/4 \\
WMT24 ja-zh & 0.20 & 12544.0 & 13133.2 & 12904.6 & 12905.7 & 13054.0 & 7/33 & 35/5 \\
WMT24 ja-zh & 0.30 & 14436.1 & 15180.5 & 14707.1 & 14707.1 & 15180.5 & 6/34 & 40/0 \\
\bottomrule
\end{tabular}
\caption{Mean-pooling acquisition at $V^+-V^-\le\lceil .02|S_0|\rceil$, $\alpha=.10$, on complete original batches. Query columns average entry counts across all 40 splits. Frontier sums predicted fractional progress toward safety and exclusion; S-only and E-only retain the corresponding single term. W/L counts Frontier strict wins/losses against Reference or Width; remaining splits tie. All methods use the same normalized scores, bound updates and certification target. Prediction-based methods share reference-fitted models; none uses quality labels. These descriptive comparisons do not establish a universally best policy.}
\label{tab:frontier-primary}
\end{table*}

\begin{table*}[t]
\centering\small
\setlength{\tabcolsep}{4pt}
\begin{tabular}{lrrrrrrrr}
\toprule
Panel & $q$ & Reference & Width & Frontier & S-only & E-only & W/L Ref. & W/L Width \\
\midrule
RewardBench 2 & 0.05 & 759.2 & 787.8 & 771.2 & 764.6 & 787.8 & 5/34 & 32/2 \\
RewardBench 2 & 0.10 & 1341.4 & 1387.1 & 1354.7 & 1353.8 & 1387.7 & 7/33 & 40/0 \\
RewardBench 2 & 0.20 & 2006.6 & 2006.4 & 1906.4 & 1906.4 & 2006.4 & 40/0 & 40/0 \\
RewardBench 2 & 0.30 & 2148.8 & 1906.5 & 1819.1 & 1819.1 & 1906.5 & 40/0 & 40/0 \\
WMT23, 13 & 0.05 & 1940.4 & 2053.1 & 1964.4 & 1953.6 & 2053.0 & 11/28 & 40/0 \\
WMT23, 13 & 0.10 & 3647.6 & 3856.5 & 3715.9 & 3715.9 & 3856.5 & 2/38 & 40/0 \\
WMT23, 13 & 0.20 & 2615.6 & 2352.5 & 2254.1 & 2254.1 & 2352.5 & 40/0 & 40/0 \\
WMT23, 13 & 0.30 & 2685.8 & 2108.1 & 2046.9 & 2046.9 & 2108.1 & 40/0 & 36/4 \\
WMT23, 12 & 0.05 & 1823.0 & 1879.4 & 1839.6 & 1833.6 & 1879.2 & 13/26 & 37/2 \\
WMT23, 12 & 0.10 & 3160.6 & 3311.7 & 3218.7 & 3217.8 & 3311.7 & 4/36 & 39/1 \\
WMT23, 12 & 0.20 & 2611.1 & 2447.8 & 2377.0 & 2377.0 & 2447.8 & 40/0 & 38/1 \\
WMT23, 12 & 0.30 & 3274.4 & 2816.3 & 2715.6 & 2715.6 & 2816.3 & 40/0 & 39/1 \\
WMT24 en-de & 0.05 & 5809.8 & 6017.4 & 5921.8 & 5917.3 & 6017.4 & 6/33 & 35/3 \\
WMT24 en-de & 0.10 & 6847.4 & 7007.2 & 6885.7 & 6884.6 & 7007.2 & 17/22 & 37/3 \\
WMT24 en-de & 0.20 & 10201.6 & 10388.9 & 10236.8 & 10236.8 & 10388.9 & 20/20 & 34/5 \\
WMT24 en-de & 0.30 & 11985.3 & 11327.7 & 10967.6 & 10967.6 & 11327.7 & 40/0 & 38/2 \\
WMT24 en-es & 0.05 & 5126.9 & 5316.6 & 5227.0 & 5211.0 & 5316.6 & 4/36 & 34/6 \\
WMT24 en-es & 0.10 & 6320.6 & 6471.5 & 6337.9 & 6337.9 & 6471.5 & 15/25 & 40/0 \\
WMT24 en-es & 0.20 & 9054.4 & 9263.8 & 9094.9 & 9094.9 & 9263.8 & 17/23 & 37/3 \\
WMT24 en-es & 0.30 & 9703.0 & 9008.4 & 8735.0 & 8735.0 & 9008.4 & 40/0 & 35/5 \\
WMT24 ja-zh & 0.05 & 4554.2 & 4672.9 & 4628.0 & 4637.5 & 4670.3 & 12/28 & 29/10 \\
WMT24 ja-zh & 0.10 & 8093.6 & 8332.5 & 8246.4 & 8251.1 & 8331.0 & 7/33 & 36/3 \\
WMT24 ja-zh & 0.20 & 13115.0 & 13549.6 & 13380.4 & 13382.7 & 13548.5 & 10/29 & 35/5 \\
WMT24 ja-zh & 0.30 & 15245.9 & 15846.2 & 15291.7 & 15291.7 & 15846.2 & 21/19 & 40/0 \\
\bottomrule
\end{tabular}
\caption{Mean-pooling acquisition at $V^+=V^-$, $\alpha=.10$, on complete original batches. Query columns average entry counts across all 40 splits. Frontier sums predicted fractional progress toward safety and exclusion; S-only and E-only retain the corresponding single term. W/L counts Frontier strict wins/losses against Reference or Width; remaining splits tie. All methods use the same normalized scores, bound updates and certification target. Prediction-based methods share reference-fitted models; none uses quality labels. These descriptive comparisons do not establish a universally best policy.}
\label{tab:frontier-exact}
\end{table*}


\subsection{Prediction accuracy and observed query cost}
Tables~\ref{tab:acquisition-median} and~\ref{tab:acquisition-mean} report every depth at the main tolerance. Across the three original mean-pooling panels, reference priority uses fewer entries on average than width at $q=.05,.10$, but more at $q=.30$. These are configuration means, not a claim about every split. On thirteen-stream WMT23 with mean pooling at $q=.05$, the mean splitwise initial rank correlation is .280. Reference priority uses 139.5 fewer queries than width on average and wins all 40 splits. At $q=.30$, the mean correlation is .747, but reference priority uses 563.2 more queries on average and loses all 40 splits. Reference uses 2,363.8 queries versus width\textquotesingle s 1,800.6. At $q=.05$ and $q=.30$, respectively, entry errors are .213 and .260; remaining-row errors at the defined diagnostic states are .138 and .139. Depth changes the cutoff, selected rows and certification target, so this comparison does not isolate a cause. Initial row-total rank correlation alone does not identify the rule with the fewest queries in these configurations.

On RewardBench with median pooling at $q=.05$, initial row-total correlation is $-.260$, yet reference beats width on 26 of 40 splits and saves 7.2 queries on average. Query choices also depend on candidate-level predicted upper-cost reductions and observed replies; stopping depends on which rows determine the count bounds. These summaries do not establish harm from better predictions or identify an optimal priority. Table~\ref{tab:query-prefix-example} remains a single-split illustration of different query choices. All three tolerances and paired records accompany the analysis code.
\begin{table*}[t]
\centering\small


\begin{tabular}{lrrrrrr}
\toprule
Panel & $q$ & Entry error & $\rho(n)$ & Remaining error & $\Delta Q$ & Wins / 40 \\
\midrule
RewardBench 2 & 0.05 & 0.508 & -0.260 (40) & 0.159 & -7.2 & 26 \\
RewardBench 2 & 0.10 & 0.501 & -0.255 (40) & 0.180 & -20.4 & 32 \\
RewardBench 2 & 0.20 & 0.490 & -0.141 (40) & 0.186 & -6.0 & 19 \\
RewardBench 2 & 0.30 & 0.491 & -0.080 (40) & 0.206 & -0.1 & 1 \\
\addlinespace
WMT23, 13 streams & 0.05 & 0.456 & 0.031 (40) & 0.138 & -212.2 & 32 \\
WMT23, 13 streams & 0.10 & 0.502 & 0.097 (40) & 0.159 & 0.0 & 0 \\
WMT23, 13 streams & 0.20 & 0.504 & 0.122 (40) & 0.136 & 0.0 & 0 \\
WMT23, 13 streams & 0.30 & 0.494 & -0.042 (40) & 0.118 & 0.0 & 0 \\
\addlinespace
WMT23, 12 streams & 0.05 & 0.516 & 0.141 (40) & 0.185 & -0.6 & 2 \\
WMT23, 12 streams & 0.10 & 0.505 & -0.018 (40) & 0.160 & -8.1 & 4 \\
WMT23, 12 streams & 0.20 & 0.506 & 0.024 (40) & 0.135 & 0.9 & 0 \\
WMT23, 12 streams & 0.30 & 0.493 & -0.125 (40) & 0.121 & 3.0 & 0 \\
\addlinespace
WMT24 en-de & 0.05 & 0.476 & 0.026 (38) & 0.197 & -18.3 & 2 \\
WMT24 en-de & 0.10 & 0.496 & -0.047 (40) & 0.174 & -13.6 & 4 \\
WMT24 en-de & 0.20 & 0.501 & 0.024 (40) & 0.152 & 0.0 & 0 \\
WMT24 en-de & 0.30 & 0.501 & 0.036 (40) & 0.150 & 0.0 & 0 \\
\addlinespace
WMT24 en-es & 0.05 & 0.483 & -0.020 (40) & 0.185 & 0.9 & 0 \\
WMT24 en-es & 0.10 & 0.498 & -0.104 (40) & 0.170 & 0.0 & 0 \\
WMT24 en-es & 0.20 & 0.504 & 0.120 (40) & 0.150 & 0.0 & 0 \\
WMT24 en-es & 0.30 & 0.507 & 0.038 (40) & 0.134 & 0.0 & 0 \\
\addlinespace
WMT24 ja-zh & 0.05 & 0.514 & -0.213 (40) & 0.199 & -13.4 & 7 \\
WMT24 ja-zh & 0.10 & 0.540 & -0.054 (40) & 0.279 & -24.2 & 4 \\
WMT24 ja-zh & 0.20 & 0.523 & 0.018 (40) & 0.259 & -5.9 & 1 \\
WMT24 ja-zh & 0.30 & 0.514 & -0.022 (40) & 0.240 & 0.0 & 0 \\
\bottomrule
\end{tabular}
\caption{Prediction and discovery diagnostics for median pooling, full-reference right ECDF, $\alpha=.10$ and $\varepsilon=\lceil.02|S_0|\rceil$. Each error statistic is defined on all 40 splits and averages them; query differences also average 40 splits. $\rho(n)$ is mean initial row-total rank correlation over the $n$ defined splits. Remaining error is normalized by the full row scale $Tp_i$ and uses $t=\max(0,Q-1)$, the initial state for zero-query stops; $\Delta Q$ is the reference-minus-largest-width entry count, not a percentage. Wins count strict query reductions against width. Zero-query stops may contribute to prediction diagnostics when active rows remain.}
\label{tab:acquisition-median}
\end{table*}

\begin{table*}[t]
\centering\small


\begin{tabular}{lrrrrrr}
\toprule
Panel & $q$ & Entry error & $\rho(n)$ & Remaining error & $\Delta Q$ & Wins / 40 \\
\midrule
RewardBench 2 & 0.05 & 0.235 & 0.297 (40) & 0.152 & -42.7 & 40 \\
RewardBench 2 & 0.10 & 0.231 & 0.581 (40) & 0.137 & -66.6 & 40 \\
RewardBench 2 & 0.20 & 0.219 & 0.718 (40) & 0.128 & -6.2 & 21 \\
RewardBench 2 & 0.30 & 0.225 & 0.746 (40) & 0.144 & 247.1 & 0 \\
\addlinespace
WMT23, 13 streams & 0.05 & 0.213 & 0.280 (40) & 0.138 & -139.5 & 40 \\
WMT23, 13 streams & 0.10 & 0.194 & 0.687 (40) & 0.158 & -263.7 & 40 \\
WMT23, 13 streams & 0.20 & 0.226 & 0.705 (40) & 0.141 & 251.5 & 1 \\
WMT23, 13 streams & 0.30 & 0.260 & 0.747 (40) & 0.139 & 563.2 & 0 \\
\addlinespace
WMT23, 12 streams & 0.05 & 0.228 & 0.591 (40) & 0.147 & -83.3 & 40 \\
WMT23, 12 streams & 0.10 & 0.191 & 0.694 (40) & 0.155 & -192.5 & 40 \\
WMT23, 12 streams & 0.20 & 0.249 & 0.720 (40) & 0.141 & 140.0 & 4 \\
WMT23, 12 streams & 0.30 & 0.264 & 0.753 (40) & 0.134 & 459.3 & 0 \\
\addlinespace
WMT24 en-de & 0.05 & 0.215 & 0.309 (40) & 0.125 & -301.0 & 40 \\
WMT24 en-de & 0.10 & 0.195 & 0.589 (40) & 0.135 & -247.0 & 36 \\
WMT24 en-de & 0.20 & 0.240 & 0.686 (40) & 0.146 & -319.9 & 32 \\
WMT24 en-de & 0.30 & 0.237 & 0.685 (40) & 0.145 & 566.2 & 0 \\
\addlinespace
WMT24 en-es & 0.05 & 0.200 & 0.522 (40) & 0.124 & -266.2 & 40 \\
WMT24 en-es & 0.10 & 0.203 & 0.619 (40) & 0.139 & -234.6 & 37 \\
WMT24 en-es & 0.20 & 0.234 & 0.711 (40) & 0.153 & -326.9 & 39 \\
WMT24 en-es & 0.30 & 0.243 & 0.710 (40) & 0.143 & 615.3 & 0 \\
\addlinespace
WMT24 ja-zh & 0.05 & 0.266 & 0.509 (40) & 0.153 & -187.5 & 40 \\
WMT24 ja-zh & 0.10 & 0.254 & 0.645 (40) & 0.161 & -347.8 & 40 \\
WMT24 ja-zh & 0.20 & 0.245 & 0.535 (40) & 0.187 & -589.2 & 36 \\
WMT24 ja-zh & 0.30 & 0.237 & 0.685 (40) & 0.202 & -744.4 & 40 \\
\bottomrule
\end{tabular}
\caption{Prediction and discovery diagnostics for mean pooling, full-reference right ECDF, $\alpha=.10$ and $\varepsilon=\lceil.02|S_0|\rceil$. Each error statistic is defined on all 40 splits and averages them; query differences also average 40 splits. $\rho(n)$ is mean initial row-total rank correlation over the $n$ defined splits. Remaining error is normalized by the full row scale $Tp_i$ and uses $t=\max(0,Q-1)$, the initial state for zero-query stops; $\Delta Q$ is the reference-minus-largest-width entry count, not a percentage. Wins count strict query reductions against width. Zero-query stops may contribute to prediction diagnostics when active rows remain.}
\label{tab:acquisition-mean}
\end{table*}

\subsection{Same-instance row and candidate interventions}
\label{app:query-interventions}
Across all 1,920 original panel--pooling--depth--split configurations, we replaced the reference policy's row estimates, candidate priorities, or both with true flip contributions. True-row priority is $\ell_i+\sum_{r\in U_i}f_{ir}=F_i$; it uses full hidden row totals. True-candidate priority orders unqueried candidates by increasing $f_{ir}$, keeping the row rule otherwise unchanged. The Both variant combines these two interventions. Ties keep the reference policy\textquotesingle s seeded row order and candidate-index order. These oracle-priority ablations use hidden test scores and are diagnostics, not deployable methods. Within each instance only the query priorities change; the score matrix, maps, cutoff, eligible rows, initial bounds, budgets and certificate conditions stay fixed.

The reference, true-row, true-candidate and combined paths run through exact stopping and record the first primary stop, giving 7,680 paths and 15,360 states. Width comparisons use the corresponding original acquisition records. The reference method reproduces all 3,840 original exact and primary query counts. An independent checker validates 21,632,445 replies, every stopping mask and first stopping state. It also reconstructs 416,920 query choices by direct row and candidate sorting on all 48 split-zero configurations.

Table~\ref{tab:query-interventions} reports mean-pooling query counts averaged over 40 splits at each depth. On thirteen-stream WMT23 with mean pooling at $q=.30$, true-row priority saves queries in all 40 splits relative to reference, but its mean remains 379.7 queries above width. True-candidate priority alone is worse on 36 splits, better on one and tied on three. At $q=.05$, true-row priority is worse on 31 and better on nine splits. True row costs rank outputs by realized churn; query value also depends on current safety deficits, achievable bound reductions and, when exclusion is active, observations ruling out larger sets. True contributions therefore do not make this fixed priority rule query-optimal. These interventions change priority within an instance without attributing each loss to a particular deficit or implying general unhelpfulness of better predictions. Complete results, including median and exact-stop gains and losses, are in \nolinkurl{outputs/query-mechanism-interventions/records.csv}.
\begin{table*}[t]
\centering\small
\begin{tabular}{lrrrrrr}
\toprule
Panel & $q$ & Reference & Width & True row & True candidate & Both \\
\midrule
RewardBench 2 & 0.05 & 727.9 & 770.6 & 712.7 & 715.1 & 712.0 \\
RewardBench 2 & 0.10 & 1289.6 & 1356.2 & 1261.0 & 1281.5 & 1260.2 \\
RewardBench 2 & 0.20 & 1919.6 & 1925.8 & 1874.5 & 1914.8 & 1874.1 \\
RewardBench 2 & 0.30 & 2028.7 & 1781.6 & 1956.5 & 2043.4 & 1956.3 \\
WMT23, 13 streams & 0.05 & 1873.7 & 2013.2 & 1892.9 & 1862.9 & 1892.1 \\
WMT23, 13 streams & 0.10 & 3518.8 & 3782.5 & 3498.5 & 3474.3 & 3497.0 \\
WMT23, 13 streams & 0.20 & 2419.2 & 2167.6 & 2349.9 & 2445.1 & 2349.7 \\
WMT23, 13 streams & 0.30 & 2363.8 & 1800.6 & 2180.3 & 2396.4 & 2180.0 \\
WMT23, 12 streams & 0.05 & 1757.7 & 1841.1 & 1744.7 & 1726.6 & 1743.8 \\
WMT23, 12 streams & 0.10 & 3047.4 & 3239.9 & 3013.7 & 3015.1 & 3012.6 \\
WMT23, 12 streams & 0.20 & 2413.8 & 2273.8 & 2341.4 & 2438.5 & 2340.9 \\
WMT23, 12 streams & 0.30 & 2979.8 & 2520.6 & 2794.5 & 3030.0 & 2793.9 \\
WMT24 en-de & 0.05 & 5619.0 & 5920.0 & 5582.0 & 5524.8 & 5581.3 \\
WMT24 en-de & 0.10 & 6538.0 & 6785.0 & 6398.8 & 6544.7 & 6398.4 \\
WMT24 en-de & 0.20 & 9591.7 & 9911.6 & 9341.0 & 9663.8 & 9340.2 \\
WMT24 en-de & 0.30 & 11135.4 & 10569.2 & 10633.1 & 11310.4 & 10633.0 \\
WMT24 en-es & 0.05 & 4962.1 & 5228.2 & 4923.6 & 4907.8 & 4922.6 \\
WMT24 en-es & 0.10 & 6038.1 & 6272.6 & 5913.7 & 6028.0 & 5913.1 \\
WMT24 en-es & 0.20 & 8509.9 & 8836.9 & 8319.1 & 8563.8 & 8318.7 \\
WMT24 en-es & 0.30 & 8924.9 & 8309.6 & 8534.1 & 9060.6 & 8533.9 \\
WMT24 ja-zh & 0.05 & 4402.5 & 4590.0 & 4399.6 & 4336.6 & 4398.6 \\
WMT24 ja-zh & 0.10 & 7834.2 & 8182.0 & 7738.2 & 7687.3 & 7736.9 \\
WMT24 ja-zh & 0.20 & 12544.0 & 13133.2 & 12318.3 & 12202.0 & 12316.7 \\
WMT24 ja-zh & 0.30 & 14436.1 & 15180.5 & 14236.8 & 14304.0 & 14235.6 \\
\bottomrule
\end{tabular}
\caption{Same-instance mean-pooling interventions on complete original batches. Entries are mean candidate--output query counts over 40 splits, computed from integer sums with decimal half-up rounding. All use $\alpha=.10$ and first stop at $V^+-V^-\le\lceil .02|S_0|\rceil$. True row replaces predicted row totals in reference priority; True candidate uses true within-row flip contributions; Both combines them. These full-information diagnostics keep maps, cutoff, initial endpoint bounds, update rules, budget and stopping target fixed. They are not deployable policies. Median and exact-stop summaries accompany the code.}
\label{tab:query-interventions}
\end{table*}

\FloatBarrier



\section{Summarization: evidence, discovery and measured scoring cost}
\label{app:summeval-confirmation}
This experiment evaluates certificate size and discovery cost on summarization, using the 100-document, sixteen-system SummEval collection \citep{fabbri-etal-2021-summeval}. Before inspecting certification outcomes, we fixed a document-disjoint 33/33/34 reference/calibration/test split with seed 20261011. We also fixed four depths $q\in\{.05,.10,.20,.30\}$, $\alpha=.10$, both pooling rules and the primary and exact stopping criteria. Calibration documents are reserved and unused. The test batch contains 544 outputs. We report all prespecified acquisition configurations without filtering on query savings. This is a new task evaluation after method development, not a claim of SummEval's absence from evaluator training.

\subsection{Frozen cached-score confirmation}
Five metric families supply BARTScore CNN source-to-summary \citep{yuan2021bartscore}, BERTScore F \citep{zhang2020bertscore}, MoverScore \citep{zhao-etal-2019-moverscore}, Prism average \citep{thompson-post-2020-automatic} and ROUGE-L F \citep{lin-2004-rouge} scores. The cached matrix is obtained from \texttt{DAMO-NLP-SG/LLM\_summeval}, revision \texttt{bd7dde15817a} (full revision in the supplement). This experiment reuses released metric scores after identity and text-alignment checks; it does not rerun the original scorers. All 1,600 document/system identities align with the original SummEval annotations. Text alignment uses case, whitespace, punctuation and Penn Treebank bracket normalization; this does not establish identical scoring preprocessing. The raw source is retained locally with its provenance.

All ten choices of three base metrics and two candidates are evaluated. Higher scores are used as better by metric definition; human labels do not orient scores, select queries or choose stops. Each composition refits reference-only right ECDF maps, cutoff and ridge priorities. Ten compositions, four depths and two pooling rules give 80 instances and 160 certification targets. Seven acquisition policies give 560 complete paths and 1,120 stopping states. The integer budget is $\lfloor2kMT\alpha\rfloor$ with $M=3,T=2$. Independent replay checks 83,886 replies and every first stopping state.

Table~\ref{tab:summeval-strata} applies the same ex post capacity stratification used for the earlier NLP panels. The full-retention stratum has $V^*=|S_0|$; the exclusion stratum has $V^*+\varepsilon<|S_0|$; the remaining near-full stratum has $V^*<|S_0|\le V^*+\varepsilon$. These conditions are computed after acquisition and are not inputs to the policy. Under mean pooling, joint frontier saves two queries on average in the exclusion stratum at either stop. It increases the mean query count in the full-retention stratum. Under median pooling, joint frontier does not lower the average query count relative to safety-only in the exclusion stratum. Across the eight near-full mean primary targets, joint frontier uses 0.75 more queries on average: one target uses fewer, four tie and three use more than safety-only. Thus the mean-pooling saving in the exclusion stratum extends to this task conditionally; it is neither a universal advantage nor a per-instance guarantee. The compositions and depths share one test partition and are not independent replications.
\begin{table*}[t]
\centering\small
\begin{tabular}{lllrrrrr}\toprule
Pooling & Stop & Stratum & Targets & Less & Equal & More & Mean query difference\\\midrule
mean & exact & exclusion & 24 & 7 & 15 & 2 & -2.000 \\
mean & exact & full & 16 & 0 & 11 & 5 & 1.188 \\
mean & primary & exclusion & 16 & 7 & 7 & 2 & -2.000 \\
mean & primary & full & 16 & 0 & 10 & 6 & 1.875 \\
mean & primary & near full & 8 & 1 & 4 & 3 & 0.750 \\
median & exact & exclusion & 22 & 5 & 11 & 6 & 0.136 \\
median & exact & full & 18 & 1 & 12 & 5 & 1.222 \\
median & primary & exclusion & 18 & 4 & 10 & 4 & 0.167 \\
median & primary & full & 18 & 2 & 11 & 5 & 1.722 \\
median & primary & near full & 4 & 2 & 0 & 2 & 1.750 \\
\bottomrule\end{tabular}
\caption{Joint frontier versus safety-only acquisition on the 160 cached-score SummEval targets: ten evaluator compositions, four depths, two pooling rules and two stopping rules at $\alpha=.10$, with primary $\varepsilon=\lceil.02|S_0|\rceil$ and exact $\varepsilon=0$. The final column is joint-frontier minus safety-only entry queries. Less/equal/more count query differences, not certificate-minimum differences. All 160 targets are included; strata use full information only after acquisition.}
\label{tab:summeval-strata}\end{table*}


For the 48 non-full-retention mean targets, each joint and best separately optimal union program closes its own incumbent--lower-bound gap numerically. Reconstructed integer witnesses satisfy both count constraints. Joint evidence is smaller on 27 targets and equal on 21; numerical bound closure is not an exact rational minimality proof. Re-solving all 48 unions with HiGHS presolve disabled reproduced every value. The 68 full-retention targets across both pooling rules admit a direct integer computation described below. Odd-median targets also use the closed-form certificates, and enumeration of all four candidate-column subsets gives the minimum column count for $T=2$. Column enumeration does not establish entry-query minimality.

\paragraph{Exact computation in the full-retention subcase.}
If the true total selected-row cost is at most $b_H$, every partial lower-cost vector permits all selected rows. On $m$ active rows, the certificate therefore requires safety for at least $m-\min(m,\varepsilon)$ rows; exclusion is automatic. For row $i$, sort the reductions $p_i-f_{ir}$ in descending order and let $c_i(a)$ be $Tp_i$ minus the sum of the first $a$ reductions. Among all $a$ observations on this row, this is the smallest attainable upper cost. A dynamic program over processed rows, dropped-row count and total queries stores the minimum sum of retained upper costs. A retained row contributes $c_i(a)$ for each $a\in\{0,\ldots,T\}$; a dropped row contributes zero cost and zero queries. The smallest query count with at most $\varepsilon$ drops and cost at most $b_H$ is the exact minimum in the rectangular completion model.

To justify the recurrence, queries on dropped rows are unnecessary, and replacing a retained row's queried set by its largest reductions cannot worsen safety. Conversely, each dynamic-program transition specifies realizable observations and a retained subset. The recurrence therefore covers an optimal certificate and reconstructs a feasible one at its reported cost. Pairing the minimum safety witness with empty exclusion yields the same best separate union. This computation covers all 68 full-retention targets; exhaustive comparison on 60 small instances checks the implementation, not the general argument. It provides no characterization of arbitrary non-full-retention instances.

\subsection{Actual scorer execution}
\label{app:real-scoring}
To measure whether reduced score acquisition lowers execution time, a second experiment recomputes scores from the source texts on the same document split. It fixes BARTScore, BERTScore and MoverScore as base metrics and Prism and ROUGE-L as candidates. This is a fresh panel, not a reproduction of the cached matrix. BERTScore uses unrescaled RoBERTa-large layer-17 F with IDF disabled; MoverScore uses the original BERT-MNLI unigram power-mean implementation. Its two IDF dictionaries are fitted separately, on reference-split system summaries and on the human summaries of reference-split documents. Prism uses the bidirectional average of model m39v1 with a single reference. ROUGE-L uses \texttt{rouge-score} 0.1.2 with stemming and a single reference. BARTScore uses \texttt{facebook/bart-large-cnn} for source-to-summary likelihood. The scorer wrappers follow \texttt{neulab/BARTScore} revision \texttt{248f511cb34a} (full revision in the supplement); the supplement records exact model revisions and dependencies.

One A40 GPU executes FP32 scoring with TF32 disabled, PyTorch 2.4.1/CUDA 12.4 and four CPU threads. Each run loads models from local weights, scores all 528 reference outputs with five metrics and all 544 test outputs with three base metrics, fits maps and priorities, and then runs the specified candidate-scoring policy. Full scoring calls each metric once per scoring phase, processing inputs internally in batches of eight. Acquisition calls the candidate scorer once per requested entry without reading a precomputed candidate-test matrix. All methods include the common preparation, including the ridge fit. Eight methods, four depths, two pooling rules and three repetitions give 192 executions in a seeded random order. Each execution records the primary stop on its path to the exact stop, so the two stopping summaries reuse the same executions. Model downloads, dependency installation and post-path rescore audits are outside the measured pipeline interval.

All 192 runs have identical mapped reference and base arrays, despite small floating-point differences in raw scores. Every queried entry is rescored in a batch after the path finishes, and all mapped scores agree. Saved traces reconstruct every first stopping state. One early full-scoring run lacks the full candidate-array archive; its timing and counters remain in the complete grid, while the other 23 full runs retain these arrays. A score-entry count, a metric function invocation and an internal neural-network forward pass are distinct units. Tables~\ref{tab:real-scoring-primary} and~\ref{tab:real-scoring-exact} report the first two along with elapsed time; detailed forward counters and every run accompany the analysis.
\begin{table*}[t]\centering\small
\begin{tabular}{lrrrrr}\toprule
Method & Candidate entries & Metric calls & Total seconds & Time/full & Allocated USD\\\midrule
Full scoring & 1088.00 & 10.00 & 108.18 & 1.000 & 0.01797 \\
Reference & 120.25 & 128.25 & 89.97 & 0.832 & 0.01495 \\
Width & 123.25 & 131.25 & 90.43 & 0.836 & 0.01502 \\
Fixed & 125.25 & 133.25 & 89.51 & 0.827 & 0.01487 \\
Random & 124.12 & 132.12 & 89.95 & 0.832 & 0.01494 \\
Joint frontier & 123.50 & 131.50 & 90.24 & 0.834 & 0.01499 \\
Safety-only & 124.00 & 132.00 & 90.26 & 0.834 & 0.01499 \\
Exclusion-only & 122.88 & 130.88 & 90.37 & 0.835 & 0.01501 \\
\bottomrule\end{tabular}
\caption{Fresh SummEval scoring at $\alpha=.10$ and primary stopping, $V^+-V^-\le\lceil.02|S_0|\rceil$. Means cover mean and three-member median pooling, $q\in\{.05,.10,.20,.30\}$ and three repetitions (24 runs per method). Total time includes model loading and reference/base scoring. Time/full averages paired ratios. Metric calls include reference/base calls and may process multiple outputs; USD is duration times USD 0.598/hour, not a per-run invoice.}
\label{tab:real-scoring-primary}\end{table*}

\begin{table*}[t]\centering\small
\begin{tabular}{lrrrrr}\toprule
Method & Candidate entries & Metric calls & Total seconds & Time/full & Allocated USD\\\midrule
Full scoring & 1088.00 & 10.00 & 108.18 & 1.000 & 0.01797 \\
Reference & 125.88 & 133.88 & 90.19 & 0.834 & 0.01498 \\
Width & 128.88 & 136.88 & 90.48 & 0.836 & 0.01503 \\
Fixed & 131.00 & 139.00 & 89.73 & 0.829 & 0.01491 \\
Random & 129.25 & 137.25 & 90.20 & 0.834 & 0.01498 \\
Joint frontier & 129.88 & 137.88 & 90.49 & 0.837 & 0.01503 \\
Safety-only & 130.25 & 138.25 & 90.47 & 0.836 & 0.01503 \\
Exclusion-only & 128.50 & 136.50 & 90.45 & 0.836 & 0.01502 \\
\bottomrule\end{tabular}
\caption{Fresh SummEval scoring at $\alpha=.10$ and exact stopping, $V^+=V^-$. Means cover mean and three-member median pooling, $q\in\{.05,.10,.20,.30\}$ and three repetitions (24 runs per method). Total time includes model loading and reference/base scoring. Time/full averages paired ratios. Metric calls include reference/base calls and may process multiple outputs; USD is duration times USD 0.598/hour, not a per-run invoice.}
\label{tab:real-scoring-exact}\end{table*}


At primary stopping, reference priority uses 120.25 candidate-test entries per run on average, versus 1,088 for full scoring. Including the shared reference and base scoring, the mean score-entry totals are 4,392.25 and 5,360, so the total entry count falls far less than candidate-test entries. Mean total time falls from 108.18 to 89.97 seconds. At the same time, metric function calls rise from 10 to 128.25 because the on-demand interface uses single-entry calls. Each of the 24 paired reference runs is faster than its corresponding full run; this is a descriptive hardware-specific comparison, not a claim about all scorer deployments. Joint frontier and safety-only have similar total times on this fresh-score panel. Fewest candidate queries do not imply lowest mean time: fixed order averages 89.51 seconds against 89.97 for reference priority. This 0.46-second gap is smaller than every within-configuration range width for these two policies (Table~\ref{tab:runtime-variation}). The smaller offline joint certificates on some cached-score targets do not by themselves establish a runtime advantage between these policies.

Allocated compute cost multiplies each pipeline duration by the RunPod A40 GPU-plus-container quote of USD 0.598 per hour, recorded on 11 October 2026. Primary full and reference runs average USD 0.01797 and USD 0.01495, respectively. These are duration-based allocations, not observed per-run invoices or total account expenditure: installation, downloads, rescore audits, process startup and idle time also consume rented resources. The complete execution queue took 19,217.75 seconds. Table~\ref{tab:runtime-variation} reports within-configuration runtime variation; the supplement retains all 128 method/pooling/depth/stop summaries. Time ratios pair runs by pooling, depth and repetition. The comparison fixes batch-eight full scoring against sequential entry acquisition; online microbatching, persistent model services and alternative hardware may change the tradeoff.

\begin{table*}[t]\centering\small
\begin{tabular}{lrrrr}\toprule
Method & Primary SD range (s) & Exact SD range (s) & Primary range width (s) & Exact range width (s)\\\midrule
Fixed & 0.34--1.61 & 0.32--1.61 & 0.66--3.13 & 0.62--3.13 \\
Exclusion-only & 0.36--1.34 & 0.36--1.34 & 0.64--2.55 & 0.64--2.56 \\
Joint frontier & 0.19--3.05 & 0.19--3.06 & 0.37--5.33 & 0.37--5.33 \\
Safety-only & 0.26--1.97 & 0.26--1.97 & 0.48--3.44 & 0.49--3.44 \\
Full & 0.33--1.39 & 0.33--1.39 & 0.59--2.76 & 0.59--2.76 \\
Width & 0.32--1.29 & 0.32--1.29 & 0.60--2.30 & 0.60--2.30 \\
Random & 0.66--1.76 & 0.66--1.76 & 1.25--3.20 & 1.26--3.20 \\
Reference & 0.31--1.79 & 0.32--1.79 & 0.61--3.56 & 0.61--3.56 \\
\bottomrule\end{tabular}
\caption{Runtime variation within fixed scoring configurations. Each sample standard deviation (SD) and range width uses three repetitions at one pooling/depth setting. The table reports their minimum and maximum across the eight settings for each method; it does not pool configurations to estimate run noise. Complete per-setting means, ranges and SDs accompany the supplement.}
\label{tab:runtime-variation}\end{table*}



\end{document}
